% notes.tex — Linear Algebra (English version)
% Compile with: xelatex -> biber -> xelatex -> xelatex
\documentclass[a4paper,11pt,twoside,openright]{book}

% Shared preamble (includes polyglossia, biblatex, etc.)
\usepackage[english]{../common/preamble}

% Bibliography file
\addbibresource{../bibliography.bib}

% ============================================================
\usepackage{tikz}
\usetikzlibrary{calc}

\begin{document}

% ============================================================
% Title page
% ============================================================
\begin{titlepage}
\begin{tikzpicture}[remember picture, overlay]
  \fill[left color=blue!15!white, right color=blue!5!white]
    (current page.south west) rectangle (current page.north east);
  \fill[color=orange!70!yellow]
    ([yshift=-2cm]current page.north west) rectangle ([yshift=-2.3cm]current page.north east);
  \fill[color=blue!80!black, rounded corners=8pt]
    ([xshift=3cm, yshift=-4cm]current page.north west) rectangle
    ([xshift=-3cm, yshift=-10cm]current page.north east);
  \node[anchor=center, text=white, font=\Huge\bfseries]
    at ([yshift=-6cm]current page.north) {Linear Algebra};
  \node[anchor=center, text=orange!80!yellow, font=\Large]
    at ([yshift=-7.5cm]current page.north) {Lecture Notes};
  \node[anchor=center, text=white!80, font=\large]
    at ([yshift=-8.8cm]current page.north) {Licence L1 --- 2025--2026};
  \node[anchor=center, text=blue!80!black, font=\Large\itshape]
    at ([yshift=-12cm]current page.north) {Ya\'e Ulrich Gaba};
  \draw[orange!70!yellow, line width=2pt]
    ([xshift=5cm, yshift=-13.5cm]current page.north west) --
    ([xshift=-5cm, yshift=-13.5cm]current page.north east);
  \node[anchor=center, text width=12cm, align=center, text=blue!60!black, font=\itshape]
    at ([yshift=-16cm]current page.north) {``The essence of mathematics lies in its freedom.''\\[4pt]--- Georg Cantor};
  \node[anchor=south, text=gray, font=\small]
    at ([yshift=1.5cm]current page.south) {\today};
\end{tikzpicture}
\end{titlepage}

\frontmatter

% ============================================================
% Table of contents
% ============================================================
\tableofcontents

% ============================================================
% Preface
% ============================================================
\chapter*{Preface}
\addcontentsline{toc}{chapter}{Preface}
\markboth{Preface}{Preface}

Linear algebra is one of the most fundamental and pervasive branches of modern mathematics. Its core concepts --- vector spaces, linear maps, matrices, determinants, eigenvalues --- form the foundation upon which vast areas of analysis, geometry, physics, computer science, and engineering are built.

These course notes are intended for first- and second-year undergraduate students in mathematics or engineering programs. They assume a solid background in high-school mathematics (calculus, functions, trigonometry) as well as elementary notions of set theory and logic.

\medskip

\noindent\textbf{Philosophy.}
We have striven to maintain a constant balance between \emph{algebraic rigor} and \emph{geometric intuition}. Every abstract concept is illustrated with concrete examples, figures, and applications. All proofs are complete: we believe that learning to read and write proofs is just as important as mastering computational techniques.

\medskip

\noindent\textbf{Organization.}
The notes consist of nine chapters, whose logical dependencies are shown below:

\begin{center}
\begin{tikzpicture}[
  node distance=1.2cm and 2cm,
  block/.style={rectangle, draw, rounded corners, minimum width=2.8cm,
    minimum height=0.8cm, text centered, font=\small},
  arrow/.style={-{Stealth[length=5pt]}, thick}
]
  \node[block] (ch0) {Ch.~0 Preliminaries};
  \node[block, below=of ch0] (ch1) {Ch.~1 Vector Spaces};
  \node[block, below left=1.2cm and 0.5cm of ch1] (ch2) {Ch.~2 Linear Maps};
  \node[block, below right=1.2cm and 0.5cm of ch1] (ch3) {Ch.~3 Linear Systems};
  \node[block, below=2.8cm of ch1] (ch4) {Ch.~4 Determinants};
  \node[block, below=of ch4] (ch5) {Ch.~5 Eigenvalues};
  \node[block, below left=1.2cm and 0.5cm of ch5] (ch6) {Ch.~6 Inner Products};
  \node[block, below right=1.2cm and 0.5cm of ch5] (ch8) {Ch.~8 Jordan Form};
  \node[block, below=2.8cm of ch5] (ch7) {Ch.~7 Spectral Theorem};

  \draw[arrow] (ch0) -- (ch1);
  \draw[arrow] (ch1) -- (ch2);
  \draw[arrow] (ch1) -- (ch3);
  \draw[arrow] (ch2) -- (ch4);
  \draw[arrow] (ch3) -- (ch4);
  \draw[arrow] (ch4) -- (ch5);
  \draw[arrow] (ch5) -- (ch6);
  \draw[arrow] (ch5) -- (ch8);
  \draw[arrow] (ch6) -- (ch7);
\end{tikzpicture}
\end{center}

Chapter~0 (Preliminaries) may be skimmed quickly by students already familiar with the basic concepts. Chapter~8 (Jordan Normal Form) is an advanced chapter, intended for students who wish to explore the subject in greater depth.

\medskip

\noindent\textbf{Exercises.}
Each chapter concludes with a set of exercises graded by difficulty:
\begin{itemize}
  \item \basic{} --- basic exercises (direct verification of definitions),
  \item \intermediate{} --- intermediate exercises (applications of theorems),
  \item \challenging{} --- challenging exercises (synthesis problems, open-ended questions).
\end{itemize}

\medskip

\noindent\textbf{Historical note.}
Linear algebra, as we know it today, developed gradually from the 18th to the 20th century. Major contributions include: Gauss's work on elimination (1809), Cauchy's theory of determinants (1815), the discovery of eigenvalues by Cauchy and Sylvester, Jordan's canonical forms (1870), and the axiomatic formalization of vector spaces by Peano (1888) and Banach (1920).

\vfill
\begin{flushright}
\textit{February 2026}
\end{flushright}

% ============================================================
% Notation
% ============================================================
\chapter*{Notation}
\addcontentsline{toc}{chapter}{Notation}
\markboth{Notation}{Notation}

We collect here the notation used throughout these notes. The reader is encouraged to refer back to this section as needed.

\bigskip

\begin{longtable}{@{}p{3cm}p{10cm}@{}}
\toprule
\textbf{Symbol} & \textbf{Meaning} \\
\midrule
\endhead

\multicolumn{2}{@{}l}{\textbf{Number sets}} \\[4pt]
$\N$ & Natural numbers $\{0,1,2,\ldots\}$ \\
$\Z$ & Integers \\
$\Q$ & Rational numbers \\
$\R$ & Real numbers \\
$\C$ & Complex numbers \\
$\K$ & Base field (usually $\R$ or $\C$) \\
$\F_p$ & Finite field with $p$ elements \\[6pt]

\multicolumn{2}{@{}l}{\textbf{Spaces and structures}} \\[4pt]
$\K^n$ & Space of $n$-tuples over $\K$ \\
$\Mat_{n,p}(\K)$ & Space of $n\times p$ matrices with entries in $\K$ \\
$\Mat_n(\K)$ & Square matrices of size $n\times n$ \\
$\GL_n(\K)$ & General linear group (invertible matrices) \\
$\cP_n(\K)$ & Space of polynomials of degree $\leq n$ \\
$\cL(E,F)$ & Space of linear maps from $E$ to $F$ \\
$\End(E)$ & Endomorphisms of $E$ ($=\cL(E,E)$) \\[6pt]

\multicolumn{2}{@{}l}{\textbf{Linear maps and matrices}} \\[4pt]
$\Ker f$ & Kernel of $f$ \\
$\im f$ & Image (range) of $f$ \\
$\rank(f)$, $\rank(A)$ & Rank of a linear map or matrix \\
$\Id$, $I_n$ & Identity map, $n\times n$ identity matrix \\
$\transpose{A}$ & Transpose of $A$ \\
$\hermit{A}$ & Conjugate transpose (adjoint) of $A$ \\
$\inv{A}$ & Inverse of $A$ \\
$\tr(A)$ & Trace of $A$ \\
$\det(A)$ & Determinant of $A$ \\[6pt]

\multicolumn{2}{@{}l}{\textbf{Vector spaces}} \\[4pt]
$\Span(v_1,\ldots,v_k)$ & Subspace spanned by $v_1,\ldots,v_k$ \\
$\dim E$ & Dimension of $E$ \\
$\codim F$ & Codimension of $F$ in $E$ \\
$E\oplus F$ & Direct sum \\
$E/F$ & Quotient space \\
$E^*$ & Dual of $E$ \\
$\cB=(e_1,\ldots,e_n)$ & Basis of $E$ \\[6pt]

\multicolumn{2}{@{}l}{\textbf{Inner products and norms}} \\[4pt]
$\inner{u,v}$ & Inner product of $u$ and $v$ \\
$\norm{v}$ & Norm of $v$ \\
$u\perp v$ & $u$ and $v$ are orthogonal \\
$F^\perp$ & Orthogonal complement of $F$ \\[6pt]

\multicolumn{2}{@{}l}{\textbf{Eigenvalues}} \\[4pt]
$\Spec(f)$ & Spectrum of $f$ (set of eigenvalues) \\
$E_\lambda$ & Eigenspace associated with $\lambda$ \\
$\chi_f(\lambda)$ & Characteristic polynomial of $f$ \\
$\mu_f(\lambda)$ & Minimal polynomial of $f$ \\[6pt]

\multicolumn{2}{@{}l}{\textbf{Miscellaneous}} \\[4pt]
$\deff$ & Equality by definition \\
$\delta_{ij}$ & Kronecker delta \\
$\sgn(\sigma)$ & Sign of a permutation $\sigma$ \\
$\mathfrak{S}_n$ & Symmetric group on $\{1,\ldots,n\}$ \\

\bottomrule
\end{longtable}

\mainmatter

% ============================================================
% Chapters
% ============================================================

% ch0-preliminaries.tex — Chapter 0: Preliminaries
% Review of sets, mappings, relations, algebraic structures, and proof techniques.

\chapter{Preliminaries}\label{ch:preliminaries}

Linear algebra does not exist in a vacuum. It rests upon a small but
essential toolkit from set theory, logic, and abstract algebra. The purpose
of this chapter is to establish that toolkit clearly and concisely so that
the rest of the course may proceed without interruption.

Most of the material here should be familiar from earlier courses. We
therefore adopt a brisk pace: definitions are stated precisely, key
examples are worked through, and proofs are included only when they
illuminate a technique that will recur later. The reader who is
comfortable with all five sections may safely skip ahead to

ef{ch:vector-spaces}; the reader who is not should work through the
exercises at the end before continuing.


% ============================================================
\section{Sets}\label{sec:sets}
% ============================================================

\begin{definition}{Set}{def:set}
A \emph{set}\index{set} is a collection of distinct objects, called its
\emph{elements}. We write $x\in A$ to mean that $x$ is an element of~$A$,
and $x\notin A$ otherwise.
\end{definition}

We recall the standard number sets that will pervade every chapter:
\[
  \N \;\subset\; \Z \;\subset\; \Q \;\subset\; \R \;\subset\; \C.
\]
Here $\N = \set{0,1,2,\ldots}$ includes zero by convention.

\begin{definition}{Subset and set equality}{def:subset}
Let $A$ and $B$ be sets.
\begin{enumerate}[label=(\roman*)]
  \item $A$ is a \emph{subset}\index{subset} of $B$, written $A\subseteq B$,
    if every element of $A$ belongs to $B$.
  \item $A = B$ if and only if $A\subseteq B$ and $B\subseteq A$.
  \item $A$ is a \emph{proper subset}\index{proper subset} of $B$, written
    $A\subsetneq B$, if $A\subseteq B$ and $A\neq B$.
\end{enumerate}
\end{definition}

\begin{definition}{Set operations}{def:set-ops}
Let $A$ and $B$ be subsets of a universal set~$U$.
\begin{enumerate}[label=(\roman*)]
  \item \textbf{Union:} $A\cup B \deff \setcond{x\in U}{x\in A \text{ or } x\in B}$.
  \item \textbf{Intersection:} $A\cap B \deff \setcond{x\in U}{x\in A \text{ and } x\in B}$.
  \item \textbf{Difference:} $A\setminus B \deff \setcond{x\in U}{x\in A \text{ and } x\notin B}$.
  \item \textbf{Complement:} $A^c \deff U\setminus A$.
  \item \textbf{Cartesian product:} $A\times B \deff \setcond{(a,b)}{a\in A,\; b\in B}$.
  \item \textbf{Power set:} $\mathcal{P}(A) \deff \setcond{S}{S\subseteq A}$.
\end{enumerate}
\end{definition}

\begin{example}{Cartesian products}{ex:cartesian}
$\R^2 = \R\times\R$ is the Euclidean plane. More generally,
$\R^n = \underbrace{\R\times\cdots\times\R}_{n}$ is the set of all
$n$-tuples of real numbers, the primary setting for Chapters~1--3.
\end{example}


% ============================================================
\section{Mappings (functions)}\label{sec:mappings}
% ============================================================

\begin{definition}{Mapping}{def:mapping}
A \emph{mapping}\index{mapping} (or \emph{function}) from a set $A$ to a
set $B$ is a rule $f\colon A\to B$ that assigns to each element $a\in A$
exactly one element $f(a)\in B$. The set $A$ is the \emph{domain} and $B$
is the \emph{codomain}.
\end{definition}

\begin{definition}{Image and preimage}{def:image-preimage}
Let $f\colon A\to B$ be a mapping.
\begin{enumerate}[label=(\roman*)]
  \item The \emph{image}\index{image} of a subset $S\subseteq A$ is
    $f(S) \deff \setcond{f(a)}{a\in S}$.
  \item The \emph{image} (or \emph{range}) of $f$ is $\im f \deff f(A)$.
  \item The \emph{preimage}\index{preimage} of a subset $T\subseteq B$ is
    $f^{-1}(T) \deff \setcond{a\in A}{f(a)\in T}$.
\end{enumerate}
\end{definition}

\begin{definition}{Injectivity, surjectivity, bijectivity}{def:inj-surj-bij}
Let $f\colon A\to B$.
\begin{enumerate}[label=(\roman*)]
  \item $f$ is \emph{injective}\index{injective} (one-to-one) if
    $f(a_1)=f(a_2)$ implies $a_1=a_2$.
  \item $f$ is \emph{surjective}\index{surjective} (onto) if
    $\im f = B$, i.e.\ for every $b\in B$ there exists $a\in A$ with $f(a)=b$.
  \item $f$ is \emph{bijective}\index{bijective} if it is both injective
    and surjective. In this case $f$ has an \emph{inverse}
    $\inv{f}\colon B\to A$ satisfying $\inv{f}\circ f = \Id_A$ and
    $f\circ\inv{f} = \Id_B$.
\end{enumerate}
\end{definition}

\begin{example}{A classical injection and surjection}{ex:inj-surj}
\begin{enumerate}[label=(\alph*)]
  \item The map $f\colon\R\to\R$, $f(x)=x^2$, is neither injective
    (since $f(-1)=f(1)=1$) nor surjective (since $-1\notin\im f$).
  \item Restricting the domain and codomain: the map
    $g\colon [0,\infty)\to [0,\infty)$, $g(x)=x^2$, is bijective, with
    inverse $\inv{g}(y)=\sqrt{y}$.
\end{enumerate}
\end{example}

\begin{definition}{Composition}{def:composition}
If $f\colon A\to B$ and $g\colon B\to C$, their \emph{composition}\index{composition}
is the mapping $g\circ f\colon A\to C$ defined by $(g\circ f)(a) \deff g(f(a))$.
Composition is associative: $h\circ(g\circ f)=(h\circ g)\circ f$.
\end{definition}

\begin{proposition}{Composition preserves injectivity and surjectivity}{prop:comp-inj-surj}
Let $f\colon A\to B$ and $g\colon B\to C$.
\begin{enumerate}[label=(\roman*)]
  \item If $f$ and $g$ are injective, then $g\circ f$ is injective.
  \item If $f$ and $g$ are surjective, then $g\circ f$ is surjective.
  \item If $g\circ f$ is injective, then $f$ is injective.
  \item If $g\circ f$ is surjective, then $g$ is surjective.
\end{enumerate}
\end{proposition}

\begin{proof}
We prove (i); the others are similar exercises (see

ef{exo:exo:comp-inj-surj}).

Suppose $f$ and $g$ are injective and $(g\circ f)(a_1) = (g\circ f)(a_2)$.
Then $g(f(a_1)) = g(f(a_2))$. Since $g$ is injective,
$f(a_1) = f(a_2)$. Since $f$ is injective, $a_1 = a_2$.
\end{proof}


% ============================================================
\section{Relations}\label{sec:relations}
% ============================================================

\begin{definition}{Binary relation}{def:relation}
A \emph{binary relation}\index{relation!binary} on a set $A$ is a subset
$\mathcal{R}\subseteq A\times A$. We write $a\;\mathcal{R}\;b$ to mean
$(a,b)\in\mathcal{R}$.
\end{definition}

\subsection{Equivalence relations}\label{subsec:equiv}

\begin{definition}{Equivalence relation}{def:equiv-rel}
A binary relation $\sim$ on $A$ is an \emph{equivalence
relation}\index{relation!equivalence} if it is:
\begin{enumerate}[label=(\roman*)]
  \item \textbf{Reflexive:} $a\sim a$ for all $a\in A$.
  \item \textbf{Symmetric:} $a\sim b$ implies $b\sim a$.
  \item \textbf{Transitive:} $a\sim b$ and $b\sim c$ imply $a\sim c$.
\end{enumerate}
\end{definition}

\begin{definition}{Equivalence class and quotient set}{def:equiv-class}
Let $\sim$ be an equivalence relation on~$A$. The \emph{equivalence
class}\index{equivalence class} of $a\in A$ is
\[
  [a] \deff \setcond{b\in A}{b\sim a}.
\]
The \emph{quotient set}\index{quotient set} is
$A/{\sim} \;\deff\; \setcond{[a]}{a\in A}$.
\end{definition}

\begin{example}{Congruence modulo $n$}{ex:congruence}
Fix an integer $n\geq 1$. On $\Z$, define $a\sim b$ if and only if
$n\mid (a-b)$ (i.e.\ $a\equiv b\pmod{n}$). This is an equivalence
relation with $n$ equivalence classes: $[0],[1],\ldots,[n-1]$. The
quotient set is denoted $\Z/n\Z$ (or simply $\Z_n$).

For $n=3$: $[0]=\set{\ldots,-6,-3,0,3,6,\ldots}$,
$[1]=\set{\ldots,-5,-2,1,4,7,\ldots}$,
$[2]=\set{\ldots,-4,-1,2,5,8,\ldots}$.
\end{example}

\begin{remark}{Equivalence classes partition the set}{rem:partition}
The equivalence classes of any equivalence relation on $A$ form a
\emph{partition} of $A$: they are nonempty, pairwise disjoint, and their
union is~$A$. Conversely, every partition of $A$ defines an equivalence
relation. This fundamental observation will reappear when we study quotient
vector spaces in 
ef{ch:linear-maps}.
\end{remark}

\subsection{Order relations}\label{subsec:order}

\begin{definition}{Partial order}{def:partial-order}
A binary relation $\leq$ on $A$ is a \emph{partial
order}\index{relation!order} if it is:
\begin{enumerate}[label=(\roman*)]
  \item \textbf{Reflexive:} $a\leq a$ for all $a\in A$.
  \item \textbf{Antisymmetric:} $a\leq b$ and $b\leq a$ imply $a = b$.
  \item \textbf{Transitive:} $a\leq b$ and $b\leq c$ imply $a\leq c$.
\end{enumerate}
If in addition every two elements are comparable ($a\leq b$ or $b\leq a$
for all $a,b\in A$), then $\leq$ is a \emph{total order}\index{total order}.
\end{definition}

\begin{example}{Orders on familiar sets}{ex:orders}
\begin{enumerate}[label=(\alph*)]
  \item $(\R,\leq)$ is totally ordered.
  \item $(\mathcal{P}(S),\subseteq)$ is partially ordered for any set $S$,
    but not totally ordered when $\card{S}\geq 2$ (e.g.\ $\set{1}$ and
    $\set{2}$ are incomparable in $\mathcal{P}(\set{1,2})$).
  \item Divisibility on $\N^*$: $a\mid b$ defines a partial order.
\end{enumerate}
\end{example}


% ============================================================
\section{Algebraic structures}\label{sec:algebraic-structures}
% ============================================================

In this section we introduce the hierarchy \emph{group} $\to$
\emph{ring} $\to$ \emph{field}. The central notion for this course is
that of a \emph{field}, since linear algebra is the study of vector
spaces \emph{over a field}. However, groups and rings appear naturally
in the theory (symmetry groups, polynomial rings, matrix rings), so a
brief treatment is warranted.

% ----- Groups -----
\subsection{Groups}\label{subsec:groups}

\begin{definition}{Group}{def:group}
A \emph{group}\index{group} is a pair $(G,\star)$ where $G$ is a nonempty
set and $\star\colon G\times G\to G$ is a binary operation satisfying:
\begin{enumerate}[label=\textbf{G\arabic*.}]
  \item \textbf{Associativity:} $(a\star b)\star c = a\star(b\star c)$
    for all $a,b,c\in G$.
  \item \textbf{Identity:} there exists $e\in G$ such that
    $e\star a = a\star e = a$ for all $a\in G$.
  \item \textbf{Inverses:} for each $a\in G$ there exists $a'\in G$ such
    that $a\star a' = a'\star a = e$.
\end{enumerate}
If in addition $a\star b = b\star a$ for all $a,b\in G$, the group is
called \emph{abelian}\index{group!abelian} (or \emph{commutative}).
\end{definition}

\begin{example}{Basic examples of groups}{ex:groups}
\begin{enumerate}[label=(\alph*)]
  \item $(\Z,+)$, $(\Q,+)$, $(\R,+)$, $(\C,+)$ are abelian groups with
    identity $0$ and inverse $-a$.
  \item $(\Q^*,\times)$, $(\R^*,\times)$, $(\C^*,\times)$ are abelian
    groups with identity $1$ and inverse $1/a$. (Here
    $\Q^*=\Q\setminus\set{0}$, etc.)
  \item $(\GL_n(\R),\times)$, the set of invertible $n\times n$ real
    matrices under matrix multiplication, is a group (non-abelian for
    $n\geq 2$). This group will play a central role starting from
    
ef{ch:determinants}.
  \item $(\Z/n\Z,+)$ is an abelian group of order~$n$.
\end{enumerate}
\end{example}

\begin{proposition}{Uniqueness of identity and inverses}{prop:unique-id}
In any group $(G,\star)$:
\begin{enumerate}[label=(\roman*)]
  \item The identity element is unique.
  \item Each element has a unique inverse.
\end{enumerate}
\end{proposition}

\begin{proof}
(i)~Suppose $e$ and $e'$ are both identities. Then $e = e\star e' = e'$.

(ii)~Suppose $a'$ and $a''$ are both inverses of~$a$. Then
$a' = a'\star e = a'\star(a\star a'') = (a'\star a)\star a'' = e\star a'' = a''$.
\end{proof}

% ----- Rings -----
\subsection{Rings}\label{subsec:rings}

\begin{definition}{Ring}{def:ring}
A \emph{ring}\index{ring} is a triple $(R,+,\cdot\,)$ where:
\begin{enumerate}[label=\textbf{R\arabic*.}]
  \item $(R,+)$ is an abelian group (with identity $0_R$).
  \item Multiplication is associative:
    $(a\cdot b)\cdot c = a\cdot(b\cdot c)$.
  \item There exists a multiplicative identity $1_R$ with
    $1_R\cdot a = a\cdot 1_R = a$ for all~$a$.
  \item Distributivity:
    $a\cdot(b+c) = a\cdot b + a\cdot c$ and
    $(a+b)\cdot c = a\cdot c + b\cdot c$.
\end{enumerate}
If $a\cdot b = b\cdot a$ for all $a,b\in R$, the ring is
\emph{commutative}\index{ring!commutative}.
\end{definition}

\begin{example}{Rings}{ex:rings}
\begin{enumerate}[label=(\alph*)]
  \item $(\Z,+,\times)$ is a commutative ring.
  \item The polynomial ring $\R[X]$ is a commutative ring.
  \item $(\Mat_n(\R),+,\times)$, the set of $n\times n$ real matrices, is
    a ring that is \emph{not} commutative for $n\geq 2$.
  \item $(\Z/n\Z,+,\times)$ is a commutative ring.
\end{enumerate}
\end{example}

% ----- Fields -----
\subsection{Fields}\label{subsec:fields}

The field axioms are the single most important algebraic prerequisite for
this course. A field is, roughly speaking, a set where one can add,
subtract, multiply, and divide (except by zero) and where all the usual
algebraic rules hold.

\begin{definition}{Field}{def:field}
A \emph{field}\index{field} is a triple $(\K,+,\cdot\,)$ satisfying the
following axioms. Let $a,b,c$ denote arbitrary elements of~$\K$.

\medskip
\noindent\textbf{Addition axioms:}
\begin{enumerate}[label=\textbf{A\arabic*.}]
  \item \textbf{Closure:} $a+b\in\K$.
  \item \textbf{Associativity:} $(a+b)+c = a+(b+c)$.
  \item \textbf{Commutativity:} $a+b = b+a$.
  \item \textbf{Identity:} there exists $0\in\K$ with $a+0=a$.
  \item \textbf{Inverses:} for each $a$ there exists $-a\in\K$ with
    $a+(-a)=0$.
\end{enumerate}

\medskip
\noindent\textbf{Multiplication axioms:}
\begin{enumerate}[label=\textbf{M\arabic*.}]
  \item \textbf{Closure:} $a\cdot b\in\K$.
  \item \textbf{Associativity:} $(a\cdot b)\cdot c = a\cdot(b\cdot c)$.
  \item \textbf{Commutativity:} $a\cdot b = b\cdot a$.
  \item \textbf{Identity:} there exists $1\in\K$, $1\neq 0$, with
    $a\cdot 1 = a$.
  \item \textbf{Inverses:} for each $a\neq 0$ there exists $a^{-1}\in\K$
    with $a\cdot a^{-1}=1$.
\end{enumerate}

\medskip
\noindent\textbf{Distributive law:}
\begin{enumerate}[label=\textbf{D\arabic*.}]
  \item $a\cdot(b+c)=a\cdot b + a\cdot c$.
\end{enumerate}
\end{definition}

\begin{remark}{Field = commutative division ring}{rem:field-division-ring}
Equivalently, a field is a commutative ring in which every nonzero element
has a multiplicative inverse. The requirement $1\neq 0$ ensures that the
field has at least two elements.
\end{remark}

\begin{proposition}{Elementary consequences of the field axioms}{prop:field-consequences}
Let $\K$ be a field and $a,b\in\K$. Then:
\begin{enumerate}[label=(\roman*)]
  \item $a\cdot 0 = 0$.
  \item $(-1)\cdot a = -a$.
  \item $a\cdot b = 0$ implies $a=0$ or $b=0$
    (a field has no zero divisors).
\end{enumerate}
\end{proposition}

\begin{proof}
(i)~$a\cdot 0 = a\cdot(0+0)=a\cdot 0 + a\cdot 0$. Adding $-(a\cdot 0)$ to
both sides gives $0 = a\cdot 0$.

(ii)~$a + (-1)\cdot a = 1\cdot a + (-1)\cdot a = (1+(-1))\cdot a = 0\cdot a
= 0$, so $(-1)\cdot a$ is the additive inverse of~$a$.

(iii)~If $a\neq 0$, multiply both sides of $ab=0$ by $a^{-1}$:
$b = 1\cdot b = (a^{-1}a)b = a^{-1}(ab) = a^{-1}\cdot 0 = 0$.
\end{proof}

We now examine the three families of fields that appear most frequently in
linear algebra.

% --- The real numbers ---
\begin{example}{The real numbers $\R$}{ex:reals}
$(\R,+,\times)$ is the prototypical field, and the default setting for
most of undergraduate linear algebra. In addition to the field axioms,
$\R$ carries a total order compatible with the field operations and
satisfies the \emph{completeness axiom} (every nonempty bounded-above
subset has a least upper bound). These extra properties are not required
by the field axioms and are not shared by all fields.
\end{example}

% --- The complex numbers ---
\begin{example}{The complex numbers $\C$}{ex:complex}
Define $\C \deff \setcond{a+bi}{a,b\in\R}$ where $i^2=-1$. Addition and
multiplication are given by:
\begin{align*}
  (a+bi)+(c+di) &= (a+c)+(b+d)i, \\
  (a+bi)(c+di)  &= (ac-bd)+(ad+bc)i.
\end{align*}
The additive identity is $0$, the multiplicative identity is $1$, and the
inverse of $a+bi\neq 0$ is
\[
  (a+bi)^{-1} = \frac{a}{a^2+b^2} - \frac{b}{a^2+b^2}\,i.
\]
A key advantage of $\C$ over $\R$ is that every polynomial of degree
$n\geq 1$ has exactly $n$ roots (counted with multiplicity) --- this is
the \emph{Fundamental Theorem of Algebra}. As a consequence, eigenvalue
theory is cleaner over $\C$, as we shall see in

ef{ch:eigenvalues,ch:spectral}.
\end{example}

% --- Finite fields ---
\begin{example}{Finite fields $\F_p$}{ex:finite-field}
Let $p$ be a prime number. On $\Z/p\Z = \set{[0],[1],\ldots,[p-1]}$,
define addition and multiplication by:
\[
  [a]+[b]\deff[a+b],\qquad [a]\cdot[b]\deff[a\cdot b],
\]
where the operations on the right are performed in $\Z$ and the result is
reduced modulo~$p$. These operations are well-defined (independent of the
choice of representatives) and make $\Z/p\Z$ into a field, denoted
$\F_p$\index{finite field}.

The key point is that $p$ being prime guarantees the existence of
multiplicative inverses: if $[a]\neq[0]$, then $\gcd(a,p)=1$, so by
B\'{e}zout's identity there exist $u,v\in\Z$ with $au+pv=1$, and hence
$[a]\cdot[u]=[1]$.

\medskip\noindent
\textbf{Multiplication table of $\F_5$:}

\smallskip
\begin{center}
\begin{tabular}{c|ccccc}
$\times$ & $[0]$ & $[1]$ & $[2]$ & $[3]$ & $[4]$ \\
\midrule
$[0]$ & $[0]$ & $[0]$ & $[0]$ & $[0]$ & $[0]$ \\
$[1]$ & $[0]$ & $[1]$ & $[2]$ & $[3]$ & $[4]$ \\
$[2]$ & $[0]$ & $[2]$ & $[4]$ & $[1]$ & $[3]$ \\
$[3]$ & $[0]$ & $[3]$ & $[1]$ & $[4]$ & $[2]$ \\
$[4]$ & $[0]$ & $[4]$ & $[3]$ & $[2]$ & $[1]$ \\
\end{tabular}
\end{center}

\smallskip
From the table we read off: $[2]^{-1}=[3]$, $[3]^{-1}=[2]$,
$[4]^{-1}=[4]$.
\end{example}

\begin{remark}{Why primality matters}{rem:why-prime}
If $n$ is not prime, $\Z/n\Z$ is \emph{not} a field. For instance, in
$\Z/6\Z$ we have $[2]\cdot[3]=[0]$ with $[2]\neq[0]$ and $[3]\neq[0]$,
violating the no-zero-divisors property. More generally, finite fields of
order $q$ exist if and only if $q$ is a prime power $p^k$; the
construction for $k\geq 2$ requires quotient rings of polynomial rings and
is beyond our scope.
\end{remark}

\begin{definition}{Characteristic of a field}{def:characteristic}
The \emph{characteristic}\index{characteristic} of a field $\K$, denoted
$\Char(\K)$, is the smallest positive integer $p$ such that
$\underbrace{1+1+\cdots+1}_{p} = 0$ in $\K$. If no such integer exists,
we set $\Char(\K)=0$.
\end{definition}

\begin{example}{Characteristics}{ex:char}
$\Char(\Q) = \Char(\R) = \Char(\C) = 0$, while $\Char(\F_p) = p$.
Throughout this course, \emph{unless otherwise stated, $\K$ denotes a
field of characteristic zero} (typically $\R$ or $\C$).
\end{example}

\begin{remark}{Convention: the symbol $\K$}{rem:convention-K}
In these notes, the letter $\K$ always denotes the base field. Most
results in Chapters~1--5 hold over an arbitrary field; we shall
explicitly note when additional hypotheses on $\K$ (such as
$\Char(\K)=0$, or algebraic closure) are needed.
\end{remark}


% ============================================================
\section{Proof techniques}\label{sec:proofs}
% ============================================================

Proofs are the backbone of mathematics. In this section we review the
principal strategies. The reader should view these not as abstract logical
rules but as practical tools that will be used repeatedly throughout the
course.

\subsection{Direct proof}\label{subsec:direct}

To prove ``if $P$ then $Q$'', assume $P$ is true and deduce $Q$ through a
chain of logical steps.

\begin{example}{A direct proof}{ex:direct-proof}
\emph{Claim:} If $n$ is an even integer, then $n^2$ is even.

\emph{Proof:} Assume $n$ is even. Then $n=2k$ for some $k\in\Z$. Hence
$n^2 = 4k^2 = 2(2k^2)$, which is even. \qed
\end{example}

\subsection{Proof by contrapositive}\label{subsec:contrapositive}

The statement ``if $P$ then $Q$'' is logically equivalent to ``if
$\lnot Q$ then $\lnot P$''. Sometimes the contrapositive is easier to
prove.

\begin{example}{A contrapositive proof}{ex:contrapositive-proof}
\emph{Claim:} If $n^2$ is odd, then $n$ is odd.

\emph{Proof (contrapositive):} We prove: if $n$ is even, then $n^2$ is
even. This is exactly the statement proved in

ef{ex:ex:direct-proof}. \qed
\end{example}

\subsection{Proof by contradiction}\label{subsec:contradiction}

To prove a statement $P$, assume $\lnot P$ and derive a logical
contradiction.

\begin{example}{Irrationality of $\sqrt{2}$}{ex:sqrt2}
\emph{Claim:} $\sqrt{2}\notin\Q$.

\emph{Proof:} Suppose for contradiction that $\sqrt{2}=a/b$ with
$a,b\in\Z$, $b\neq 0$, and $\gcd(a,b)=1$. Then $2b^2=a^2$, so $a^2$ is
even, hence $a$ is even: $a=2k$. Then $2b^2=4k^2$, so $b^2=2k^2$, hence
$b$ is even. But then $\gcd(a,b)\geq 2$, contradicting
$\gcd(a,b)=1$. \qed
\end{example}

\subsection{Proof by induction}\label{subsec:induction}

\begin{definition}{Principle of mathematical induction}{def:induction}
Let $P(n)$ be a statement depending on $n\in\N$. If
\begin{enumerate}[label=(\roman*)]
  \item \textbf{Base case:} $P(n_0)$ is true for some $n_0\in\N$, and
  \item \textbf{Inductive step:} for every $n\geq n_0$, $P(n)$ implies
    $P(n+1)$,
\end{enumerate}
then $P(n)$ is true for all $n\geq n_0$.
\end{definition}

\begin{remark}{Strong induction}{rem:strong-induction}
In the \emph{strong} form of induction, the inductive hypothesis in step
(ii) is strengthened: one assumes $P(k)$ for all $n_0\leq k\leq n$ and
deduces $P(n+1)$. The two forms are logically equivalent.
\end{remark}

\begin{example}{Sum of the first $n$ integers}{ex:sum-integers}
\emph{Claim:} For all $n\geq 1$,
$\displaystyle\sum_{k=1}^{n}k = \frac{n(n+1)}{2}$.

\emph{Proof by induction.}

\textbf{Base case ($n=1$):}
$\displaystyle\sum_{k=1}^{1}k = 1 = \frac{1\cdot 2}{2}$. \checkmark

\textbf{Inductive step:} Assume the formula holds for some $n\geq 1$.
Then
\[
  \sum_{k=1}^{n+1}k = \left(\sum_{k=1}^{n}k\right) + (n+1)
  = \frac{n(n+1)}{2} + (n+1)
  = \frac{n(n+1)+2(n+1)}{2}
  = \frac{(n+1)(n+2)}{2}.
\]
This is the formula with $n$ replaced by $n+1$. \qed
\end{example}

\begin{example}{A preview: dimension of a subspace}{ex:induction-linalg}
Induction is the primary tool for proving results about finite-dimensional
vector spaces. For instance, one proves by induction on $\dim V$ that
every subspace of a finite-dimensional vector space is itself
finite-dimensional (see 
ef{ch:vector-spaces}). The base case is
$\dim V = 0$ (the trivial space), and the inductive step reduces a space
of dimension $n+1$ to one of dimension~$n$.
\end{example}

\subsection{Proof by double inclusion}\label{subsec:double-inclusion}

To show that two sets $A$ and $B$ are equal, one proves $A\subseteq B$
and $B\subseteq A$ separately. This technique is used constantly in
linear algebra (e.g.\ showing that $\Ker f = \set{\vzero}$).

\subsection{Proof by double counting / dimension argument}\label{subsec:dim-argument}

Later in the course, a powerful and characteristically linear-algebraic
proof technique will emerge: the \emph{dimension argument}. To show that
two subspaces $U$ and $W$ are equal, one often shows $U\subseteq W$ (or
$W\subseteq U$) and $\dim U = \dim W$. This will be formalized in

ef{ch:vector-spaces}.


% ============================================================
\section{Exercises}\label{sec:ch0-exercises}
% ============================================================

\begin{exercise}{Set operations \basic}{exo:set-ops}
Let $A=\set{1,2,3,4,5}$ and $B=\set{3,4,5,6,7}$.
Compute $A\cup B$, $A\cap B$, $A\setminus B$, $B\setminus A$, and
$A\times\set{0,1}$.
\end{exercise}

\begin{exercise}{Injectivity and surjectivity \basic}{exo:inj-surj}
For each of the following mappings, determine whether it is injective,
surjective, bijective, or none.
\begin{enumerate}[label=(\alph*)]
  \item $f\colon\Z\to\Z$, $f(n)=2n$.
  \item $g\colon\Z\to\Z$, $g(n)=n+3$.
  \item $h\colon\R\to\R$, $h(x)=x^3-x$.
  \item $\vphi\colon\R\to\R_{>0}$, $\vphi(x)=e^x$.
\end{enumerate}
\end{exercise}

\begin{exercise}{Composition and invertibility \intermediate}{exo:comp-inj-surj}
Prove parts (ii)--(iv) of 
ef{prop:prop:comp-inj-surj}.
\end{exercise}

\begin{exercise}{Equivalence relations \basic}{exo:equiv-rel}
On $\Z\times(\Z\setminus\set{0})$, define $(a,b)\sim(c,d)$ if and only if
$ad=bc$. Verify that $\sim$ is an equivalence relation. What familiar set
does the quotient $(\Z\times(\Z\setminus\set{0}))/{\sim}$ represent?
\end{exercise}

\begin{exercise}{Verifying group axioms \basic}{exo:group-axioms}
Verify that $(\Z/5\Z\setminus\set{[0]},\times)$ is an abelian group. What
is the inverse of $[3]$?
\end{exercise}

\begin{exercise}{Why $\Z$ is not a field \basic}{exo:Z-not-field}
Which field axiom fails for $(\Z,+,\times)$? Give a specific
counterexample.
\end{exercise}

\begin{exercise}{Arithmetic in $\F_7$ \intermediate}{exo:F7}
\begin{enumerate}[label=(\alph*)]
  \item Write out the complete multiplication table of $\F_7$.
  \item Find $[3]^{-1}$, $[5]^{-1}$, and $[6]^{-1}$ in $\F_7$.
  \item Solve the equation $[3]x + [2] = [5]$ in $\F_7$.
\end{enumerate}
\end{exercise}

\begin{exercise}{Complex arithmetic \basic}{exo:complex-arith}
Let $z=3+4i$ and $w=1-2i$. Compute $z+w$, $zw$, $\abs{z}$, $z^{-1}$, and
$z/w$.
\end{exercise}

\begin{exercise}{Field of four elements \challenging}{exo:F4}
We know $\F_p$ exists for every prime $p$. A field of order $4$ also
exists, but it is \emph{not} $\Z/4\Z$ (which has zero divisors). Consider
the set $\F_4=\set{0,1,\alpha,\alpha+1}$ with the rule $\alpha^2+\alpha+1=0$
(i.e.\ $\alpha^2=\alpha+1$) and characteristic~$2$ (so $1+1=0$, and
$x=-x$ for all~$x$).
\begin{enumerate}[label=(\alph*)]
  \item Write out the complete addition and multiplication tables for
    $\F_4$.
  \item Verify that $\F_4$ is indeed a field.
\end{enumerate}
\end{exercise}

\begin{exercise}{Induction: sum of squares \intermediate}{exo:sum-squares}
Prove by induction that for all $n\geq 1$:
\[
  \sum_{k=1}^{n}k^2 = \frac{n(n+1)(2n+1)}{6}.
\]
\end{exercise}

\begin{exercise}{Induction: divisibility \intermediate}{exo:induction-div}
Prove by induction that $3^{2n}-1$ is divisible by $8$ for all $n\geq 1$.
\end{exercise}

\begin{exercise}{Proof strategies \challenging}{exo:proof-strategies}
Prove each of the following statements using the most appropriate proof
technique. State which technique you use.
\begin{enumerate}[label=(\alph*)]
  \item If $a,b\in\R$ and $a+b$ is irrational, then at least one of
    $a,b$ is irrational.
  \item For every $n\in\N$, $n^2+n$ is even.
  \item There is no smallest positive rational number.
  \item For all $n\geq 1$, $\displaystyle\sum_{k=0}^{n}2^k = 2^{n+1}-1$.
\end{enumerate}
\end{exercise}


% ============================================================
\section*{Chapter summary}\label{sec:ch0-summary}
\addcontentsline{toc}{section}{Chapter summary}
% ============================================================

\begin{itemize}
  \item A \textbf{set} is a collection of distinct objects. The standard
    operations are union, intersection, difference, and Cartesian product.
    The number sets $\N\subset\Z\subset\Q\subset\R\subset\C$ form the
    foundation for all that follows.

  \item A \textbf{mapping} $f\colon A\to B$ assigns to each element of
    $A$ exactly one element of $B$. The fundamental classification is into
    injective, surjective, and bijective maps.

  \item An \textbf{equivalence relation} partitions a set into
    equivalence classes. An \textbf{order relation} provides a way to
    compare elements.

  \item A \textbf{group} $(G,\star)$ captures the idea of a reversible
    operation. A \textbf{ring} $(R,+,\cdot\,)$ has two compatible
    operations. A \textbf{field} $(\K,+,\cdot\,)$ is a commutative ring
    in which every nonzero element has a multiplicative inverse.

  \item The principal fields in this course are $\R$ (the reals), $\C$
    (the complex numbers), and $\F_p$ (finite fields of prime order).

  \item The main proof techniques are: direct proof, contrapositive,
    contradiction, mathematical induction, and double inclusion. Mastering
    these is essential for the remainder of the course.
\end{itemize}

% ch1-vector-spaces.tex — Chapter 1: Vector Spaces and Subspaces
% The foundational chapter of the course: vector space axioms, subspaces,
% linear combinations, bases, dimension, and the Grassmann formula.

\chapter{Vector Spaces and Subspaces}\label{ch:vector-spaces}

The concept of a \emph{vector space}\index{vector space} is the central
object of linear algebra. It abstracts the familiar geometry of arrows in
the plane and in space into a framework that encompasses polynomials,
matrices, functions, and solutions of differential equations --- all
governed by the same structural rules.

In the previous chapter we established the notion of a field~$\K$. A
vector space is, in essence, a set whose elements (called \emph{vectors})
can be added together and scaled by elements of~$\K$, subject to axioms
that generalize the intuitive properties of $\R^2$ and $\R^3$.

\medskip
\noindent\textbf{Why abstract?}\ \ The reader may wonder why we bother
with axioms when we could simply work in $\R^n$. The answer is twofold.
First, many naturally occurring mathematical objects---polynomial spaces,
function spaces, solution sets of linear systems---share the same
algebraic structure as $\R^n$, and it is both economical and illuminating
to develop the theory once, in the abstract, rather than repeat arguments
for each concrete example. Second, the abstract framework reveals
\emph{which properties matter}: the dimension of a space, for instance,
is a far more fundamental invariant than the specific nature of its
elements.

We begin with the motivating example of $\R^2$ and $\R^3$, then give the
full axiomatic definition and a rich collection of examples. The bulk of
the chapter develops the core structural theory: subspaces, linear
combinations, linear independence, bases, and dimension, culminating in
the Grassmann formula.

Throughout this chapter, $\K$ denotes a field (typically $\R$ or $\C$).


% ============================================================
\section{Geometric motivation}\label{sec:geometric-motivation}
% ============================================================

Before plunging into axioms, let us recall the geometric picture that
motivates the entire theory.

In the Euclidean plane $\R^2$, a vector $\vec{v} = (v_1, v_2)$ can be
visualized as an arrow from the origin to the point $(v_1, v_2)$. Two
fundamental operations are available:
\begin{itemize}
  \item \textbf{Addition:} $(v_1,v_2)+(w_1,w_2) = (v_1+w_1, v_2+w_2)$,
    which corresponds to the parallelogram rule.
  \item \textbf{Scalar multiplication:}
    $\lambda(v_1,v_2) = (\lambda v_1, \lambda v_2)$, which stretches or
    compresses the arrow (and reverses it if $\lambda < 0$).
\end{itemize}

These operations satisfy many natural properties: addition is commutative
and associative, there is a zero vector $(0,0)$, every vector has an
additive inverse, scalar multiplication distributes over addition, and so
on. The exact list of these properties constitutes the axioms of a vector
space.

\begin{figure}[ht]
\centering
\begin{tikzpicture}[>=Stealth, scale=1.2]
  % Grid
  \draw[gray!30, very thin, step=0.5] (-0.5,-0.5) grid (4.5,3.5);
  \draw[->] (-0.5,0) -- (4.5,0) node[right]{$x_1$};
  \draw[->] (0,-0.5) -- (0,3.5) node[above]{$x_2$};
  % Vectors
  \draw[->, thick, blue!70!black] (0,0) -- (2,1)
    node[midway, below right]{$\vec{v}$};
  \draw[->, thick, red!70!black] (0,0) -- (1,2)
    node[midway, above left]{$\vec{w}$};
  % Parallelogram
  \draw[dashed, blue!50] (2,1) -- (3,3);
  \draw[dashed, red!50] (1,2) -- (3,3);
  \draw[->, thick, purple!70!black] (0,0) -- (3,3)
    node[midway, above left]{$\vec{v}+\vec{w}$};
  % Scalar multiple
  \draw[->, thick, blue!40, dashed] (0,0) -- (3,1.5)
    node[right]{$\frac{3}{2}\vec{v}$};
\end{tikzpicture}
\caption{Vector addition (parallelogram rule) and scalar multiplication
in $\R^2$.}
\label{fig:vector-addition-R2}
\end{figure}

The same picture extends to $\R^3$, where vectors are triples
$(v_1,v_2,v_3)$ and the operations are defined componentwise. But what
about $\R^n$ for $n\geq 4$, where geometric visualization fails? And
what about objects like polynomials or matrices, which can also be added
and scaled? The axioms of a vector space capture exactly the common
structure.


% ============================================================
\section{Definition of a vector space}\label{sec:vs-definition}
% ============================================================

\begin{definition}{Vector space}{def:vector-space}
\index{vector space!definition}
A \emph{vector space over $\K$} (or \emph{$\K$-vector space}) is a set
$E$ equipped with two operations:
\begin{itemize}
  \item \textbf{Addition:}\index{addition!vector} a map
    $+\colon E\times E\to E$, $(u,v)\mapsto u+v$.
  \item \textbf{Scalar multiplication:}\index{scalar multiplication}
    a map $\cdot\colon\K\times E\to E$, $(\lambda,v)\mapsto \lambda\cdot v$
    (usually written $\lambda v$).
\end{itemize}
These operations satisfy the following axioms for all
$u,v,w\in E$ and all $\lambda,\mu\in\K$:

\medskip
\noindent\textbf{Addition axioms:}
\begin{enumerate}[label=\textbf{(V\arabic*)}]
  \item\label{ax:assoc} \textbf{Associativity:}
    $(u+v)+w = u+(v+w)$.
  \item\label{ax:comm} \textbf{Commutativity:}
    $u+v = v+u$.
  \item\label{ax:zero} \textbf{Zero vector:}\index{zero vector}
    there exists $\vzero\in E$ such that $v+\vzero = v$ for all $v\in E$.
  \item\label{ax:inverse} \textbf{Additive inverse:}
    for each $v\in E$ there exists $-v\in E$ such that $v+(-v)=\vzero$.
\end{enumerate}

\medskip
\noindent\textbf{Scalar multiplication axioms:}
\begin{enumerate}[label=\textbf{(V\arabic*)},resume]
  \item\label{ax:compat} \textbf{Compatibility:}
    $\lambda(\mu v) = (\lambda\mu) v$.
  \item\label{ax:unit} \textbf{Unit:}
    $1_\K\, v = v$.
  \item\label{ax:dist-scalar} \textbf{Distributivity over vector addition:}
    $\lambda(u+v) = \lambda u + \lambda v$.
  \item\label{ax:dist-field} \textbf{Distributivity over scalar addition:}
    $(\lambda+\mu)v = \lambda v + \mu v$.
\end{enumerate}
The elements of $E$ are called \emph{vectors}\index{vector} and the
elements of $\K$ are called \emph{scalars}\index{scalar}.
\end{definition}

\begin{remark}{Notation}{rem:notation-vs}
When the field $\K$ is clear from context we simply say ``vector space''
rather than ``$\K$-vector space''. We write $\vzero$ for the zero vector
and $0$ for the zero scalar, relying on context to distinguish them.
\end{remark}

\begin{proposition}{Elementary properties}{prop:elementary-vs}
\index{vector space!elementary properties}
Let $E$ be a $\K$-vector space. For all $v\in E$ and $\lambda\in\K$:
\begin{enumerate}[label=(\roman*)]
  \item The zero vector $\vzero$ is unique.
  \item The additive inverse $-v$ is unique.
  \item $0_\K\, v = \vzero$.
  \item $\lambda\,\vzero = \vzero$.
  \item $(-1)v = -v$.
  \item $\lambda v = \vzero$ implies $\lambda = 0$ or $v = \vzero$.
\end{enumerate}
\end{proposition}

\begin{proof}
\begin{enumerate}[label=(\roman*)]
  \item Suppose $\vzero$ and $\vzero'$ are both zero vectors. Then
    $\vzero = \vzero + \vzero' = \vzero'$, using the zero-vector
    property of $\vzero'$ and then of $\vzero$.

  \item Suppose $w$ and $w'$ are both additive inverses of $v$. Then
    $w = w + \vzero = w + (v + w') = (w + v) + w' = \vzero + w' = w'$.

  \item $0_\K\,v = (0_\K + 0_\K)v = 0_\K\,v + 0_\K\,v$ by~\ref{ax:dist-field}.
    Adding $-(0_\K\,v)$ to both sides gives $\vzero = 0_\K\,v$.

  \item $\lambda\,\vzero = \lambda(\vzero + \vzero) =
    \lambda\,\vzero + \lambda\,\vzero$ by~\ref{ax:dist-scalar}. Adding
    $-(\lambda\,\vzero)$ yields $\vzero = \lambda\,\vzero$.

  \item $v + (-1)v = 1\cdot v + (-1)v = (1+(-1))v = 0_\K\,v = \vzero$
    by~\ref{ax:dist-field} and (iii). By uniqueness of inverses,
    $(-1)v = -v$.

  \item If $\lambda\neq 0$, then $\lambda^{-1}$ exists in $\K$ and
    $v = 1\cdot v = (\lambda^{-1}\lambda)v = \lambda^{-1}(\lambda v) =
    \lambda^{-1}\vzero = \vzero$ by (iv). \qedhere
\end{enumerate}
\end{proof}


% ============================================================
\section{Key examples of vector spaces}\label{sec:vs-examples}
% ============================================================

We now present a gallery of examples that will accoSpany us throughout
the course. The reader is encouraged to verify the axioms in each case
(at least mentally).

% --- K^n ---
\begin{example}{The space $\K^n$}{ex:Kn}
\index{Kn@$\K^n$}
For any positive integer $n$, the set
\[
  \K^n = \setcond{(x_1,\ldots,x_n)}{x_1,\ldots,x_n\in\K}
\]
equipped with componentwise addition and scalar multiplication:
\begin{align*}
  (x_1,\ldots,x_n) + (y_1,\ldots,y_n) &= (x_1+y_1,\ldots,x_n+y_n), \\
  \lambda(x_1,\ldots,x_n) &= (\lambda x_1,\ldots,\lambda x_n),
\end{align*}
is a $\K$-vector space. The zero vector is $\vzero = (0,\ldots,0)$.

When $\K = \R$, the cases $n=1,2,3$ correspond to the real line, the
plane, and three-dimensional space, respectively.
\end{example}

% --- Matrices ---
\begin{example}{The space $\Mat_{n,p}(\K)$}{ex:matrix-space}
\index{matrix space}
The set $\Mat_{n,p}(\K)$ of all $n\times p$ matrices with entries in $\K$
is a $\K$-vector space under the usual matrix addition and scalar
multiplication. The zero vector is the $n\times p$ zero matrix.

By identifying an $n\times p$ matrix with an element of $\K^{np}$ (by
listing entries row by row), we see that $\Mat_{n,p}(\K)$ and $\K^{np}$
have the ``same'' structure; this will be made precise by the notion of
\emph{isomorphism} in 
ef{ch:linear-maps}.
\end{example}

% --- Polynomials ---
\begin{example}{Polynomial spaces}{ex:polynomial-space}
\index{polynomial space}
The set $\K[X]$ of all polynomials with coefficients in $\K$ is a
$\K$-vector space under the usual addition and scalar multiplication of
polynomials. The zero vector is the zero polynomial.

For each $n\in\N$, the subset
\[
  \K_n[X] \deff \setcond{P\in\K[X]}{\deg P\leq n}\cup\set{0}
\]
of polynomials of degree at most $n$ (together with the zero polynomial)
is also a $\K$-vector space.
\end{example}

% --- Function spaces ---
\begin{example}{Function spaces}{ex:function-space}
\index{function space}
Let $S$ be a nonempty set. The set $\K^S \deff \setcond{f}{f\colon S\to\K}$
of all functions from $S$ to $\K$ is a $\K$-vector space under pointwise
operations:
\begin{align*}
  (f+g)(s) &\deff f(s) + g(s), \\
  (\lambda f)(s) &\deff \lambda\cdot f(s),
\end{align*}
for all $s\in S$, $\lambda\in\K$. The zero vector is the constant
function $s\mapsto 0$.

Important special cases include:
\begin{itemize}
  \item $\cC^0([a,b],\R)$: the space of continuous real-valued functions
    on $[a,b]$.
  \item $\cC^\infty(\R,\R)$: the space of infinitely differentiable
    functions.
  \item The space of sequences $(\K^\N)$: functions from $\N$ to $\K$.
\end{itemize}
\end{example}

% --- Solution spaces ---
\begin{example}{Solution space of a homogeneous linear system}{ex:homogeneous-system}
\index{solution space}\index{homogeneous system}
Let $A\in\Mat_{n,p}(\K)$. The set of solutions
\[
  \cS = \setcond{X\in\K^p}{AX = \vzero}
\]
of the homogeneous system $AX = \vzero$ is a $\K$-vector space (a
subspace of $\K^p$, as we shall see in~
ef{sec:subspaces}).

Indeed, if $AX_1 = \vzero$ and $AX_2 = \vzero$, then
$A(X_1+X_2) = AX_1 + AX_2 = \vzero$ and
$A(\lambda X_1) = \lambda AX_1 = \lambda\vzero = \vzero$.

\textbf{Warning:} The solution set of a \emph{non}-homogeneous system
$AX = B$ (with $B\neq\vzero$) is \emph{not} a vector space, since it
does not contain $\vzero$ (unless $B=\vzero$).
\end{example}

\begin{example}{The trivial vector space}{ex:trivial-space}
\index{trivial vector space}
The set $\set{\vzero}$, consisting of a single element, is a
$\K$-vector space (the only possibility for the operations being
$\vzero+\vzero = \vzero$ and $\lambda\,\vzero = \vzero$). It is called
the \emph{trivial} or \emph{zero} vector space.
\end{example}

\begin{remark}{The field $\K$ is a vector space over itself}{rem:K-over-K}
The field $\K$ is itself a $\K$-vector space (with field multiplication
playing the role of scalar multiplication). More generally, $\C$ is an
$\R$-vector space (of dimension~$2$, as we shall see).
\end{remark}


% ============================================================
\section{Subspaces}\label{sec:subspaces}
% ============================================================

A subspace is a subset of a vector space that is itself a vector space
(with the inherited operations). Rather than checking all eight axioms,
the following characterization provides a convenient shortcut.

\begin{definition}{Vector subspace}{def:subspace}
\index{subspace!definition}
Let $E$ be a $\K$-vector space. A subset $F\subseteq E$ is a
\emph{vector subspace} (or simply \emph{subspace}) of $E$ if:
\begin{enumerate}[label=(\roman*)]
  \item $F$ is nonempty (equivalently, $\vzero\in F$).
  \item $F$ is closed under addition: $u+v\in F$ for all $u,v\in F$.
  \item $F$ is closed under scalar multiplication: $\lambda v\in F$ for
    all $\lambda\in\K$ and $v\in F$.
\end{enumerate}
\end{definition}

\begin{proposition}{Subspace characterization}{prop:subspace-char}
\index{subspace!characterization}
A subset $F\subseteq E$ is a subspace of $E$ if and only if
\begin{enumerate}[label=(\alph*)]
  \item $F\neq\varnothing$ (or equivalently, $\vzero\in F$), and
  \item for all $\lambda\in\K$ and all $u,v\in F$:
    $\lambda u + v\in F$.
\end{enumerate}
Equivalently, $F$ is nonempty and closed under linear combinations: for
all $\lambda,\mu\in\K$ and $u,v\in F$, $\lambda u + \mu v\in F$.
\end{proposition}

\begin{proof}
If $F$ is a subspace, then (a) and (b) are immediate from the definition.
Conversely, assume (a) and (b). Taking $\lambda = 0$ in (b) shows
$\vzero + v = v\in F$ for any $v\in F$, confirming closure under
addition (take $\lambda=1$). Taking $v = \vzero$ in (b) shows
$\lambda u\in F$, confirming closure under scalar multiplication. Thus
$F$ satisfies all three conditions of the definition.
\end{proof}

\begin{remark}{The nonemptiness condition}{rem:nonempty}
Condition~(a) cannot be dropped. The empty set $\varnothing$ is trivially
closed under addition and scalar multiplication (vacuously) but is not a
vector space.
\end{remark}

\begin{example}{Subspaces of $\R^2$}{ex:subspaces-R2}
\index{subspace!of $\R^2$}
The subspaces of $\R^2$ are precisely:
\begin{enumerate}[label=(\roman*)]
  \item $\set{\vzero}$ (the origin).
  \item Lines through the origin: $\setcond{t\vec{v}}{t\in\R}$ for
    some $\vec{v}\neq\vzero$.
  \item $\R^2$ itself.
\end{enumerate}
Note that a line \emph{not} passing through the origin (e.g.\
$\setcond{(x,y)\in\R^2}{x+y=1}$) is \emph{not} a subspace.
\end{example}

\begin{figure}[ht]
\centering
\begin{tikzpicture}[>=Stealth, scale=1.0]
  % Axes
  \draw[gray!30, very thin, step=1] (-3.2,-2.2) grid (3.2,2.7);
  \draw[->] (-3.2,0) -- (3.2,0) node[right]{$x_1$};
  \draw[->] (0,-2.2) -- (0,2.7) node[above]{$x_2$};
  % Subspace: line through origin
  \draw[thick, blue!70!black] (-3,-1.5) -- (3,1.5);
  \node[blue!70!black, above left] at (2.4,1.2)
    {$F = \Span(\vec{v})$};
  \draw[->, very thick, blue!70!black] (0,0) -- (2,1)
    node[below right]{$\vec{v}$};
  % NOT a subspace: translated line
  \draw[thick, red!70!black, dashed] (-2.5,2.5+0.75) -- (2.5,-2.5+0.75);
  \node[red!70!black, right] at (1.8,-0.75)
    {$x_1+x_2=1$ (not a subspace)};
  % Origin
  \fill (0,0) circle (2pt) node[below left]{$\vzero$};
\end{tikzpicture}
\caption{A subspace of $\R^2$ (a line through the origin) versus a
non-subspace (a translated line).}
\label{fig:subspace-R2}
\end{figure}

\begin{example}{Subspaces of $\R^3$}{ex:subspaces-R3}
\index{subspace!of $\R^3$}
The subspaces of $\R^3$ are:
\begin{enumerate}[label=(\roman*)]
  \item $\set{\vzero}$.
  \item Lines through the origin.
  \item Planes through the origin (e.g.\
    $\setcond{(x,y,z)\in\R^3}{ax+by+cz=0}$ with $(a,b,c)\neq\vzero$).
  \item $\R^3$ itself.
\end{enumerate}
This classification will follow from the theory of dimension developed
later in this chapter.
\end{example}

\begin{example}{Further examples of subspaces}{ex:further-subspaces}
\begin{enumerate}[label=(\alph*)]
  \item The set of symmetric $n\times n$ matrices
    $\mathrm{Sym}_n(\K) = \setcond{A\in\Mat_n(\K)}{\transpose{A}=A}$ is a
    subspace of $\Mat_n(\K)$.
  \item The set of skew-symmetric matrices
    $\mathrm{Alt}_n(\K) = \setcond{A\in\Mat_n(\K)}{\transpose{A}=-A}$
    is a subspace of $\Mat_n(\K)$.
  \item The set of polynomials of degree at most $n$, $\K_n[X]$, is a
    subspace of $\K[X]$.
  \item The set of even polynomials
    $\setcond{P\in\K[X]}{P(-X)=P(X)}$ is a subspace of $\K[X]$.
  \item The solution set of the differential equation $y'' + y = 0$ (a
    subset of $\cC^\infty(\R,\R)$) is a subspace (the zero function is a
    solution, and linear combinations of solutions are solutions).
\end{enumerate}
\end{example}


% ============================================================
\section{Intersection and sum of subspaces}\label{sec:intersection-sum}
% ============================================================

\begin{proposition}{Intersection of subspaces}{prop:intersection-subspace}
\index{intersection of subspaces}
Let $(F_i)_{i\in I}$ be a (possibly infinite) family of subspaces of a
$\K$-vector space $E$. Then $\displaystyle\bigcap_{i\in I} F_i$ is a
subspace of $E$.
\end{proposition}

\begin{proof}
Let $F = \bigcap_{i\in I} F_i$.
\begin{itemize}
  \item \textbf{Nonempty:} $\vzero\in F_i$ for every $i$, so
    $\vzero\in F$.
  \item \textbf{Closed under linear combinations:} Let $u,v\in F$ and
    $\lambda,\mu\in\K$. For each $i\in I$, $u,v\in F_i$ and $F_i$ is a
    subspace, so $\lambda u+\mu v\in F_i$. Hence
    $\lambda u+\mu v\in F$. \qedhere
\end{itemize}
\end{proof}

\begin{remark}{Union is not a subspace in general}{rem:union-not-subspace}
\index{union of subspaces}
In contrast, the \emph{union} of two subspaces is generally \emph{not}
a subspace. For instance, in $\R^2$, let $F_1$ be the $x$-axis and $F_2$
the $y$-axis. Then $(1,0)\in F_1$ and $(0,1)\in F_2$, but
$(1,0)+(0,1)=(1,1)\notin F_1\cup F_2$.

In fact, $F_1\cup F_2$ is a subspace if and only if $F_1\subseteq F_2$
or $F_2\subseteq F_1$.
\end{remark}

The correct replacement for union is the \emph{sum} of subspaces.

\begin{definition}{Sum of subspaces}{def:sum-subspaces}
\index{sum of subspaces}
Let $F_1,\ldots,F_r$ be subspaces of $E$. Their \emph{sum} is
\[
  F_1 + F_2 + \cdots + F_r \deff
  \setcond{v_1 + v_2 + \cdots + v_r}{v_i\in F_i \text{ for each } i}.
\]
For two subspaces, $F+G = \setcond{u+v}{u\in F,\;v\in G}$.
\end{definition}

\begin{proposition}{Sum is the smallest containing subspace}{prop:sum-smallest}
The sum $F_1+\cdots+F_r$ is a subspace of $E$. It is the smallest
subspace of $E$ containing each $F_i$, i.e.\
$F_1+\cdots+F_r = \bigcap\setcond{H}{H \text{ subspace of } E,\;
F_i\subseteq H \text{ for all } i}$.
\end{proposition}

\begin{proof}
We prove the case $r=2$; the general case follows by induction.

Let $S = F + G$. First, $\vzero = \vzero + \vzero\in S$, so $S$ is
nonempty. Let $u_1+v_1, u_2+v_2\in S$ (with $u_i\in F$, $v_i\in G$) and
$\lambda\in\K$. Then
\[
  \lambda(u_1+v_1)+(u_2+v_2) = (\lambda u_1+u_2)+(\lambda v_1+v_2)\in S,
\]
since $\lambda u_1+u_2\in F$ and $\lambda v_1+v_2\in G$. Hence $S$ is a
subspace.

For the minimality claim: $F\subseteq S$ (take $v=\vzero$) and
$G\subseteq S$ (take $u=\vzero$), so $S$ contains both. If $H$ is any
subspace containing $F$ and $G$, then for any $u+v\in S$ with $u\in F$
and $v\in G$, we have $u,v\in H$, so $u+v\in H$; thus $S\subseteq H$.
\end{proof}


% ============================================================
\section{Direct sum and complementary subspaces}\label{sec:direct-sum}
% ============================================================

\begin{definition}{Direct sum (internal)}{def:direct-sum}
\index{direct sum}\index{direct sum!internal}
Let $F$ and $G$ be subspaces of $E$. We say that $E$ is the
\emph{direct sum} of $F$ and $G$, written $E = F\oplus G$, if:
\begin{enumerate}[label=(\roman*)]
  \item $E = F + G$ (every vector in $E$ can be written as $u+v$ with
    $u\in F$, $v\in G$), and
  \item $F\cap G = \set{\vzero}$.
\end{enumerate}
Equivalently, every $x\in E$ can be written \emph{uniquely} as $x=u+v$
with $u\in F$ and $v\in G$.
\end{definition}

\begin{proposition}{Characterization of direct sum}{prop:direct-sum-char}
\index{direct sum!characterization}
Let $F$ and $G$ be subspaces of $E$ with $E = F+G$. The following are
equivalent:
\begin{enumerate}[label=(\roman*)]
  \item $E = F\oplus G$.
  \item $F\cap G = \set{\vzero}$.
  \item Every $x\in E$ has a \emph{unique} decomposition $x = u + v$
    with $u\in F$ and $v\in G$.
\end{enumerate}
\end{proposition}

\begin{proof}
\textbf{(i)$\Leftrightarrow$(ii)} is immediate from the definition.

\textbf{(ii)$\Rightarrow$(iii):} Suppose $x = u_1+v_1 = u_2+v_2$ with
$u_i\in F$ and $v_i\in G$. Then $u_1-u_2 = v_2-v_1\in F\cap G = \set{\vzero}$,
so $u_1=u_2$ and $v_1=v_2$.

\textbf{(iii)$\Rightarrow$(ii):} The zero vector can be written as
$\vzero = \vzero + \vzero$. If $w\in F\cap G$, then
$\vzero = w + (-w)$ with $w\in F$ and $-w\in G$. By uniqueness,
$w = \vzero$.
\end{proof}

\begin{definition}{Complementary subspace}{def:complement}
\index{complementary subspace}
If $E = F\oplus G$, we say that $G$ is a \emph{complement} of $F$ in
$E$ (and $F$ is a complement of $G$). In this case, $F$ and $G$ are
called \emph{complementary subspaces}.
\end{definition}

\begin{remark}{Non-uniqueness of complements}{rem:complement-nonunique}
A subspace may have many different complements. For instance, in $\R^2$,
the $x$-axis has any non-horizontal line through the origin as a
complement.
\end{remark}

The direct sum extends naturally to more than two subspaces.

\begin{definition}{Direct sum of $r$ subspaces}{def:direct-sum-r}
\index{direct sum!of multiple subspaces}
Let $F_1,\ldots,F_r$ be subspaces of $E$. We say that
$E = F_1\oplus\cdots\oplus F_r$ if:
\begin{enumerate}[label=(\roman*)]
  \item $E = F_1 + \cdots + F_r$, and
  \item for each $j\in\set{1,\ldots,r}$,
    $F_j \cap (F_1+\cdots+F_{j-1}+F_{j+1}+\cdots+F_r) = \set{\vzero}$.
\end{enumerate}
Equivalently, every $x\in E$ has a unique decomposition
$x = v_1 + \cdots + v_r$ with $v_i\in F_i$.
\end{definition}

\begin{example}{Direct sum decomposition in $\R^3$}{ex:direct-sum-R3}
\index{direct sum!example in $\R^3$}
Let
\begin{align*}
  F &= \setcond{(x,0,0)}{x\in\R} &&\text{(the $x$-axis)}, \\
  G &= \setcond{(0,y,z)}{y,z\in\R} &&\text{(the $yz$-plane)}.
\end{align*}
Then $\R^3 = F\oplus G$, since every $(x,y,z) = (x,0,0)+(0,y,z)$
uniquely and $F\cap G = \set{\vzero}$.
\end{example}

\begin{figure}[ht]
\centering
\begin{tikzpicture}[>=Stealth, scale=1.1,
  x={(-0.4cm,-0.3cm)}, y={(1cm,0cm)}, z={(0cm,1cm)}]
  % Plane G (yz-plane)
  \fill[blue!10, opacity=0.6] (0,-2.2,-1.8) -- (0,2.2,-1.8) --
    (0,2.2,1.8) -- (0,-2.2,1.8) -- cycle;
  \node[blue!70!black] at (0,1.8,1.5) {$G$};
  % Line F (x-axis)
  \draw[thick, red!70!black] (2.8,0,0) -- (-2.8,0,0);
  \node[red!70!black, above] at (-2.2,0,0) {$F$};
  % Axes
  \draw[->] (0,0,0) -- (-3,0,0) node[left]{$x$};
  \draw[->] (0,0,0) -- (0,2.5,0) node[right]{$y$};
  \draw[->] (0,0,0) -- (0,0,2.2) node[above]{$z$};
  % Point and decomposition
  \coordinate (P) at (-1.5,1.2,0.8);
  \coordinate (Pf) at (-1.5,0,0);
  \coordinate (Pg) at (0,1.2,0.8);
  \fill[black] (P) circle (1.5pt) node[above right]{$(x,y,z)$};
  \draw[dashed, red!50] (P) -- (Pg);
  \draw[dashed, blue!50] (P) -- (Pf);
  \draw[->, thick, red!70!black] (0,0,0) -- (Pf)
    node[below]{$(x,0,0)$};
  \draw[->, thick, blue!70!black] (0,0,0) -- (Pg)
    node[right]{$(0,y,z)$};
  \fill (0,0,0) circle (1.5pt);
\end{tikzpicture}
\caption{Direct sum decomposition $\R^3 = F\oplus G$ where $F$ is the
$x$-axis and $G$ is the $yz$-plane.}
\label{fig:direct-sum-R3}
\end{figure}


% ============================================================
\section{Linear combinations and Span}\label{sec:Span}
% ============================================================

\begin{definition}{Linear combination}{def:linear-combination}
\index{linear combination}
Let $E$ be a $\K$-vector space and $v_1,\ldots,v_n\in E$. A
\emph{linear combination} of $v_1,\ldots,v_n$ is any vector of the form
\[
  \lambda_1 v_1 + \lambda_2 v_2 + \cdots + \lambda_n v_n,
  \qquad \lambda_1,\ldots,\lambda_n\in\K.
\]
The scalars $\lambda_1,\ldots,\lambda_n$ are called the
\emph{coefficients} of the linear combination.
\end{definition}

\begin{definition}{Span (generating set)}{def:Span}
\index{Span}\index{generating family}
Let $\cF = (v_1,\ldots,v_n)$ be a family of vectors in $E$. The
\emph{Span} of $\cF$ is
\[
  \Span(v_1,\ldots,v_n) \deff
  \setcond{\lambda_1 v_1 + \cdots + \lambda_n v_n}
    {\lambda_1,\ldots,\lambda_n\in\K}.
\]
It is the set of all linear combinations of $v_1,\ldots,v_n$.

A family $\cF$ is a \emph{generating family} (or \emph{Spanning set})
for $E$ if $\Span(\cF) = E$. In this case we say that $\cF$
\emph{Spans} (or \emph{generates}) $E$.
\end{definition}

\begin{proposition}{Span is a subspace}{prop:Span-subspace}
$\Span(v_1,\ldots,v_n)$ is a subspace of $E$. Moreover, it is the
smallest subspace containing $v_1,\ldots,v_n$.
\end{proposition}

\begin{proof}
Let $W = \Span(v_1,\ldots,v_n)$.
\begin{itemize}
  \item $\vzero = 0v_1+\cdots+0v_n\in W$, so $W\neq\varnothing$.
  \item If $u = \sum_i \alpha_i v_i$ and $w = \sum_i\beta_i v_i$ are in
    $W$ and $\lambda\in\K$, then
    $\lambda u + w = \sum_i(\lambda\alpha_i+\beta_i)v_i\in W$.
\end{itemize}
Hence $W$ is a subspace. It contains each $v_j$ (take
$\lambda_j=1$ and $\lambda_i=0$ for $i\neq j$). If $H$ is any subspace
containing every $v_j$, then $H$ contains all linear combinations of
$v_1,\ldots,v_n$, so $W\subseteq H$.
\end{proof}

\begin{example}{Spanning $\R^3$}{ex:Span-R3}
The canonical vectors $\vece{1}=(1,0,0)$, $\vece{2}=(0,1,0)$,
$\vece{3}=(0,0,1)$ Span $\R^3$, since every $(a,b,c)\in\R^3$ equals
$a\,\vece{1}+b\,\vece{2}+c\,\vece{3}$.
\end{example}

\begin{example}{Spanning a plane in $\R^3$}{ex:Span-plane}
The vectors $\vec{u}=(1,1,0)$ and $\vec{v}=(0,1,1)$ Span the plane
\[
  \Span(\vec{u},\vec{v}) = \setcond{(s,s+t,t)}{s,t\in\R}
  = \setcond{(x,y,z)\in\R^3}{x-y+z=0}.
\]
\end{example}


% ============================================================
\section{Linear independence}\label{sec:linear-independence}
% ============================================================

\begin{definition}{Linear independence}{def:linear-independence}
\index{linear independence}
A family $(v_1,\ldots,v_n)$ of vectors in $E$ is \emph{linearly
independent} (or \emph{free}) if
\[
  \lambda_1 v_1 + \cdots + \lambda_n v_n = \vzero
  \implies
  \lambda_1 = \lambda_2 = \cdots = \lambda_n = 0.
\]
A family that is not linearly independent is called \emph{linearly
dependent}\index{linear dependence}.
\end{definition}

\begin{remark}{Interpretation}{rem:li-interpretation}
Linear independence means that no vector in the family can be written as
a linear combination of the others. Equivalently, the only way to
represent $\vzero$ as a linear combination of $v_1,\ldots,v_n$ is the
\emph{trivial} one (all coefficients zero).
\end{remark}

\begin{proposition}{Characterization of linear dependence}{prop:ld-char}
\index{linear dependence!characterization}
A family $(v_1,\ldots,v_n)$ with $n\geq 2$ is linearly dependent if and
only if at least one vector $v_j$ can be written as a linear combination
of the others:
\[
  v_j = \sum_{\substack{i=1\\i\neq j}}^n \mu_i\, v_i.
\]
\end{proposition}

\begin{proof}
$(\Rightarrow)$~If the family is dependent, there exist
$\lambda_1,\ldots,\lambda_n$, not all zero, with
$\sum_i\lambda_i v_i = \vzero$. Pick $j$ with $\lambda_j\neq 0$. Then
$v_j = -\lambda_j^{-1}\sum_{i\neq j}\lambda_i v_i$.

$(\Leftarrow)$~If $v_j = \sum_{i\neq j}\mu_i v_i$, then
$\sum_{i\neq j}\mu_i v_i + (-1)v_j = \vzero$ is a nontrivial relation.
\end{proof}

\begin{proposition}{Properties of linear independence}{prop:li-properties}
\index{linear independence!properties}
Let $E$ be a $\K$-vector space.
\begin{enumerate}[label=(\roman*)]
  \item A single nonzero vector $(v)$ is linearly independent.
  \item Any subfamily of a linearly independent family is linearly
    independent.
  \item If $(v_1,\ldots,v_n)$ is linearly independent and $w\notin
    \Span(v_1,\ldots,v_n)$, then $(v_1,\ldots,v_n,w)$ is linearly
    independent.
  \item Any family containing $\vzero$ is linearly dependent.
  \item Any family containing a repeated vector is linearly dependent.
\end{enumerate}
\end{proposition}

\begin{proof}
\begin{enumerate}[label=(\roman*)]
  \item If $\lambda v = \vzero$ with $v\neq\vzero$, then $\lambda = 0$
    by 
ef{prop:prop:elementary-vs}(vi).

  \item Suppose $(v_{i_1},\ldots,v_{i_k})$ is a subfamily and
    $\sum_{j=1}^k \lambda_j v_{i_j} = \vzero$. Extend to a combination
    of the full family by setting the remaining coefficients to zero;
    independence of the full family forces all $\lambda_j = 0$.

  \item Suppose $\lambda_1 v_1+\cdots+\lambda_n v_n + \mu w = \vzero$.
    If $\mu\neq 0$, then
    $w = -\mu^{-1}(\lambda_1 v_1+\cdots+\lambda_n v_n)\in
    \Span(v_1,\ldots,v_n)$, contradicting the hypothesis. So $\mu=0$,
    and then $\lambda_1=\cdots=\lambda_n=0$ by independence.

  \item $1\cdot\vzero = \vzero$ is a nontrivial relation.

  \item If $v_i = v_j$ for $i\neq j$, then
    $1\cdot v_i + (-1)\cdot v_j = \vzero$ is nontrivial. \qedhere
\end{enumerate}
\end{proof}

\begin{example}{Independence in $\R^3$}{ex:independence-R3}
\begin{enumerate}[label=(\alph*)]
  \item The canonical basis vectors $\vece{1}, \vece{2}, \vece{3}$ are
    linearly independent: if $a\,\vece{1}+b\,\vece{2}+c\,\vece{3}=\vzero$,
    then $(a,b,c)=(0,0,0)$.
  \item The vectors $(1,1,0)$, $(1,0,1)$, $(0,1,1)$ are linearly
    independent. Indeed,
    $\alpha(1,1,0)+\beta(1,0,1)+\gamma(0,1,1)=\vzero$ gives
    \[
      \alpha+\beta=0,\quad \alpha+\gamma=0,\quad \beta+\gamma=0,
    \]
    which has the unique solution $\alpha=\beta=\gamma=0$.
  \item The vectors $(1,2,3)$, $(4,5,6)$, $(7,8,9)$ are linearly
    \emph{dependent}: $(1,2,3) - 2(4,5,6) + (7,8,9) = (0,0,0)$.
\end{enumerate}
\end{example}

\begin{example}{Independence in $\K_n[X]$}{ex:independence-polynomials}
The polynomials $1,X,X^2,\ldots,X^n$ are linearly independent in
$\K[X]$: if $\alpha_0 + \alpha_1 X + \cdots + \alpha_n X^n = 0$ (the
zero polynomial), then all $\alpha_i = 0$ by comparing coefficients.
\end{example}


% ============================================================
\section{Bases}\label{sec:bases}
% ============================================================

\begin{definition}{Basis}{def:basis}
\index{basis!definition}
A family $\cB = (v_1,\ldots,v_n)$ of vectors in $E$ is a \emph{basis}
of $E$ if it is both linearly independent and Spanning: $\cB$ is
linearly independent and $\Span(\cB) = E$.
\end{definition}

\begin{theorem}{Characterization of a basis}{thm:basis-char}
\index{basis!characterization}
Let $\cB = (v_1,\ldots,v_n)$ be a family in $E$. The following are
equivalent:
\begin{enumerate}[label=(\roman*)]
  \item $\cB$ is a basis of $E$.
  \item Every $v\in E$ can be written \emph{uniquely} as
    $v = \lambda_1 v_1 + \cdots + \lambda_n v_n$.
  \item $\cB$ is a maximal linearly independent family.
  \item $\cB$ is a minimal Spanning family.
\end{enumerate}
\end{theorem}

\begin{proof}
\textbf{(i)$\Rightarrow$(ii):} Since $\cB$ Spans $E$, at least one
representation exists. Suppose
$v = \sum_i\lambda_i v_i = \sum_i\mu_i v_i$. Then
$\sum_i(\lambda_i-\mu_i)v_i = \vzero$, and independence gives
$\lambda_i = \mu_i$ for all $i$.

\textbf{(ii)$\Rightarrow$(i):} Existence of a representation means $\cB$
Spans $E$. If $\sum_i\lambda_i v_i = \vzero$, then
$\sum_i\lambda_i v_i = \sum_i 0\cdot v_i$ are two representations of
$\vzero$; uniqueness gives $\lambda_i = 0$ for all $i$.

\textbf{(i)$\Rightarrow$(iii):} If $\cB$ is a basis and we add any
$w\in E$, then $w = \sum_i\lambda_i v_i$, so the enlarged family
contains $w - \sum_i\lambda_i v_i = \vzero$ as a nontrivial relation;
hence it is dependent. So $\cB$ is maximal independent.

\textbf{(iii)$\Rightarrow$(i):} If $\cB$ does not Span $E$, there exists
$w\notin\Span(\cB)$, and $(v_1,\ldots,v_n,w)$ is independent by

ef{prop:prop:li-properties}(iii), contradicting maximality.

\textbf{(i)$\Rightarrow$(iv):} If we remove $v_j$, we claim the
remaining family does not Span $E$. Indeed, $v_j$ cannot be written as a
linear combination of the other $v_i$ (by independence), so $v_j$ is not
in the Span of the reduced family.

\textbf{(iv)$\Rightarrow$(i):} $\cB$ Spans $E$ by assumption. If $\cB$
is dependent, some $v_j$ is a combination of the others (by

ef{prop:prop:ld-char}), so removing $v_j$ still gives a Spanning family,
contradicting minimality.
\end{proof}

\begin{definition}{Coordinates}{def:coordinates}
\index{coordinates}
If $\cB = (v_1,\ldots,v_n)$ is a basis of $E$ and $v = \sum_i\lambda_iv_i$,
the uniquely determined scalars $\lambda_1,\ldots,\lambda_n$ are the
\emph{coordinates} (or \emph{components}) of $v$ relative to $\cB$. The
column vector
\[
  [v]_\cB = \begin{pmatrix}\lambda_1\\\vdots\\\lambda_n\end{pmatrix}
  \in\K^n
\]
is called the \emph{coordinate vector} of $v$ in basis $\cB$.
\end{definition}

\begin{example}{Canonical basis of $\K^n$}{ex:canonical-basis}
\index{canonical basis!of $\K^n$}
The \emph{canonical} (or \emph{standard}) basis of $\K^n$ is
$\cE = (\vece{1},\ldots,\vece{n})$ where $\vece{i}$ is the $n$-tuple
with a $1$ in position $i$ and $0$'s elsewhere. The coordinates of
$(x_1,\ldots,x_n)$ in $\cE$ are simply $x_1,\ldots,x_n$.
\end{example}

\begin{example}{Canonical basis of $\K_n[X]$}{ex:canonical-basis-poly}
\index{canonical basis!of polynomial space}
The family $(1,X,X^2,\ldots,X^n)$ is a basis of $\K_n[X]$: every
polynomial of degree $\leq n$ is a unique linear combination of these
monomials. In this basis, the coordinates of a polynomial are its
coefficients.
\end{example}

\begin{example}{Canonical basis of $\Mat_{n,p}(\K)$}{ex:canonical-basis-matrix}
\index{canonical basis!of matrix space}
The family $(E_{ij})_{1\leq i\leq n,\,1\leq j\leq p}$, where $E_{ij}$
is the matrix with $1$ in position $(i,j)$ and $0$'s elsewhere, is a
basis of $\Mat_{n,p}(\K)$.
\end{example}


% ============================================================
\section{Dimension}\label{sec:dimension}
% ============================================================

The notion of dimension captures the ``size'' of a vector space. Its
definition rests on a fundamental theorem about the relationship between
independent families and Spanning families.

\begin{theorem}{Steinitz exchange lemma}{thm:steinitz}
\index{Steinitz exchange lemma}
Let $E$ be a $\K$-vector space. If $(v_1,\ldots,v_p)$ is linearly
independent and $(w_1,\ldots,w_q)$ Spans $E$, then $p\leq q$.
\end{theorem}

\begin{proof}
We proceed by induction on $p$.

\textbf{Base case ($p=0$):} The empty family is independent and $0\leq q$.

\textbf{Inductive step:} Assume the result holds for all independent
families of size $p-1$. Suppose $(v_1,\ldots,v_p)$ is independent and
$(w_1,\ldots,w_q)$ Spans $E$.

Since $(w_1,\ldots,w_q)$ Spans $E$, we can write
\[
  v_p = \alpha_1 w_1 + \cdots + \alpha_q w_q
\]
for some $\alpha_i\in\K$. Not all $\alpha_i$ can be zero (since
$v_p\neq\vzero$). By reindexing, assume $\alpha_q\neq 0$. Then
\[
  w_q = \alpha_q^{-1}\Bigl(v_p - \sum_{i=1}^{q-1}\alpha_i w_i\Bigr).
\]
Therefore every vector in $E$ that was a linear combination of
$w_1,\ldots,w_q$ can also be expressed using $w_1,\ldots,w_{q-1},v_p$.
That is, $(w_1,\ldots,w_{q-1},v_p)$ still Spans $E$.

Now $(v_1,\ldots,v_{p-1})$ is independent (as a subfamily of an
independent family), and $(w_1,\ldots,w_{q-1},v_p)$ Spans $E$. We claim
$(v_1,\ldots,v_{p-1})$ remains independent in $E$, so by the induction
hypothesis, $p-1\leq q-1+1 = q$... but we need to be more careful.

We refine the argument. Since $(w_1,\ldots,w_{q-1},v_p)$ Spans $E$,
each $v_i$ ($1\leq i\leq p-1$) can be written as a combination of
$w_1,\ldots,w_{q-1},v_p$. Since $(v_1,\ldots,v_p)$ is independent,
$v_1,\ldots,v_{p-1}$ are not in $\Span(v_p)$ and hence depend
nontrivially on $w_1,\ldots,w_{q-1}$. We may thus view
$(v_1,\ldots,v_{p-1})$ as a linearly independent family in the space
$\Span(w_1,\ldots,w_{q-1},v_p)$.

Apply the exchange process again: write $v_{p-1}$ as a combination of
$w_1,\ldots,w_{q-1},v_p$. The coefficient of some $w_j$ must be nonzero
(otherwise $v_{p-1}\in\Span(v_p)$, contradicting independence). Swap
$w_j$ for $v_{p-1}$.

Continuing this process $p$ times, we replace $p$ of the $w_i$'s by
$v_1,\ldots,v_p$. At each step, at least one $w_i$ must be available
for replacement (its coefficient must be nonzero). This is possible only
if $p\leq q$.
\end{proof}

\begin{corollary}{Invariance of basis size}{cor:basis-size}
\index{basis!invariance of size}
If $E$ has a finite basis, then every basis of $E$ has the same number of
elements.
\end{corollary}

\begin{proof}
Let $\cB = (v_1,\ldots,v_n)$ and $\cB' = (w_1,\ldots,w_m)$ be two
bases. Since $\cB$ is independent and $\cB'$ Spans $E$, the Steinitz
lemma gives $n\leq m$. Symmetrically, $m\leq n$. Hence $n=m$.
\end{proof}

\begin{definition}{Dimension}{def:dimension}
\index{dimension}
A $\K$-vector space $E$ is \emph{finite-dimensional} if it has a finite
Spanning set. In this case, the \emph{dimension} of $E$, denoted
$\dim_\K E$ (or simply $\dim E$), is the common cardinality of all its
bases.

By convention, $\dim\set{\vzero} = 0$.

If $E$ has no finite Spanning set, it is \emph{infinite-dimensional}.
\end{definition}

\begin{example}{Dimensions of standard spaces}{ex:dimensions}
\begin{enumerate}[label=(\alph*)]
  \item $\dim_\K \K^n = n$ (the canonical basis has $n$ elements).
  \item $\dim_\K \K_n[X] = n+1$ (the basis $1,X,\ldots,X^n$ has $n+1$
    elements).
  \item $\dim_\K \Mat_{n,p}(\K) = np$.
  \item $\K[X]$ is infinite-dimensional: the family
    $(1,X,X^2,\ldots)$ is independent and Spans $\K[X]$, but it is
    not finite.
\end{enumerate}
\end{example}

\begin{theorem}{Existence of a basis (finite-dimensional case)}{thm:basis-existence}
\index{basis!existence}
Every finite-dimensional vector space $E\neq\set{\vzero}$ has a basis.
More precisely:
\begin{enumerate}[label=(\roman*)]
  \item Every Spanning family contains a basis (one can extract a basis
    from it).
  \item Every linearly independent family can be extended to a basis.
\end{enumerate}
\end{theorem}

\begin{proof}
\begin{enumerate}[label=(\roman*)]
  \item Let $\cF = (w_1,\ldots,w_q)$ Span $E$. If $\cF$ is independent,
    it is a basis. If not, some $w_j$ is a combination of the others;
    remove it. The resulting family still Spans $E$ (any combination
    involving $w_j$ can be rewritten using the others). Repeat until the
    family is independent. The process terminates since we cannot reduce
    below a single vector (assuming $E\neq\set{\vzero}$).

  \item Let $(v_1,\ldots,v_p)$ be independent. If it Spans $E$, it is a
    basis. If not, there exists $w\notin\Span(v_1,\ldots,v_p)$, and
    $(v_1,\ldots,v_p,w)$ is independent by
    
ef{prop:prop:li-properties}(iii). Repeat. By the Steinitz lemma, the
    size of an independent family is bounded by any Spanning family's
    size, so the process terminates. \qedhere
\end{enumerate}
\end{proof}

\begin{theorem}{Existence of complements}{thm:complement-existence}
\index{complementary subspace!existence}
Let $E$ be a finite-dimensional $\K$-vector space and $F$ a subspace of
$E$. Then $F$ has at least one complement in $E$: there exists a
subspace $G$ such that $E = F\oplus G$.
\end{theorem}

\begin{proof}
Let $(v_1,\ldots,v_p)$ be a basis of $F$. This is a linearly independent
family in $E$; by 
ef{thm:thm:basis-existence}(ii), extend it to a basis
$(v_1,\ldots,v_p,w_1,\ldots,w_q)$ of $E$. Set
$G = \Span(w_1,\ldots,w_q)$.

\textbf{$E = F+G$:} Any $x\in E$ can be written as
$x = \sum_i\alpha_iv_i + \sum_j\beta_jw_j\in F+G$.

\textbf{$F\cap G = \set{\vzero}$:} If $x\in F\cap G$, write
$x = \sum_i\alpha_iv_i = \sum_j\beta_jw_j$. Then
$\sum_i\alpha_iv_i - \sum_j\beta_jw_j = \vzero$, and since
$(v_1,\ldots,v_p,w_1,\ldots,w_q)$ is a basis (hence independent), all
coefficients vanish: $x = \vzero$.
\end{proof}


% ============================================================
\section{Dimension formulas}\label{sec:dimension-formulas}
% ============================================================

\begin{proposition}{Dimension of a subspace}{prop:dim-subspace}
\index{dimension!of a subspace}
Let $E$ be a finite-dimensional $\K$-vector space and $F$ a subspace of
$E$. Then:
\begin{enumerate}[label=(\roman*)]
  \item $F$ is finite-dimensional and $\dim F\leq\dim E$.
  \item $\dim F = \dim E$ if and only if $F = E$.
\end{enumerate}
\end{proposition}

\begin{proof}
\begin{enumerate}[label=(\roman*)]
  \item Any independent family in $F$ is also independent in $E$, so its
    size is at most $\dim E$ (by the Steinitz lemma). Take a maximal
    independent family in $F$; it is a basis of $F$ (if it did not Span
    $F$, it could be enlarged, contradicting maximality). Hence $F$ is
    finite-dimensional and $\dim F\leq\dim E$.

  \item If $\dim F = \dim E = n$, let $(v_1,\ldots,v_n)$ be a basis
    of $F$. This is an independent family of size $n$ in $E$, hence a
    basis of $E$ (a maximal independent family, since no independent
    family in $E$ has more than $\dim E = n$ elements). Thus
    $F = \Span(v_1,\ldots,v_n) = E$.

    Conversely, $F = E$ trivially implies $\dim F = \dim E$. \qedhere
\end{enumerate}
\end{proof}

\begin{proposition}{Dimension of a direct sum}{prop:dim-direct-sum}
\index{dimension!of a direct sum}
If $E = F\oplus G$, then $\dim E = \dim F + \dim G$.
\end{proposition}

\begin{proof}
Let $(v_1,\ldots,v_p)$ be a basis of $F$ and $(w_1,\ldots,w_q)$ a basis
of $G$. We show $(v_1,\ldots,v_p,w_1,\ldots,w_q)$ is a basis of $E$.

\textbf{Spanning:} Any $x\in E$ can be written $x = u + v$ with $u\in F$,
$v\in G$. Expanding $u$ and $v$ in their respective bases gives $x$ as a
combination of $v_1,\ldots,v_p,w_1,\ldots,w_q$.

\textbf{Independence:} Suppose
$\sum_i\alpha_i v_i + \sum_j\beta_j w_j = \vzero$. Then
$\sum_i\alpha_iv_i = -\sum_j\beta_jw_j\in F\cap G = \set{\vzero}$.
Hence $\sum_i\alpha_iv_i = \vzero$ and $\sum_j\beta_jw_j = \vzero$,
which by independence of each basis gives all $\alpha_i = 0$ and all
$\beta_j = 0$.

Therefore the combined family is a basis of $E$ with $p+q$ elements.
\end{proof}

\begin{corollary}{Dimension of a complement}{cor:dim-complement}
If $F$ is a subspace of a finite-dimensional space $E$, then any
complement $G$ of $F$ satisfies $\dim G = \dim E - \dim F$.
\end{corollary}


% ============================================================
\section{The Grassmann formula}\label{sec:grassmann}
% ============================================================

\begin{theorem}{Grassmann formula (dimension of a sum)}{thm:grassmann}
\index{Grassmann formula}
Let $E$ be a finite-dimensional $\K$-vector space and let $F,G$ be
subspaces of $E$. Then
\[
  \dim(F+G) = \dim F + \dim G - \dim(F\cap G).
\]
\end{theorem}

\begin{proof}
Let $p = \dim F$, $q = \dim G$, and $r = \dim(F\cap G)$.

\textbf{Step 1.} Let $(u_1,\ldots,u_r)$ be a basis of $F\cap G$.

\textbf{Step 2.} Since $F\cap G\subseteq F$, we can extend
$(u_1,\ldots,u_r)$ to a basis of $F$:
\[
  (u_1,\ldots,u_r,\,v_1,\ldots,v_s) \quad\text{is a basis of } F,
\]
where $s = p - r$.

\textbf{Step 3.} Similarly, extend $(u_1,\ldots,u_r)$ to a basis of $G$:
\[
  (u_1,\ldots,u_r,\,w_1,\ldots,w_t) \quad\text{is a basis of } G,
\]
where $t = q - r$.

\textbf{Step 4.} We claim that
\[
  \cB = (u_1,\ldots,u_r,\,v_1,\ldots,v_s,\,w_1,\ldots,w_t)
\]
is a basis of $F+G$.

\textit{Spanning:} Any $x\in F+G$ is $x = f+g$ with $f\in F$ and
$g\in G$. Writing $f$ in the basis of $F$ and $g$ in the basis of $G$:
\[
  x = \Bigl(\sum_{i=1}^r \alpha_i u_i + \sum_{j=1}^s \beta_j v_j\Bigr)
    + \Bigl(\sum_{i=1}^r \gamma_i u_i + \sum_{k=1}^t \delta_k w_k\Bigr)
  = \sum_{i=1}^r(\alpha_i+\gamma_i)u_i
    + \sum_{j=1}^s\beta_j v_j
    + \sum_{k=1}^t\delta_k w_k.
\]
So $x\in\Span(\cB)$.

\textit{Independence:} Suppose
\begin{equation}\label{eq:grassmann-indep}
  \sum_{i=1}^r \alpha_i u_i + \sum_{j=1}^s \beta_j v_j
  + \sum_{k=1}^t \delta_k w_k = \vzero.
\end{equation}
Set $h \deff \sum_{k=1}^t\delta_k w_k =
-\sum_{i=1}^r\alpha_i u_i - \sum_{j=1}^s\beta_j v_j$.

The right-hand side lies in $F$ (since $u_i,v_j\in F$). The left-hand
side lies in $G$ (since $w_k\in G$). Hence $h\in F\cap G$.

Since $(u_1,\ldots,u_r)$ is a basis of $F\cap G$, we can write
$h = \sum_{i=1}^r\mu_i u_i$ for some $\mu_i\in\K$. But also
$h = \sum_{k=1}^t\delta_k w_k$, so
\[
  \sum_{k=1}^t\delta_k w_k - \sum_{i=1}^r\mu_i u_i = \vzero.
\]
Since $(u_1,\ldots,u_r,w_1,\ldots,w_t)$ is a basis of $G$ (hence
independent), all $\delta_k = 0$ and all $\mu_i = 0$. In particular,
$h = \vzero$.

Returning to~\eqref{eq:grassmann-indep} with all $\delta_k=0$:
$\sum_i\alpha_iu_i + \sum_j\beta_jv_j = \vzero$. Since
$(u_1,\ldots,u_r,v_1,\ldots,v_s)$ is a basis of $F$ (hence independent),
all $\alpha_i = 0$ and all $\beta_j = 0$.

\textbf{Conclusion.} $\cB$ is a basis of $F+G$ with
$r + s + t = r + (p-r) + (q-r) = p+q-r$ elements, so
\[
  \dim(F+G) = p + q - r = \dim F + \dim G - \dim(F\cap G). \qedhere
\]
\end{proof}

\begin{corollary}{Criterion for direct sum via dimension}{cor:direct-sum-dim}
\index{direct sum!dimension criterion}
Let $F$ and $G$ be subspaces of a finite-dimensional space $E$. Then
$E = F\oplus G$ if and only if $E = F+G$ and $\dim E = \dim F+\dim G$.
Equivalently, any two of the following three conditions imply the third:
\begin{enumerate}[label=(\roman*)]
  \item $E = F+G$,
  \item $F\cap G = \set{\vzero}$,
  \item $\dim E = \dim F + \dim G$.
\end{enumerate}
\end{corollary}

\begin{proof}
By the Grassmann formula, $\dim(F+G) = \dim F + \dim G - \dim(F\cap G)$.
Condition~(ii) says $\dim(F\cap G)=0$. Condition~(i) says
$\dim(F+G)=\dim E$. Condition~(iii) is
$\dim E = \dim F + \dim G$.

If (i) and (ii) hold: $\dim E = \dim(F+G) = \dim F+\dim G$, giving (iii).

If (i) and (iii) hold: $\dim F+\dim G = \dim E = \dim(F+G) =
\dim F+\dim G-\dim(F\cap G)$, so $\dim(F\cap G)=0$, giving (ii).

If (ii) and (iii) hold: $\dim(F+G) = \dim F+\dim G = \dim E$, and since
$F+G\subseteq E$ is a subspace with the same dimension as $E$,

ef{prop:prop:dim-subspace}(ii) gives $F+G=E$, i.e.\ (i).
\end{proof}


% ============================================================
\section{Rank of a family of vectors}\label{sec:rank}
% ============================================================

\begin{definition}{Rank of a family}{def:rank-family}
\index{rank!of a family of vectors}
Let $(v_1,\ldots,v_p)$ be a family of vectors in a $\K$-vector space
$E$. The \emph{rank} of this family is
\[
  \rank(v_1,\ldots,v_p) \deff \dim\Span(v_1,\ldots,v_p).
\]
\end{definition}

\begin{proposition}{Properties of rank}{prop:rank-properties}
Let $\cF = (v_1,\ldots,v_p)$ be a family of vectors in $E$.
\begin{enumerate}[label=(\roman*)]
  \item $0\leq\rank(\cF)\leq p$.
  \item $\rank(\cF) = p$ if and only if $\cF$ is linearly independent.
  \item If $E$ is finite-dimensional, $\cF$ Spans $E$ if and only if
    $\rank(\cF) = \dim E$.
  \item $\cF$ is a basis of $E$ if and only if
    $\rank(\cF) = p = \dim E$.
\end{enumerate}
\end{proposition}

\begin{proof}
\begin{enumerate}[label=(\roman*)]
  \item $\Span(\cF)$ is a subspace Spanned by $p$ vectors, so
    $\dim\Span(\cF)\leq p$. The lower bound is clear.

  \item $\rank(\cF) = p$ means $\Span(\cF)$ has a basis of size $p$.
    Since $\cF$ Spans $\Span(\cF)$, we can extract a basis of size
    $\dim\Span(\cF)=p$ from $\cF$; but $\cF$ already has $p$ elements,
    so it must be the basis itself---hence $\cF$ is independent.
    Conversely, if $\cF$ is independent, it is a basis of $\Span(\cF)$,
    so $\rank(\cF)=p$.

  \item $\cF$ Spans $E$ iff $\Span(\cF) = E$ iff
    $\dim\Span(\cF)=\dim E$ (by 
ef{prop:prop:dim-subspace}(ii)).

  \item Combines (ii) and (iii). \qedhere
\end{enumerate}
\end{proof}

\begin{remark}{Rank of the columns of a matrix}{rem:rank-matrix}
\index{rank!of a matrix}
If $A\in\Mat_{n,p}(\K)$ has columns $C_1,\ldots,C_p\in\K^n$, then
$\rank(C_1,\ldots,C_p) = \dim\Span(C_1,\ldots,C_p)$. This number is
called the \emph{column rank} of $A$. It equals the \emph{row rank} and
is simply called the \emph{rank} of $A$, denoted $\rank(A)$. The proof
of this equality will be given in 
ef{ch:linear-maps}.
\end{remark}


% ============================================================
\section{Applications}\label{sec:vs-applications}
% ============================================================

\subsection{Solution spaces of homogeneous linear
systems}\label{subsec:solution-spaces}

\begin{proposition}{Dimension of the solution space}{prop:solution-space-dim}
\index{solution space!dimension}
Let $A\in\Mat_{n,p}(\K)$ with $\rank(A)=r$. The solution space
$\cS = \setcond{X\in\K^p}{AX=\vzero}$ is a subspace of $\K^p$ of
dimension $p - r$.
\end{proposition}

\begin{proof}
That $\cS$ is a subspace was shown in 
ef{ex:ex:homogeneous-system}. The
dimension formula $\dim\cS = p - r$ is the \emph{rank--nullity theorem},
which will be proved rigorously in 
ef{ch:linear-maps} once the
necessary language of linear maps is available. For now, we note that
Gaussian elimination produces exactly $p - r$ free variables, each of
which gives rise to one basis vector of $\cS$.
\end{proof}

\begin{example}{A concrete system}{ex:concrete-system}
Consider the homogeneous system in $\R^4$:
\[
  \begin{cases}
    x_1 + 2x_2 - x_3 + x_4 = 0, \\
    2x_1 + 4x_2 + x_3 - 2x_4 = 0.
  \end{cases}
\]
The coefficient matrix has rank $2$ (since the two rows are not
proportional after reduction). Hence $\dim\cS = 4-2 = 2$.

Row reduction gives $x_1 = -2x_2 + x_3 - x_4$ and $x_3 = \frac{4}{3}x_4$
(after back-substitution). Setting the free variables $x_2$ and $x_4$ to
$(1,0)$ and $(0,1)$ respectively yields a basis of $\cS$.
\end{example}

\subsection{Polynomial spaces}\label{subsec:polynomial-applications}

\begin{example}{Even and odd polynomials}{ex:even-odd-poly}
\index{polynomial space!even and odd}
Let $E_n = \K_n[X]$ (polynomials of degree $\leq n$). Define
\begin{align*}
  F &= \setcond{P\in E_n}{P(-X)=P(X)} &&\text{(even polynomials)}, \\
  G &= \setcond{P\in E_n}{P(-X)=-P(X)} &&\text{(odd polynomials)}.
\end{align*}
Then $F$ and $G$ are subspaces of $E_n$ and $E_n = F\oplus G$.

Indeed, every polynomial $P$ can be uniquely decomposed as
\[
  P(X) = \underbrace{\frac{P(X)+P(-X)}{2}}_{\in F}
        + \underbrace{\frac{P(X)-P(-X)}{2}}_{\in G}.
\]
A basis of $F$ is $(1,X^2,X^4,\ldots)$ (even powers up to $n$) and a
basis of $G$ is $(X,X^3,X^5,\ldots)$ (odd powers up to $n$). The
dimensions satisfy
\[
  \dim F + \dim G = \left\lfloor\tfrac{n}{2}\right\rfloor+1
  + \left\lfloor\tfrac{n+1}{2}\right\rfloor = n+1 = \dim E_n. \qedhere
\]
\end{example}

\begin{example}{Symmetric and skew-symmetric matrices}{ex:sym-skew}
\index{symmetric matrix}\index{skew-symmetric matrix}
Let $E = \Mat_n(\K)$ with $\Char(\K)\neq 2$. Define
\begin{align*}
  S &= \setcond{A\in E}{\transpose{A}=A}, \\
  A &= \setcond{A\in E}{\transpose{A}=-A}.
\end{align*}
Then $E = S\oplus A$, since every matrix $M$ decomposes uniquely as
\[
  M = \underbrace{\frac{M+\transpose{M}}{2}}_{\in S}
    + \underbrace{\frac{M-\transpose{M}}{2}}_{\in A}.
\]
We have $\dim S = \frac{n(n+1)}{2}$ and $\dim A = \frac{n(n-1)}{2}$,
and indeed $\frac{n(n+1)}{2}+\frac{n(n-1)}{2}=n^2=\dim E$.
\end{example}


% ============================================================
\section{Exercises}\label{sec:ch1-exercises}
% ============================================================

\begin{exercise}{Verifying vector space axioms \basic}{exo:verify-axioms}
Verify that the set $\R^+ = (0,\infty)$ with the operations
\[
  x\oplus y \deff xy, \qquad \lambda\odot x \deff x^\lambda,
\]
is an $\R$-vector space. What is the zero vector? What is the additive
inverse of $x$?
\end{exercise}

\begin{exercise}{Subspace verification \basic}{exo:subspace-check}
Determine which of the following subsets of $\R^3$ are subspaces. Justify
your answers.
\begin{enumerate}[label=(\alph*)]
  \item $S_1 = \setcond{(x,y,z)\in\R^3}{x+2y-z=0}$.
  \item $S_2 = \setcond{(x,y,z)\in\R^3}{x^2+y^2+z^2\leq 1}$.
  \item $S_3 = \setcond{(x,y,z)\in\R^3}{x=y=z}$.
  \item $S_4 = \setcond{(x,y,z)\in\R^3}{xy=0}$.
  \item $S_5 = \setcond{(x,y,z)\in\R^3}{x+y+z=1}$.
\end{enumerate}
\end{exercise}

\begin{exercise}{Subspaces of function spaces \basic}{exo:function-subspace}
Let $E = \R^\R$ be the vector space of all functions $\R\to\R$. Prove
that the following subsets are subspaces:
\begin{enumerate}[label=(\alph*)]
  \item The set of even functions $\setcond{f\in E}{f(-x)=f(x)
    \text{ for all } x}$.
  \item The set of functions vanishing at $0$:
    $\setcond{f\in E}{f(0)=0}$.
  \item The set of periodic functions with period $T>0$:
    $\setcond{f\in E}{f(x+T)=f(x) \text{ for all } x}$.
\end{enumerate}
\end{exercise}

\begin{exercise}{Linear independence in $\R^4$ \basic}{exo:independence-R4}
Determine whether the following vectors in $\R^4$ are linearly
independent:
\[
  v_1 = (1,0,1,0),\quad v_2 = (0,1,0,1),\quad v_3 = (1,1,1,1),\quad
  v_4 = (1,0,0,1).
\]
If they are dependent, find an explicit nontrivial linear relation.
\end{exercise}

\begin{exercise}{Extending to a basis \intermediate}{exo:extend-basis}
In $\R^4$, the vectors $v_1=(1,1,0,0)$ and $v_2=(0,1,1,0)$ are linearly
independent. Extend $\set{v_1,v_2}$ to a basis of $\R^4$ by adding
suitable vectors from the canonical basis.
\end{exercise}

\begin{exercise}{Sum and intersection \intermediate}{exo:sum-intersection}
In $\R^4$, let
\begin{align*}
  F &= \Span\bigl((1,0,1,0),\,(0,1,0,1)\bigr), \\
  G &= \Span\bigl((1,1,0,0),\,(0,0,1,1)\bigr).
\end{align*}
\begin{enumerate}[label=(\alph*)]
  \item Find $\dim F$, $\dim G$, and $\dim(F+G)$.
  \item Determine $F\cap G$ and verify the Grassmann formula.
  \item Is $F+G$ a direct sum?
\end{enumerate}
\end{exercise}

\begin{exercise}{Direct sum decomposition \intermediate}{exo:direct-sum-decomp}
Let $E = \Mat_2(\R)$. Define
\[
  F = \setcond{A\in E}{\transpose{A}=A}, \qquad
  G = \setcond{A\in E}{\transpose{A}=-A}.
\]
\begin{enumerate}[label=(\alph*)]
  \item Show that $E = F\oplus G$.
  \item Find bases for $F$ and $G$, and verify that
    $\dim F+\dim G=\dim E$.
\end{enumerate}
\end{exercise}

\begin{exercise}{Polynomial subspace \intermediate}{exo:poly-subspace}
In $\R_3[X]$ (polynomials of degree $\leq 3$), let
\[
  F = \setcond{P\in\R_3[X]}{P(0)=0 \text{ and } P(1)=0}.
\]
\begin{enumerate}[label=(\alph*)]
  \item Show that $F$ is a subspace of $\R_3[X]$.
  \item Find a basis of $F$ and determine $\dim F$.
  \item Find a complement of $F$ in $\R_3[X]$.
\end{enumerate}
\end{exercise}

\begin{exercise}{The Grassmann formula in practice \intermediate}{exo:grassmann-practice}
Let $E=\R^5$ and
\begin{align*}
  F &= \setcond{(x_1,x_2,x_3,x_4,x_5)\in\R^5}
    {x_1-x_2=0,\; x_3-x_4+x_5=0}, \\
  G &= \setcond{(x_1,x_2,x_3,x_4,x_5)\in\R^5}
    {x_1=0,\; x_2-x_3=0,\; x_4=0}.
\end{align*}
\begin{enumerate}[label=(\alph*)]
  \item Find bases for $F$ and $G$.
  \item Compute $\dim F$, $\dim G$, $\dim(F\cap G)$, and $\dim(F+G)$.
  \item Verify the Grassmann formula.
\end{enumerate}
\end{exercise}

\begin{exercise}{Abstract dimension argument \intermediate}{exo:abstract-dim}
Let $E$ be a vector space of dimension $n$ and let $F,G$ be subspaces
with $\dim F + \dim G > n$. Prove that $F\cap G\neq\set{\vzero}$.
\end{exercise}

\begin{exercise}{Characterizing bases via cardinality \challenging}{exo:basis-card}
Let $E$ be a $\K$-vector space with $\dim E = n$. Let
$\cF = (v_1,\ldots,v_n)$ be a family of $n$ vectors. Prove that the
following are equivalent:
\begin{enumerate}[label=(\roman*)]
  \item $\cF$ is a basis of $E$.
  \item $\cF$ is linearly independent.
  \item $\cF$ Spans $E$.
\end{enumerate}
(Note: In general, one needs \emph{both} independence and Spanning to
guarantee a basis; the point of this exercise is that when the cardinality
equals $\dim E$, either condition alone suffices.)
\end{exercise}

\begin{exercise}{Direct sum of three subspaces \challenging}{exo:direct-sum-three}
In $\R^3$, let
\begin{align*}
  F_1 &= \Span\bigl((1,0,0)\bigr), \\
  F_2 &= \Span\bigl((0,1,0)\bigr), \\
  F_3 &= \Span\bigl((1,1,1)\bigr).
\end{align*}
\begin{enumerate}[label=(\alph*)]
  \item Show that $\R^3 = F_1\oplus F_2\oplus F_3$.
  \item Exhibit three subspaces $G_1,G_2,G_3$ of $\R^3$ such that
    $G_i\cap G_j = \set{\vzero}$ for all $i\neq j$ and
    $G_1+G_2+G_3=\R^3$, but $\R^3\neq G_1\oplus G_2\oplus G_3$.

    (\emph{Hint:} pairwise trivial intersection does not imply direct sum
    for three or more subspaces.)
\end{enumerate}
\end{exercise}

\begin{exercise}{Infinite-dimensional spaces \challenging}{exo:infinite-dim}
\begin{enumerate}[label=(\alph*)]
  \item Prove that $\K[X]$ is infinite-dimensional by showing that the
    family $(1,X,X^2,\ldots)$ is linearly independent.
  \item Show that the function space $\R^\R$ is infinite-dimensional.
    (\emph{Hint:} Consider the family $(\exp_\lambda)_{\lambda\in\R}$
    where $\exp_\lambda(x)=e^{\lambda x}$, or use the inclusion
    $\K[X]\hookrightarrow\R^\R$ via the evaluation map.)
  \item Prove that every vector space has a basis (assuming the Axiom of
    Choice, or equivalently, Zorn's Lemma).
\end{enumerate}
\end{exercise}

\begin{exercise}{Grassmann formula for three subspaces \challenging}{exo:grassmann-three}
Let $F_1,F_2,F_3$ be subspaces of a finite-dimensional space $E$. Prove
the inequality
\[
  \dim(F_1+F_2+F_3) \geq \dim F_1+\dim F_2+\dim F_3
  -\dim(F_1\cap F_2)-\dim(F_1\cap F_3)-\dim(F_2\cap F_3).
\]
Give an example where strict inequality holds.
\end{exercise}


% ============================================================
\section*{Chapter summary}\label{sec:ch1-summary}
\addcontentsline{toc}{section}{Chapter summary}
% ============================================================

\begin{itemize}
  \item A \textbf{vector space} over a field $\K$ is a set $E$ with
    addition and scalar multiplication satisfying eight axioms
    \ref{ax:assoc}--\ref{ax:dist-field}. The central examples are
    $\K^n$, $\Mat_{n,p}(\K)$, polynomial spaces $\K_n[X]$ and $\K[X]$,
    and function spaces.

  \item A \textbf{subspace} is a nonempty subset closed under addition
    and scalar multiplication. The intersection of subspaces is a
    subspace; their union is generally not, and is replaced by their
    \textbf{sum}.

  \item The sum $E = F\oplus G$ is \textbf{direct} when
    $F\cap G = \set{\vzero}$, equivalently when every vector has a
    unique $F$-plus-$G$ decomposition. Every subspace of a
    finite-dimensional space admits a complement.

  \item A \textbf{linear combination} of $v_1,\ldots,v_n$ is a vector
    $\sum\lambda_iv_i$. The \textbf{Span} of a family is the subspace of
    all such combinations. A family is \textbf{linearly independent} if
    the only combination yielding $\vzero$ is the trivial one.

  \item A \textbf{basis} is a family that is both independent and
    Spanning, or equivalently, a family giving a \emph{unique}
    representation for every vector. The \textbf{Steinitz exchange lemma}
    ensures that all bases have the same size; this common size is the
    \textbf{dimension}.

  \item Key dimension formulas:
    \begin{itemize}
      \item $\dim(F\oplus G) = \dim F + \dim G$.
      \item \textbf{Grassmann formula:}
        $\dim(F+G)=\dim F+\dim G-\dim(F\cap G)$.
      \item $F\subseteq E$ and $\dim F=\dim E$ imply $F = E$.
    \end{itemize}

  \item The \textbf{rank} of a family of vectors is the dimension of its
    Span. It equals the maximal number of independent vectors in the
    family.

  \item \textbf{Applications:} The solution space of a homogeneous system
    $AX=\vzero$ is a subspace of $\K^p$ of dimension $p-\rank(A)$.
    Polynomial spaces and matrix spaces admit natural direct sum
    decompositions (e.g.\ even/odd polynomials, symmetric/skew-symmetric
    matrices).
\end{itemize}

% ch2-linear-maps.tex — Chapter 2: Linear Maps and Matrices
% Definitions, kernel/image, rank-nullity theorem, matrix representations,
% change of basis, operations on matrices, dual spaces, applications.

\chapter{Linear Maps and Matrices}\label{ch:linear-maps}

% ============================================================
% Motivating introduction
% ============================================================

The study of vector spaces becomes truly powerful only when we understand
the \emph{maps} between them that respect their algebraic structure.
Consider the Euclidean plane~$\R^2$. Many natural geometric
transformations---rotations, reflections, projections onto a line,
scalings---share a remarkable property: they preserve the operations of
vector addition and scalar multiplication. A rotation through angle
$\theta$ sends the sum of two vectors to the sum of their rotated
images; scaling by a factor~$\lambda$ commutes with addition.

\begin{center}
\begin{tikzpicture}[>=Stealth, scale=1.2]
  % Original vectors
  \draw[->, thick, blue!70!black] (0,0) -- (2,0.5) node[right] {$\vec{v}$};
  \draw[->, thick, red!70!black] (0,0) -- (0.8,1.8) node[above] {$\vec{w}$};
  \draw[->, thick, black!60] (0,0) -- (2.8,2.3) node[above right] {$\vec{v}+\vec{w}$};

  % Rotation arrow
  \draw[->, thick, dashed, gray] (3.5,1.2) to[bend left=20]
    node[midway,above] {\small rotation $R_\theta$} (5.5,1.2);

  % Rotated vectors (approx 45 degrees)
  \begin{scope}[shift={(7,0)}, rotate=40]
    \draw[->, thick, blue!70!black] (0,0) -- (2,0.5)
      node[right] {$R_\theta(\vec{v})$};
    \draw[->, thick, red!70!black] (0,0) -- (0.8,1.8)
      node[above] {$R_\theta(\vec{w})$};
    \draw[->, thick, black!60] (0,0) -- (2.8,2.3)
      node[above right] {$R_\theta(\vec{v}+\vec{w})$};
  \end{scope}
\end{tikzpicture}
\captionof{figure}{A rotation maps the parallelogram law to a rotated
  parallelogram: the image of a sum is the sum of the images.}
\label{fig:rotation-intro}
\end{center}

These structure-preserving maps are called \emph{linear maps}. They form
the bridge between the abstract theory of vector spaces and the concrete
computations of matrix algebra. In this chapter we develop their theory
systematically: we define linear maps, study their kernel and image,
prove the fundamental rank--nullity theorem, represent linear maps by
matrices, and learn to change bases. Along the way we shall see that
matrix multiplication is nothing but the composition of linear maps in
disguise.

Throughout this chapter, $\K$ denotes a field (typically $\R$ or $\C$),
and $E$, $F$, $G$ denote $\K$-vector spaces unless stated otherwise.
We freely use results from 
ef{ch:vector-spaces}, especially the
notions of subspace, basis, and dimension.


% ============================================================
\section{Definition and first properties}\label{sec:linear-map-def}
% ============================================================

\begin{definition}{Linear map}{def:linear-map}
\index{linear map}
Let $E$ and $F$ be vector spaces over~$\K$. A map
$f\colon E \to F$ is \emph{linear} (or is a \emph{$\K$-linear map},
or a \emph{homomorphism of vector spaces}) if it satisfies the
following two conditions:
\begin{enumerate}[label=(\roman*)]
  \item \textbf{Additivity:}
    $f(u+v) = f(u) + f(v)$\quad for all $u,v\in E$.
  \item \textbf{Homogeneity:}
    $f(\lambda u) = \lambda\, f(u)$\quad for all $\lambda\in\K$,
    $u\in E$.
\end{enumerate}
Equivalently, $f$ is linear if and only if
$f(\lambda u + \mu v) = \lambda\, f(u) + \mu\, f(v)$
for all $\lambda,\mu\in\K$ and all $u,v\in E$.
\end{definition}

\begin{remark}{Terminology}{rem:terminology}
\index{endomorphism}\index{isomorphism}\index{automorphism}
Several special names are used:
\begin{itemize}
  \item A linear map $f\colon E\to E$ (same domain and codomain) is
    called an \emph{endomorphism} of~$E$.
  \item A bijective linear map is called an \emph{isomorphism}.
  \item A bijective endomorphism is called an \emph{automorphism}.
  \item A linear map $f\colon E\to\K$ (codomain is the base field)
    is called a \emph{linear form}\index{linear form}.
\end{itemize}
\end{remark}

\begin{proposition}{Elementary properties}{prop:linear-elementary}
Let $f\colon E\to F$ be a linear map. Then:
\begin{enumerate}[label=(\roman*)]
  \item $f(\vzero_E) = \vzero_F$.
  \item $f(-v) = -f(v)$ for all $v\in E$.
  \item $f\!\left(\sum_{i=1}^{n} \lambda_i v_i\right)
    = \sum_{i=1}^{n} \lambda_i\, f(v_i)$
    for all $\lambda_i\in\K$, $v_i\in E$.
\end{enumerate}
\end{proposition}

\begin{proof}
(i) Setting $\lambda = 0$ in homogeneity: $f(\vzero_E) = f(0\cdot v) =
0\cdot f(v) = \vzero_F$.

(ii) $f(-v) = f((-1)\cdot v) = (-1)\cdot f(v) = -f(v)$.

(iii) By induction on~$n$, using additivity and homogeneity at each step.
The base case $n=1$ is homogeneity, and the inductive step follows from
$f\!\left(\sum_{i=1}^{n+1}\lambda_i v_i\right)
= f\!\left(\sum_{i=1}^{n}\lambda_i v_i + \lambda_{n+1}v_{n+1}\right)
= f\!\left(\sum_{i=1}^{n}\lambda_i v_i\right) + \lambda_{n+1}\,f(v_{n+1})$.
\end{proof}

\begin{proposition}{A linear map is determined by its values on a basis}%
{prop:linear-determined-by-basis}
\index{linear map!determined by basis}
Let $\cB = (\vece{1},\ldots,\vece{n})$ be a basis of~$E$, and let
$w_1,\ldots,w_n$ be arbitrary vectors in~$F$. Then there exists a
\emph{unique} linear map $f\colon E\to F$ such that
$f(\vece{i}) = w_i$ for all $i\in\set{1,\ldots,n}$.
\end{proposition}

\begin{proof}
\textit{Existence.}\; Every $v\in E$ has a unique decomposition
$v = \sum_{i=1}^{n}\lambda_i\vece{i}$. Define
$f(v) \deff \sum_{i=1}^{n}\lambda_i w_i$.
One verifies directly that $f$ is linear: for
$v = \sum \lambda_i \vece{i}$ and $v' = \sum \mu_i \vece{i}$ and
$\alpha,\beta\in\K$,
\[
  f(\alpha v + \beta v')
  = f\!\Bigl(\sum_{i}(\alpha\lambda_i+\beta\mu_i)\vece{i}\Bigr)
  = \sum_{i}(\alpha\lambda_i+\beta\mu_i)w_i
  = \alpha\sum_{i}\lambda_i w_i + \beta\sum_{i}\mu_i w_i
  = \alpha f(v)+\beta f(v').
\]
Moreover $f(\vece{i}) = w_i$ since
$\vece{i} = 0\cdot\vece{1}+\cdots+1\cdot\vece{i}+\cdots+0\cdot\vece{n}$.

\textit{Uniqueness.}\; Suppose $g\colon E\to F$ is linear with
$g(\vece{i}) = w_i$. For any $v = \sum \lambda_i\vece{i}$,
$g(v) = \sum \lambda_i\, g(\vece{i}) = \sum \lambda_i w_i = f(v)$.
\end{proof}


% ============================================================
\section{Examples of linear maps}\label{sec:linear-map-examples}
% ============================================================

\begin{example}{Rotation in $\R^2$}{ex:rotation}
\index{rotation}
For a fixed angle $\theta$, the map $R_\theta\colon\R^2\to\R^2$
defined by
\[
  R_\theta\!\begin{pmatrix}x\\y\end{pmatrix}
  = \begin{pmatrix}\cos\theta & -\sin\theta\\
    \sin\theta & \cos\theta\end{pmatrix}
    \begin{pmatrix}x\\y\end{pmatrix}
  = \begin{pmatrix}x\cos\theta - y\sin\theta\\
    x\sin\theta + y\cos\theta\end{pmatrix}
\]
is a linear map (an automorphism, in fact, with inverse~$R_{-\theta}$).
\end{example}

\begin{example}{Orthogonal projection onto a line}{ex:projection-line}
\index{projection!orthogonal}
Let $\ell\subset\R^2$ be the line through the origin with unit
direction vector $\vec{u} = (\cos\alpha,\sin\alpha)^{\mathsf{T}}$.
The orthogonal projection onto $\ell$ is
\[
  \proj_\ell(\vec{v}) = \inner{\vec{v},\vec{u}}\,\vec{u}
  = \begin{pmatrix}\cos^2\alpha & \cos\alpha\sin\alpha\\
    \cos\alpha\sin\alpha & \sin^2\alpha\end{pmatrix}
  \begin{pmatrix}x\\y\end{pmatrix}.
\]
This is linear. Note $\proj_\ell \circ \proj_\ell = \proj_\ell$
(idempotent).
\end{example}

\begin{example}{Reflection}{ex:reflection}
\index{reflection}
The reflection of $\R^2$ across the line $\ell$ through the origin
with direction angle~$\alpha$ is given by
\[
  S_\alpha\!\begin{pmatrix}x\\y\end{pmatrix}
  = \begin{pmatrix}\cos 2\alpha & \sin 2\alpha\\
    \sin 2\alpha & -\cos 2\alpha\end{pmatrix}
  \begin{pmatrix}x\\y\end{pmatrix}.
\]
This is a linear map satisfying $S_\alpha^2 = \Id$.
\end{example}

\begin{example}{Scaling (homothety)}{ex:scaling}
\index{scaling}\index{homothety}
For $\lambda\in\K$, the map $h_\lambda\colon E\to E$ defined by
$h_\lambda(v) = \lambda v$ is linear. It contracts ($|\lambda|<1$),
dilates ($|\lambda|>1$), or reflects ($\lambda<0$) vectors.
\end{example}

\begin{example}{Differentiation}{ex:differentiation}
\index{differentiation}
Let $\cP_n(\R)$ be the space of real polynomials of degree at most~$n$.
The derivative map $D\colon\cP_n(\R)\to\cP_{n-1}(\R)$ defined by
$D(p) = p'$ is linear, since $(p+q)' = p'+q'$ and $(\lambda p)' =
\lambda p'$.
\end{example}

\begin{example}{Integration}{ex:integration}
\index{integration}
The map $I\colon\mathcal{C}([0,1],\R)\to\R$ defined by
$I(f) = \int_0^1 f(t)\,dt$ is a linear form on the vector space of
continuous real-valued functions on $[0,1]$.
\end{example}

\begin{example}{Matrix--vector multiplication}{ex:matrix-vector}
\index{matrix!multiplication}
Let $A\in\Mat_{m,n}(\K)$. The map $f_A\colon\K^n\to\K^m$ defined by
$f_A(x) = Ax$ is linear. Conversely, every linear map
$\K^n\to\K^m$ arises this way (see 
ef{sec:matrix-of-linear-map}).
\end{example}


% ============================================================
\section{Geometric illustrations}\label{sec:geometric-illustrations}
% ============================================================

The following figures illustrate how several linear maps transform
the unit square in~$\R^2$.

\begin{figure}[ht]
\centering
\begin{tikzpicture}[>=Stealth, scale=1.3]
  % --- Rotation ---
  \begin{scope}
    \fill[blue!10] (0,0) -- (1,0) -- (1,1) -- (0,1) -- cycle;
    \draw[blue!50, thick] (0,0) -- (1,0) -- (1,1) -- (0,1) -- cycle;
    \node[below] at (0.5,-0.1) {\small original};
  \end{scope}
  \begin{scope}[shift={(3,0)}, rotate=30]
    \fill[red!15] (0,0) -- (1,0) -- (1,1) -- (0,1) -- cycle;
    \draw[red!60, thick] (0,0) -- (1,0) -- (1,1) -- (0,1) -- cycle;
  \end{scope}
  \node[below] at (3.3,-0.2) {\small rotation ($30^\circ$)};

  % --- Projection ---
  \begin{scope}[shift={(6.5,0)}]
    \fill[blue!10] (0,0) -- (1,0) -- (1,1) -- (0,1) -- cycle;
    \draw[blue!50, thick] (0,0) -- (1,0) -- (1,1) -- (0,1) -- cycle;
    \node[below] at (0.5,-0.1) {\small original};
  \end{scope}
  \begin{scope}[shift={(9.5,0)}]
    \draw[green!60!black, very thick] (0,0) -- (1,0);
    \draw[green!60!black, very thick] (0,0) -- (1,0)
      node[midway, below=4pt] {\small proj.\ onto $x$-axis};
    \fill[green!20, opacity=0.5] (0,0) -- (1,0) -- (1,0) -- (0,0);
  \end{scope}
\end{tikzpicture}

\bigskip

\begin{tikzpicture}[>=Stealth, scale=1.3]
  % --- Shear ---
  \begin{scope}
    \fill[blue!10] (0,0) -- (1,0) -- (1,1) -- (0,1) -- cycle;
    \draw[blue!50, thick] (0,0) -- (1,0) -- (1,1) -- (0,1) -- cycle;
    \node[below] at (0.5,-0.1) {\small original};
  \end{scope}
  \begin{scope}[shift={(3,0)}]
    \fill[orange!15] (0,0) -- (1,0) -- (1.5,1) -- (0.5,1) -- cycle;
    \draw[orange!70, thick] (0,0) -- (1,0) -- (1.5,1) -- (0.5,1) -- cycle;
    \node[below] at (0.75,-0.1) {\small shear ($k=0.5$)};
  \end{scope}

  % --- Scaling ---
  \begin{scope}[shift={(6.5,0)}]
    \fill[blue!10] (0,0) -- (1,0) -- (1,1) -- (0,1) -- cycle;
    \draw[blue!50, thick] (0,0) -- (1,0) -- (1,1) -- (0,1) -- cycle;
    \node[below] at (0.5,-0.1) {\small original};
  \end{scope}
  \begin{scope}[shift={(9.5,0)}, scale=0.6]
    \fill[violet!15] (0,0) -- (1,0) -- (1,1) -- (0,1) -- cycle;
    \draw[violet!60, thick] (0,0) -- (1,0) -- (1,1) -- (0,1) -- cycle;
  \end{scope}
  \node[below] at (9.8,-0.1) {\small scaling ($\lambda=0.6$)};
\end{tikzpicture}

\caption{Effect of linear maps on the unit square: rotation, projection
  onto the $x$-axis, horizontal shear, and uniform scaling.}
\label{fig:linear-map-effects}
\end{figure}


% ============================================================
\section{Kernel and image}\label{sec:kernel-image}
% ============================================================

\begin{definition}{Kernel}{def:kernel}
\index{kernel}
Let $f\colon E\to F$ be a linear map. The \emph{kernel} (or
\emph{null space}) of~$f$ is
\[
  \Ker f \deff \setcond{v\in E}{f(v) = \vzero_F}.
\]
\end{definition}

\begin{definition}{Image}{def:image-linear}
\index{image!of a linear map}
The \emph{image} (or \emph{range}) of~$f$ is
\[
  \im f \deff \setcond{f(v)}{v\in E}
  = \setcond{w\in F}{\exists\, v\in E,\; f(v) = w}.
\]
\end{definition}

\begin{theorem}{Kernel and image are subspaces}{thm:ker-im-subspaces}
Let $f\colon E\to F$ be linear. Then:
\begin{enumerate}[label=(\roman*)]
  \item $\Ker f$ is a subspace of~$E$.
  \item $\im f$ is a subspace of~$F$.
\end{enumerate}
\end{theorem}

\begin{proof}
(i) We verify the three subspace criteria.
\emph{Non-empty:} $f(\vzero_E) = \vzero_F$, so
$\vzero_E\in\Ker f$.
\emph{Closed under addition:} if $u,v\in\Ker f$, then
$f(u+v) = f(u)+f(v) = \vzero_F+\vzero_F = \vzero_F$, so $u+v\in\Ker f$.
\emph{Closed under scalar multiplication:} if $v\in\Ker f$ and
$\lambda\in\K$, then
$f(\lambda v) = \lambda f(v) = \lambda\cdot\vzero_F = \vzero_F$,
so $\lambda v\in\Ker f$.

(ii) \emph{Non-empty:} $\vzero_F = f(\vzero_E)\in\im f$.
\emph{Closed under addition:} if $w_1,w_2\in\im f$, write
$w_i = f(v_i)$; then $w_1+w_2 = f(v_1)+f(v_2) = f(v_1+v_2)\in\im f$.
\emph{Scalar multiplication:}
$\lambda w_1 = \lambda f(v_1) = f(\lambda v_1)\in\im f$.
\end{proof}

\begin{figure}[ht]
\centering
\begin{tikzpicture}[>=Stealth, scale=1.0]
  % E ellipse
  \draw[thick, blue!60] (0,0) ellipse (2cm and 3cm);
  \node[above, blue!80!black, font=\large] at (0,3.2) {$E$};

  % Ker f inside E
  \fill[red!15, opacity=0.7] (0,-0.5) ellipse (1cm and 1.5cm);
  \draw[thick, red!60] (0,-0.5) ellipse (1cm and 1.5cm);
  \node[red!70!black, font=\small] at (0,-0.5) {$\Ker f$};

  % F ellipse
  \draw[thick, green!50!black] (7,0) ellipse (2cm and 3cm);
  \node[above, green!60!black, font=\large] at (7,3.2) {$F$};

  % Im f inside F
  \fill[orange!20, opacity=0.7] (7,0.3) ellipse (1.2cm and 1.8cm);
  \draw[thick, orange!70] (7,0.3) ellipse (1.2cm and 1.8cm);
  \node[orange!70!black, font=\small] at (7,0.3) {$\im f$};

  % zero in F
  \fill (7,-1.2) circle (2pt) node[right, font=\small] {$\vzero_F$};

  % Arrow from Ker f to 0_F
  \draw[->, thick, red!60, dashed] (1,-0.5) to[bend left=15]
    node[midway, above, font=\footnotesize] {$f$} (6.3,-1.2);

  % Arrow from E to Im f
  \draw[->, thick, blue!60] (1.5,1.5) to[bend left=20]
    node[midway, above, font=\footnotesize] {$f$} (5.8,1.5);
\end{tikzpicture}
\caption{The kernel is a subspace of~$E$ mapping to~$\vzero_F$; the
  image is a subspace of~$F$.}
\label{fig:kernel-image}
\end{figure}

\begin{proposition}{Image of a Spanning set}{prop:image-Spanning}
If $E = \Span(v_1,\ldots,v_p)$, then
$\im f = \Span(f(v_1),\ldots,f(v_p))$.
In particular, if $\cB = (\vece{1},\ldots,\vece{n})$ is a basis
of~$E$, then $\im f = \Span(f(\vece{1}),\ldots,f(\vece{n}))$.
\end{proposition}

\begin{proof}
Every $w\in\im f$ has the form $w = f(v)$ where
$v = \sum_{i=1}^{p}\lambda_i v_i$. By linearity,
$w = \sum \lambda_i f(v_i) \in \Span(f(v_1),\ldots,f(v_p))$.
Conversely, each $f(v_i)\in\im f$, and $\im f$ is a subspace, so it
contains $\Span(f(v_1),\ldots,f(v_p))$.
\end{proof}

\begin{definition}{Rank and nullity}{def:rank-nullity}
\index{rank}\index{nullity}
Let $f\colon E\to F$ be a linear map between finite-dimensional spaces.
\begin{itemize}
  \item The \emph{rank} of $f$ is $\rank(f) \deff \dim(\im f)$.
  \item The \emph{nullity} of $f$ is $\dim(\Ker f)$.
\end{itemize}
\end{definition}


% ============================================================
\section{The rank--nullity theorem}\label{sec:rank-nullity}
% ============================================================

\begin{theorem}{Rank--nullity theorem (dimension theorem)}%
{thm:rank-nullity}
\index{rank--nullity theorem}\index{dimension theorem}
Let $E$ be a finite-dimensional vector space and let $f\colon E\to F$
be a linear map. Then $\im f$ is finite-dimensional and
\[
  \boxed{\dim E \;=\; \dim(\Ker f) \;+\; \dim(\im f)
    \;=\; \dim(\Ker f) \;+\; \rank(f).}
\]
\end{theorem}

\begin{proof}
Set $n = \dim E$ and $r = \dim(\Ker f)$.
Let $(u_1,\ldots,u_r)$ be a basis of $\Ker f$. Since $\Ker f$ is a
subspace of the finite-dimensional space~$E$, we may extend this to a
basis of~$E$:
\[
  (u_1,\ldots,u_r,\, v_1,\ldots,v_s)
\]
where $s = n-r$. We claim that $(f(v_1),\ldots,f(v_s))$ is a basis
of $\im f$, which will prove $\dim(\im f) = s = n - r$.

\medskip
\textit{Step 1: $(f(v_1),\ldots,f(v_s))$ Spans $\im f$.}\;
Let $w\in\im f$, say $w = f(x)$ with
$x = \alpha_1 u_1 + \cdots + \alpha_r u_r +
\beta_1 v_1 + \cdots + \beta_s v_s$.
Then
\[
  w = f(x) = \alpha_1 \underbrace{f(u_1)}_{=\,\vzero}
  + \cdots + \alpha_r \underbrace{f(u_r)}_{=\,\vzero}
  + \beta_1 f(v_1) + \cdots + \beta_s f(v_s)
  = \beta_1 f(v_1) + \cdots + \beta_s f(v_s),
\]
so $w\in\Span(f(v_1),\ldots,f(v_s))$.

\medskip
\textit{Step 2: $(f(v_1),\ldots,f(v_s))$ is linearly independent.}\;
Suppose $\beta_1 f(v_1) + \cdots + \beta_s f(v_s) = \vzero_F$.
By linearity, $f(\beta_1 v_1 + \cdots + \beta_s v_s) = \vzero_F$,
so $\beta_1 v_1 + \cdots + \beta_s v_s \in \Ker f$.
Since $(u_1,\ldots,u_r)$ is a basis of $\Ker f$, there exist
$\gamma_1,\ldots,\gamma_r\in\K$ such that
\[
  \beta_1 v_1 + \cdots + \beta_s v_s
  = \gamma_1 u_1 + \cdots + \gamma_r u_r.
\]
Rearranging:
$\gamma_1 u_1 + \cdots + \gamma_r u_r
- \beta_1 v_1 - \cdots - \beta_s v_s = \vzero_E$.
Since $(u_1,\ldots,u_r,v_1,\ldots,v_s)$ is a basis of~$E$, all
coefficients are zero. In particular, $\beta_1 = \cdots = \beta_s = 0$.

\medskip
We conclude that $(f(v_1),\ldots,f(v_s))$ is a basis of $\im f$,
so $\dim(\im f) = s = n - r$, i.e.,
$\dim E = \dim(\Ker f) + \dim(\im f)$.
\end{proof}

\begin{corollary}{Dimension bound on rank}{cor:rank-bound}
For any linear map $f\colon E\to F$ with $E$ finite-dimensional,
\[
  \rank(f) \leq \min\!\bigl(\dim E,\, \dim F\bigr).
\]
\end{corollary}

\begin{proof}
$\rank(f) = \dim(\im f) \leq \dim F$ since $\im f\subseteq F$, and
$\rank(f) = \dim E - \dim(\Ker f) \leq \dim E$.
\end{proof}


% ============================================================
\section{Injectivity, surjectivity, bijectivity}%
\label{sec:inject-surject-biject}
% ============================================================

\begin{theorem}{Injectivity criterion}{thm:inject-kernel}
\index{injectivity!of linear maps}
A linear map $f\colon E\to F$ is injective if and only if
$\Ker f = \set{\vzero_E}$.
\end{theorem}

\begin{proof}
$(\Rightarrow)$\; Suppose $f$ is injective. If $v\in\Ker f$, then
$f(v) = \vzero_F = f(\vzero_E)$. Injectivity gives $v = \vzero_E$.

$(\Leftarrow)$\; Suppose $\Ker f = \set{\vzero_E}$. If
$f(u) = f(v)$, then $f(u-v) = f(u)-f(v) = \vzero_F$, so
$u - v\in\Ker f = \set{\vzero_E}$, hence $u = v$.
\end{proof}

\begin{proposition}{Characterizations in finite dimension}%
{prop:iso-finite-dim}
Let $f\colon E\to F$ be linear with $\dim E = \dim F = n < \infty$.
Then the following are equivalent:
\begin{enumerate}[label=(\roman*)]
  \item $f$ is injective.
  \item $f$ is surjective.
  \item $f$ is bijective (i.e., an isomorphism).
  \item $\rank(f) = n$.
\end{enumerate}
\end{proposition}

\begin{proof}
By the rank--nullity theorem, $\rank(f) = n - \dim(\Ker f)$.

(i) $\Leftrightarrow$ (iv):
$f$ is injective $\Leftrightarrow$ $\Ker f = \set{\vzero}$
$\Leftrightarrow$ $\dim(\Ker f) = 0$ $\Leftrightarrow$ $\rank(f) = n$.

(ii) $\Leftrightarrow$ (iv):
$f$ is surjective $\Leftrightarrow$ $\im f = F$
$\Leftrightarrow$ $\dim(\im f) = \dim F = n$ $\Leftrightarrow$
$\rank(f) = n$.

(iii) $\Leftrightarrow$ (i) $\wedge$ (ii): immediate from the
equivalences above.
\end{proof}

\begin{remark}{Same dimension is essential}{rem:same-dim}
The equivalence of injectivity and surjectivity \emph{fails} when
$\dim E\neq\dim F$. For instance, the inclusion
$\R^2\hookrightarrow\R^3$ (sending $(x,y)\mapsto(x,y,0)$) is
injective but not surjective.
\end{remark}

\begin{definition}{Isomorphism of vector spaces}{def:isomorphism}
\index{isomorphism}
Two vector spaces $E$ and $F$ over $\K$ are \emph{isomorphic},
written $E\cong F$, if there exists an isomorphism $f\colon E\to F$.
\end{definition}

\begin{theorem}{Classification by dimension}{thm:classification-dim}
\index{isomorphism!classification}
Two finite-dimensional $\K$-vector spaces are isomorphic if and only
if they have the same dimension. In particular, every
$n$-dimensional $\K$-vector space is isomorphic to~$\K^n$.
\end{theorem}

\begin{proof}
$(\Rightarrow)$\; If $f\colon E\isoto F$ is an isomorphism and
$(\vece{1},\ldots,\vece{n})$ is a basis of~$E$, then
$(f(\vece{1}),\ldots,f(\vece{n}))$ is a basis of~$F$
(an injective linear map sends linearly independent sets to linearly
independent sets, and a surjective map ensures they Span~$F$). Hence
$\dim F = n = \dim E$.

$(\Leftarrow)$\; Let $\dim E = \dim F = n$. Choose bases
$(\vece{1},\ldots,\vece{n})$ of~$E$ and
$(\vecf{1},\ldots,\vecf{n})$ of~$F$. By

ef{prop:prop:linear-determined-by-basis}, there is a unique linear map
$f\colon E\to F$ with $f(\vece{i}) = \vecf{i}$. Since $f$ maps
a basis to a basis, it is bijective.
\end{proof}


% ============================================================
\section{Composition of linear maps}\label{sec:composition}
% ============================================================

\begin{proposition}{Composition is linear}{prop:composition-linear}
\index{composition!of linear maps}
If $f\colon E\to F$ and $g\colon F\to G$ are linear maps, then the
composition $g\circ f\colon E\to G$ is linear.
\end{proposition}

\begin{proof}
For all $\lambda,\mu\in\K$ and $u,v\in E$:
\[
  (g\circ f)(\lambda u+\mu v)
  = g\!\bigl(f(\lambda u+\mu v)\bigr)
  = g\!\bigl(\lambda f(u)+\mu f(v)\bigr)
  = \lambda\, g(f(u))+\mu\, g(f(v))
  = \lambda\,(g\circ f)(u)+\mu\,(g\circ f)(v).\qedhere
\]
\end{proof}

\begin{proposition}{Properties of composition}{prop:composition-props}
Let $f\colon E\to F$, $g\colon F\to G$, $h\colon G\to H$ be linear.
\begin{enumerate}[label=(\roman*)]
  \item \textbf{Associativity:} $h\circ(g\circ f) = (h\circ g)\circ f$.
  \item \textbf{Identity:} $f\circ\Id_E = f = \Id_F\circ f$.
  \item If $f$ is an isomorphism, its inverse $\inv{f}\colon F\to E$
    is also linear, with $\inv{f}\circ f = \Id_E$ and
    $f\circ\inv{f} = \Id_F$.
\end{enumerate}
\end{proposition}

\begin{proof}
(i) and (ii) follow from the corresponding properties of function
composition. For (iii), let $w_1,w_2\in F$ and $\lambda\in\K$. Write
$w_i = f(v_i)$. Then $\inv{f}(\lambda w_1 + w_2) =
\inv{f}(\lambda f(v_1)+f(v_2)) = \inv{f}(f(\lambda v_1+v_2)) =
\lambda v_1+v_2 = \lambda\inv{f}(w_1)+\inv{f}(w_2)$.
\end{proof}

\begin{proposition}{Kernel and image under composition}%
{prop:ker-im-composition}
Let $f\colon E\to F$ and $g\colon F\to G$ be linear. Then:
\begin{enumerate}[label=(\roman*)]
  \item $\Ker f \subseteq \Ker(g\circ f)$.
  \item $\im(g\circ f) \subseteq \im g$.
  \item If $g$ is injective, then $\Ker(g\circ f) = \Ker f$.
  \item If $f$ is surjective, then $\im(g\circ f) = \im g$.
\end{enumerate}
\end{proposition}

\begin{proof}
(i) If $v\in\Ker f$, then $(g\circ f)(v) = g(\vzero_F) = \vzero_G$.

(ii) Every element of $\im(g\circ f)$ has the form $g(f(v))\in\im g$.

(iii) $v\in\Ker(g\circ f)$ means $g(f(v))=\vzero_G$. If $g$ is
injective, $f(v) = \vzero_F$, so $v\in\Ker f$.

(iv) If $f$ is surjective, every $w\in F$ is $f(v)$ for some $v$,
so $g(w) = g(f(v)) \in \im(g\circ f)$. Hence $\im g\subseteq\im(g\circ f)$.
\end{proof}


% ============================================================
\section{The vector space of linear maps}\label{sec:space-of-linear-maps}
% ============================================================

\begin{theorem}{$\cL(E,F)$ is a vector space}{thm:L-E-F}
\index{linear map!space of}
Let $E$ and $F$ be vector spaces over~$\K$. Define operations on the
set $\cL(E,F)$ of all linear maps from $E$ to $F$:
\begin{itemize}
  \item \textbf{Addition:} $(f+g)(v) \deff f(v)+g(v)$.
  \item \textbf{Scalar multiplication:} $(\lambda f)(v) \deff \lambda\,f(v)$.
\end{itemize}
Then $\cL(E,F)$ is a $\K$-vector space, with zero element the zero map
$\vzero\colon v\mapsto \vzero_F$.
\end{theorem}

\begin{proof}
One must verify the vector space axioms. We check that $f+g$ and
$\lambda f$ are indeed linear (so that $\cL(E,F)$ is closed under
these operations), and the remaining axioms are inherited from those
of~$F$ applied pointwise.

\emph{$f+g$ is linear:}
$(f+g)(\alpha u + \beta v)
= f(\alpha u+\beta v) + g(\alpha u+\beta v)
= \alpha f(u)+\beta f(v)+\alpha g(u)+\beta g(v)
= \alpha(f+g)(u)+\beta(f+g)(v)$.

\emph{$\lambda f$ is linear:}
$(\lambda f)(\alpha u+\beta v)
= \lambda f(\alpha u+\beta v)
= \lambda(\alpha f(u)+\beta f(v))
= \alpha(\lambda f)(u)+\beta(\lambda f)(v)$.

Associativity of addition, commutativity, etc., all follow from the
corresponding properties in~$F$.
\end{proof}

\begin{proposition}{Dimension of $\cL(E,F)$}{prop:dim-L-E-F}
If $\dim E = n$ and $\dim F = m$, then $\dim\cL(E,F) = nm$.
\end{proposition}

\begin{proof}
This follows from the isomorphism $\cL(E,F)\cong\Mat_{m,n}(\K)$
established in 
ef{sec:matrix-of-linear-map}: the space of
$m\times n$ matrices has dimension~$mn$.
\end{proof}

\begin{remark}{The endomorphism algebra}{rem:endomorphism-algebra}
\index{endomorphism!algebra}
The space $\End(E) = \cL(E,E)$ carries an additional operation:
composition. This makes $\End(E)$ a \emph{$\K$-algebra} (a vector
space equipped with an associative bilinear multiplication and a
unit element~$\Id_E$). Note that composition is not commutative
in general.
\end{remark}


% ============================================================
\section{Matrix of a linear map}\label{sec:matrix-of-linear-map}
% ============================================================

\begin{definition}{Matrix representation}{def:matrix-representation}
\index{matrix!of a linear map}
Let $E$ and $F$ be finite-dimensional with ordered bases
$\cB = (\vece{1},\ldots,\vece{n})$ and
$\cC = (\vecf{1},\ldots,\vecf{m})$ respectively.
Let $f\colon E\to F$ be linear. For each $j\in\set{1,\ldots,n}$, write
\[
  f(\vece{j}) = \sum_{i=1}^{m} a_{ij}\,\vecf{i}.
\]
The $m\times n$ matrix
\[
  \Mat_{\cC,\cB}(f) \deff (a_{ij})_{\substack{1\le i\le m\\1\le j\le n}}
  = \begin{pmatrix}
    a_{11} & a_{12} & \cdots & a_{1n}\\
    a_{21} & a_{22} & \cdots & a_{2n}\\
    \vdots & \vdots & \ddots & \vdots\\
    a_{m1} & a_{m2} & \cdots & a_{mn}
  \end{pmatrix}
\]
is called the \emph{matrix of~$f$ relative to the bases $\cB$
and~$\cC$}. The $j$-th column is the coordinate vector of $f(\vece{j})$
in basis~$\cC$.
\end{definition}

\begin{remark}{Convention}{rem:matrix-convention}
When $E = F$ and $\cB = \cC$, we simply write $\Mat_\cB(f)$.
When working in $\K^n$ with the canonical basis~$\cE$, we identify
$f$ with the matrix $A = \Mat_\cE(f)$ and write $f(x) = Ax$.
\end{remark}

\begin{example}{Matrix of a rotation}{ex:matrix-rotation}
The rotation $R_\theta\colon\R^2\to\R^2$ of

ef{ex:ex:rotation} has matrix in the canonical basis:
\[
  \Mat_\cE(R_\theta)
  = \begin{pmatrix}\cos\theta & -\sin\theta\\
    \sin\theta & \cos\theta\end{pmatrix}.
\]
The first column is $R_\theta(\vece{1}) = (\cos\theta,\sin\theta)^{\mathsf{T}}$
and the second column is $R_\theta(\vece{2}) = (-\sin\theta,\cos\theta)^{\mathsf{T}}$.
\end{example}

\begin{example}{Matrix of differentiation}{ex:matrix-diff}
Consider $D\colon\cP_3(\R)\to\cP_3(\R)$, $D(p) = p'$, with basis
$\cB = (1,t,t^2,t^3)$. Then $D(1)=0$, $D(t)=1$, $D(t^2)=2t$,
$D(t^3)=3t^2$, so
\[
  \Mat_\cB(D) = \begin{pmatrix}
    0 & 1 & 0 & 0\\
    0 & 0 & 2 & 0\\
    0 & 0 & 0 & 3\\
    0 & 0 & 0 & 0
  \end{pmatrix}.
\]
\end{example}

\begin{theorem}{The matrix representation is an isomorphism}%
{thm:matrix-iso}
The map
\[
  \Phi\colon \cL(E,F) \to \Mat_{m,n}(\K), \qquad
  f \mapsto \Mat_{\cC,\cB}(f)
\]
is an isomorphism of vector spaces. Moreover:
\begin{enumerate}[label=(\roman*)]
  \item $\Mat_{\cC,\cB}(f+g) = \Mat_{\cC,\cB}(f) + \Mat_{\cC,\cB}(g)$.
  \item $\Mat_{\cC,\cB}(\lambda f) = \lambda\,\Mat_{\cC,\cB}(f)$.
  \item If $g\colon F\to G$ is linear with $G$ having basis~$\cD$, then
    \[
      \Mat_{\cD,\cB}(g\circ f)
      = \Mat_{\cD,\cC}(g)\cdot\Mat_{\cC,\cB}(f).
    \]
\end{enumerate}
\end{theorem}

\begin{proof}
Linearity of $\Phi$ (parts (i) and (ii)) follows directly from the
definitions. Bijectivity follows from

ef{prop:prop:linear-determined-by-basis}: for every matrix
$A\in\Mat_{m,n}(\K)$, there is a unique $f\in\cL(E,F)$ with
$\Phi(f) = A$ (define $f$ by $f(\vece{j}) = \sum_i a_{ij}\vecf{i}$).

(iii) Let $A = \Mat_{\cC,\cB}(f)$ and $B = \Mat_{\cD,\cC}(g)$.
The $j$-th column of $\Mat_{\cD,\cB}(g\circ f)$ is the coordinate
vector of $(g\circ f)(\vece{j}) = g(f(\vece{j}))$ in basis~$\cD$.
Now $f(\vece{j}) = \sum_{k=1}^m a_{kj}\vecf{k}$, so
\[
  g(f(\vece{j})) = \sum_{k=1}^m a_{kj}\,g(\vecf{k})
  = \sum_{k=1}^m a_{kj}\sum_{i=1}^\ell b_{ik}\,\vec{d}_i
  = \sum_{i=1}^\ell\Bigl(\sum_{k=1}^m b_{ik}\,a_{kj}\Bigr)\vec{d}_i.
\]
The coefficient of $\vec{d}_i$ is $\sum_k b_{ik}a_{kj} = (BA)_{ij}$,
proving $\Mat_{\cD,\cB}(g\circ f) = BA$.
\end{proof}

\begin{remark}{Matrix multiplication = composition}{rem:mult-comp}
\index{matrix!multiplication!and composition}

ef{thm:thm:matrix-iso}(iii) reveals the true meaning of matrix
multiplication: it is the algebraic encoding of the composition of
linear maps. The product $BA$ of matrices is defined precisely so that
the matrix of $g\circ f$ equals the product of the matrices of $g$
and~$f$.
\end{remark}


% ============================================================
\section{Change of basis}\label{sec:change-of-basis}
% ============================================================

\begin{definition}{Transition matrix (change-of-basis matrix)}%
{def:transition-matrix}
\index{transition matrix}\index{change of basis}
Let $\cB = (\vece{1},\ldots,\vece{n})$ and
$\cB' = (\vece{1}',\ldots,\vece{n}')$ be two bases of~$E$.
The \emph{transition matrix} from $\cB$ to $\cB'$ is
\[
  P_{\cB\to\cB'} \deff \Mat_{\cB}(\Id_E,\cB',\cB)
  = \bigl[\,[\vece{1}']_\cB \;\; [\vece{2}']_\cB \;\; \cdots \;\;
    [\vece{n}']_\cB\,\bigr],
\]
where $[\vece{j}']_\cB$ is the column of coordinates of $\vece{j}'$
in basis~$\cB$.
\end{definition}

\begin{proposition}{Properties of transition matrices}%
{prop:transition-props}
\begin{enumerate}[label=(\roman*)]
  \item $P_{\cB\to\cB'}$ is invertible, with
    $\inv{(P_{\cB\to\cB'})} = P_{\cB'\to\cB}$.
  \item If $[v]_\cB$ and $[v]_{\cB'}$ denote the coordinate
    vectors of $v\in E$ in the two bases, then
    \[
      [v]_\cB = P_{\cB\to\cB'}\,[v]_{\cB'}.
    \]
  \item For three bases $\cB,\cB',\cB''$:
    $P_{\cB\to\cB''} = P_{\cB\to\cB'}\cdot P_{\cB'\to\cB''}$.
\end{enumerate}
\end{proposition}

\begin{proof}
(i) The transition matrix is the matrix of $\Id_E$ relative to
$\cB'$ (source) and $\cB$ (target). Since $\Id_E$ is an isomorphism,
its matrix is invertible. The inverse is the matrix of $\Id_E$ relative
to $\cB$ and $\cB'$, which is $P_{\cB'\to\cB}$.

(ii) Let $v = \sum_j x_j'\,\vece{j}' = \sum_i x_i\,\vece{i}$. Then
$[v]_\cB = (x_1,\ldots,x_n)^{\mathsf{T}}$ and
$[v]_{\cB'} = (x_1',\ldots,x_n')^{\mathsf{T}}$. Since
$\vece{j}' = \sum_i p_{ij}\vece{i}$ where $P = (p_{ij})$,
$v = \sum_j x_j' \sum_i p_{ij}\vece{i} = \sum_i(\sum_j p_{ij}x_j')\vece{i}$,
giving $x_i = \sum_j p_{ij}x_j'$, i.e., $[v]_\cB = P[v]_{\cB'}$.

(iii) Follows from (i) and the multiplicativity of matrix
representations under composition:
$\Id_E = \Id_E\circ\Id_E$.
\end{proof}

\begin{theorem}{Change-of-basis formula for linear maps}%
{thm:change-of-basis}
\index{change of basis!formula}
Let $f\in\End(E)$. Let $\cB$ and $\cB'$ be bases of~$E$, and let
$P = P_{\cB\to\cB'}$. If $A = \Mat_\cB(f)$ and $A' = \Mat_{\cB'}(f)$,
then
\[
  \boxed{A' = \inv{P}\, A\, P.}
\]
More generally, if $f\colon E\to F$ with bases $\cB,\cB'$ of~$E$ and
$\cC,\cC'$ of~$F$, and $P = P_{\cB\to\cB'}$, $Q = P_{\cC\to\cC'}$, then
\[
  \Mat_{\cC',\cB'}(f) = \inv{Q}\,\Mat_{\cC,\cB}(f)\,P.
\]
\end{theorem}

\begin{proof}
We have $f = \Id_F\circ f\circ\Id_E$. Using

ef{thm:thm:matrix-iso}(iii):
\begin{align*}
  \Mat_{\cC',\cB'}(f)
  &= \Mat_{\cC',\cC'}(\Id_F)\cdot\Mat_{\cC',\cB'}(f) \\
  &= \Mat_{\cC',\cC}(\Id_F)\cdot\Mat_{\cC,\cB}(f)\cdot
    \Mat_{\cB,\cB'}(\Id_E) \\
  &= \inv{Q}\cdot A \cdot P.
\end{align*}
Here we used $\Mat_{\cC',\cC}(\Id_F) = \inv{(P_{\cC\to\cC'})} = \inv{Q}$
and $\Mat_{\cB,\cB'}(\Id_E) = P_{\cB\to\cB'} = P$.

For the special case $E = F$, $\cB = \cC$, $\cB' = \cC'$, $P = Q$:
$A' = \inv{P}AP$.
\end{proof}

\begin{definition}{Similar matrices}{def:similar}
\index{similar matrices}
Two matrices $A,A'\in\Mat_n(\K)$ are \emph{similar} if there exists an
invertible matrix $P\in\GL_n(\K)$ such that $A' = \inv{P}AP$. Similar
matrices represent the same endomorphism in different bases.
\end{definition}

\begin{example}{Change of basis for a projection}{ex:change-basis-proj}
Consider the projection $\pi\colon\R^2\to\R^2$ onto the $x$-axis,
with canonical matrix $A = \begin{psmallmatrix}1&0\\0&0\end{psmallmatrix}$.
In the basis $\cB' = \bigl((1,1)^{\mathsf{T}},\,(1,-1)^{\mathsf{T}}\bigr)$,
the transition matrix is
$P = \begin{psmallmatrix}1&1\\1&-1\end{psmallmatrix}$, with
$\inv{P} = \tfrac{1}{2}\begin{psmallmatrix}1&1\\1&-1\end{psmallmatrix}$.
Then
\[
  A' = \inv{P}AP
  = \frac{1}{2}\begin{pmatrix}1&1\\1&-1\end{pmatrix}
    \begin{pmatrix}1&0\\0&0\end{pmatrix}
    \begin{pmatrix}1&1\\1&-1\end{pmatrix}
  = \frac{1}{2}\begin{pmatrix}1&1\\1&-1\end{pmatrix}
    \begin{pmatrix}1&1\\0&0\end{pmatrix}
  = \begin{pmatrix}\tfrac{1}{2}&\tfrac{1}{2}\\[3pt]
    \tfrac{1}{2}&\tfrac{1}{2}\end{pmatrix}.
\]
\end{example}


% ============================================================
\section{Operations on matrices}\label{sec:matrix-operations}
% ============================================================

We now consolidate the algebraic operations on matrices, all of which
are motivated by the corresponding operations on linear maps.

\begin{definition}{Matrix addition and scalar multiplication}%
{def:matrix-add-scalar}
\index{matrix!addition}\index{matrix!scalar multiplication}
Let $A = (a_{ij})$ and $B = (b_{ij})$ be $m\times n$ matrices over~$\K$.
\begin{enumerate}[label=(\roman*)]
  \item $A + B \deff (a_{ij}+b_{ij})$.
  \item $\lambda A \deff (\lambda\,a_{ij})$ for $\lambda\in\K$.
\end{enumerate}
With these operations, $\Mat_{m,n}(\K)$ is a vector space of
dimension~$mn$.
\end{definition}

\begin{definition}{Matrix multiplication}{def:matrix-mult}
\index{matrix!multiplication}
Let $A = (a_{ij})\in\Mat_{m,n}(\K)$ and
$B = (b_{jk})\in\Mat_{n,p}(\K)$. Their product is the matrix
$AB = (c_{ik})\in\Mat_{m,p}(\K)$ where
\[
  c_{ik} = \sum_{j=1}^{n} a_{ij}\,b_{jk}.
\]
This is the ``row-by-column'' rule: the $(i,k)$-entry of $AB$ is the
dot product of the $i$-th row of $A$ with the $k$-th column of~$B$.
\end{definition}

\begin{proposition}{Properties of matrix multiplication}%
{prop:matrix-mult-props}
\begin{enumerate}[label=(\roman*)]
  \item \textbf{Associativity:} $A(BC) = (AB)C$ (when dimensions match).
  \item \textbf{Distributivity:} $A(B+C) = AB+AC$ and $(A+B)C = AC+BC$.
  \item \textbf{Scalar compatibility:} $\lambda(AB) = (\lambda A)B = A(\lambda B)$.
  \item \textbf{Identity:} $I_m A = A = A I_n$ for
    $A\in\Mat_{m,n}(\K)$.
  \item \textbf{Non-commutativity:} In general, $AB\neq BA$.
\end{enumerate}
\end{proposition}

\begin{definition}{Transpose}{def:transpose}
\index{transpose}
The \emph{transpose} of $A = (a_{ij})\in\Mat_{m,n}(\K)$ is
$\transpose{A} = (a_{ji})\in\Mat_{n,m}(\K)$.
\end{definition}

\begin{proposition}{Properties of the transpose}{prop:transpose-props}
\begin{enumerate}[label=(\roman*)]
  \item $\transpose{(A+B)} = \transpose{A}+\transpose{B}$.
  \item $\transpose{(\lambda A)} = \lambda\,\transpose{A}$.
  \item $\transpose{(\transpose{A})} = A$.
  \item $\transpose{(AB)} = \transpose{B}\,\transpose{A}$
    \quad (reversal of order).
\end{enumerate}
\end{proposition}

\begin{proof}
We prove (iv); the others are immediate.
Let $A\in\Mat_{m,n}(\K)$, $B\in\Mat_{n,p}(\K)$. The $(k,i)$-entry of
$\transpose{(AB)}$ is $(AB)_{ik} = \sum_j a_{ij}b_{jk}$.
The $(k,i)$-entry of $\transpose{B}\,\transpose{A}$ is
$\sum_j (\transpose{B})_{kj}(\transpose{A})_{ji}
= \sum_j b_{jk}\,a_{ij} = \sum_j a_{ij}b_{jk}$.
\end{proof}


% ============================================================
\section{Inverse matrices}\label{sec:inverse-matrices}
% ============================================================

\begin{definition}{Invertible matrix}{def:invertible}
\index{matrix!invertible}\index{invertible matrix}
A square matrix $A\in\Mat_n(\K)$ is \emph{invertible} (or
\emph{non-singular}) if there exists a matrix $B\in\Mat_n(\K)$ such
that $AB = BA = I_n$. The matrix $B$ is unique and is denoted
$\inv{A}$.
\end{definition}

\begin{proposition}{Properties of inverses}{prop:inverse-props}
Let $A,B\in\GL_n(\K)$ (i.e., $A$ and $B$ are invertible).
\begin{enumerate}[label=(\roman*)]
  \item $\inv{(\inv{A})} = A$.
  \item $\inv{(AB)} = \inv{B}\,\inv{A}$ \quad (reversal of order).
  \item $\inv{(\transpose{A})} = \transpose{(\inv{A})}$.
  \item $\inv{(\lambda A)} = \lambda^{-1}\inv{A}$ for $\lambda\neq 0$.
\end{enumerate}
\end{proposition}

\begin{proof}
(ii) We verify: $(AB)(\inv{B}\inv{A}) = A(B\inv{B})\inv{A}
= AI_n\inv{A} = A\inv{A} = I_n$, and similarly
$(\inv{B}\inv{A})(AB) = I_n$.
The remaining parts are proved similarly.
\end{proof}

\begin{theorem}{Invertibility criteria}{thm:invertibility}
\index{matrix!invertibility criteria}
For $A\in\Mat_n(\K)$, the following are equivalent:
\begin{enumerate}[label=(\roman*)]
  \item $A$ is invertible.
  \item The linear map $x\mapsto Ax$ is an isomorphism of $\K^n$.
  \item $\Ker(A) = \set{\vzero}$ (i.e., $Ax = 0$ implies $x = 0$).
  \item The columns of $A$ form a basis of $\K^n$.
  \item $\rank(A) = n$.
\end{enumerate}
\end{theorem}

\begin{proof}
This follows from 
ef{prop:prop:iso-finite-dim} applied to the linear map
$f_A\colon\K^n\to\K^n$, $f_A(x) = Ax$. The matrix $A$ is invertible if
and only if $f_A$ is bijective (i.e., an isomorphism). The equivalences
(ii)--(v) were established in 
ef{prop:prop:iso-finite-dim} and the
preceding discussion.
\end{proof}

\begin{example}{Inverse of a $2\times 2$ matrix}{ex:inverse-2x2}
\index{matrix!inverse!$2\times 2$}
For $A = \begin{psmallmatrix}a&b\\c&d\end{psmallmatrix}$ with
$ad-bc\neq 0$:
\[
  \inv{A} = \frac{1}{ad-bc}\begin{pmatrix}d & -b\\-c & a\end{pmatrix}.
\]
\end{example}


% ============================================================
\section{Dual spaces}\label{sec:dual-spaces}
% ============================================================

\begin{definition}{Dual space}{def:dual-space}
\index{dual space}\index{linear form}
Let $E$ be a vector space over~$\K$. The \emph{dual space} of $E$ is
\[
  E^* \deff \cL(E,\K) = \setcond{\vphi\colon E\to\K}{\vphi
    \text{ is linear}}.
\]
Elements of $E^*$ are called \emph{linear forms} (or \emph{covectors}
or \emph{linear functionals}).
\end{definition}

\begin{example}{Linear forms on $\R^n$}{ex:linear-form-Rn}
Every linear form on $\R^n$ has the form
$\vphi(\vec{x}) = a_1 x_1 + a_2 x_2 + \cdots + a_n x_n$
for some fixed scalars $a_1,\ldots,a_n\in\R$.
\end{example}

\begin{definition}{Dual basis}{def:dual-basis}
\index{dual basis}
Let $\cB = (\vece{1},\ldots,\vece{n})$ be a basis of the
finite-dimensional space~$E$. For each $i\in\set{1,\ldots,n}$, define
the linear form $\vece{i}^*\in E^*$ by
\[
  \vece{i}^*(\vece{j}) = \delta_{ij}
  = \begin{cases}1 & \text{if } i=j,\\0 & \text{if } i\neq j.\end{cases}
\]
The family $\cB^* = (\vece{1}^*,\ldots,\vece{n}^*)$ is called the
\emph{dual basis} of~$\cB$.
\end{definition}

\begin{theorem}{The dual basis is a basis of $E^*$}{thm:dual-basis}
If $\dim E = n$, then $\cB^*$ is a basis of $E^*$, and
$\dim E^* = n$. Every linear form $\vphi\in E^*$ can be written
uniquely as
\[
  \vphi = \sum_{i=1}^{n}\vphi(\vece{i})\,\vece{i}^*.
\]
\end{theorem}

\begin{proof}
\textit{Spanning.}\; Let $\vphi\in E^*$ and set
$\psi = \sum_{i=1}^n \vphi(\vece{i})\,\vece{i}^*$. For each basis
vector $\vece{j}$:
$\psi(\vece{j}) = \sum_i \vphi(\vece{i})\,\vece{i}^*(\vece{j})
= \sum_i \vphi(\vece{i})\,\delta_{ij} = \vphi(\vece{j})$.
Since $\vphi$ and $\psi$ agree on a basis, they are equal
(
ef{prop:prop:linear-determined-by-basis}).

\textit{Linear independence.}\; Suppose
$\sum_i \lambda_i\,\vece{i}^* = 0$ (the zero form). Evaluating at
$\vece{j}$: $\sum_i \lambda_i\,\delta_{ij} = \lambda_j = 0$ for
each~$j$.
\end{proof}

\begin{remark}{Finite dimension is essential}{rem:dual-infinite}
In infinite dimension, $E^*$ is generally much larger than~$E$.
For instance, if $E = \K[t]$ (polynomials of all degrees), then
$E^*$ is not isomorphic to~$E$. The theory of dual spaces becomes
significantly richer (and subtler) in that setting.
\end{remark}


% ============================================================
\section{Applications}\label{sec:applications-ch2}
% ============================================================

\subsection{Computer graphics transformations}%
\label{subsec:computer-graphics}
\index{computer graphics}

Every 2D transformation used in computer graphics---translation,
rotation, scaling, shearing---can be expressed as a linear map (or, in
the case of translations, an \emph{affine} map made linear by
\emph{homogeneous coordinates}).

In homogeneous coordinates, a point $(x,y)$ is represented as the
triple $(x,y,1)^{\mathsf{T}}$, and transformations become
$3\times 3$ matrices:
\[
  \underbrace{\begin{pmatrix}
    \cos\theta & -\sin\theta & t_x\\
    \sin\theta & \cos\theta & t_y\\
    0 & 0 & 1
  \end{pmatrix}}_{\text{rotate by $\theta$, translate by $(t_x,t_y)$}}
  \begin{pmatrix}x\\y\\1\end{pmatrix}
  = \begin{pmatrix}
    x\cos\theta - y\sin\theta + t_x\\
    x\sin\theta + y\cos\theta + t_y\\
    1
  \end{pmatrix}.
\]

Composing transformations corresponds to multiplying their matrices.
This is the foundation of the rendering pipeline in OpenGL, DirectX,
and every modern graphics system: the entire scene geometry is
transformed by chains of matrix multiplications.

\subsection{Preview: Markov chains}\label{subsec:markov-preview}
\index{Markov chain}

A \emph{Markov chain} models a system that transitions between finitely
many states with fixed probabilities. The transition probabilities are
encoded in a \emph{stochastic matrix}
$P = (p_{ij})\in\Mat_n(\R)$ where $p_{ij}\geq 0$ and
$\sum_j p_{ij} = 1$ for each row~$i$.

If $\vec{x}^{(k)}\in\R^n$ is the probability distribution at time
step~$k$ (a row vector), then
\[
  \vec{x}^{(k+1)} = \vec{x}^{(k)}\cdot P.
\]
The long-term behavior of the chain is governed by the eigenvalues and
eigenvectors of~$P$---a topic we will study in detail in

ef{ch:eigenvalues}. For now, observe that computing the state after
$k$ steps amounts to computing $\vec{x}^{(0)}\cdot P^k$, a purely
linear-algebraic problem.


% ============================================================
\section{Exercises}\label{sec:exercises-ch2}
% ============================================================

\begin{exercise}{Verifying linearity \basic}{exo:verify-linearity}
Which of the following maps are linear? Justify each answer.
\begin{enumerate}[label=(\alph*)]
  \item $f\colon\R^2\to\R$, $f(x,y) = 3x - 2y$.
  \item $g\colon\R^2\to\R$, $g(x,y) = x^2 + y$.
  \item $h\colon\R^3\to\R^2$, $h(x,y,z) = (x+z,\, 2y-x)$.
  \item $k\colon\R^2\to\R^2$, $k(x,y) = (x+1,\, y)$.
\end{enumerate}
\end{exercise}

\begin{exercise}{Kernel and image computation \basic}{exo:ker-im-compute}
Let $f\colon\R^3\to\R^2$ be defined by $f(x,y,z) = (x+y-z,\; 2x+y+z)$.
\begin{enumerate}[label=(\alph*)]
  \item Find $\Ker f$ and give a basis.
  \item Find $\im f$ and determine $\rank(f)$.
  \item Verify the rank--nullity theorem for this map.
\end{enumerate}
\end{exercise}

\begin{exercise}{Matrix of a linear map \basic}{exo:matrix-linear-map}
Let $f\colon\R^3\to\R^2$ be defined by
$f(\vece{1}) = (1,2)^{\mathsf{T}}$,
$f(\vece{2}) = (0,-1)^{\mathsf{T}}$,
$f(\vece{3}) = (3,0)^{\mathsf{T}}$
(with respect to the canonical bases). Write down $\Mat(f)$ and
compute $f(1,1,1)^{\mathsf{T}}$.
\end{exercise}

\begin{exercise}{Composition and matrix multiplication \intermediate}%
{exo:composition-matrices}
Let $f\colon\R^2\to\R^3$ and $g\colon\R^3\to\R^2$ have matrices
\[
  A = \begin{pmatrix}1&0\\2&1\\0&3\end{pmatrix}, \qquad
  B = \begin{pmatrix}1&-1&2\\0&1&1\end{pmatrix}.
\]
\begin{enumerate}[label=(\alph*)]
  \item Compute $BA$ and $AB$.
  \item Which product represents $g\circ f$? Which represents $f\circ g$?
  \item Determine $\Ker(g\circ f)$ and $\rank(g\circ f)$.
\end{enumerate}
\end{exercise}

\begin{exercise}{Change of basis \intermediate}{exo:change-basis}
Let $f\colon\R^2\to\R^2$ be defined by $f(x,y) = (3x+y,\; x+3y)$.
\begin{enumerate}[label=(\alph*)]
  \item Write $\Mat_\cE(f)$ in the canonical basis $\cE$.
  \item Let $\cB' = \bigl((1,1)^{\mathsf{T}},(1,-1)^{\mathsf{T}}\bigr)$.
    Find the transition matrix $P$ from $\cE$ to~$\cB'$.
  \item Compute $\Mat_{\cB'}(f) = \inv{P}\Mat_\cE(f)\,P$.
  \item Interpret the result geometrically.
\end{enumerate}
\end{exercise}

\begin{exercise}{Rank--nullity applications \intermediate}%
{exo:rank-nullity-app}
\begin{enumerate}[label=(\alph*)]
  \item Let $f\colon\R^5\to\R^3$ be linear with $\dim(\Ker f) = 2$.
    What is $\rank(f)$?
  \item Let $g\colon\R^4\to\R^4$ be linear with $\rank(g) = 3$.
    What is $\dim(\Ker g)$? Is $g$ surjective?
  \item Prove that there is no surjective linear map from $\R^3$ to $\R^4$.
  \item Prove that there is no injective linear map from $\R^4$ to $\R^3$.
\end{enumerate}
\end{exercise}

\begin{exercise}{Projections \intermediate}{exo:projections}
\index{projection!idempotent}
A linear map $p\colon E\to E$ is called a \emph{projection} if
$p\circ p = p$ (idempotent).
\begin{enumerate}[label=(\alph*)]
  \item Show that $\im p = \Ker(\Id_E - p)$.
  \item Show that $\Ker p = \im(\Id_E - p)$.
  \item Deduce that $E = \Ker p \oplus \im p$.
  \item If $\dim E = n$ and $\rank(p) = r$, show there exists a basis
    in which $\Mat(p) = \diag(1,\ldots,1,0,\ldots,0)$ ($r$ ones).
\end{enumerate}
\end{exercise}

\begin{exercise}{Dual basis computation \intermediate}{exo:dual-basis}
Let $\cB = (\vece{1},\vece{2})$ be the basis of $\R^2$ with
$\vece{1} = (1,1)^{\mathsf{T}}$, $\vece{2} = (1,-1)^{\mathsf{T}}$.
Find explicit formulas for $\vece{1}^*$ and $\vece{2}^*$ (as functions
of the canonical coordinates $x_1,x_2$).
\end{exercise}

\begin{exercise}{Injectivity and surjectivity \challenging}%
{exo:inject-surject}
Let $f\colon E\to F$ and $g\colon F\to G$ be linear maps between
finite-dimensional spaces.
\begin{enumerate}[label=(\alph*)]
  \item Prove: $\rank(g\circ f) \leq \min(\rank f,\, \rank g)$.
  \item Prove the \emph{Sylvester rank inequality}\index{Sylvester rank inequality}:
    $\rank f + \rank g - \dim F \leq \rank(g\circ f)$.
    \textit{Hint:} Apply the rank--nullity theorem to the restriction
    $g\restrict{\im f}$.
  \item Deduce: if $g\circ f = 0$, then
    $\rank f + \rank g \leq \dim F$.
\end{enumerate}
\end{exercise}

\begin{exercise}{Nilpotent endomorphisms \challenging}{exo:nilpotent}
\index{nilpotent}
An endomorphism $f\in\End(E)$ is \emph{nilpotent} of index~$k$ if
$f^k = 0$ but $f^{k-1}\neq 0$.
\begin{enumerate}[label=(\alph*)]
  \item Show that if $f$ is nilpotent, then $f$ is not invertible.
  \item Show that the chain $\set{\vzero}\subseteq\Ker f\subseteq
    \Ker f^2\subseteq\cdots$ is strictly increasing until it
    stabilizes at~$E$.
  \item Prove that the nilpotency index satisfies $k\leq\dim E$.
  \item Prove that $\Id + f$ is invertible, and find a formula for
    $(\Id + f)^{-1}$ in terms of powers of $f$.
    \textit{Hint:} geometric series.
\end{enumerate}
\end{exercise}

\begin{exercise}{Linear maps on polynomial spaces \challenging}%
{exo:polynomial-maps}
Let $E = \cP_3(\R)$ (polynomials of degree $\leq 3$) with basis
$\cB = (1,t,t^2,t^3)$. Define $T\colon E\to E$ by $T(p)(t) = p(t+1)$.
\begin{enumerate}[label=(\alph*)]
  \item Verify that $T$ is linear.
  \item Compute the matrix $\Mat_\cB(T)$.
  \item Determine $\Ker T$ and $\im T$. Is $T$ an isomorphism?
  \item Compute $T^{-1}$ and interpret it.
\end{enumerate}
\end{exercise}

\begin{exercise}{The trace as a linear form \basic}{exo:trace-linear}
\index{trace}
Show that the trace map $\tr\colon\Mat_n(\K)\to\K$, defined by
$\tr(A) = \sum_{i=1}^n a_{ii}$, is a linear form.
Determine $\Ker(\tr)$ and its dimension.
\end{exercise}


% ============================================================
\section*{Chapter summary}\label{sec:summary-ch2}
\addcontentsline{toc}{section}{Chapter summary}
% ============================================================

\begin{itemize}
  \item A \textbf{linear map}\index{linear map} $f\colon E\to F$
    preserves addition and scalar multiplication. It is completely
    determined by its values on a basis
    (
ef{prop:prop:linear-determined-by-basis}).

  \item The \textbf{kernel} $\Ker f$ and \textbf{image} $\im f$ are
    subspaces (
ef{thm:thm:ker-im-subspaces}). The kernel measures
    ``information lost'' by $f$; the image measures ``information
    reached.''

  \item The \textbf{rank--nullity theorem}
    (
ef{thm:thm:rank-nullity}):
    $\dim E = \dim(\Ker f) + \rank(f)$.

  \item A linear map between spaces of equal finite dimension is
    injective if and only if it is surjective if and only if it is
    bijective (
ef{prop:prop:iso-finite-dim}).

  \item The set $\cL(E,F)$ of all linear maps is itself a vector space
    of dimension $\dim E\cdot\dim F$
    (
ef{thm:thm:L-E-F,prop:prop:dim-L-E-F}).

  \item Choosing bases allows us to represent every linear map as a
    \textbf{matrix} (
ef{def:def:matrix-representation}).
    \textbf{Composition} of linear maps corresponds to
    \textbf{matrix multiplication}
    (
ef{thm:thm:matrix-iso}).

  \item Under a \textbf{change of basis} with transition matrix~$P$,
    the matrix of an endomorphism transforms as
    $A' = \inv{P}AP$ (
ef{thm:thm:change-of-basis}).

  \item A square matrix is \textbf{invertible} if and only if its
    kernel is trivial, equivalently if and only if its rank equals
    its size (
ef{thm:thm:invertibility}).

  \item The \textbf{dual space} $E^* = \cL(E,\K)$ has the same
    dimension as~$E$ in finite dimension, and admits a canonical
    \textbf{dual basis} (
ef{thm:thm:dual-basis}).
\end{itemize}

% ch3-linear-systems.tex — Chapter 3: Systems of Linear Equations and Gaussian Elimination
% Systems of linear equations, row reduction, echelon forms, rank,
% Rouché–Capelli theorem, structure of solutions, applications.

\chapter{Systems of Linear Equations and Gaussian Elimination}%
\label{ch:linear-systems}

% ============================================================
% Motivating introduction
% ============================================================

Systems of linear equations are among the oldest and most ubiquitous
objects in mathematics. They appear whenever several quantities are
linked by proportional relationships: currents in an electrical
network must satisfy Kirchhoff's laws, prices in a multi-market
economic equilibrium are determined by supply-and-demand balances,
forces on a bridge truss must cancel for the structure to remain
in static equilibrium. Even the modern theory of machine learning
ultimately reduces the training of linear models to the solution of
(often very large) linear systems.

The algorithmic heart of the subject is \emph{Gaussian elimination}%
\index{Gaussian elimination}, a systematic procedure that transforms
any system into an equivalent one whose solutions can be read off by
back-substitution. Despite its simplicity, the algorithm is
remarkably powerful: it simultaneously tells us whether a system has
solutions, how many free parameters the solution set carries, and
provides an explicit parametric description.

In this chapter we formalise systems of linear equations, introduce
the matrix language~$Ax=b$, develop row reduction to (reduced) echelon
form, define the rank of a matrix, and prove the
Rouch\'{e}--Capelli theorem that characterises consistency. We then
establish the fundamental structure theorem: the general solution of a
consistent system is a translate of the solution space of the
associated homogeneous system. Throughout, we provide geometric
interpretations and worked examples.

\medskip
Throughout this chapter, $\K$ denotes a field (typically $\R$ or $\C$),
and we freely use results from 
ef{ch:vector-spaces} and

ef{ch:linear-maps}.


% ============================================================
\section{Systems of linear equations}\label{sec:linear-systems}
% ============================================================

\begin{definition}{System of linear equations}{def:linear-system}
\index{system of linear equations}\index{linear system}
A \emph{system of $m$ linear equations in $n$ unknowns} over a
field~$\K$ is a collection of equations
\[
  \left\{
  \begin{aligned}
    a_{11}\,x_1 + a_{12}\,x_2 + \cdots + a_{1n}\,x_n &= b_1,\\
    a_{21}\,x_1 + a_{22}\,x_2 + \cdots + a_{2n}\,x_n &= b_2,\\
    &\;\;\vdots \\
    a_{m1}\,x_1 + a_{m2}\,x_2 + \cdots + a_{mn}\,x_n &= b_m,
  \end{aligned}
  \right.
\]
where $a_{ij}\in\K$ are the \emph{coefficients}\index{coefficient}
and $b_i\in\K$ are the \emph{right-hand sides} (or \emph{constant
terms}). A vector $(s_1,\dots,s_n)\in\K^n$ is a \emph{solution} if
substituting $x_j = s_j$ satisfies every equation simultaneously. The
\emph{solution set}\index{solution set} is
\[
  \cS \deff \setcond{x\in\K^n}{Ax = b}.
\]
\end{definition}

\subsection{Matrix form}\label{subsec:matrix-form}
\index{matrix form of a linear system}

The system above can be written compactly as a single matrix equation
\begin{equation}\label{eq:Axb}
  Ax = b,
\end{equation}
where
\[
  A = \begin{pmatrix}
    a_{11} & a_{12} & \cdots & a_{1n}\\
    a_{21} & a_{22} & \cdots & a_{2n}\\
    \vdots & \vdots & \ddots & \vdots\\
    a_{m1} & a_{m2} & \cdots & a_{mn}
  \end{pmatrix}
  \in \Mat_{m,n}(\K),
  \qquad
  x = \begin{pmatrix} x_1\\ x_2\\ \vdots\\ x_n \end{pmatrix}
  \in \K^n,
  \qquad
  b = \begin{pmatrix} b_1\\ b_2\\ \vdots\\ b_m \end{pmatrix}
  \in \K^m.
\]
The matrix $A$ is called the \emph{coefficient matrix}%
\index{coefficient matrix}. The \emph{augmented matrix}%
\index{augmented matrix} is the $m\times(n+1)$ matrix
\begin{equation}\label{eq:augmented}
  (A \mid b) \deff
  \left(\begin{array}{cccc|c}
    a_{11} & a_{12} & \cdots & a_{1n} & b_1\\
    a_{21} & a_{22} & \cdots & a_{2n} & b_2\\
    \vdots & \vdots & \ddots & \vdots & \vdots\\
    a_{m1} & a_{m2} & \cdots & a_{mn} & b_m
  \end{array}\right).
\end{equation}

\begin{remark}{Column picture}{rem:column-picture}
\index{column picture}
The equation $Ax = b$ can also be read as asking whether $b$ lies in
the column space of~$A$: writing $A = (\vec{a}_1 \mid \vec{a}_2 \mid
\cdots \mid \vec{a}_n)$, the system is consistent if and only if
\[
  b = x_1\,\vec{a}_1 + x_2\,\vec{a}_2 + \cdots + x_n\,\vec{a}_n
\]
for some scalars $x_j\in\K$; that is, $b\in\im(A)$ where we view $A$
as a linear map $\K^n\to\K^m$.
\end{remark}


% ============================================================
\section{Homogeneous and non-homogeneous systems}%
\label{sec:homogeneous}
% ============================================================

\begin{definition}{Homogeneous system}{def:homogeneous}
\index{homogeneous system}\index{non-homogeneous system}
The system $Ax = b$ is called \emph{homogeneous} if $b = \vzero$, and
\emph{non-homogeneous} (or \emph{inhomogeneous}) otherwise.
Given a (possibly non-homogeneous) system $Ax = b$, the system
$Ax = \vzero$ is called its \emph{associated homogeneous system}.
\end{definition}

\begin{proposition}{The solution set of a homogeneous system is a subspace}%
{prop:homog-subspace}
\index{kernel!of a matrix}
Let $A\in\Mat_{m,n}(\K)$. The solution set of $Ax=\vzero$ is a
vector subspace of~$\K^n$. It coincides with $\Ker(A)$, the kernel
of the linear map $x\mapsto Ax$.
\end{proposition}

\begin{proof}
The zero vector satisfies $A\vzero = \vzero$, so $\cS_0\neq\emptyset$.
If $u,v\in\cS_0$ and $\lambda\in\K$, then
$A(\lambda u + v) = \lambda Au + Av = \lambda\vzero + \vzero = \vzero$,
so $\lambda u + v\in\cS_0$.
By the subspace criterion (
ef{ch:vector-spaces}), $\cS_0$ is a
subspace of~$\K^n$. By definition,
$\cS_0 = \setcond{x\in\K^n}{Ax = \vzero} = \Ker(A)$.
\end{proof}

\begin{proposition}{Affine structure of the solution set}%
{prop:affine-structure}
\index{affine subspace}
Let $Ax = b$ be a consistent system (i.e.\ $\cS\neq\emptyset$), and
let $x_0\in\cS$ be any particular solution. Then
\begin{equation}\label{eq:affine-sol}
  \cS = x_0 + \Ker(A) \deff \setcond{x_0 + h}{h \in \Ker(A)}.
\end{equation}
In particular, $\cS$ is an \emph{affine subspace}\index{affine subspace}
of~$\K^n$ of dimension $\dim\Ker(A)$.
\end{proposition}

\begin{proof}
($\supseteq$) If $h\in\Ker(A)$, then $A(x_0+h) = Ax_0 + Ah = b +
\vzero = b$, so $x_0+h\in\cS$.

($\subseteq$) If $x\in\cS$, set $h\deff x - x_0$. Then
$Ah = Ax - Ax_0 = b - b = \vzero$, so $h\in\Ker(A)$ and
$x = x_0 + h$.
\end{proof}


% ============================================================
\section{Elementary row operations}\label{sec:row-operations}
% ============================================================

\begin{definition}{Elementary row operations}{def:ero}
\index{elementary row operations}\index{row operations}
The following three operations on the rows of a matrix are called
\emph{elementary row operations} (EROs):
\begin{enumerate}[label=(E\arabic*)]
  \item\label{ero:swap}
    \textbf{Row interchange:} Swap rows $i$ and~$j$
    \quad ($R_i \leftrightarrow R_j$).
  \item\label{ero:scale}
    \textbf{Row scaling:} Multiply row $i$ by a non-zero scalar
    $\lambda\in\K^*$ \quad ($R_i \leftarrow \lambda\,R_i$).
  \item\label{ero:add}
    \textbf{Row replacement:} Add $\lambda$ times row $j$ to row $i$
    \quad ($R_i \leftarrow R_i + \lambda\,R_j$), where $i\neq j$.
\end{enumerate}
\end{definition}

\begin{proposition}{EROs preserve the solution set}{prop:ero-preserve}
If the augmented matrix $(A'\mid b')$ is obtained from $(A\mid b)$ by
an elementary row operation, then the systems $Ax=b$ and $A'x=b'$
have the same solution set.
\end{proposition}

\begin{proof}
Each ERO is reversible: \ref{ero:swap} is its own inverse;
\ref{ero:scale} is reversed by scaling by $\lambda^{-1}$;
\ref{ero:add} is reversed by subtracting $\lambda$ times row~$j$.
Hence any solution of one system is a solution of the other.
More precisely, each ERO amounts to left-multiplication by an
invertible \emph{elementary matrix}\index{elementary matrix}
$E\in\GL_m(\K)$. Since $Ax = b$ if and only if $E(Ax) = Eb$, i.e.\
$(EA)x = Eb$, the solution set is unchanged.
\end{proof}

\begin{definition}{Row equivalence}{def:row-equiv}
\index{row equivalence}
Two matrices are \emph{row equivalent} if one can be obtained from
the other by a finite sequence of elementary row operations. We write
$A \sim A'$.
\end{definition}


% ============================================================
\section{Echelon forms}\label{sec:echelon}
% ============================================================

\begin{definition}{Row echelon form}{def:ref}
\index{row echelon form}\index{echelon form}
A matrix is in \emph{row echelon form} (REF) if:
\begin{enumerate}[label=(\roman*)]
  \item All zero rows are at the bottom.
  \item The first non-zero entry in each non-zero row (called the
    \emph{pivot}\index{pivot}) is strictly to the right of the pivot
    in the row above.
\end{enumerate}
\end{definition}

\begin{definition}{Reduced row echelon form}{def:rref}
\index{reduced row echelon form}\index{RREF}
A matrix is in \emph{reduced row echelon form} (RREF) if it is in REF
and, additionally:
\begin{enumerate}[label=(\roman*),resume]
  \item Every pivot is equal to~$1$.
  \item Each pivot is the only non-zero entry in its column.
\end{enumerate}
\end{definition}

\begin{example}{Echelon forms}{ex:echelon-shapes}
The following matrix is in row echelon form (pivots are boxed):
\[
  \begin{pmatrix}
    \boxed{2} & 3 & -1 & 7\\
    0 & \boxed{5} & 4 & 2\\
    0 & 0 & 0 & \boxed{1}\\
    0 & 0 & 0 & 0
  \end{pmatrix}.
\]
Its reduced row echelon form is
\[
  \begin{pmatrix}
    \boxed{1} & 0 & * & 0\\
    0 & \boxed{1} & * & 0\\
    0 & 0 & 0 & \boxed{1}\\
    0 & 0 & 0 & 0
  \end{pmatrix},
\]
where the $*$ entries are determined by the original entries.
\end{example}

\begin{theorem}{Existence and uniqueness of RREF}{thm:rref-unique}
\index{RREF!uniqueness}
Every matrix $A\in\Mat_{m,n}(\K)$ is row equivalent to a unique
matrix in reduced row echelon form.
\end{theorem}

\begin{proof}
\emph{Existence} follows from the Gaussian elimination algorithm
described in 
ef{sec:gauss-elim}, which produces a REF, followed
by Gauss--Jordan elimination (
ef{sec:gauss-jordan}), which
produces the RREF.

\emph{Uniqueness.}
Suppose $R$ and $R'$ are two RREFs row equivalent to~$A$. Then
$R$ and $R'$ are row equivalent to each other, so they define the
same system of homogeneous equations $Rx=\vzero$ and $R'x=\vzero$
(by 
ef{prop:prop:ero-preserve}). In particular, their solution sets
coincide: $\Ker(R) = \Ker(R')$.

We show $R = R'$ column by column. Let $j$ be any column index.
If $j$ is a pivot column of~$R$, say with pivot in row~$i$, then
the $i$-th equation of $Rx = \vzero$ reads $x_j = -\sum_{k}
r_{ik}\,x_k$ where the sum is over the non-pivot (free) variables.
Since $\Ker(R) = \Ker(R')$, column~$j$ must also be a pivot column
of~$R'$ in the same row, and the coefficients must agree. Similarly
for non-pivot columns. Hence $R = R'$.
\end{proof}


% ============================================================
\section{Gaussian elimination}\label{sec:gauss-elim}
% ============================================================

\index{Gaussian elimination!algorithm}
Gaussian elimination transforms a matrix into row echelon form using
elementary row operations. The algorithm proceeds as follows.

\medskip
\noindent\textbf{Algorithm} (Gaussian elimination).
\begin{enumerate}[label=\textbf{Step \arabic*.}]
  \item Start with column $j=1$ and row $i=1$.
  \item \emph{Pivot selection.} In column $j$, find a non-zero entry
    in rows $i, i+1, \dots, m$. If none exists, increment $j$ and
    repeat. If $j > n$, stop.
  \item \emph{Swap.} If the non-zero entry is in row $k\neq i$,
    swap $R_i \leftrightarrow R_k$.
  \item \emph{Elimination.} For each row $\ell = i+1, \dots, m$,
    perform $R_\ell \leftarrow R_\ell -
    \frac{a_{\ell j}}{a_{ij}}\,R_i$ to create zeros below the pivot.
  \item Increment $i$ and $j$; go to Step~2.
\end{enumerate}

\begin{example}{Gaussian elimination: a $3\times 4$ system}%
{ex:gauss-3x4}
Solve the system
\[
  \left\{
  \begin{aligned}
    x_1 + 2x_2 - x_3 + 3x_4 &= 5,\\
    2x_1 + 5x_2 + x_3 + 2x_4 &= 8,\\
    x_1 + 3x_2 + 2x_3 - x_4 &= 1.
  \end{aligned}
  \right.
\]
The augmented matrix and its row reduction:
\begin{align*}
  &\left(\begin{array}{cccc|c}
    1 & 2 & -1 & 3 & 5\\
    2 & 5 & 1 & 2 & 8\\
    1 & 3 & 2 & -1 & 1
  \end{array}\right)
  \xrightarrow[\substack{R_3 \leftarrow R_3 - R_1}]{R_2 \leftarrow R_2 - 2R_1}
  \left(\begin{array}{cccc|c}
    1 & 2 & -1 & 3 & 5\\
    0 & 1 & 3 & -4 & -2\\
    0 & 1 & 3 & -4 & -4
  \end{array}\right)\\[6pt]
  &\xrightarrow{R_3 \leftarrow R_3 - R_2}
  \left(\begin{array}{cccc|c}
    1 & 2 & -1 & 3 & 5\\
    0 & 1 & 3 & -4 & -2\\
    0 & 0 & 0 & 0 & -2
  \end{array}\right).
\end{align*}
The last row reads $0 = -2$, a contradiction. Hence the system is
\textbf{inconsistent}\index{inconsistent system}: $\cS = \emptyset$.
\end{example}

\begin{example}{A consistent system with free variables}%
{ex:gauss-free}
Solve the system
\[
  \left\{
  \begin{aligned}
    x_1 + 2x_2 + x_3 &= 4,\\
    2x_1 + 5x_2 + 4x_3 &= 11,\\
    x_1 + 3x_2 + 3x_3 &= 7.
  \end{aligned}
  \right.
\]
Row reduction of the augmented matrix:
\begin{align*}
  &\left(\begin{array}{ccc|c}
    1 & 2 & 1 & 4\\
    2 & 5 & 4 & 11\\
    1 & 3 & 3 & 7
  \end{array}\right)
  \xrightarrow[\substack{R_3 \leftarrow R_3 - R_1}]{R_2 \leftarrow R_2 - 2R_1}
  \left(\begin{array}{ccc|c}
    1 & 2 & 1 & 4\\
    0 & 1 & 2 & 3\\
    0 & 1 & 2 & 3
  \end{array}\right)
  \xrightarrow{R_3 \leftarrow R_3 - R_2}
  \left(\begin{array}{ccc|c}
    1 & 2 & 1 & 4\\
    0 & 1 & 2 & 3\\
    0 & 0 & 0 & 0
  \end{array}\right).
\end{align*}
The pivot columns are $1$ and $2$; the variable $x_3$ is free
\index{free variable}. Back-substitution gives:
\[
  x_2 = 3 - 2x_3, \qquad x_1 = 4 - 2x_2 - x_3 = -2 + 3x_3.
\]
Setting $x_3 = t \in \K$, the solution set is
\[
  \cS = \set{\begin{pmatrix} -2+3t\\ 3-2t\\ t \end{pmatrix}
  : t\in\K}
  = \begin{pmatrix} -2\\3\\0 \end{pmatrix}
  + t\begin{pmatrix} 3\\-2\\1 \end{pmatrix},
  \quad t\in\K.
\]
This is a line in $\K^3$: a particular solution plus the kernel
$\Ker(A) = \Span\!\left(\begin{pmatrix}3\\-2\\1\end{pmatrix}\right)$.
\end{example}


% ============================================================
\section{Gauss--Jordan elimination}\label{sec:gauss-jordan}
% ============================================================

\index{Gauss--Jordan elimination}
After reaching row echelon form, \emph{Gauss--Jordan elimination}
continues with back-elimination to reach the reduced row echelon form.

\medskip
\noindent\textbf{Algorithm} (Gauss--Jordan: from REF to RREF).
\begin{enumerate}[label=\textbf{Step \arabic*.}]
  \item Starting from the \emph{bottom-most} pivot, scale its row so
    the pivot becomes~$1$.
  \item Use this pivot to eliminate all entries \emph{above} it in the
    same column.
  \item Move to the next pivot above and repeat.
\end{enumerate}

\begin{example}{Gauss--Jordan continuation}{ex:gauss-jordan}
Continuing from 
ef{ex:ex:gauss-free}, we had the REF
\[
  \left(\begin{array}{ccc|c}
    1 & 2 & 1 & 4\\
    0 & 1 & 2 & 3\\
    0 & 0 & 0 & 0
  \end{array}\right).
\]
The pivots are already~$1$. Eliminate above the second pivot:
\[
  \xrightarrow{R_1 \leftarrow R_1 - 2R_2}
  \left(\begin{array}{ccc|c}
    1 & 0 & -3 & -2\\
    0 & 1 & 2 & 3\\
    0 & 0 & 0 & 0
  \end{array}\right).
\]
This is the RREF. The solution can be read off directly:
$x_1 = -2 + 3x_3$, $x_2 = 3 - 2x_3$, confirming the result of

ef{ex:ex:gauss-free}.
\end{example}


% ============================================================
\section{Rank of a matrix}\label{sec:matrix-rank}
% ============================================================

\begin{definition}{Rank via echelon form}{def:rank-echelon}
\index{rank!of a matrix}
The \emph{rank} of a matrix $A\in\Mat_{m,n}(\K)$, denoted
$\rank(A)$, is the number of pivots in any row echelon form of~$A$.
\end{definition}

\begin{proposition}{Rank is well-defined}{prop:rank-well-defined}
The rank does not depend on which row echelon form is chosen: it
equals the number of pivots in the unique RREF of~$A$.
\end{proposition}

\begin{proof}
By 
ef{thm:thm:rref-unique}, the RREF is unique. Since the number of
pivots in any REF equals the number of pivots in the RREF (the
Gauss--Jordan steps do not create or destroy pivots), the rank is
well-defined.
\end{proof}

\begin{proposition}{Rank equals column rank equals row rank}%
{prop:rank-equiv}
\index{rank!column rank}\index{rank!row rank}
For any matrix $A\in\Mat_{m,n}(\K)$, the following are equal:
\begin{enumerate}[label=(\roman*)]
  \item The number of pivots in an echelon form of~$A$.
  \item $\dim(\im A)$, i.e.\ the dimension of the column space
    of~$A$.
  \item The dimension of the row space of~$A$.
\end{enumerate}
In particular, $\rank(A) = \rank(\transpose{A})$.
\end{proposition}

\begin{proof}
(i)$=$\,(ii): The pivot columns of~$A$ form a basis for the column
space $\im(A)$ (the corresponding columns of the
\emph{original}~$A$, selected via the pivot positions of the RREF,
are linearly independent and Span~$\im(A)$). Thus
$\dim(\im A)$ equals the number of pivots.

(i)$=$\,(iii): Row operations do not change the row space (each new
row is a linear combination of the old rows, and the operation is
reversible). The non-zero rows of a REF are linearly independent
(by the staircase structure), so they form a basis for the row space.
Their count equals the number of pivots.
\end{proof}

\begin{remark}{Rank and the rank--nullity theorem}{rem:rank-nullity-systems}
\index{rank--nullity theorem}
The rank--nullity theorem (
ef{ch:linear-maps}) applied to the
linear map $x\mapsto Ax$ gives
\begin{equation}\label{eq:rank-nullity}
  \rank(A) + \dim\Ker(A) = n,
\end{equation}
where $n$ is the number of columns (i.e.\ unknowns). Thus the number
of free variables in the general solution equals $n - \rank(A)$.
\end{remark}


% ============================================================
\section{The Rouch\'{e}--Capelli theorem}\label{sec:rouche-capelli}
% ============================================================

The central consistency criterion is the following classical result.

\begin{theorem}{Rouch\'{e}--Capelli theorem}%
{thm:rouche-capelli}
\index{Rouche-Capelli theorem@Rouch\'{e}--Capelli theorem}
\index{consistency criterion}
The system $Ax = b$ is consistent (i.e.\ has at least one solution)
if and only if
\begin{equation}\label{eq:rouche-capelli}
  \rank(A) = \rank(A\mid b).
\end{equation}
When this condition holds, the dimension of the solution set is
$n - \rank(A)$, where $n$ is the number of unknowns. In particular:
\begin{itemize}
  \item If $\rank(A) = n$, the solution is unique.
  \item If $\rank(A) < n$, there are infinitely many solutions
    (assuming $\K$ is infinite), parametrised by $n - \rank(A)$ free
    variables.
\end{itemize}
\end{theorem}

\begin{proof}
Write $r = \rank(A)$ and $r' = \rank(A\mid b)$. Since every column
of~$A$ is also a column of $(A\mid b)$, we always have $r \leq r'
\leq r+1$.

\emph{($\Rightarrow$)}
Suppose $\cS\neq\emptyset$, say $x_0\in\cS$. Then $b = Ax_0$ is a
linear combination of the columns of~$A$, so $b$ lies in the column
space of~$A$. Adding the column~$b$ to~$A$ does not increase the
dimension of the column space:
$\rank(A\mid b) = \dim\Span(\text{columns of } A, b) =
\dim\Span(\text{columns of } A) = \rank(A)$.
Hence $r' = r$.

\emph{($\Leftarrow$)}
Suppose $\rank(A) = \rank(A\mid b)$. Then the column spaces satisfy
$\im(A\mid b) = \im(A)$, which means $b$ is a linear combination of
the columns of~$A$. That is, there exists $x_0\in\K^n$ with
$Ax_0 = b$, so the system is consistent.

The dimension statement follows from 
ef{prop:prop:affine-structure} and

ef{eq:rank-nullity}: $\dim\cS = \dim\Ker(A) = n - \rank(A)$.
\end{proof}

\begin{corollary}{Homogeneous systems always have non-trivial solutions
when $n > m$}{cor:homog-nontrivial}
\index{homogeneous system!non-trivial solution}
If $A\in\Mat_{m,n}(\K)$ with $n > m$ (more unknowns than equations),
then $Ax = \vzero$ has a non-trivial solution.
\end{corollary}

\begin{proof}
We have $\rank(A) \leq m < n$, so $\dim\Ker(A) = n - \rank(A) \geq 1$.
\end{proof}


% ============================================================
\section{Structure of the solution set}\label{sec:solution-structure}
% ============================================================

We now state and prove the complete structure theorem, synthesising
the results of the preceding sections.

\begin{theorem}{Structure of the general solution}{thm:solution-structure}
\index{general solution}\index{solution set!structure}
Let $A\in\Mat_{m,n}(\K)$ and $b\in\K^m$, and suppose $Ax=b$ is
consistent. Let $x_p\in\K^n$ be any particular solution. Then:
\begin{enumerate}[label=(\roman*)]
  \item $\Ker(A) = \setcond{x\in\K^n}{Ax = \vzero}$ is a vector
    subspace of $\K^n$ of dimension $n - \rank(A)$.
  \item The general solution of $Ax = b$ is
    \begin{equation}\label{eq:gen-sol}
      \cS = \setcond{x_p + h}{h\in\Ker(A)} = x_p + \Ker(A).
    \end{equation}
  \item If $\set{h_1, \dots, h_k}$ is a basis of $\Ker(A)$ (where
    $k = n - \rank(A)$), then every solution can be written uniquely as
    \begin{equation}\label{eq:parametric-sol}
      x = x_p + t_1 h_1 + t_2 h_2 + \cdots + t_k h_k,
      \qquad t_1, \dots, t_k \in \K.
    \end{equation}
\end{enumerate}
\end{theorem}

\begin{proof}
(i) follows from 
ef{prop:prop:homog-subspace} and

ef{eq:rank-nullity}.

(ii) is exactly 
ef{prop:prop:affine-structure}.

(iii) Since $\set{h_1,\dots,h_k}$ is a basis of $\Ker(A)$, every
$h\in\Ker(A)$ can be written uniquely as $h = t_1 h_1 + \cdots +
t_k h_k$. Combined with~(ii), every $x\in\cS$ has the unique
representation $x = x_p + t_1 h_1 + \cdots + t_k h_k$.
The uniqueness of the $t_i$ follows from the linear independence of
the $h_i$: if $x_p + \sum t_i h_i = x_p + \sum t_i' h_i$, then
$\sum (t_i - t_i') h_i = \vzero$, forcing $t_i = t_i'$ for all~$i$.
\end{proof}


% ============================================================
\section{Parametric representation of solutions}%
\label{sec:parametric}
% ============================================================

\index{parametric form}
When solving a system $Ax = b$ by Gaussian elimination, the RREF
identifies which variables are \emph{pivot variables}\index{pivot variable}
(corresponding to pivot columns) and which are \emph{free variables}%
\index{free variable} (corresponding to non-pivot columns). A
parametric representation is obtained by:
\begin{enumerate}
  \item Expressing each pivot variable in terms of the free variables.
  \item Assigning arbitrary parameters $t_1, t_2, \dots$ to the free
    variables.
  \item Writing the general solution as a vector.
\end{enumerate}

\begin{example}{Parametric form with two free variables}%
{ex:parametric-two-free}
Consider the system with RREF
\[
  \left(\begin{array}{ccccc|c}
    1 & 3 & 0 & -2 & 0 & 4\\
    0 & 0 & 1 & 5 & 0 & -1\\
    0 & 0 & 0 & 0 & 1 & 2\\
    0 & 0 & 0 & 0 & 0 & 0
  \end{array}\right).
\]
Pivot variables: $x_1, x_3, x_5$.\quad Free variables: $x_2 = s$,
$x_4 = t$. Reading off:
\[
  x_1 = 4 - 3s + 2t, \quad x_3 = -1 - 5t, \quad x_5 = 2.
\]
The parametric form is
\[
  \begin{pmatrix} x_1\\ x_2\\ x_3\\ x_4\\ x_5 \end{pmatrix}
  =
  \underbrace{\begin{pmatrix} 4\\ 0\\ -1\\ 0\\ 2 \end{pmatrix}}_{x_p}
  + s\underbrace{\begin{pmatrix} -3\\ 1\\ 0\\ 0\\ 0 \end{pmatrix}}_{h_1}
  + t\underbrace{\begin{pmatrix} 2\\ 0\\ -5\\ 1\\ 0 \end{pmatrix}}_{h_2},
  \qquad s,t\in\K.
\]
Here $\rank(A) = 3$, $n = 5$, and $\dim\Ker(A) = 5 - 3 = 2$.
\end{example}


% ============================================================
\section{Geometric interpretation}\label{sec:geometric-interp}
% ============================================================

\index{geometric interpretation!of linear systems}

\subsection{Two equations in two unknowns}

A single linear equation $a_1 x_1 + a_2 x_2 = b$ in two unknowns
describes a \emph{line} in~$\R^2$. A system of two such equations
corresponds to two lines, and the solution set is their intersection.

\begin{center}
\begin{tikzpicture}[>=Stealth, scale=1.0]
  % Unique solution
  \begin{scope}[shift={(0,0)}]
    \draw[->] (-0.5,0) -- (3.5,0) node[right] {$x_1$};
    \draw[->] (0,-0.5) -- (0,3) node[above] {$x_2$};
    \draw[thick, blue!70!black, domain=-0.3:3.2] plot (\x, {2.5 - 0.7*\x});
    \draw[thick, red!70!black, domain=-0.3:3.2] plot (\x, {0.3 + 0.6*\x});
    \fill[black] (1.692, 1.315) circle (2pt);
    \node[below] at (1.75, -0.8) {\small Unique solution};
    \node[below] at (1.75, -1.2) {\small $\rank(A) = 2$};
  \end{scope}
  % No solution
  \begin{scope}[shift={(5.5,0)}]
    \draw[->] (-0.5,0) -- (3.5,0) node[right] {$x_1$};
    \draw[->] (0,-0.5) -- (0,3) node[above] {$x_2$};
    \draw[thick, blue!70!black, domain=-0.3:3.2] plot (\x, {2.5 - 0.7*\x});
    \draw[thick, red!70!black, domain=-0.3:3.2] plot (\x, {1.2 - 0.7*\x});
    \node[below] at (1.75, -0.8) {\small No solution};
    \node[below] at (1.75, -1.2) {\small (parallel lines)};
  \end{scope}
  % Infinitely many
  \begin{scope}[shift={(11,0)}]
    \draw[->] (-0.5,0) -- (3.5,0) node[right] {$x_1$};
    \draw[->] (0,-0.5) -- (0,3) node[above] {$x_2$};
    \draw[thick, blue!70!black, domain=-0.3:3.2] plot (\x, {2.5 - 0.7*\x});
    \draw[thick, red!70!black, dashed, domain=-0.3:3.2]
      plot (\x, {2.5 - 0.7*\x});
    \node[below] at (1.75, -0.8) {\small Infinitely many};
    \node[below] at (1.75, -1.2) {\small (coincident lines)};
  \end{scope}
\end{tikzpicture}
\captionof{figure}{Three possibilities for a $2\times 2$ linear
  system over $\R$.}
\label{fig:2x2-systems}
\end{center}

\subsection{Three equations in three unknowns}

Each equation $a_1 x_1 + a_2 x_2 + a_3 x_3 = b$ describes a
\emph{plane}\index{plane} in~$\R^3$. A system of three such equations
asks for the common intersection of three planes. The possibilities
are richer than in $\R^2$.

\begin{center}
\begin{tikzpicture}[>=Stealth, scale=0.95]
  % --- Unique point ---
  \begin{scope}[shift={(0,0)}]
    % Three planes meeting at a point (schematic)
    \fill[blue!15, opacity=0.7] (-1.2,-0.8) -- (1.5,-0.5)
      -- (1.8,1.3) -- (-0.9,1.0) -- cycle;
    \fill[red!15, opacity=0.6] (-0.5,-1.0) -- (1.8,-0.3)
      -- (1.2,1.5) -- (-1.1,0.8) -- cycle;
    \fill[green!15, opacity=0.5] (-1.0,-0.4) -- (1.0,-1.0)
      -- (1.6,0.9) -- (-0.4,1.5) -- cycle;
    \fill[black] (0.5,0.3) circle (2.5pt);
    \node[below] at (0.4,-1.3) {\small Unique point};
  \end{scope}
  % --- Line ---
  \begin{scope}[shift={(4.5,0)}]
    \fill[blue!15, opacity=0.7] (-1.2,-1.0) -- (1.5,-0.8)
      -- (1.8,1.0) -- (-0.9,0.8) -- cycle;
    \fill[red!15, opacity=0.6] (-1.0,-0.8) -- (1.8,-0.3)
      -- (1.3,1.2) -- (-1.5,0.7) -- cycle;
    \fill[green!15, opacity=0.5] (-0.8,-1.0) -- (1.6,-0.5)
      -- (1.0,1.3) -- (-1.4,0.8) -- cycle;
    \draw[thick, black] (-0.6,-0.6) -- (1.2,0.9);
    \node[below] at (0.4,-1.3) {\small Line};
  \end{scope}
  % --- Plane ---
  \begin{scope}[shift={(9,0)}]
    \fill[blue!20, opacity=0.8] (-1.4,-1.0) -- (1.8,-0.8)
      -- (1.6,1.2) -- (-1.6,1.0) -- cycle;
    \node at (0.1,0.1) {\small (3 coincident)};
    \node[below] at (0.2,-1.3) {\small Plane};
  \end{scope}
  % --- Empty ---
  \begin{scope}[shift={(13.5,0)}]
    \fill[blue!15, opacity=0.6] (-1.2,-0.2) -- (1.5,0.0)
      -- (1.5,0.8) -- (-1.2,0.6) -- cycle;
    \fill[red!15, opacity=0.6] (-1.2,-0.7) -- (1.5,-0.5)
      -- (1.5,0.3) -- (-1.2,0.1) -- cycle;
    \fill[green!15, opacity=0.5] (-1.2,0.4) -- (1.5,0.6)
      -- (1.5,1.4) -- (-1.2,1.2) -- cycle;
    \node[below] at (0.2,-1.3) {\small No solution};
    \node[below] at (0.2,-1.7) {\small (parallel planes)};
  \end{scope}
\end{tikzpicture}
\captionof{figure}{Selected configurations for three planes
  in~$\R^3$: unique intersection point, intersection along a line,
  three coincident planes, and three parallel planes.}
\label{fig:3x3-systems}
\end{center}


% ============================================================
\section{Applications}\label{sec:applications-systems}
% ============================================================

\subsection{Network flows}\label{subsec:network-flows}
\index{network flow}

Consider a network of one-way streets or pipelines with known flow
rates entering and leaving the network. At each junction
(node)\index{node}, conservation of flow requires that the total
flow in equals the total flow out. This yields one linear equation per
node.

\begin{example}{Traffic flow}{ex:traffic-flow}
Consider the network with four intersections where external flows are
given and the internal flows $x_1, x_2, x_3, x_4$ are unknown:

\begin{center}
\begin{tikzpicture}[>=Stealth, node distance=2.5cm,
  every node/.style={circle, draw, minimum size=8mm, inner sep=1pt}]
  \node (A) at (0,2) {$A$};
  \node (B) at (3,2) {$B$};
  \node (C) at (3,0) {$C$};
  \node (D) at (0,0) {$D$};
  % Internal edges
  \draw[->, thick] (A) -- node[above] {$x_1$} (B);
  \draw[->, thick] (B) -- node[right] {$x_2$} (C);
  \draw[->, thick] (C) -- node[below] {$x_3$} (D);
  \draw[->, thick] (D) -- node[left]  {$x_4$} (A);
  % External flows
  \draw[->, thick, blue!70!black] (-1.5,2) -- (A)
    node[midway, above, draw=none, fill=none] {$100$};
  \draw[->, thick, blue!70!black] (B) -- (4.5,2)
    node[midway, above, draw=none, fill=none] {$80$};
  \draw[->, thick, blue!70!black] (4.5,0) -- (C)
    node[midway, above, draw=none, fill=none] {$50$};
  \draw[->, thick, blue!70!black] (D) -- (-1.5,0)
    node[midway, above, draw=none, fill=none] {$70$};
\end{tikzpicture}
\end{center}

Flow conservation at each node:
\[
  \begin{aligned}
    A&: & 100 + x_4 &= x_1, \\
    B&: & x_1 &= 80 + x_2, \\
    C&: & x_2 + 50 &= x_3, \\
    D&: & x_3 &= 70 + x_4.
  \end{aligned}
  \qquad\Longleftrightarrow\qquad
  \begin{aligned}
    x_1 - x_4 &= 100,\\
    x_1 - x_2 &= 80,\\
    x_2 - x_3 &= -50,\\
    x_3 - x_4 &= 70.
  \end{aligned}
\]
Row reduction shows $\rank(A) = 3$ and $n=4$, giving one free
variable. Setting $x_4 = t$:
\[
  x_1 = 100+t, \quad x_2 = 20+t, \quad x_3 = 70+t, \quad x_4 = t.
\]
Physical constraints ($x_i \geq 0$) require $t \geq 0$.
\end{example}


\subsection{Electrical circuits}\label{subsec:circuits}
\index{Kirchhoff's laws}\index{electrical circuit}

In an electrical circuit, Kirchhoff's current law (KCL) states that
the algebraic sum of currents at any node is zero, and Kirchhoff's
voltage law (KVL) states that the algebraic sum of voltage drops
around any closed loop is zero. Combined with Ohm's law
($V = IR$)\index{Ohm's law}, these yield a linear system for the
unknown currents.

\begin{example}{A simple resistive circuit}{ex:circuit}
Consider a circuit with two loops sharing a branch. Let $I_1$, $I_2$,
$I_3$ be the branch currents and $R_1 = 2\,\Omega$,
$R_2 = 4\,\Omega$, $R_3 = 6\,\Omega$ the resistances, with a
voltage source $V = 12$\,V. Applying KCL and KVL:
\[
  \left\{
  \begin{aligned}
    I_1 - I_2 - I_3 &= 0 && (\text{KCL at node}),\\
    2I_1 + 6I_3 &= 12 && (\text{KVL, loop 1}),\\
    4I_2 - 6I_3 &= 0 && (\text{KVL, loop 2}).
  \end{aligned}
  \right.
\]
Solving by Gaussian elimination gives
$I_1 = \frac{36}{11}$\,A,\;
$I_2 = \frac{18}{11}$\,A,\;
$I_3 = \frac{18}{11} \cdot \frac{2}{3} = \frac{12}{11}$\,A.
\end{example}


\subsection{Preview: least squares}\label{subsec:least-squares-preview}
\index{least squares!preview}

When a system $Ax = b$ is \emph{overdetermined} (more equations than
unknowns, typically inconsistent), one seeks the vector $\hat{x}$
that minimises $\norm{Ax - b}^2$. A fundamental result (to be proved
in the chapter on inner product spaces) states that $\hat{x}$
satisfies the \emph{normal equations}\index{normal equations}
\begin{equation}\label{eq:normal-equations}
  \transpose{A}A\,\hat{x} = \transpose{A}b.
\end{equation}
This is a (consistent) square linear system, solvable by Gaussian
elimination. Least squares is the mathematical foundation of linear
regression, signal processing, and data fitting.


% ============================================================
\section{Exercises}\label{sec:exercises-systems}
% ============================================================

% --- Exercise 1 ---
\begin{exercise}{Row reduction practice \basic}{exo:row-reduce-1}
Reduce the following matrix to RREF and find $\rank(A)$:
\[
  A = \begin{pmatrix}
    1 & 2 & -1 & 0\\
    2 & 4 & 1 & 3\\
    3 & 6 & 2 & 5
  \end{pmatrix}.
\]
\end{exercise}

% --- Exercise 2 ---
\begin{exercise}{Solving a $3\times 3$ system \basic}{exo:solve-3x3}
Solve the system
\[
  \left\{
  \begin{aligned}
    x + y + z &= 6,\\
    2x + 3y + z &= 14,\\
    x + 2y - z &= 2.
  \end{aligned}
  \right.
\]
\end{exercise}

% --- Exercise 3 ---
\begin{exercise}{Homogeneous system \basic}{exo:homog-system}
Find the general solution of $Ax = \vzero$ where
\[
  A = \begin{pmatrix}
    1 & 2 & 3\\
    2 & 4 & 6\\
    1 & 2 & 3
  \end{pmatrix}.
\]
What is $\dim\Ker(A)$?
\end{exercise}

% --- Exercise 4 ---
\begin{exercise}{Consistency conditions \basic}{exo:consistency}
For which values of $a\in\R$ is the system
\[
  \left\{
  \begin{aligned}
    x + y &= 1,\\
    2x + 2y &= a,
  \end{aligned}
  \right.
\]
consistent? When consistent, describe the solution set geometrically.
\end{exercise}

% --- Exercise 5 ---
\begin{exercise}{Parametric solutions \intermediate}{exo:parametric-sol}
Solve the system and express the solution in parametric vector form:
\[
  \left\{
  \begin{aligned}
    x_1 + 3x_2 - 2x_3 + x_4 &= 4,\\
    2x_1 + 6x_2 - 3x_3 + 5x_4 &= 11,\\
    x_1 + 3x_2 + x_4 &= 5.
  \end{aligned}
  \right.
\]
Identify the particular solution and a basis for the kernel.
\end{exercise}

% --- Exercise 6 ---
\begin{exercise}{Rank and the Rouch\'{e}--Capelli theorem
\intermediate}{exo:rouche-capelli-app}
Let
\[
  A = \begin{pmatrix}
    1 & 2 & 1\\
    2 & 4 & 3\\
    3 & 6 & 4
  \end{pmatrix},
  \qquad
  b = \begin{pmatrix} 1\\3\\4 \end{pmatrix}.
\]
Compute $\rank(A)$ and $\rank(A\mid b)$. Is $Ax = b$ consistent?
If so, find the solution set.
\end{exercise}

% --- Exercise 7 ---
\begin{exercise}{Family of systems \intermediate}{exo:family-systems}
For which values of $\lambda\in\R$ does the system
\[
  \left\{
  \begin{aligned}
    x + y + \lambda z &= 1,\\
    x + \lambda y + z &= 1,\\
    \lambda x + y + z &= 1,
  \end{aligned}
  \right.
\]
have (a) a unique solution, (b) infinitely many solutions,
(c) no solution?
\end{exercise}

% --- Exercise 8 ---
\begin{exercise}{Network flow \intermediate}{exo:network-flow}
In a road network with five nodes, the flow conservation equations
yield the system
\[
  \left\{
  \begin{aligned}
    x_1 + x_2 &= 200,\\
    x_2 + x_3 &= 150,\\
    x_3 + x_4 &= 180,\\
    x_4 + x_5 &= 120.
  \end{aligned}
  \right.
\]
Find the general solution. If all flows must be non-negative,
determine the feasible range for~$x_1$.
\end{exercise}

% --- Exercise 9 ---
\begin{exercise}{Computing a kernel basis \intermediate}{exo:kernel-basis}
Find a basis for $\Ker(A)$ where
\[
  A = \begin{pmatrix}
    1 & 0 & -2 & 1 & 3\\
    0 & 1 & 3 & -1 & 0\\
    2 & 1 & -1 & 1 & 6\\
    1 & 1 & 1 & 0 & 3
  \end{pmatrix}.
\]
Verify that $\rank(A) + \dim\Ker(A) = 5$.
\end{exercise}

% --- Exercise 10 ---
\begin{exercise}{Rank inequalities \challenging}{exo:rank-ineq}
Let $A\in\Mat_{m,n}(\K)$ and $B\in\Mat_{n,p}(\K)$. Prove:
\begin{enumerate}[label=(\alph*)]
  \item $\rank(AB) \leq \min(\rank(A), \rank(B))$.
  \item \emph{Sylvester's rank inequality:}\index{Sylvester's rank inequality}
    $\rank(AB) \geq \rank(A) + \rank(B) - n$.
\end{enumerate}
\textit{Hint for (b):} use the rank--nullity theorem and show
$\Ker(B) \subseteq \Ker(AB)$.
\end{exercise}

% --- Exercise 11 ---
\begin{exercise}{Normal equations \challenging}{exo:normal-eq}
Consider the inconsistent system (three equations, two unknowns):
\[
  \begin{pmatrix} 1 & 1\\ 1 & 2\\ 1 & 3 \end{pmatrix}
  \begin{pmatrix} x_1\\ x_2 \end{pmatrix}
  = \begin{pmatrix} 1\\ 2\\ 4 \end{pmatrix}.
\]
\begin{enumerate}[label=(\alph*)]
  \item Form the normal equations $\transpose{A}A\,\hat{x} =
    \transpose{A}b$.
  \item Solve for $\hat{x}$ by Gaussian elimination.
  \item Interpret $\hat{x}$ as the coefficients of the best-fit
    line $y = x_1 + x_2 t$ through the data points
    $(1,1)$, $(2,2)$, $(3,4)$.
\end{enumerate}
\end{exercise}

% --- Exercise 12 ---
\begin{exercise}{Intersection of subspaces via systems
\challenging}{exo:intersect-subspaces}
Let $U = \Span\!\left(\begin{pmatrix}1\\1\\0\\1\end{pmatrix},
\begin{pmatrix}0\\1\\1\\0\end{pmatrix}\right)$ and
$W = \Span\!\left(\begin{pmatrix}1\\0\\1\\1\end{pmatrix},
\begin{pmatrix}1\\1\\1\\0\end{pmatrix}\right)$ be subspaces
of~$\R^4$.
\begin{enumerate}[label=(\alph*)]
  \item Set up a linear system whose solutions describe $U\cap W$.
  \item Solve the system and find a basis for $U\cap W$.
  \item Verify the Grassmann formula:
    $\dim(U+W) = \dim U + \dim W - \dim(U\cap W)$.
\end{enumerate}
\end{exercise}

% --- Exercise 13 ---
\begin{exercise}{Proving row equivalence preserves solution sets
\intermediate}{exo:ero-proof}
Give a detailed proof that each elementary row operation on the
augmented matrix $(A\mid b)$ preserves the solution set of $Ax = b$.
Explicitly construct the inverse operation in each case.
\end{exercise}

% --- Exercise 14 ---
\begin{exercise}{RREF uniqueness \challenging}{exo:rref-unique-detail}
Provide a complete proof that two row-equivalent matrices in reduced
row echelon form must be identical.
\textit{Hint:} proceed by induction on the number of columns, using
the fact that the two matrices have identical null spaces.
\end{exercise}


% ============================================================
\section{Chapter summary}\label{sec:summary-ch3}
% ============================================================

\begin{itemize}
  \item A \textbf{system of linear equations}\index{system of linear equations}
    $Ax = b$ involves a coefficient matrix $A\in\Mat_{m,n}(\K)$, an
    unknown vector $x\in\K^n$, and a right-hand side $b\in\K^m$.

  \item \textbf{Elementary row operations}\index{elementary row operations}
    (row swap, row scaling, row replacement) transform a system into a
    row-equivalent one with the same solution set.

  \item \textbf{Gaussian elimination}\index{Gaussian elimination}
    reduces $A$ to \textbf{row echelon form}\index{row echelon form}
    (REF); \textbf{Gauss--Jordan elimination}%
    \index{Gauss--Jordan elimination} continues to the unique
    \textbf{reduced row echelon form}\index{reduced row echelon form}
    (RREF).

  \item The \textbf{rank}\index{rank} $\rank(A)$ is the number of
    pivots, equal to the dimension of both the column space and the
    row space.

  \item \textbf{Rouch\'{e}--Capelli theorem}%
    \index{Rouche-Capelli theorem@Rouch\'{e}--Capelli theorem}:
    $Ax = b$ is consistent $\iff$ $\rank(A) = \rank(A\mid b)$.

  \item The \textbf{general solution}\index{general solution} of a
    consistent system is
    $\cS = x_p + \Ker(A)$,
    an affine subspace of dimension $n - \rank(A)$.

  \item \textbf{Rank--nullity:}\index{rank--nullity theorem}
    $\rank(A) + \dim\Ker(A) = n$ (number of columns).

  \item Geometrically, a linear equation in $\R^2$ defines a line, in
    $\R^3$ a plane; the solution set of a system is the intersection
    of these geometric objects.

  \item Applications include network flows, electrical circuits
    (Kirchhoff's laws), and least squares (normal equations).
\end{itemize}

% ch4-determinants.tex — Chapter 4: Determinants
% Permutations, Leibniz formula, cofactor expansion, adjugate, Cramer's rule,
% geometric interpretation, applications.

\chapter{Determinants}\label{ch:determinants}

% ============================================================
% Motivating introduction
% ============================================================

How can one tell, by a single number, whether a square matrix is
invertible? How can one compute the area of a parallelogram Spanned by
two vectors, or the volume of a parallelepiped in~$\R^3$? These
seemingly different questions are all answered by the same object:
the \emph{determinant}\index{determinant}.

Consider two vectors $\vec{u}=(a,b)$ and $\vec{v}=(c,d)$ in~$\R^2$.
The signed area of the parallelogram they Span is
$ad - bc$---precisely the determinant of the matrix whose rows (or
columns) are $\vec{u}$ and~$\vec{v}$. This quantity is zero exactly
when the two vectors are collinear, i.e.\ when the matrix
$\begin{psmallmatrix} a & b \\ c & d \end{psmallmatrix}$ fails to be
invertible. The sign records the \emph{orientation}: positive if
$\vec{v}$ lies to the left of~$\vec{u}$, negative otherwise.

\begin{center}
\begin{tikzpicture}[>=Stealth, scale=1.1]
  % Parallelogram
  \coordinate (O) at (0,0);
  \coordinate (U) at (3,0.8);
  \coordinate (V) at (1,2.2);
  \coordinate (UV) at ($(U)+(V)$);

  \fill[blue!10] (O) -- (U) -- (UV) -- (V) -- cycle;
  \draw[thick, blue!70!black, ->] (O) -- (U) node[midway, below right] {$\vec{u}$};
  \draw[thick, red!70!black, ->] (O) -- (V) node[midway, above left] {$\vec{v}$};
  \draw[dashed, gray] (U) -- (UV) -- (V);

  \node at ($(O)!0.5!(UV)$) {$\abs{\det}$};

  % Orientation arc
  \draw[->, thick, green!50!black] (0.9,0.24) arc[start angle=15, end angle=55, radius=0.9]
    node[midway, right] {\small $+$};
\end{tikzpicture}
\captionof{figure}{The absolute value of $\det\begin{psmallmatrix}a&b\\c&d\end{psmallmatrix}$
  gives the area of the parallelogram; the sign encodes orientation.}
\label{fig:det-parallelogram-intro}
\end{center}

In this chapter we develop the theory of determinants rigorously. We
begin with the combinatorial machinery of permutations and their
signatures, define the determinant via the Leibniz formula, establish
its fundamental properties, prove cofactor expansion, derive Cramer's
rule, and explore geometric interpretations and applications.

Throughout this chapter, $\K$ denotes a field (typically $\R$ or~$\C$),
and all matrices are square unless stated otherwise. We write
$\Mat_n(\K)$ for the space of $n\times n$ matrices with entries
in~$\K$, and we freely use notation and results from

ef{ch:vector-spaces,ch:linear-maps}.


% ============================================================
\section{Permutations and the symmetric group}\label{sec:permutations}
% ============================================================

\begin{definition}{Permutation}{def:permutation}
\index{permutation}\index{symmetric group}
A \emph{permutation} of the set $\set{1,2,\ldots,n}$ is a bijection
$\sigma\colon\set{1,\ldots,n}\to\set{1,\ldots,n}$.
The set of all such permutations, equipped with composition, is the
\emph{symmetric group}\index{symmetric group} $\mathfrak{S}_n$.
Its cardinality is $\card{\mathfrak{S}_n} = n!$.
\end{definition}

We use two-line notation and cycle notation interchangeably. For
instance, the permutation $\sigma\in\mathfrak{S}_4$ with
$\sigma(1)=2$, $\sigma(2)=4$, $\sigma(3)=3$, $\sigma(4)=1$ is written
\[
  \sigma = \begin{pmatrix} 1 & 2 & 3 & 4 \\ 2 & 4 & 3 & 1 \end{pmatrix}
  = (1\;2\;4).
\]

\begin{definition}{Transposition}{def:transposition}
\index{transposition}
A \emph{transposition} is a permutation that exchanges exactly two
elements and fixes all others. We write $\tau_{ij}$ for the
transposition that swaps $i$ and~$j$ (with $i\neq j$).
\end{definition}

\begin{proposition}{Generation by transpositions}{prop:gen-transpositions}
Every permutation $\sigma\in\mathfrak{S}_n$ can be written as a
product (composition) of transpositions.
\end{proposition}

\begin{proof}
It suffices to show every cycle decomposes into transpositions, since
every permutation is a product of disjoint cycles. Indeed,
$(a_1\;a_2\;\cdots\;a_k) = \tau_{a_1 a_k}\circ\tau_{a_1 a_{k-1}}
\circ\cdots\circ\tau_{a_1 a_2}$.
\end{proof}

\begin{definition}{Inversions and signature}{def:signature}
\index{inversion}\index{signature!of a permutation}
An \emph{inversion} of $\sigma\in\mathfrak{S}_n$ is a pair $(i,j)$
with $1\leq i < j \leq n$ and $\sigma(i)>\sigma(j)$.
The \emph{signature} (or \emph{sign}) of~$\sigma$ is
\[
  \sgn(\sigma) \deff (-1)^{N(\sigma)},
\]
where $N(\sigma)$ denotes the number of inversions of~$\sigma$.
Equivalently,
\[
  \sgn(\sigma) = \prod_{1\leq i < j \leq n}
    \frac{\sigma(j)-\sigma(i)}{j-i}.
\]
A permutation is called \emph{even}\index{even permutation} if
$\sgn(\sigma)=+1$ and \emph{odd}\index{odd permutation} if
$\sgn(\sigma)=-1$.
\end{definition}

\begin{proposition}{Properties of the signature}{prop:sgn-properties}
For all $\sigma,\tau\in\mathfrak{S}_n$:
\begin{enumerate}[label=(\roman*)]
  \item $\sgn(\sigma\circ\tau) = \sgn(\sigma)\cdot\sgn(\tau)$.
  \item $\sgn(\sigma^{-1}) = \sgn(\sigma)$.
  \item Every transposition has signature $-1$.
  \item If $\sigma = \tau_1\circ\cdots\circ\tau_r$ is a product of
    transpositions, then $\sgn(\sigma)=(-1)^r$. In particular, the
    parity of~$r$ depends only on~$\sigma$, not on the decomposition.
\end{enumerate}
\end{proposition}

\begin{proof}
Consider the product formula
$\sgn(\sigma) = \prod_{i<j}\frac{\sigma(j)-\sigma(i)}{j-i}$.

(iii) Let $\tau = \tau_{pq}$ with $p < q$. In the product
$\prod_{i<j}(\tau(j)-\tau(i))$, compared to $\prod_{i<j}(j-i)$,
the factor $(q-p)$ changes sign (to $(p-q)$), and for each $k$
strictly between $p$ and $q$, the two factors involving $k$ and $p$,
and $k$ and $q$, exchange roles, contributing no net sign change.
A careful count shows the total sign is $-1$.

(i) From the product formula,
$\sgn(\sigma\circ\tau) = \prod_{i<j}
\frac{(\sigma\circ\tau)(j)-(\sigma\circ\tau)(i)}{j-i}$.
Write $j-i = \frac{j-i}{\tau(j)-\tau(i)}\cdot(\tau(j)-\tau(i))$
(since $\tau$ is a bijection, the denominator is nonzero when $i<j$),
and rearrange to factor the product as $\sgn(\sigma)\cdot\sgn(\tau)$.

(ii) Follows from (i) applied to $\sigma\circ\sigma^{-1}=\Id$
and $\sgn(\Id)=+1$.

(iv) Follows from (i) and (iii) by induction.
\end{proof}

\begin{example}{Signatures in $\mathfrak{S}_3$}{ex:S3-signatures}
The group $\mathfrak{S}_3$ has $3!=6$ elements:
\[
\begin{array}{c|cccccc}
  \sigma & \Id & (1\;2) & (1\;3) & (2\;3) & (1\;2\;3) & (1\;3\;2) \\
  \hline
  \sgn(\sigma) & +1 & -1 & -1 & -1 & +1 & +1
\end{array}
\]
The three-cycles are even (each is a product of two transpositions:
$(1\;2\;3)=(1\;3)(1\;2)$), while the transpositions are odd.
\end{example}

\begin{example}{Computing a signature}{ex:signature-computation}
Let $\sigma = (1\;3\;5\;2)\in\mathfrak{S}_5$, so $\sigma(1)=3$,
$\sigma(3)=5$, $\sigma(5)=2$, $\sigma(2)=1$, $\sigma(4)=4$.
The inversions are $(1,2),(2,5),(3,5)$, giving $N(\sigma)=3$.
Alternatively, $(1\;3\;5\;2)=(1\;2)(1\;5)(1\;3)$ is a product of
3~transpositions, so $\sgn(\sigma)=(-1)^3 = -1$.
\end{example}


% ============================================================
\section{Definition of the determinant}\label{sec:det-definition}
% ============================================================

\begin{definition}{Determinant (Leibniz formula)}{def:determinant}
\index{determinant!Leibniz formula}\index{Leibniz formula}
Let $A = (a_{ij})_{1\leq i,j\leq n}\in\Mat_n(\K)$. The
\emph{determinant} of~$A$ is
\[
  \det(A) \deff \sum_{\sigma\in\mathfrak{S}_n}
    \sgn(\sigma)\,\prod_{i=1}^n a_{i,\sigma(i)}.
\]
\end{definition}

\begin{remark}{Small cases}{rem:small-det}
For $n=1$: $\det(a) = a$.

For $n=2$:
$\det\begin{pmatrix} a & b \\ c & d \end{pmatrix} = ad - bc$.

For $n=3$:
$\det\begin{pmatrix} a_1 & a_2 & a_3 \\ b_1 & b_2 & b_3 \\
c_1 & c_2 & c_3 \end{pmatrix}
= a_1 b_2 c_3 + a_2 b_3 c_1 + a_3 b_1 c_2
  - a_3 b_2 c_1 - a_2 b_1 c_3 - a_1 b_3 c_2$,
which is the \emph{Sarrus rule}\index{Sarrus rule}.
\end{remark}


% ============================================================
\section{Alternating multilinear characterization}\label{sec:multilinear}
% ============================================================

We now give an axiomatic characterization of the determinant. Denote
the rows of $A$ by $R_1,\ldots,R_n\in\K^n$, so that we may write
$\det(A)$ as a function of these row vectors:
$\det(A) = D(R_1,\ldots,R_n)$.

\begin{definition}{Multilinear and alternating}{def:multilinear-alternating}
\index{multilinear form}\index{alternating form}
A function $D\colon (\K^n)^n \to \K$ is:
\begin{enumerate}[label=(\roman*)]
  \item \emph{multilinear} (in the rows) if for each $i$,
    $D$ is linear in the $i$-th argument with all other arguments fixed:
    \[
      D(R_1,\ldots,\alpha R_i + \beta R_i',\ldots,R_n)
      = \alpha\,D(R_1,\ldots,R_i,\ldots,R_n)
        + \beta\,D(R_1,\ldots,R_i',\ldots,R_n).
    \]
  \item \emph{alternating} if $D(R_1,\ldots,R_n)=0$ whenever two rows
    are equal: $R_i = R_j$ for some $i\neq j$.
\end{enumerate}
\end{definition}

\begin{lemma}{Alternating implies antisymmetric}{lem:alt-antisymmetric}
If $D$ is multilinear and alternating, then swapping two rows changes
the sign:
\[
  D(R_1,\ldots,R_i,\ldots,R_j,\ldots,R_n)
  = -D(R_1,\ldots,R_j,\ldots,R_i,\ldots,R_n).
\]
\end{lemma}

\begin{proof}
Expand $D(\ldots,R_i+R_j,\ldots,R_i+R_j,\ldots) = 0$ by multilinearity.
Two of the four terms vanish (identical rows), leaving
$D(\ldots,R_i,\ldots,R_j,\ldots) + D(\ldots,R_j,\ldots,R_i,\ldots) = 0$.
\end{proof}

\begin{theorem}{Uniqueness of the determinant}{thm:det-unique}
\index{determinant!axiomatic characterization}
The determinant is the unique function
$D\colon (\K^n)^n \to \K$ that is:
\begin{enumerate}[label=(\roman*)]
  \item multilinear in the rows,
  \item alternating,
  \item normalized: $D(I_n)=1$.
\end{enumerate}
\end{theorem}

\begin{proof}
\textbf{Existence.} We verify that the Leibniz formula defines a
function satisfying (i)--(iii).

Multilinearity: each term
$\sgn(\sigma)\prod_{i=1}^n a_{i,\sigma(i)}$ is linear in each
row~$R_k$ (the factor $a_{k,\sigma(k)}$ appears linearly), and the
sum of linear functions is linear.

Alternating: suppose $R_p = R_q$ for $p\neq q$. Pair each
$\sigma\in\mathfrak{S}_n$ with $\sigma'=\sigma\circ\tau_{pq}$.
Then $\sgn(\sigma')=-\sgn(\sigma)$ while
$\prod_i a_{i,\sigma'(i)}=\prod_i a_{i,\sigma(i)}$ (since
swapping rows $p,q$ with identical entries has no effect on the
product). Hence the terms cancel pairwise, giving $D=0$.

Normalization: for $A=I_n$, $a_{i,\sigma(i)}=\delta_{i,\sigma(i)}$,
so the only surviving permutation is $\sigma=\Id$, giving
$D(I_n)=\sgn(\Id)\cdot 1 = 1$.

\textbf{Uniqueness.} Let $D$ satisfy (i)--(iii). Write
$R_i = \sum_{j=1}^n a_{ij}\,\vece{j}$. By multilinearity,
\[
  D(R_1,\ldots,R_n)
  = \sum_{j_1,\ldots,j_n=1}^{n}
    a_{1,j_1}\cdots a_{n,j_n}\,D(\vece{j_1},\ldots,\vece{j_n}).
\]
By the alternating property, $D(\vece{j_1},\ldots,\vece{j_n})=0$
unless $j_1,\ldots,j_n$ are all distinct, i.e.\ $(j_1,\ldots,j_n)$
is a permutation $\sigma$. In that case, repeated application of

ef{lem:lem:alt-antisymmetric} gives
$D(\vece{\sigma(1)},\ldots,\vece{\sigma(n)})
= \sgn(\sigma)\,D(\vece{1},\ldots,\vece{n})
= \sgn(\sigma)$.
Therefore $D(R_1,\ldots,R_n) = \sum_{\sigma\in\mathfrak{S}_n}
\sgn(\sigma)\prod_{i=1}^n a_{i,\sigma(i)}$, which is the Leibniz
formula.
\end{proof}


% ============================================================
\section{Properties of the determinant}\label{sec:det-properties}
% ============================================================

\begin{proposition}{Effect of row operations}{prop:row-ops}
\index{determinant!row operations}
Let $A\in\Mat_n(\K)$ with rows $R_1,\ldots,R_n$.
\begin{enumerate}[label=(\roman*)]
  \item \textbf{Scaling:} Multiplying row $R_i$ by $\lambda\in\K$
    multiplies $\det(A)$ by $\lambda$.
  \item \textbf{Swap:} Swapping two rows multiplies $\det(A)$ by $-1$.
  \item \textbf{Row addition:} Adding a scalar multiple of one row to
    another does not change $\det(A)$:
    replacing $R_i$ by $R_i + \lambda R_j$ (with $j\neq i$) leaves
    $\det$ unchanged.
\end{enumerate}
In particular, $\det(\lambda A)=\lambda^n\det(A)$ for all
$\lambda\in\K$.
\end{proposition}

\begin{proof}
(i) and (ii) follow immediately from multilinearity and

ef{lem:lem:alt-antisymmetric}.

(iii) By multilinearity in row~$i$,
$D(R_1,\ldots,R_i+\lambda R_j,\ldots,R_n)
= D(R_1,\ldots,R_i,\ldots,R_n)
+ \lambda\,D(R_1,\ldots,R_j,\ldots,R_n)$,
where in the second term row~$j$ appears in both positions $i$
and~$j$, so the alternating property gives $0$.
\end{proof}

\begin{theorem}{Determinant of a transpose}{thm:det-transpose}
\index{determinant!of transpose}
For all $A\in\Mat_n(\K)$,
$\det(\transpose{A}) = \det(A)$.
\end{theorem}

\begin{proof}
Denote $B = \transpose{A}$, so $b_{ij}=a_{ji}$. Then
\begin{align*}
  \det(B) &= \sum_{\sigma\in\mathfrak{S}_n}\sgn(\sigma)
    \prod_{i=1}^n b_{i,\sigma(i)}
  = \sum_{\sigma\in\mathfrak{S}_n}\sgn(\sigma)
    \prod_{i=1}^n a_{\sigma(i),i}.
\end{align*}
Substituting $\tau = \sigma^{-1}$ (a bijection on $\mathfrak{S}_n$
with $\sgn(\tau)=\sgn(\sigma)$), and writing $j=\sigma(i)$
so $i=\tau(j)$:
\[
  \det(B) = \sum_{\tau\in\mathfrak{S}_n}\sgn(\tau)
    \prod_{j=1}^n a_{j,\tau(j)} = \det(A). \qedhere
\]
\end{proof}

\begin{remark}{Columns behave like rows}{rem:columns}
Since $\det(\transpose{A})=\det(A)$, every property of $\det$ stated
for rows also holds for columns: the determinant is also multilinear
and alternating in the columns.
\end{remark}

\begin{theorem}{Determinant of a product}{thm:det-product}
\index{determinant!of product}
For all $A,B\in\Mat_n(\K)$,
\[
  \det(AB) = \det(A)\det(B).
\]
\end{theorem}

\begin{proof}
Fix $A\in\Mat_n(\K)$ and define $D\colon\Mat_n(\K)\to\K$ by
$D(B) = \det(AB)$, viewed as a function of the rows of~$B$.

If $B$ has rows $R_1,\ldots,R_n$, then $AB$ has rows
$AR_1^{\mathsf{T}},\ldots,AR_n^{\mathsf{T}}$ (where we view each
$R_i$ as a row vector and $A$ acts by left-multiplication after
transposing). More precisely, the $i$-th row of $AB$ is
$\sum_{k=1}^n b_{ik}\,C_k$, where $C_k$ is the $k$-th row of~$A$
viewed appropriately. Rewriting: the $i$-th row of $AB$ is a linear
function of the $i$-th row of~$B$.

Hence $B\mapsto\det(AB)$ is multilinear and alternating in the
rows of~$B$ (since $\det$ itself is). By the uniqueness

ef{thm:thm:det-unique},
$D(B) = D(I_n)\cdot\det(B) = \det(A)\cdot\det(B)$.
\end{proof}

\begin{corollary}{Determinant of an inverse}{cor:det-inverse}
\index{determinant!of inverse}
If $A\in\GL_n(\K)$, then $\det(\inv{A}) = \frac{1}{\det(A)}$.
\end{corollary}

\begin{proof}
$\det(A)\det(\inv{A}) = \det(A\inv{A}) = \det(I_n)=1$.
\end{proof}

\begin{proposition}{Triangular matrices}{prop:det-triangular}
\index{determinant!triangular matrix}\index{triangular matrix!determinant}
If $A$ is upper or lower triangular, then
$\det(A) = \prod_{i=1}^n a_{ii}$.
In particular, $\det(I_n)=1$ and
$\det(\diag(\lambda_1,\ldots,\lambda_n)) = \lambda_1\cdots\lambda_n$.
\end{proposition}

\begin{proof}
In the Leibniz formula, the product $\prod_{i=1}^n a_{i,\sigma(i)}$
vanishes unless $\sigma(i)\geq i$ for all~$i$ (upper triangular case).
The only permutation satisfying this is $\sigma=\Id$, so
$\det(A) = \prod_{i=1}^n a_{ii}$.
\end{proof}

\begin{proposition}{Block triangular matrices}{prop:det-block}
\index{determinant!block matrix}
Let $A$ be a block upper (or lower) triangular matrix:
\[
  A = \begin{pmatrix} A_{11} & A_{12} \\ 0 & A_{22} \end{pmatrix},
\]
where $A_{11}\in\Mat_p(\K)$ and $A_{22}\in\Mat_q(\K)$ with $p+q=n$.
Then $\det(A) = \det(A_{11})\cdot\det(A_{22})$.
\end{proposition}

\begin{proof}
Fix $A_{11}$ and $A_{12}$, and consider
$D(A_{22}) \deff \det(A)$ as a function of the rows of $A_{22}$
(these are the last $q$ rows of $A$). Since the block has zeros below
$A_{11}$, the last $q$ rows of $A$ are
$(0,\ldots,0,r_{q,1},\ldots,r_{q,q})$ where $(r_{q,1},\ldots,r_{q,q})$
is a row of~$A_{22}$. By multilinearity and the alternating property
of $\det$, $D$ is multilinear and alternating in the rows of $A_{22}$.
Hence $D(A_{22}) = D(I_q)\cdot\det(A_{22})$. But when $A_{22}=I_q$,
the matrix $A$ is block upper triangular with $A_{22}=I_q$, and we
can row-reduce $A_{12}$ to zero using the last $q$ rows without
changing the determinant, obtaining
$\det\begin{psmallmatrix} A_{11} & 0 \\ 0 & I_q\end{psmallmatrix}
= \det(A_{11})$.
Therefore $D(I_q)=\det(A_{11})$.
\end{proof}


% ============================================================
\section{Cofactor expansion (Laplace expansion)}\label{sec:cofactor}
% ============================================================

\begin{definition}{Minor and cofactor}{def:cofactor}
\index{minor}\index{cofactor}
Let $A=(a_{ij})\in\Mat_n(\K)$. The \emph{$(i,j)$-minor} of~$A$,
denoted $M_{ij}$, is the determinant of the $(n-1)\times(n-1)$
submatrix obtained by deleting row~$i$ and column~$j$. The
\emph{$(i,j)$-cofactor} is
\[
  C_{ij} \deff (-1)^{i+j}\,M_{ij}.
\]
\end{definition}

\begin{theorem}{Laplace expansion}{thm:laplace}
\index{Laplace expansion}\index{cofactor expansion}
For any $A\in\Mat_n(\K)$:
\begin{enumerate}[label=(\roman*)]
  \item \textbf{Expansion along row~$i$:}\quad
    $\displaystyle\det(A) = \sum_{j=1}^n a_{ij}\,C_{ij}
     = \sum_{j=1}^n (-1)^{i+j}\,a_{ij}\,M_{ij}$.
  \item \textbf{Expansion along column~$j$:}\quad
    $\displaystyle\det(A) = \sum_{i=1}^n a_{ij}\,C_{ij}
     = \sum_{i=1}^n (-1)^{i+j}\,a_{ij}\,M_{ij}$.
\end{enumerate}
Moreover, expanding along a ``foreign'' row or column gives zero:
\[
  \sum_{j=1}^n a_{ij}\,C_{kj} = 0 \quad\text{for } i\neq k,
  \qquad
  \sum_{i=1}^n a_{ij}\,C_{ik} = 0 \quad\text{for } j\neq k.
\]
\end{theorem}

\begin{proof}
\textbf{(i) Expansion along row~$i$.}
In the Leibniz formula $\det(A)=\sum_{\sigma}\sgn(\sigma)
\prod_{\ell=1}^n a_{\ell,\sigma(\ell)}$, group the terms according
to the value $\sigma(i)=j$:
\[
  \det(A) = \sum_{j=1}^n a_{ij}
    \sum_{\substack{\sigma\in\mathfrak{S}_n\\\sigma(i)=j}}
    \sgn(\sigma)\prod_{\ell\neq i}a_{\ell,\sigma(\ell)}.
\]
We claim that
$\sum_{\sigma:\sigma(i)=j}\sgn(\sigma)\prod_{\ell\neq i}
a_{\ell,\sigma(\ell)} = (-1)^{i+j}M_{ij}$.

To see this, consider the permutations $\sigma\in\mathfrak{S}_n$
with $\sigma(i)=j$. Any such $\sigma$ is determined by its restriction
to $\set{1,\ldots,n}\setminus\set{i}$, which is a bijection onto
$\set{1,\ldots,n}\setminus\set{j}$. We can identify this restriction
with a permutation $\sigma'\in\mathfrak{S}_{n-1}$ after relabelling:
the rows $1,\ldots,i-1,i+1,\ldots,n$ become $1,\ldots,n-1$ and
similarly for columns. The relabelling involves moving index~$i$ to
position~$n$ (via $i-1$ adjacent transpositions in the row indices)
and index~$j$ to position~$n$ (via $j-1$ adjacent transpositions
in the column indices), contributing a sign of
$(-1)^{(n-i)+(n-j)}=(-1)^{i+j}$ (since $(-1)^{2n}=1$).
Thus $\sgn(\sigma) = (-1)^{i+j}\sgn(\sigma')$, and the inner sum
becomes $(-1)^{i+j}\det(A_{ij})=(-1)^{i+j}M_{ij}=C_{ij}$.

\textbf{(ii)} follows from (i) applied to $\transpose{A}$, using
$\det(\transpose{A})=\det(A)$.

\textbf{Foreign expansion.}
Consider $\sum_{j=1}^n a_{ij}C_{kj}$ with $i\neq k$. This equals the
determinant of the matrix $A'$ obtained from $A$ by replacing row~$k$
with row~$R_i$ (since the cofactors $C_{kj}$ depend only on the rows
other than row~$k$, and the formula gives the expansion of $\det(A')$
along row~$k$). But $A'$ has two identical rows ($R_i$ appears in both
positions $i$ and~$k$), so $\det(A')=0$ by the alternating property.
\end{proof}

\begin{example}{$3\times 3$ determinant by cofactor expansion}{ex:3x3-cofactor}
Compute $\det\begin{pmatrix} 2 & 1 & 3 \\ 0 & -1 & 4 \\ 5 & 2 & 1
\end{pmatrix}$ by expanding along the first column:
\begin{align*}
  \det(A) &= 2\cdot(-1)^{1+1}\det\begin{pmatrix}-1&4\\2&1\end{pmatrix}
  + 0\cdot(-1)^{2+1}\det\begin{pmatrix}1&3\\2&1\end{pmatrix}
  + 5\cdot(-1)^{3+1}\det\begin{pmatrix}1&3\\-1&4\end{pmatrix}\\
  &= 2(-1-8) + 0 + 5(4+3) = 2(-9)+5\cdot 7 = -18+35 = 17.
\end{align*}
\end{example}


% ============================================================
\section{The adjugate matrix and Cramer's rule}\label{sec:adjugate}
% ============================================================

\begin{definition}{Adjugate (classical adjoint)}{def:adjugate}
\index{adjugate matrix}\index{comatrix}
The \emph{adjugate} (or \emph{classical adjoint}, or
\emph{comatrix}\index{comatrix|see{adjugate matrix}}) of
$A\in\Mat_n(\K)$ is the $n\times n$ matrix
\[
  \operatorname{adj}(A) \deff \bigl(C_{ji}\bigr)_{1\leq i,j\leq n}
  = \transpose{(\text{cofactor matrix})},
\]
where $C_{ij}=(-1)^{i+j}M_{ij}$ is the $(i,j)$-cofactor.
\end{definition}

\begin{theorem}{Adjugate identity}{thm:adjugate}
\index{adjugate matrix!identity}
For all $A\in\Mat_n(\K)$,
\[
  A\cdot\operatorname{adj}(A) = \operatorname{adj}(A)\cdot A
  = \det(A)\cdot I_n.
\]
\end{theorem}

\begin{proof}
The $(i,k)$-entry of $A\cdot\operatorname{adj}(A)$ is
$\sum_{j=1}^n a_{ij}\,[\operatorname{adj}(A)]_{jk}
= \sum_{j=1}^n a_{ij}\,C_{kj}$.
By 
ef{thm:thm:laplace}, this equals $\det(A)$ if $i=k$ and $0$ if
$i\neq k$ (foreign expansion). Thus
$A\cdot\operatorname{adj}(A)=\det(A)\,I_n$. The identity
$\operatorname{adj}(A)\cdot A = \det(A)\,I_n$ is proved by
considering columns instead of rows (or by applying the row result
to $\transpose{A}$ and transposing).
\end{proof}

\begin{corollary}{Inverse via adjugate}{cor:inverse-adjugate}
\index{inverse matrix!via adjugate}
If $\det(A)\neq 0$, then $A$ is invertible and
\[
  \inv{A} = \frac{1}{\det(A)}\operatorname{adj}(A).
\]
\end{corollary}

\begin{theorem}{Cramer's rule}{thm:cramer}
\index{Cramer's rule}
Let $A\in\Mat_n(\K)$ with $\det(A)\neq 0$, and let $b\in\K^n$.
The unique solution $x=(x_1,\ldots,x_n)^{\mathsf{T}}$ of $Ax=b$ is
given by
\[
  x_j = \frac{\det(A_j)}{\det(A)},\qquad j=1,\ldots,n,
\]
where $A_j$ is the matrix obtained from~$A$ by replacing column~$j$
with~$b$.
\end{theorem}

\begin{proof}
Since $\det(A)\neq 0$, the system has a unique solution
$x = \inv{A}b = \frac{1}{\det(A)}\operatorname{adj}(A)\,b$.
The $j$-th component is
\[
  x_j = \frac{1}{\det(A)}\sum_{i=1}^n C_{ij}\,b_i
  = \frac{1}{\det(A)}\sum_{i=1}^n (-1)^{i+j}M_{ij}\,b_i.
\]
But $\sum_{i=1}^n (-1)^{i+j}b_i\,M_{ij}$ is precisely the cofactor
expansion of $\det(A_j)$ along column~$j$ (since column~$j$ of $A_j$
is~$b$, and the minors $M_{ij}$ are unchanged). Hence
$x_j = \det(A_j)/\det(A)$.
\end{proof}

\begin{example}{Cramer's rule in dimension 2}{ex:cramer-2d}
Solve $\begin{cases} 3x+2y=7 \\ x-y=1\end{cases}$.

Here $A=\begin{psmallmatrix}3&2\\1&-1\end{psmallmatrix}$,
$\det(A)=-3-2=-5$,
$\det(A_1)=\det\begin{psmallmatrix}7&2\\1&-1\end{psmallmatrix}=-9$,
$\det(A_2)=\det\begin{psmallmatrix}3&7\\1&1\end{psmallmatrix}=-4$.
So $x = \frac{-9}{-5}=\frac{9}{5}$,
$y=\frac{-4}{-5}=\frac{4}{5}$.
\end{example}


% ============================================================
\section{Determinant and invertibility}\label{sec:det-invertibility}
% ============================================================

\begin{theorem}{Invertibility criterion}{thm:det-invertibility}
\index{determinant!and invertibility}\index{invertible matrix!determinant criterion}
A matrix $A\in\Mat_n(\K)$ is invertible if and only if $\det(A)\neq 0$.
\end{theorem}

\begin{proof}
$(\Rightarrow)$\; If $A$ is invertible, then
$\det(A)\det(\inv{A})=\det(I_n)=1$, so $\det(A)\neq 0$.

$(\Leftarrow)$\; If $\det(A)\neq 0$, then by 
ef{thm:thm:adjugate},
$A\cdot\frac{1}{\det(A)}\operatorname{adj}(A) = I_n$,
so $A$ is invertible (with explicit inverse given by

ef{cor:cor:inverse-adjugate}).

Alternatively, if $A$ is not invertible, its columns are linearly
dependent: some column is a linear combination of the others. By
multilinearity and the alternating property in columns,
$\det(A)=0$.
\end{proof}

\begin{corollary}{The general linear group}{cor:GL-det}
\index{general linear group}
$\GL_n(\K) = \setcond{A\in\Mat_n(\K)}{\det(A)\neq 0}$.
The determinant defines a group homomorphism
$\det\colon\GL_n(\K)\to\K^*$ (where $\K^*=\K\setminus\set{0}$
is the multiplicative group of~$\K$).
\end{corollary}

\begin{proof}
The first statement is 
ef{thm:thm:det-invertibility}. The
multiplicativity $\det(AB)=\det(A)\det(B)$
(
ef{thm:thm:det-product}) shows that $\det$ is a group homomorphism.
Its kernel is the \emph{special linear group}
$\SL_n(\K)=\setcond{A\in\GL_n(\K)}{\det(A)=1}$.
\end{proof}


% ============================================================
\section{Determinant of a linear map}\label{sec:det-linear-map}
% ============================================================

\begin{definition}{Determinant of an endomorphism}{def:det-endo}
\index{determinant!of a linear map}
Let $E$ be a finite-dimensional $\K$-vector space and
$f\in\End(E)$. Choose any basis $\cB$ of~$E$ and define
\[
  \det(f) \deff \det\bigl(\Mat_\cB(f)\bigr).
\]
\end{definition}

\begin{proposition}{Well-definedness}{prop:det-well-defined}
The determinant $\det(f)$ does not depend on the choice of basis.
\end{proposition}

\begin{proof}
Let $\cB$ and $\cB'$ be two bases of~$E$, and let $P$ be the
change-of-basis matrix from $\cB$ to $\cB'$. Then
$\Mat_{\cB'}(f) = \inv{P}\,\Mat_\cB(f)\,P$, so
\[
  \det\bigl(\Mat_{\cB'}(f)\bigr)
  = \det(\inv{P})\det\bigl(\Mat_\cB(f)\bigr)\det(P)
  = \frac{1}{\det(P)}\det\bigl(\Mat_\cB(f)\bigr)\det(P)
  = \det\bigl(\Mat_\cB(f)\bigr). \qedhere
\]
\end{proof}

\begin{proposition}{Properties of $\det$ for endomorphisms}{prop:det-endo-props}
Let $E$ be a finite-dimensional $\K$-vector space and
$f,g\in\End(E)$.
\begin{enumerate}[label=(\roman*)]
  \item $\det(g\circ f) = \det(g)\det(f)$.
  \item $f$ is an automorphism if and only if $\det(f)\neq 0$.
  \item $\det(\Id_E)=1$.
  \item $\det(\lambda f)=\lambda^n\det(f)$ where $n=\dim E$.
\end{enumerate}
\end{proposition}

\begin{proof}
All properties follow immediately from the corresponding properties
of matrix determinants and the identity
$\Mat_\cB(g\circ f)=\Mat_\cB(g)\cdot\Mat_\cB(f)$.
\end{proof}


% ============================================================
\section{Geometric interpretation}\label{sec:det-geometry}
% ============================================================

The determinant has a rich geometric meaning in~$\R^n$.

\begin{theorem}{Determinant as signed volume}{thm:signed-volume}
\index{determinant!geometric interpretation}\index{signed volume}
Let $v_1,\ldots,v_n\in\R^n$ be the column vectors of a matrix
$A\in\Mat_n(\R)$. Then:
\begin{enumerate}[label=(\roman*)]
  \item $\abs{\det(A)}$ equals the $n$-dimensional volume of the
    parallelepiped
    $P = \setcond{t_1 v_1+\cdots+t_n v_n}{0\leq t_i\leq 1}$.
  \item $\det(A)>0$ if and only if the ordered basis
    $(v_1,\ldots,v_n)$ has the same orientation as the standard basis.
\end{enumerate}
\end{theorem}

\begin{proof}[Proof sketch]
(i) Both $\abs{\det}$ and the volume function are multiplicative
under composition ($\abs{\det(AB)}=\abs{\det A}\abs{\det B}$), agree
on elementary matrices (which correspond to elementary row
operations with known geometric effects), and every invertible matrix
is a product of elementary matrices. For the singular case, both
sides equal zero.

(ii) The determinant is a continuous function of the columns.
The set of invertible matrices $\GL_n(\R)$ has exactly two connected
components, distinguished by the sign of $\det$. The identity matrix
$I_n$ has $\det=1>0$ and corresponds to the standard orientation.
\end{proof}

\begin{center}
\begin{tikzpicture}[>=Stealth, scale=1.0]
  % 2D parallelogram
  \begin{scope}
    \coordinate (O) at (0,0);
    \coordinate (U) at (2.5,0.5);
    \coordinate (V) at (0.8,2);
    \coordinate (UV) at ($(U)+(V)$);

    \fill[blue!15] (O) -- (U) -- (UV) -- (V) -- cycle;
    \draw[thick, blue!70!black, ->] (O) -- (U)
      node[below right] {$\vec{u}$};
    \draw[thick, red!70!black, ->] (O) -- (V)
      node[above left] {$\vec{v}$};
    \draw[dashed, gray] (U) -- (UV) -- (V);
    \node at (1.6,1.2) {Area $= \abs{\det}$};
    \node[below] at (1.6,-0.5) {\textbf{2D:} parallelogram};
  \end{scope}

  % 3D parallelepiped
  \begin{scope}[shift={(7,0)}, x={(1,0)}, y={(0.4,0.5)}, z={(0,1)}]
    \coordinate (O) at (0,0,0);
    \coordinate (A) at (2,0,0);
    \coordinate (B) at (0,1.5,0);
    \coordinate (C) at (0,0,1.8);
    \coordinate (AB) at (2,1.5,0);
    \coordinate (AC) at (2,0,1.8);
    \coordinate (BC) at (0,1.5,1.8);
    \coordinate (ABC) at (2,1.5,1.8);

    % Back faces
    \fill[red!8] (O) -- (B) -- (AB) -- (A) -- cycle;
    \fill[blue!8] (O) -- (B) -- (BC) -- (C) -- cycle;
    \fill[green!8] (O) -- (A) -- (AC) -- (C) -- cycle;

    % Front faces
    \fill[red!15] (C) -- (BC) -- (ABC) -- (AC) -- cycle;
    \fill[blue!15] (A) -- (AB) -- (ABC) -- (AC) -- cycle;
    \fill[green!15] (B) -- (AB) -- (ABC) -- (BC) -- cycle;

    % Edges
    \draw[thick] (O) -- (A) -- (AB) -- (B) -- cycle;
    \draw[thick] (O) -- (C) -- (AC) -- (A);
    \draw[thick] (C) -- (BC) -- (B);
    \draw[thick] (AC) -- (ABC) -- (AB);
    \draw[thick] (ABC) -- (BC);

    % Vectors
    \draw[->, very thick, blue!70!black] (O) -- (A)
      node[below right] {$\vec{u}$};
    \draw[->, very thick, red!70!black] (O) -- (B)
      node[left] {$\vec{v}$};
    \draw[->, very thick, green!60!black] (O) -- (C)
      node[left] {$\vec{w}$};

    \node at (1.2,0.7,0.9) {\small Vol};
    \node[below] at (1,-0.2,-0.5) {\textbf{3D:} parallelepiped};
  \end{scope}
\end{tikzpicture}
\captionof{figure}{The determinant computes signed areas and volumes.}
\label{fig:det-area-volume}
\end{center}

\begin{center}
\begin{tikzpicture}[>=Stealth, scale=1.2]
  % Positive orientation
  \begin{scope}
    \draw[->] (-0.3,0) -- (2.5,0) node[right] {$x$};
    \draw[->] (0,-0.3) -- (0,2.2) node[above] {$y$};
    \draw[->, very thick, blue!70!black] (0,0) -- (2,0.3)
      node[below right] {$\vec{u}$};
    \draw[->, very thick, red!70!black] (0,0) -- (-0.2,1.8)
      node[above left] {$\vec{v}$};
    \draw[->, thick, green!50!black] (0.7,0.1) arc[start angle=8,
      end angle=84, radius=0.7];
    \node[green!50!black] at (0.7,0.9) {\small $+$};
    \node[below] at (1,-0.6) {$\det > 0$};
    \node[below] at (1,-1.0) {(positive orientation)};
  \end{scope}

  % Negative orientation
  \begin{scope}[shift={(5.5,0)}]
    \draw[->] (-0.3,0) -- (2.5,0) node[right] {$x$};
    \draw[->] (0,-0.3) -- (0,2.2) node[above] {$y$};
    \draw[->, very thick, blue!70!black] (0,0) -- (-0.2,1.8)
      node[above left] {$\vec{u}$};
    \draw[->, very thick, red!70!black] (0,0) -- (2,0.3)
      node[below right] {$\vec{v}$};
    \draw[->, thick, orange!70!red] (0.05,0.65) arc[start angle=84,
      end angle=8, radius=0.7];
    \node[orange!70!red] at (0.7,0.9) {\small $-$};
    \node[below] at (1,-0.6) {$\det < 0$};
    \node[below] at (1,-1.0) {(negative orientation)};
  \end{scope}
\end{tikzpicture}
\captionof{figure}{The sign of the determinant encodes orientation.}
\label{fig:det-orientation}
\end{center}


% ============================================================
\section{Applications}\label{sec:det-applications}
% ============================================================

\subsection{Cross product in $\R^3$}\label{subsec:cross-product}

\begin{definition}{Cross product}{def:cross-product}
\index{cross product}
For $\vec{u}=(u_1,u_2,u_3)$ and $\vec{v}=(v_1,v_2,v_3)$ in~$\R^3$,
the \emph{cross product} is defined (formally) by
\[
  \vec{u}\times\vec{v} \deff
  \det\begin{pmatrix}
    \vece{1} & \vece{2} & \vece{3} \\
    u_1 & u_2 & u_3 \\
    v_1 & v_2 & v_3
  \end{pmatrix}
  = \begin{pmatrix}
    u_2 v_3 - u_3 v_2 \\
    u_3 v_1 - u_1 v_3 \\
    u_1 v_2 - u_2 v_1
  \end{pmatrix}.
\]
\end{definition}

\begin{proposition}{Properties of the cross product}{prop:cross-product}
For $\vec{u},\vec{v},\vec{w}\in\R^3$ and $\alpha\in\R$:
\begin{enumerate}[label=(\roman*)]
  \item \textbf{Antisymmetry:}
    $\vec{v}\times\vec{u} = -(\vec{u}\times\vec{v})$.
  \item \textbf{Bilinearity:}
    $(\alpha\vec{u}+\vec{w})\times\vec{v}
    = \alpha(\vec{u}\times\vec{v})+\vec{w}\times\vec{v}$.
  \item \textbf{Orthogonality:}
    $(\vec{u}\times\vec{v})\cdot\vec{u}
    = (\vec{u}\times\vec{v})\cdot\vec{v} = 0$.
  \item \textbf{Norm:}
    $\norm{\vec{u}\times\vec{v}}$ equals the area of the
    parallelogram Spanned by $\vec{u}$ and~$\vec{v}$.
  \item \textbf{Scalar triple product:}
    $(\vec{u}\times\vec{v})\cdot\vec{w}
    = \det(\vec{u},\vec{v},\vec{w})$ (columns).
\end{enumerate}
\end{proposition}

\begin{proof}
(i)--(ii) follow from properties of the determinant.
(iii) If we substitute $\vec{u}$ for $\vec{w}$ in the triple product
$\det(\vec{u},\vec{v},\vec{w})$, the matrix has two identical columns,
so $\det=0$.
(iv) follows from the identity
$\norm{\vec{u}\times\vec{v}}^2 = \norm{\vec{u}}^2\norm{\vec{v}}^2
- (\vec{u}\cdot\vec{v})^2$ (Lagrange identity).
(v) is verified by direct computation: both sides equal
$u_1(v_2 w_3-v_3 w_2) - u_2(v_1 w_3-v_3 w_1) + u_3(v_1 w_2-v_2 w_1)$.
\end{proof}

\subsection{Area and volume formulas}\label{subsec:area-volume}

\begin{proposition}{Area of a triangle}{prop:triangle-area}
\index{area!triangle}
The area of a triangle with vertices $P_1=(x_1,y_1)$,
$P_2=(x_2,y_2)$, $P_3=(x_3,y_3)$ in $\R^2$ is
\[
  \text{Area} = \frac{1}{2}\abs*{\det\begin{pmatrix}
    x_2-x_1 & x_3-x_1 \\
    y_2-y_1 & y_3-y_1
  \end{pmatrix}}
  = \frac{1}{2}\abs*{\det\begin{pmatrix}
    x_1 & y_1 & 1 \\
    x_2 & y_2 & 1 \\
    x_3 & y_3 & 1
  \end{pmatrix}}.
\]
\end{proposition}

\begin{proof}
The triangle is half the parallelogram Spanned by
$\overrightarrow{P_1 P_2}$ and $\overrightarrow{P_1 P_3}$.
The second formula follows by cofactor expansion along the third
column of the $3\times 3$ matrix.
\end{proof}

\subsection{Vandermonde determinant}\label{subsec:vandermonde}

\begin{theorem}{Vandermonde determinant}{thm:vandermonde}
\index{Vandermonde determinant}
For $x_1,\ldots,x_n\in\K$,
\[
  V(x_1,\ldots,x_n) \deff
  \det\begin{pmatrix}
    1 & x_1 & x_1^2 & \cdots & x_1^{n-1} \\
    1 & x_2 & x_2^2 & \cdots & x_2^{n-1} \\
    \vdots & \vdots & \vdots & \ddots & \vdots \\
    1 & x_n & x_n^2 & \cdots & x_n^{n-1}
  \end{pmatrix}
  = \prod_{1\leq i < j \leq n} (x_j - x_i).
\]
In particular, $V(x_1,\ldots,x_n)\neq 0$ if and only if the
$x_i$ are pairwise distinct.
\end{theorem}

\begin{proof}
We proceed by induction on~$n$. For $n=1$, $V(x_1)=1$ and the empty
product is~$1$. For $n=2$, $V(x_1,x_2) = x_2 - x_1$.

For the inductive step, subtract $x_1$ times column $j-1$ from
column~$j$, for $j=n,n-1,\ldots,2$ (right to left). This does not
change the determinant. Row~$1$ becomes $(1,0,0,\ldots,0)$, and for
$i\geq 2$, the entry in row~$i$, column~$j$ becomes
$x_i^{j-1}-x_1 x_i^{j-2} = x_i^{j-2}(x_i - x_1)$.
Expanding along row~$1$ and factoring out $(x_i-x_1)$ from row~$i$
(for $i=2,\ldots,n$):
\[
  V(x_1,\ldots,x_n) = \prod_{i=2}^n (x_i-x_1)\cdot V(x_2,\ldots,x_n).
\]
By the induction hypothesis,
$V(x_2,\ldots,x_n) = \prod_{2\leq i<j\leq n}(x_j-x_i)$, so
\[
  V(x_1,\ldots,x_n)
  = \prod_{i=2}^n (x_i-x_1)\cdot\prod_{2\leq i<j\leq n}(x_j-x_i)
  = \prod_{1\leq i<j\leq n}(x_j-x_i). \qedhere
\]
\end{proof}

\begin{example}{Vandermonde for $n=3$}{ex:vandermonde-3}
\[
  \det\begin{pmatrix}1 & a & a^2 \\ 1 & b & b^2 \\ 1 & c & c^2
  \end{pmatrix}
  = (b-a)(c-a)(c-b).
\]
\end{example}

\begin{remark}{Application of Vandermonde}{rem:vandermonde-app}
The Vandermonde determinant appears in polynomial interpolation:
the system expressing a polynomial of degree $\leq n-1$ through
$n$ points $(x_i,y_i)$ has a unique solution if and only if
the $x_i$ are distinct, since the coefficient matrix is the
Vandermonde matrix.
\end{remark}


% ============================================================
\section{Exercises}\label{sec:det-exercises}
% ============================================================

\begin{exercise}{Determinants of small matrices \basic}{exo:det-small}
Compute the following determinants:
\begin{enumerate}[label=(\alph*)]
  \item $\det\begin{pmatrix}3 & 1 \\ 5 & 2\end{pmatrix}$.
  \item $\det\begin{pmatrix}2 & -1 & 3 \\ 0 & 4 & 1 \\ 1 & 2 & -1\end{pmatrix}$.
  \item $\det\begin{pmatrix}1 & 2 & 3 & 4 \\ 0 & 1 & 2 & 3 \\
    0 & 0 & 1 & 2 \\ 0 & 0 & 0 & 1\end{pmatrix}$.
\end{enumerate}
\end{exercise}

\begin{exercise}{Row operations and determinant \basic}{exo:row-ops-det}
Let $A\in\Mat_3(\R)$ with $\det(A)=6$. Compute $\det(B)$ where $B$ is
obtained from $A$ by:
\begin{enumerate}[label=(\alph*)]
  \item swapping rows 1 and 3,
  \item multiplying row 2 by $-4$,
  \item adding 5 times row 1 to row 2,
  \item performing all three operations in sequence.
\end{enumerate}
\end{exercise}

\begin{exercise}{Determinant and scalar multiplication \basic}{exo:det-scalar}
\index{determinant!scalar multiplication}
Let $A\in\Mat_n(\K)$ with $\det(A)=d$. Express the following in terms
of $d$ and~$n$:
\begin{enumerate}[label=(\alph*)]
  \item $\det(3A)$.
  \item $\det(-A)$.
  \item $\det(\transpose{A})$.
  \item $\det(A^2)$.
\end{enumerate}
\end{exercise}

\begin{exercise}{Determinant of a rank-1 update \intermediate}{exo:rank-1-update}
\index{matrix determinant lemma}
Let $A\in\GL_n(\K)$ and let $u,v\in\K^n$ (column vectors). Prove
the \emph{matrix determinant lemma}:
\[
  \det(A + uv^{\mathsf{T}}) = (1 + v^{\mathsf{T}}\inv{A}u)\det(A).
\]
\textit{Hint:} Factor $A+uv^{\mathsf{T}} = A(I+\inv{A}uv^{\mathsf{T}})$
and compute the determinant of $I+wv^{\mathsf{T}}$ for a column
vector $w=\inv{A}u$.
\end{exercise}

\begin{exercise}{Cofactor expansion practice \basic}{exo:cofactor-practice}
Compute the determinant of
\[
  A = \begin{pmatrix}
    2 & 0 & 1 & 3 \\
    1 & -1 & 0 & 2 \\
    3 & 2 & 0 & 1 \\
    0 & 4 & -1 & 0
  \end{pmatrix}
\]
by expanding along a row or column of your choice (choose wisely to
minimise work).
\end{exercise}

\begin{exercise}{Adjugate of a $2\times 2$ matrix \basic}{exo:adjugate-2x2}
For $A=\begin{psmallmatrix}a&b\\c&d\end{psmallmatrix}$, compute
$\operatorname{adj}(A)$ and verify that
$A\cdot\operatorname{adj}(A)=\det(A)\,I_2$.
\end{exercise}

\begin{exercise}{Cramer's rule application \intermediate}{exo:cramer-app}
Use Cramer's rule to solve the system
\[
  \begin{cases}
    x + 2y - z = 3,\\
    2x - y + 3z = 1,\\
    3x + y + 2z = 7.
  \end{cases}
\]
\end{exercise}

\begin{exercise}{Vandermonde determinant \intermediate}{exo:vandermonde}
\begin{enumerate}[label=(\alph*)]
  \item Compute the Vandermonde determinant $V(1,2,3,4)$.
  \item For which values of $a\in\R$ is the Vandermonde matrix
    with nodes $1,a,a^2$ singular?
\end{enumerate}
\end{exercise}

\begin{exercise}{Determinant and eigenvalues \intermediate}{exo:det-eigenvalues}
\index{eigenvalue!and determinant}
Let $A\in\Mat_n(\K)$.
\begin{enumerate}[label=(\alph*)]
  \item Show that $\lambda$ is an eigenvalue of $A$ if and only if
    $\det(A-\lambda I_n)=0$.
  \item If $A$ has eigenvalues $\lambda_1,\ldots,\lambda_n$
    (counted with algebraic multiplicity), prove that
    $\det(A)=\lambda_1\cdots\lambda_n$.
\end{enumerate}
\end{exercise}

\begin{exercise}{Determinant of a block diagonal matrix \basic}{exo:block-diag}
Let $A=\diag(A_1,A_2,\ldots,A_k)$ be a block diagonal matrix where
$A_i\in\Mat_{n_i}(\K)$. Prove that
$\det(A)=\det(A_1)\det(A_2)\cdots\det(A_k)$.
\end{exercise}

\begin{exercise}{Antisymmetric matrices \intermediate}{exo:antisymmetric}
\index{antisymmetric matrix!determinant}
Let $A\in\Mat_n(\R)$ be \emph{antisymmetric}, i.e.\
$\transpose{A}=-A$.
\begin{enumerate}[label=(\alph*)]
  \item Show that if $n$ is odd, then $\det(A)=0$.
  \item Give an example of a $2\times 2$ antisymmetric matrix
    with $\det(A)\neq 0$.
  \item Show that for even~$n$, $\det(A)\geq 0$.
    \textit{Hint:} Consider $\det(A)=\det(\transpose{A})=\det(-A)$.
\end{enumerate}
\end{exercise}

\begin{exercise}{Derivative of the determinant \challenging}{exo:det-derivative}
\index{determinant!derivative}
Let $A(t)=(a_{ij}(t))$ be an $n\times n$ matrix whose entries are
differentiable functions of $t\in\R$. Prove \emph{Jacobi's formula}:
\[
  \frac{d}{dt}\det(A(t))
  = \sum_{i=1}^n \det\bigl(A_1(t),\ldots,A_i'(t),\ldots,A_n(t)\bigr),
\]
where $A_i(t)$ denotes the $i$-th row of $A(t)$ and $A_i'(t)$ its
derivative.
\textit{Hint:} Use multilinearity of $\det$ in the rows and the
product rule.
\end{exercise}

\begin{exercise}{Cauchy determinant \challenging}{exo:cauchy-det}
\index{Cauchy determinant}
Let $x_1,\ldots,x_n$ and $y_1,\ldots,y_n$ be elements of a
field~$\K$ with $x_i+y_j\neq 0$ for all $i,j$. The \emph{Cauchy
matrix} is $C=(c_{ij})$ with $c_{ij}=\frac{1}{x_i+y_j}$.
Prove the \emph{Cauchy determinant formula}:
\[
  \det(C) = \frac{\prod_{1\leq i<j\leq n}(x_j-x_i)(y_j-y_i)}%
  {\prod_{i,j=1}^{n}(x_i+y_j)}.
\]
\textit{Hint:} Proceed by induction, reducing to a Vandermonde-type
computation.
\end{exercise}

\begin{exercise}{Determinant of the commutator \challenging}{exo:commutator-det}
Let $A,B\in\Mat_n(\K)$.
\begin{enumerate}[label=(\alph*)]
  \item Show that if $n=2$ and $\tr(A)=\tr(B)=0$, then
    $\det(AB-BA)\leq 0$ when $\K=\R$.
  \item Construct $A,B\in\Mat_3(\R)$ such that $\det(AB-BA)\neq 0$.
  \item Show that for any $n$, $\tr(AB-BA)=0$.
\end{enumerate}
\end{exercise}

\begin{exercise}{Polynomial identity for determinants \intermediate}{exo:det-poly}
Let $A,B\in\Mat_n(\K)$. Show that
\[
  \det(A+B) + \det(A-B) = 2\sum_{k\text{ even}}D_k(A,B),
\]
where $D_k(A,B)$ is the sum of all determinants obtained by choosing
$k$ rows from $B$ and the remaining $n-k$ rows from $A$ (in their
original positions). Verify this for $n=2$ explicitly.
\end{exercise}


% ============================================================
\section*{Chapter summary}\label{sec:summary-ch4}
\addcontentsline{toc}{section}{Chapter summary}
% ============================================================

\begin{itemize}
  \item A \textbf{permutation}\index{permutation} of $\set{1,\ldots,n}$
    is a bijection; the set of all permutations forms the
    \textbf{symmetric group}~$\mathfrak{S}_n$ of order $n!$.
    Every permutation has a well-defined
    \textbf{signature}\index{signature!of a permutation}
    $\sgn(\sigma)=\pm 1$ (
ef{def:def:signature}).

  \item The \textbf{determinant}\index{determinant} of
    $A\in\Mat_n(\K)$ is defined by the \textbf{Leibniz formula}
    (
ef{def:def:determinant}):
    $\det(A) = \sum_{\sigma\in\mathfrak{S}_n}\sgn(\sigma)
    \prod_{i=1}^n a_{i,\sigma(i)}$.

  \item The determinant is the \emph{unique} function that is
    \textbf{multilinear}, \textbf{alternating} in the rows, and
    satisfies $\det(I_n)=1$ (
ef{thm:thm:det-unique}).

  \item Key \textbf{properties}: $\det(\transpose{A})=\det(A)$
    (
ef{thm:thm:det-transpose}), $\det(AB)=\det(A)\det(B)$
    (
ef{thm:thm:det-product}), triangular matrices have
    $\det=\prod a_{ii}$ (
ef{prop:prop:det-triangular}).

  \item \textbf{Cofactor (Laplace) expansion}\index{Laplace expansion}
    computes $\det$ by expanding along any row or column
    (
ef{thm:thm:laplace}).

  \item The \textbf{adjugate matrix}\index{adjugate matrix} satisfies
    $A\cdot\operatorname{adj}(A) = \det(A)\,I_n$
    (
ef{thm:thm:adjugate}), yielding the formula
    $\inv{A}=\frac{1}{\det(A)}\operatorname{adj}(A)$.

  \item \textbf{Cramer's rule}\index{Cramer's rule} gives an explicit
    formula for the solution of $Ax=b$ when $\det(A)\neq 0$
    (
ef{thm:thm:cramer}).

  \item $A$ is \textbf{invertible} if and only if $\det(A)\neq 0$
    (
ef{thm:thm:det-invertibility}).

  \item The determinant of a \textbf{linear map} is well-defined
    (independent of the choice of basis, 
ef{prop:prop:det-well-defined}).

  \item \textbf{Geometrically}, $\abs{\det(A)}$ is the volume of the
    parallelepiped Spanned by the columns, and $\sgn(\det)$ encodes
    orientation (
ef{thm:thm:signed-volume}).

  \item \textbf{Applications}: the cross product in~$\R^3$, area and
    volume formulas, and the \textbf{Vandermonde
    determinant}\index{Vandermonde determinant}
    $\prod_{i<j}(x_j-x_i)$ (
ef{thm:thm:vandermonde}).
\end{itemize}

% ch5-eigenvalues.tex — Chapter 5: Eigenvalues, Eigenvectors, and Diagonalization
% Definitions, characteristic polynomial, Cayley–Hamilton, diagonalizability,
% minimal polynomial, applications (A^n, Fibonacci, Markov chains, ODEs).

\chapter{Eigenvalues, Eigenvectors, and Diagonalization}\label{ch:eigenvalues}

% ============================================================
% Motivating introduction
% ============================================================

Among all the questions one can ask about a linear map, perhaps the
most fundamental is this: \emph{does there exist a basis in which the
map takes a particularly simple form?}  The simplest form one can hope
for is a diagonal matrix, and the search for such a basis leads
directly to the theory of eigenvalues and eigenvectors.

This theory is not merely an exercise in abstraction.  Eigenvalues
appear throughout mathematics and its applications:
\begin{itemize}
  \item \textbf{Differential equations.}  The solutions of a linear
    system $\vec{x}'(t) = A\vec{x}(t)$ are determined by the
    eigenvalues of~$A$; exponential growth, decay, and oscillation all
    correspond to different eigenvalue configurations.
  \item \textbf{Markov chains.}  The long-run behaviour of a stochastic
    process is governed by the eigenvalue~$1$ and its eigenvector (the
    steady-state distribution).
  \item \textbf{Vibrations and resonance.}  The natural frequencies of a
    vibrating string, membrane, or mechanical system are the square
    roots of the eigenvalues of the associated stiffness matrix.
  \item \textbf{Principal component analysis (PCA).}  In statistics and
    data science, the principal components are the eigenvectors of the
    covariance matrix, and the eigenvalues measure the variance along
    each component.
\end{itemize}

\begin{center}
\begin{tikzpicture}[>=Stealth, scale=1.4]
  % --- Eigenvector illustration: stretching and flipping ---
  % Original basis vectors
  \fill[blue!10] (-0.3,-0.3) rectangle (2.5,2.5);
  \draw[->, thick, blue!70!black] (0,0) -- (2,1)
    node[right] {$\vec{v}$};
  \draw[->, thick, red!70!black] (0,0) -- (0.5,2)
    node[above] {$\vec{w}$};

  % Arrow for the map
  \draw[->, thick, dashed, gray] (3,1) to[bend left=15]
    node[midway,above] {\small $f$} (5,1);

  % Images
  \begin{scope}[shift={(6,0)}]
    \fill[blue!10] (-1.3,-0.8) rectangle (4.3,2.5);
    % Eigenvector stretched by 2
    \draw[->, thick, blue!70!black] (0,0) -- (4,2)
      node[right] {$f(\vec{v}) = 2\vec{v}$};
    % Eigenvector flipped (eigenvalue -1)
    \draw[->, thick, red!70!black] (0,0) -- (-0.5,-2)
      node[below] {$f(\vec{w}) = -\vec{w}$};
    \draw[->, thick, red!70!black, dashed, opacity=0.35] (0,0) -- (0.5,2)
      node[above, opacity=0.5] {\small (original $\vec{w}$)};
  \end{scope}

  % Labels
  \node[below] at (1.2,-0.5) {\small Eigenvector $\vec{v}$: eigenvalue $\lambda=2$};
  \node[below] at (7.5,-0.5) {\small Eigenvector $\vec{w}$: eigenvalue $\lambda=-1$};
\end{tikzpicture}
\captionof{figure}{An endomorphism acts on eigenvectors by scaling:
  $\vec{v}$ is stretched by factor~$2$, and $\vec{w}$ is flipped
  (eigenvalue~$-1$).}
\label{fig:eigenvector-intro}
\end{center}

Throughout this chapter, $\K$ denotes a field (typically $\R$ or~$\C$),
$E$ denotes a finite-dimensional $\K$-vector space with $\dim E = n$,
and $f\in\End(E)$ is an endomorphism of~$E$.  We shall use results
from 
ef{ch:linear-maps} freely, especially the notions of kernel,
image, determinant, and matrix representation.


% ============================================================
\section{Eigenvalues and eigenvectors}\label{sec:eigenvalues-def}
% ============================================================

\begin{definition}{Eigenvalue and eigenvector of an endomorphism}%
{def:eigenvalue-endomorphism}
\index{eigenvalue}\index{eigenvector}
Let $f\in\End(E)$.  A scalar $\lambda\in\K$ is called an
\emph{eigenvalue} of~$f$ if there exists a nonzero vector $v\in E$
such that
\[
  f(v) = \lambda v.
\]
Any such nonzero vector~$v$ is called an \emph{eigenvector} of~$f$
associated with (or belonging to) the eigenvalue~$\lambda$.
\end{definition}

\begin{remark}{The zero vector is never an eigenvector}{rem:zero-not-eigen}
By convention, the zero vector is \emph{not} an eigenvector, even
though $f(\vzero) = \lambda\vzero$ holds for every~$\lambda$.  This
convention ensures that the set of eigenvectors associated with a
given eigenvalue, together with~$\vzero$, forms a subspace.
\end{remark}

\begin{definition}{Eigenvalue and eigenvector of a matrix}%
{def:eigenvalue-matrix}
\index{eigenvalue!of a matrix}\index{eigenvector!of a matrix}
Let $A\in\Mat_n(\K)$.  A scalar $\lambda\in\K$ is an eigenvalue
of~$A$ if there exists a nonzero column vector
$X\in\K^n$ such that $AX = \lambda X$.  Such an~$X$ is an
eigenvector of~$A$ for~$\lambda$.
\end{definition}

\begin{proposition}{Reformulation via the kernel}%
{prop:eigenvalue-kernel}
Let $f\in\End(E)$ and $\lambda\in\K$.  Then $\lambda$ is an
eigenvalue of~$f$ if and only if
\[
  \Ker(f - \lambda\Id) \neq \set{\vzero}.
\]
Equivalently, for $A\in\Mat_n(\K)$, $\lambda$ is an eigenvalue of~$A$
if and only if $\det(A - \lambda I_n) = 0$.
\end{proposition}

\begin{proof}
We have $f(v) = \lambda v$ if and only if $(f - \lambda\Id)(v) = \vzero$,
i.e.\ $v\in\Ker(f - \lambda\Id)$.  A nonzero such~$v$ exists precisely
when $\Ker(f - \lambda\Id) \neq \set{\vzero}$, which is equivalent to
$f - \lambda\Id$ being non-injective, hence non-invertible (in finite
dimension).  For a matrix~$A$ this means $\det(A - \lambda I_n) = 0$.
\end{proof}

\begin{example}{Eigenvalues of a $2\times 2$ matrix}{ex:eigen-2x2}
Let $A = \begin{pmatrix} 3 & 1 \\ 0 & 2 \end{pmatrix}$.  We solve
$\det(A - \lambda I_2) = 0$:
\[
  \det\begin{pmatrix} 3-\lambda & 1 \\ 0 & 2-\lambda \end{pmatrix}
  = (3-\lambda)(2-\lambda) = 0,
\]
so $\lambda_1 = 3$ and $\lambda_2 = 2$.

For $\lambda_1 = 3$: $(A - 3I_2)X = \vzero$ gives
$\begin{pmatrix} 0 & 1 \\ 0 & -1 \end{pmatrix}
\begin{pmatrix} x \\ y \end{pmatrix} = \vzero$,
hence $y = 0$ and $x$ is free.  An eigenvector is
$\begin{psmallmatrix} 1 \\ 0 \end{psmallmatrix}$.

For $\lambda_2 = 2$: $(A - 2I_2)X = \vzero$ gives
$\begin{pmatrix} 1 & 1 \\ 0 & 0 \end{pmatrix}
\begin{pmatrix} x \\ y \end{pmatrix} = \vzero$,
hence $x = -y$.  An eigenvector is
$\begin{psmallmatrix} -1 \\ 1 \end{psmallmatrix}$.
\end{example}

\begin{proposition}{Eigenvectors for distinct eigenvalues are linearly
independent}{prop:distinct-eigenvalues-li}
\index{eigenvector!linear independence}
Let $f\in\End(E)$ and let $\lambda_1,\ldots,\lambda_r$ be pairwise
distinct eigenvalues of~$f$, with associated eigenvectors
$v_1,\ldots,v_r$.  Then $(v_1,\ldots,v_r)$ is linearly independent.
\end{proposition}

\begin{proof}
By induction on~$r$.  The case $r=1$ is clear since $v_1\neq\vzero$.

Suppose the result holds for $r-1$ eigenvalues and consider a
relation $\sum_{i=1}^{r}\alpha_i v_i = \vzero$.  Apply $f$:
\[
  \sum_{i=1}^{r}\alpha_i \lambda_i v_i = \vzero.
\]
Subtract $\lambda_r$ times the original relation:
\[
  \sum_{i=1}^{r-1}\alpha_i(\lambda_i - \lambda_r) v_i = \vzero.
\]
By the induction hypothesis, $\alpha_i(\lambda_i - \lambda_r) = 0$
for all $i = 1,\ldots,r-1$.  Since the eigenvalues are distinct,
$\lambda_i - \lambda_r \neq 0$, so $\alpha_i = 0$ for
$i = 1,\ldots,r-1$.  Substituting back gives $\alpha_r v_r = \vzero$,
whence $\alpha_r = 0$.
\end{proof}

\begin{corollary}{Bound on the number of eigenvalues}%
{cor:eigenvalue-bound}
An endomorphism of an $n$-dimensional space has at most $n$ distinct
eigenvalues.
\end{corollary}

\begin{proof}
If there were $n+1$ distinct eigenvalues, the corresponding
eigenvectors would form a linearly independent family of $n+1$ vectors
in an $n$-dimensional space, contradicting the definition of dimension.
\end{proof}


% ============================================================
\section{Eigenspaces}\label{sec:eigenspaces}
% ============================================================

\begin{definition}{Eigenspace}{def:eigenspace}
\index{eigenspace}
Let $f\in\End(E)$ and $\lambda\in\K$ be an eigenvalue of~$f$.  The
\emph{eigenspace} of~$f$ associated with~$\lambda$ is
\[
  E_\lambda(f) \deff \Ker(f - \lambda\Id)
  = \setcond{v\in E}{f(v) = \lambda v}.
\]
For a matrix $A\in\Mat_n(\K)$, we write $E_\lambda(A) = \Ker(A-\lambda I_n)$.
\end{definition}

\begin{proposition}{Eigenspaces are subspaces}{prop:eigenspace-subspace}
\index{eigenspace!is a subspace}
For every eigenvalue~$\lambda$ of~$f$, the eigenspace $E_\lambda(f)$
is a vector subspace of~$E$ of dimension at least~$1$.
\end{proposition}

\begin{proof}
Since $f - \lambda\Id$ is a linear map, its kernel $\Ker(f - \lambda\Id)$
is a subspace of~$E$ (the kernel of any linear map is a subspace).
It contains at least one nonzero vector (the eigenvector guaranteed by
the definition of eigenvalue), so $\dim E_\lambda(f) \geq 1$.
\end{proof}

\begin{proposition}{Sum of eigenspaces is direct}{prop:eigenspace-direct-sum}
\index{eigenspace!direct sum}
Let $\lambda_1,\ldots,\lambda_r$ be pairwise distinct eigenvalues
of~$f$.  Then the sum $E_{\lambda_1}(f) + \cdots + E_{\lambda_r}(f)$
is a direct sum:
\[
  E_{\lambda_1}(f) + \cdots + E_{\lambda_r}(f)
  = E_{\lambda_1}(f) \oplus \cdots \oplus E_{\lambda_r}(f).
\]
\end{proposition}

\begin{proof}
Suppose $v_1 + \cdots + v_r = \vzero$ with $v_i\in E_{\lambda_i}(f)$.
We must show each $v_i = \vzero$.  If some $v_i \neq \vzero$, let
$S = \setcond{i}{v_i \neq \vzero}$.  Then the nonzero vectors
$(v_i)_{i\in S}$ are eigenvectors for pairwise distinct eigenvalues,
hence linearly independent by

ef{prop:prop:distinct-eigenvalues-li}.  But
$\sum_{i\in S} v_i = \vzero$ (the terms with $v_i = \vzero$ contribute
nothing), contradicting linear independence.  Therefore $S = \emptyset$
and all $v_i = \vzero$.
\end{proof}


% ============================================================
\section{Characteristic polynomial}\label{sec:char-poly}
% ============================================================

\begin{definition}{Characteristic polynomial of a matrix}%
{def:char-poly-matrix}
\index{characteristic polynomial!of a matrix}
Let $A\in\Mat_n(\K)$.  The \emph{characteristic polynomial} of~$A$ is
\[
  \chi_A(\lambda) \deff \det(A - \lambda I_n) \in \K[\lambda].
\]
It is a polynomial of degree~$n$ in~$\lambda$.
\end{definition}

\begin{remark}{Sign convention}{rem:sign-convention}
Some authors define $\chi_A(\lambda) = \det(\lambda I_n - A)$
instead, which differs from our convention by a factor of~$(-1)^n$.
With our convention, the leading term is $(-1)^n\lambda^n$; with the
other convention it is $\lambda^n$.  Both conventions are standard.
The eigenvalues (roots of~$\chi_A$) are the same either way.
\end{remark}

\begin{proposition}{Properties of the characteristic polynomial}%
{prop:char-poly-properties}
Let $A\in\Mat_n(\K)$ with characteristic polynomial
$\chi_A(\lambda) = (-1)^n\lambda^n + c_{n-1}\lambda^{n-1} + \cdots + c_0$.
Then:
\begin{enumerate}[label=(\roman*)]
  \item $\deg\chi_A = n$ and the leading coefficient is $(-1)^n$.
  \item $c_{n-1} = (-1)^{n-1}\tr(A)$.
  \item $c_0 = \chi_A(0) = \det(A)$.
  \item The eigenvalues of~$A$ are exactly the roots of~$\chi_A$
    in~$\K$.
  \item Similar matrices have the same characteristic polynomial:
    if $B = \inv{P}AP$ then $\chi_B = \chi_A$.
\end{enumerate}
\end{proposition}

\begin{proof}
(i) The determinant of $A - \lambda I_n$ is a sum over permutations
$\sigma\in S_n$ of products $\prod_{i=1}^n (A - \lambda I_n)_{i,\sigma(i)}$.
The only permutation contributing a term of degree~$n$ is the identity
$\sigma = \Id$, which gives $\prod_{i=1}^n (a_{ii} - \lambda)$.  The
leading term is $(-\lambda)^n = (-1)^n\lambda^n$.

(ii) In the expansion of $\prod_{i=1}^n(a_{ii} - \lambda)$, the coefficient
of $\lambda^{n-1}$ is $(-1)^{n-1}\sum_{i=1}^n a_{ii} = (-1)^{n-1}\tr(A)$.
No other permutation contributes a term of degree~$n-1$ (such a
permutation would need at least two non-diagonal factors, each of degree~$0$
in~$\lambda$, leaving degree at most $n-2$).

(iii) $\chi_A(0) = \det(A - 0\cdot I_n) = \det(A)$.

(iv) $\lambda$ is an eigenvalue if and only if $\det(A - \lambda I_n) = 0$,
i.e.\ $\chi_A(\lambda) = 0$.

(v) $\chi_B(\lambda) = \det(\inv{P}AP - \lambda I_n)
= \det(\inv{P}(A - \lambda I_n)P)
= \det(\inv{P})\det(A-\lambda I_n)\det(P)
= \chi_A(\lambda)$.
\end{proof}

\begin{definition}{Characteristic polynomial of an endomorphism}%
{def:char-poly-endo}
\index{characteristic polynomial!of an endomorphism}
Let $f\in\End(E)$ and let $A = \Mat_\cB(f)$ be the matrix of~$f$ in
some basis~$\cB$ of~$E$.  The \emph{characteristic polynomial}
of~$f$ is
\[
  \chi_f(\lambda) \deff \chi_A(\lambda) = \det(A - \lambda I_n).
\]
By property~(v) above, this is independent of the choice of
basis~$\cB$.
\end{definition}

\begin{example}{Characteristic polynomial of a $2\times 2$ matrix}%
{ex:char-poly-2x2}
Let $A = \begin{pmatrix} a & b \\ c & d \end{pmatrix}$.  Then
\[
  \chi_A(\lambda) = \det\begin{pmatrix} a-\lambda & b \\ c & d-\lambda
  \end{pmatrix}
  = (a-\lambda)(d-\lambda) - bc
  = \lambda^2 - (a+d)\lambda + (ad-bc)
  = \lambda^2 - \tr(A)\lambda + \det(A).
\]
\end{example}

\begin{example}{Characteristic polynomial of a $3\times 3$ matrix}%
{ex:char-poly-3x3}
Let $A = \begin{pmatrix} 2 & 1 & 0 \\ 0 & 3 & 1 \\ 0 & 0 & 2
\end{pmatrix}$.  Then
\begin{align*}
  \chi_A(\lambda) &= \det\begin{pmatrix} 2-\lambda & 1 & 0 \\
  0 & 3-\lambda & 1 \\ 0 & 0 & 2-\lambda \end{pmatrix} \\
  &= (2-\lambda)(3-\lambda)(2-\lambda) \\
  &= (2-\lambda)^2(3-\lambda)
  = -(\lambda-2)^2(\lambda-3).
\end{align*}
The eigenvalues are $\lambda_1 = 2$ (a double root) and $\lambda_2 = 3$
(a simple root).
\end{example}


% ============================================================
\section{The Cayley--Hamilton theorem}\label{sec:cayley-hamilton}
% ============================================================

\begin{theorem}{Cayley--Hamilton}{thm:cayley-hamilton}
\index{Cayley--Hamilton theorem}
Let $A\in\Mat_n(\K)$ with characteristic polynomial~$\chi_A$.  Then
\[
  \chi_A(A) = 0_{n\times n}.
\]
That is, every square matrix satisfies its own characteristic
polynomial.  Equivalently, if $f\in\End(E)$, then $\chi_f(f) = 0$
(the zero endomorphism).
\end{theorem}

\begin{proof}
Write $B(\lambda) \deff \operatorname{adj}(A - \lambda I_n)$ for the
classical adjoint (adjugate) of $A - \lambda I_n$.  By the adjugate
formula for determinants,
\[
  (A - \lambda I_n)\, B(\lambda) = \det(A - \lambda I_n)\, I_n
  = \chi_A(\lambda)\, I_n.
\]
Each entry of $B(\lambda)$ is a cofactor of $A - \lambda I_n$, hence a
polynomial in~$\lambda$ of degree at most $n-1$.  We can therefore
write
\[
  B(\lambda) = B_0 + B_1\lambda + B_2\lambda^2 + \cdots
    + B_{n-1}\lambda^{n-1},
\]
where $B_0,B_1,\ldots,B_{n-1}\in\Mat_n(\K)$ are constant matrices.
Similarly, write the characteristic polynomial as
\[
  \chi_A(\lambda) = c_0 + c_1\lambda + c_2\lambda^2 + \cdots
    + c_n\lambda^n,
\]
where $c_n = (-1)^n$.

Now expand the identity $(A - \lambda I_n)\,B(\lambda) = \chi_A(\lambda)\,I_n$:
\begin{align*}
  A B_0 &= c_0 I_n, \\
  A B_1 - B_0 &= c_1 I_n, \\
  A B_2 - B_1 &= c_2 I_n, \\
  &\;\;\vdots \\
  A B_{n-1} - B_{n-2} &= c_{n-1} I_n, \\
  -B_{n-1} &= c_n I_n.
\end{align*}
Multiply the $k$-th equation (from $k=0$ to $k=n$) on the left by
$A^k$ and sum:
\begin{align*}
  &\phantom{{}={}} A^0(AB_0) + A^1(AB_1 - B_0) + A^2(AB_2 - B_1)
    + \cdots + A^n(-B_{n-1}) \\
  &= c_0 I_n + c_1 A + c_2 A^2 + \cdots + c_n A^n \\
  &= \chi_A(A).
\end{align*}
On the left-hand side, all terms cancel telescopically:
\[
  AB_0 + A^2 B_1 - AB_0 + A^3 B_2 - A^2 B_1 + \cdots - A^n B_{n-1}
  = 0_{n\times n}.
\]
Therefore $\chi_A(A) = 0_{n\times n}$.
\end{proof}

\begin{remark}{Warning about ``substituting $A$ for $\lambda$''}%
{rem:cayley-hamilton-warning}
One cannot prove Cayley--Hamilton by simply writing
``$\chi_A(A) = \det(A - A\cdot I_n) = \det(0) = 0$.''
This is \emph{fallacious}: $\chi_A(\lambda)$ is a polynomial
in the \emph{scalar}~$\lambda$, while the substitution $\lambda = A$
places a \emph{matrix} inside a determinant in an invalid way.  The
correct proof, as given above, requires careful bookkeeping with matrix
polynomials.
\end{remark}

\begin{example}{Cayley--Hamilton for a $2\times 2$ matrix}{ex:cayley-hamilton-2x2}
Take $A = \begin{pmatrix} 1 & 2 \\ 3 & 4 \end{pmatrix}$.  Then
$\chi_A(\lambda) = \lambda^2 - 5\lambda - 2$.  We verify:
\[
  A^2 - 5A - 2I_2
  = \begin{pmatrix}7 & 10 \\ 15 & 22\end{pmatrix}
  - \begin{pmatrix}5 & 10 \\ 15 & 20\end{pmatrix}
  - \begin{pmatrix}2 & 0 \\ 0 & 2\end{pmatrix}
  = \begin{pmatrix}0 & 0 \\ 0 & 0\end{pmatrix}.\qedhere
\]
\end{example}


% ============================================================
\section{Spectrum of an endomorphism}\label{sec:spectrum}
% ============================================================

\begin{definition}{Spectrum}{def:spectrum}
\index{spectrum}
The \emph{spectrum} of an endomorphism $f\in\End(E)$ (or of a matrix
$A\in\Mat_n(\K)$) is the set of all its eigenvalues:
\[
  \Spec(f) \deff \setcond{\lambda\in\K}{\lambda \text{ is an eigenvalue of } f}.
\]
\end{definition}

\begin{remark}{Dependence on the base field}{rem:spectrum-field}
The spectrum depends on the base field.  For example, the rotation
matrix $R_{\pi/2} = \begin{psmallmatrix} 0 & -1 \\ 1 & 0
\end{psmallmatrix}$ has $\chi_{R_{\pi/2}}(\lambda) = \lambda^2 + 1$,
which has no real roots, so $\Spec_\R(R_{\pi/2}) = \emptyset$.
Over~$\C$, $\Spec_\C(R_{\pi/2}) = \set{i,-i}$.
\end{remark}


% ============================================================
\section{Algebraic and geometric multiplicity}\label{sec:multiplicities}
% ============================================================

\begin{definition}{Algebraic multiplicity}{def:alg-mult}
\index{multiplicity!algebraic}
Let $\lambda$ be an eigenvalue of~$f$.  The \emph{algebraic
multiplicity} of~$\lambda$, denoted $m_a(\lambda)$, is the
multiplicity of~$\lambda$ as a root of the characteristic
polynomial~$\chi_f$.  In other words, $m_a(\lambda)$ is the largest
integer $k$ such that $(\lambda_0 - \lambda)^k$ divides
$\chi_f(\lambda_0)$, where we view $\chi_f$ as a polynomial in the
variable~$\lambda_0$.
\end{definition}

\begin{definition}{Geometric multiplicity}{def:geom-mult}
\index{multiplicity!geometric}
The \emph{geometric multiplicity} of an eigenvalue~$\lambda$,
denoted $m_g(\lambda)$, is the dimension of the eigenspace:
\[
  m_g(\lambda) \deff \dim E_\lambda(f) = \dim\Ker(f - \lambda\Id).
\]
\end{definition}

\begin{theorem}{Inequality between multiplicities}%
{thm:mult-inequality}
\index{multiplicity!inequality}
For every eigenvalue~$\lambda$ of~$f$,
\[
  1 \leq m_g(\lambda) \leq m_a(\lambda).
\]
\end{theorem}

\begin{proof}
The inequality $m_g(\lambda) \geq 1$ holds because $E_\lambda(f)$
contains at least one nonzero eigenvector.

For the inequality $m_g(\lambda) \leq m_a(\lambda)$, let
$d = m_g(\lambda) = \dim E_\lambda(f)$ and choose a basis
$(v_1,\ldots,v_d)$ of~$E_\lambda(f)$.  Extend it to a basis
$\cB = (v_1,\ldots,v_d,v_{d+1},\ldots,v_n)$ of~$E$.  In this basis,
the matrix of~$f$ has the block form
\[
  \Mat_\cB(f) = \begin{pmatrix} \lambda I_d & C \\ 0 & D \end{pmatrix},
\]
where $C\in\Mat_{d,n-d}(\K)$ and $D\in\Mat_{n-d}(\K)$.  This is
because $f(v_i) = \lambda v_i$ for $i = 1,\ldots,d$.  Then
\[
  \chi_f(\mu) = \det\begin{pmatrix} \lambda I_d - \mu I_d & C \\
    0 & D - \mu I_{n-d} \end{pmatrix}
  = (\lambda - \mu)^d \det(D - \mu I_{n-d}).
\]
Therefore $(\lambda - \mu)^d$ divides $\chi_f(\mu)$, which means
$m_a(\lambda) \geq d = m_g(\lambda)$.
\end{proof}

\begin{example}{Strict inequality between multiplicities}%
{ex:strict-mult-inequality}
The matrix $A = \begin{pmatrix} 2 & 1 \\ 0 & 2 \end{pmatrix}$ has
$\chi_A(\lambda) = (2-\lambda)^2$, so $\lambda = 2$ has algebraic
multiplicity $m_a(2) = 2$.  However,
$E_2(A) = \Ker\begin{psmallmatrix} 0 & 1 \\ 0 & 0 \end{psmallmatrix}
= \Span\!\left(\begin{psmallmatrix} 1 \\ 0 \end{psmallmatrix}\right)$,
which has dimension~$1$.  So $m_g(2) = 1 < 2 = m_a(2)$.
\end{example}


% ============================================================
\section{Diagonalizability}\label{sec:diagonalizability}
% ============================================================

\begin{definition}{Diagonalizable endomorphism}{def:diagonalizable-endo}
\index{diagonalizable!endomorphism}
An endomorphism $f\in\End(E)$ is \emph{diagonalizable} if there
exists a basis of~$E$ consisting entirely of eigenvectors of~$f$.
Equivalently, $f$ is diagonalizable if its matrix in some basis is
diagonal.
\end{definition}

\begin{definition}{Diagonalizable matrix}{def:diagonalizable-matrix}
\index{diagonalizable!matrix}
A matrix $A\in\Mat_n(\K)$ is \emph{diagonalizable} if it is similar
to a diagonal matrix, i.e.\ if there exists an invertible matrix
$P\in\GL_n(\K)$ and a diagonal matrix
$D = \diag(\lambda_1,\ldots,\lambda_n)$ such that
\[
  A = P D \inv{P}.
\]
The columns of~$P$ are eigenvectors of~$A$, and the diagonal entries
of~$D$ are the corresponding eigenvalues.
\end{definition}

\begin{center}
\begin{tikzpicture}[>=Stealth, node distance=3.5cm]
  \node (E1) {$E$};
  \node[right of=E1, node distance=4cm] (E2) {$E$};
  \node[below of=E1, node distance=2.5cm] (K1) {$\K^n$};
  \node[below of=E2, node distance=2.5cm] (K2) {$\K^n$};

  \draw[->] (E1) -- node[above] {$f$} (E2);
  \draw[->] (K1) -- node[below] {$D = \diag(\lambda_1,\ldots,\lambda_n)$} (K2);
  \draw[->] (E1) -- node[left] {$[\cdot]_\cB$} (K1);
  \draw[->] (E2) -- node[right] {$[\cdot]_\cB$} (K2);

  \node[below of=K1, node distance=1.5cm, text width=8cm, align=center]
    {\small In the eigenbasis $\cB$, the map $f$ acts as
     multiplication by $\lambda_i$ on each coordinate axis.};
\end{tikzpicture}
\captionof{figure}{Diagonalization: in the eigenbasis $\cB$, the
  endomorphism~$f$ is represented by the diagonal matrix~$D$.}
\label{fig:diag-diagram}
\end{center}

\begin{theorem}{Diagonalization criterion}{thm:diag-criterion}
\index{diagonalizable!criterion}
Let $f\in\End(E)$ with $\dim E = n$, and let
$\Spec(f) = \set{\lambda_1,\ldots,\lambda_r}$.  The following
conditions are equivalent:
\begin{enumerate}[label=(\roman*)]
  \item $f$ is diagonalizable.
  \item $E = E_{\lambda_1}(f) \oplus E_{\lambda_2}(f) \oplus \cdots
    \oplus E_{\lambda_r}(f)$.
  \item $\sum_{i=1}^{r} \dim E_{\lambda_i}(f) = n$.
  \item The characteristic polynomial $\chi_f$ splits over~$\K$
    (i.e.\ has all its roots in~$\K$), \emph{and}
    $m_g(\lambda_i) = m_a(\lambda_i)$ for every eigenvalue
    $\lambda_i$.
\end{enumerate}
\end{theorem}

\begin{proof}
\textbf{(i) $\Rightarrow$ (ii).}\;
If $f$ is diagonalizable, there is a basis $(v_1,\ldots,v_n)$ of~$E$
with $f(v_j) = \mu_j v_j$ for some $\mu_j\in\K$.  Each $v_j$ belongs
to $E_{\mu_j}(f)$.  Since the $v_j$ Span~$E$, we have
$E = \sum_{i=1}^r E_{\lambda_i}(f)$.  This sum is direct by

ef{prop:prop:eigenspace-direct-sum}.

\textbf{(ii) $\Rightarrow$ (iii).}\;
If the direct sum decomposition holds, then
$n = \dim E = \sum_{i=1}^r \dim E_{\lambda_i}(f)$ by the dimension
formula for direct sums.

\textbf{(iii) $\Rightarrow$ (iv).}\;
From $\sum m_g(\lambda_i) = n$ and the inequality
$m_g(\lambda_i) \leq m_a(\lambda_i)$ for each~$i$, together with
$\sum m_a(\lambda_i) \leq n$ (the sum of algebraic multiplicities is
at most~$\deg\chi_f = n$), we get
\[
  n = \sum_{i=1}^r m_g(\lambda_i) \leq \sum_{i=1}^r m_a(\lambda_i) \leq n.
\]
All inequalities are equalities.  In particular,
$m_g(\lambda_i) = m_a(\lambda_i)$ for each~$i$, and
$\sum m_a(\lambda_i) = n$, which means $\chi_f$ splits over~$\K$
(every root is in~$\K$ and the multiplicities account for the full
degree~$n$).

\textbf{(iv) $\Rightarrow$ (i).}\;
If $\chi_f$ splits and $m_g(\lambda_i) = m_a(\lambda_i)$ for each~$i$,
then
\[
  \sum_{i=1}^r \dim E_{\lambda_i}(f)
  = \sum_{i=1}^r m_g(\lambda_i)
  = \sum_{i=1}^r m_a(\lambda_i) = n.
\]
Since the sum of eigenspaces is direct (
ef{prop:prop:eigenspace-direct-sum}),
we have $E = \bigoplus_{i=1}^r E_{\lambda_i}(f)$.  Choosing a basis
for each $E_{\lambda_i}(f)$ and concatenating gives a basis of~$E$
consisting of eigenvectors, so $f$ is diagonalizable.
\end{proof}

\begin{corollary}{Endomorphism with $n$ distinct eigenvalues}%
{cor:n-distinct-eigenvalues}
If $f\in\End(E)$ has $n = \dim E$ distinct eigenvalues, then $f$ is
diagonalizable.
\end{corollary}

\begin{proof}
Each eigenvalue has algebraic multiplicity~$1$ (since there are $n$
eigenvalues and $\deg\chi_f = n$), and the geometric multiplicity
satisfies $1 \leq m_g(\lambda) \leq m_a(\lambda) = 1$.  Thus
$m_g = m_a$ for all eigenvalues and $\chi_f$ splits, so $f$ is
diagonalizable by 
ef{thm:thm:diag-criterion}.
\end{proof}

\begin{remark}{The converse is false}{rem:diag-distinct-converse}
Having $n$ distinct eigenvalues is sufficient but not necessary for
diagonalizability.  For instance, the identity matrix $I_n$ is
diagonal (hence diagonalizable) but has a single eigenvalue
$\lambda = 1$ with multiplicity~$n$.
\end{remark}


% ============================================================
\section{The diagonalization algorithm}\label{sec:diag-algorithm}
% ============================================================

Given $A\in\Mat_n(\K)$, the following procedure determines whether
$A$ is diagonalizable and, if so, produces $P$ and $D$ such that
$A = P D \inv{P}$.

\begin{enumerate}[label=\textbf{Step \arabic*.}]
  \item \textbf{Compute the characteristic polynomial}
    $\chi_A(\lambda) = \det(A - \lambda I_n)$.
  \item \textbf{Find the eigenvalues} by solving $\chi_A(\lambda) = 0$.
    If $\chi_A$ does not split over~$\K$, then $A$ is \emph{not}
    diagonalizable over~$\K$.
  \item \textbf{For each eigenvalue~$\lambda_i$}, compute the eigenspace
    $E_{\lambda_i} = \Ker(A - \lambda_i I_n)$ by solving the
    homogeneous system $(A - \lambda_i I_n)X = \vzero$.
  \item \textbf{Check the multiplicity condition:}
    $\dim E_{\lambda_i} = m_a(\lambda_i)$ for each~$i$.
    If this fails for any eigenvalue, $A$ is \emph{not} diagonalizable.
  \item \textbf{Form} $P$ by placing the eigenvectors (bases of each
    eigenspace) as columns, and $D = \diag(\lambda_1,\ldots,\lambda_1,
    \lambda_2,\ldots,\lambda_2,\ldots)$ with each $\lambda_i$ repeated
    $m_a(\lambda_i)$ times.
\end{enumerate}

\begin{example}{Full diagonalization of a $3\times 3$ matrix}%
{ex:diag-3x3}
\index{diagonalization!example}
Diagonalize $A = \begin{pmatrix} 4 & 0 & 1 \\ 2 & 3 & 2 \\ 1 & 0 & 4
\end{pmatrix}$.

\textbf{Step 1.}  The characteristic polynomial:
\begin{align*}
  \chi_A(\lambda) &= \det\begin{pmatrix} 4-\lambda & 0 & 1 \\
    2 & 3-\lambda & 2 \\ 1 & 0 & 4-\lambda \end{pmatrix} \\
  &= (4-\lambda)\bigl[(3-\lambda)(4-\lambda) - 0\bigr]
    - 0 + 1\bigl[0 - (3-\lambda)\bigr] \\
  &= (4-\lambda)(3-\lambda)(4-\lambda) - (3-\lambda) \\
  &= (3-\lambda)\bigl[(4-\lambda)^2 - 1\bigr] \\
  &= (3-\lambda)(16 - 8\lambda + \lambda^2 - 1) \\
  &= (3-\lambda)(\lambda^2 - 8\lambda + 15) \\
  &= (3-\lambda)(\lambda - 3)(\lambda - 5) \\
  &= -(3-\lambda)^2(\lambda - 5) = -(\lambda-3)^2(\lambda-5).
\end{align*}

\textbf{Step 2.}  Eigenvalues: $\lambda_1 = 3$ (algebraic
multiplicity~$2$) and $\lambda_2 = 5$ (algebraic multiplicity~$1$).

\textbf{Step 3.}  Eigenspaces:

For $\lambda_1 = 3$: solve $(A - 3I_3)X = \vzero$:
\[
  A - 3I_3 = \begin{pmatrix} 1 & 0 & 1 \\ 2 & 0 & 2 \\ 1 & 0 & 1
  \end{pmatrix} \longrightarrow \begin{pmatrix} 1 & 0 & 1 \\ 0 & 0 & 0
  \\ 0 & 0 & 0 \end{pmatrix}.
\]
So $x_1 = -x_3$, $x_2$ free, $x_3$ free.  A basis of $E_3(A)$ is
\[
  \left\{ \begin{pmatrix} 0\\1\\0 \end{pmatrix},\;
  \begin{pmatrix} -1\\0\\1 \end{pmatrix} \right\},
  \quad \dim E_3(A) = 2 = m_a(3).
\]

For $\lambda_2 = 5$: solve $(A - 5I_3)X = \vzero$:
\[
  A - 5I_3 = \begin{pmatrix} -1 & 0 & 1 \\ 2 & -2 & 2 \\ 1 & 0 & -1
  \end{pmatrix} \longrightarrow \begin{pmatrix} 1 & 0 & -1 \\ 0 & 1 & -2
  \\ 0 & 0 & 0 \end{pmatrix}.
\]
So $x_1 = x_3$, $x_2 = 2x_3$, $x_3$ free.  A basis of $E_5(A)$ is
$\left\{ \begin{psmallmatrix} 1\\2\\1 \end{psmallmatrix} \right\}$,
and $\dim E_5(A) = 1 = m_a(5)$.

\textbf{Step 4.}  Since $m_g = m_a$ for both eigenvalues, $A$ is
diagonalizable.

\textbf{Step 5.}  We obtain
\[
  P = \begin{pmatrix} 0 & -1 & 1 \\ 1 & 0 & 2 \\ 0 & 1 & 1
  \end{pmatrix}, \qquad
  D = \begin{pmatrix} 3 & 0 & 0 \\ 0 & 3 & 0 \\ 0 & 0 & 5
  \end{pmatrix}, \qquad
  A = P D \inv{P}.
\]
\end{example}

\begin{example}{A non-diagonalizable matrix}{ex:non-diag}
\index{diagonalizable!non-example}
Consider $A = \begin{pmatrix} 2 & 1 \\ 0 & 2 \end{pmatrix}$.  Then
$\chi_A(\lambda) = (2-\lambda)^2$, so $\lambda = 2$ is the only
eigenvalue with $m_a(2) = 2$.  But
$E_2(A) = \Ker\begin{psmallmatrix} 0 & 1 \\ 0 & 0 \end{psmallmatrix}
= \Span\!\left(\begin{psmallmatrix} 1 \\ 0 \end{psmallmatrix}\right)$,
so $m_g(2) = 1 < 2 = m_a(2)$.  By

ef{thm:thm:diag-criterion}, $A$ is \emph{not} diagonalizable.
\end{example}


% ============================================================
\section{Trigonalization}\label{sec:trigonalization}
% ============================================================

\begin{definition}{Trigonalizable endomorphism}{def:trigonalizable}
\index{trigonalizable}\index{triangularizable|see{trigonalizable}}
An endomorphism $f\in\End(E)$ is \emph{trigonalizable} (over~$\K$)
if there exists a basis of~$E$ in which the matrix of~$f$ is upper
triangular.  A matrix $A\in\Mat_n(\K)$ is trigonalizable if it is
similar to an upper triangular matrix.
\end{definition}

\begin{theorem}{Trigonalization over an algebraically closed field}%
{thm:trigonalization}
\index{trigonalization theorem}
Let $f\in\End(E)$ with $\dim E = n$.  Then $f$ is trigonalizable
over~$\K$ if and only if $\chi_f$ splits over~$\K$.

In particular, every endomorphism of a finite-dimensional complex
vector space is trigonalizable (since $\C$ is algebraically closed).
\end{theorem}

\begin{proof}
$(\Leftarrow)$\; We proceed by induction on $n = \dim E$.
The case $n = 1$ is trivial.

Since $\chi_f$ splits, $f$ has at least one eigenvalue
$\lambda_1\in\K$.  Let $v_1$ be a corresponding eigenvector and set
$F = \Span(v_1)$.  The quotient space $\bar{E} = E/F$ has
dimension~$n-1$, and $f$ induces a well-defined endomorphism
$\bar{f}\colon\bar{E}\to\bar{E}$ (since $F$ is $f$-invariant).
Moreover, $\chi_{\bar{f}}$ divides $\chi_f$ (up to a constant),
so $\chi_{\bar{f}}$ also splits over~$\K$.

By the induction hypothesis, $\bar{f}$ is trigonalizable: there
exists a basis $(\bar{v}_2,\ldots,\bar{v}_n)$ of~$\bar{E}$ in which
$\bar{f}$ is upper triangular.  Choosing representatives
$v_2,\ldots,v_n\in E$ with $\bar{v}_i = v_i + F$, the family
$(v_1,v_2,\ldots,v_n)$ is a basis of~$E$ and the matrix of~$f$ in
this basis is upper triangular with $\lambda_1$ in the $(1,1)$
position.

$(\Rightarrow)$\; If $f$ is trigonalizable with upper triangular
matrix $T$, then $\chi_f(\lambda) = \chi_T(\lambda) = \prod_{i=1}^n
(t_{ii} - \lambda)$, which splits over~$\K$.
\end{proof}

\begin{corollary}{Trigonalization over $\C$}{cor:trig-C}
Every matrix $A\in\Mat_n(\C)$ is trigonalizable over~$\C$.
\end{corollary}


% ============================================================
\section{Minimal polynomial}\label{sec:minimal-poly}
% ============================================================

\begin{definition}{Annihilating polynomial}{def:annihilating-poly}
\index{annihilating polynomial}
A polynomial $p\in\K[\lambda]$ is an \emph{annihilating polynomial}
(or \emph{annihilator}) of $f\in\End(E)$ if $p(f) = 0$ (the zero
endomorphism).  Similarly, $p$ annihilates $A\in\Mat_n(\K)$ if
$p(A) = 0_{n\times n}$.
\end{definition}

By the Cayley--Hamilton theorem (
ef{thm:thm:cayley-hamilton}), $\chi_f$
is always an annihilating polynomial of~$f$.  The minimal polynomial
is the ``smallest'' such polynomial.

\begin{definition}{Minimal polynomial}{def:minimal-poly}
\index{minimal polynomial}
The \emph{minimal polynomial} of $f\in\End(E)$, denoted $\mu_f$, is
the unique monic polynomial of smallest degree that annihilates~$f$.
\end{definition}

\begin{proposition}{Existence and uniqueness of the minimal polynomial}%
{prop:minimal-poly-unique}
The minimal polynomial $\mu_f$ exists and is unique.  Moreover, $\mu_f$
divides every annihilating polynomial of~$f$.
\end{proposition}

\begin{proof}
\textit{Existence.}\; The set of annihilating polynomials is nonempty
(it contains $\chi_f$ by Cayley--Hamilton) and every polynomial has a
nonnegative degree, so there is one of minimal degree.  Normalising
to be monic gives~$\mu_f$.

\textit{Divides every annihilator.}\; Let $p$ be an annihilating
polynomial.  Perform Euclidean division: $p = q\mu_f + r$ with
$\deg r < \deg\mu_f$.  Then $r(f) = p(f) - q(f)\mu_f(f) = 0 - 0 = 0$.
If $r \neq 0$, then $r/\operatorname{lc}(r)$ would be a monic
annihilator of smaller degree than~$\mu_f$, a contradiction.  Hence
$r = 0$ and $\mu_f \mid p$.

\textit{Uniqueness.}\; If $\mu$ and $\mu'$ are both monic annihilating
polynomials of minimal degree, then $\mu \mid \mu'$ and $\mu' \mid \mu$
(by the above), so $\mu = \mu'$ (both being monic).
\end{proof}

\begin{proposition}{Properties of the minimal polynomial}%
{prop:minimal-poly-properties}
Let $f\in\End(E)$.  Then:
\begin{enumerate}[label=(\roman*)]
  \item $\mu_f$ divides $\chi_f$.
  \item $\mu_f$ and $\chi_f$ have the same roots (i.e.\ the same
    eigenvalues), though possibly with different multiplicities.
  \item If $A$ and $B$ are similar matrices, they have the same
    minimal polynomial.
\end{enumerate}
\end{proposition}

\begin{proof}
(i) Cayley--Hamilton gives $\chi_f(f) = 0$, so $\mu_f \mid \chi_f$.

(ii) Since $\mu_f \mid \chi_f$, every root of $\mu_f$ is a root of
$\chi_f$.  Conversely, if $\lambda$ is an eigenvalue of~$f$ with
eigenvector~$v$, then $\mu_f(f)(v) = \mu_f(\lambda)v = \vzero$,
so $\mu_f(\lambda) = 0$ (since $v \neq \vzero$).

(iii) If $B = \inv{P}AP$, then $p(B) = \inv{P}p(A)P$ for any
polynomial~$p$.  Thus $p(A) = 0$ if and only if $p(B) = 0$.
\end{proof}

\begin{theorem}{Diagonalizability via the minimal polynomial}%
{thm:diag-minimal-poly}
\index{diagonalizable!minimal polynomial criterion}
An endomorphism $f\in\End(E)$ is diagonalizable if and only if its
minimal polynomial $\mu_f$ splits into distinct linear factors
over~$\K$:
\[
  \mu_f(\lambda) = (\lambda - \lambda_1)(\lambda - \lambda_2)\cdots
    (\lambda - \lambda_r)
\]
where $\lambda_1,\ldots,\lambda_r$ are pairwise distinct.
\end{theorem}

\begin{proof}
$(\Rightarrow)$\; Suppose $f$ is diagonalizable with eigenvalues
$\lambda_1,\ldots,\lambda_r$.  Set
$p(\lambda) = (\lambda-\lambda_1)\cdots(\lambda-\lambda_r)$.  Since $f$
is diagonalizable, there is a basis of eigenvectors.  For any
eigenvector~$v$ with eigenvalue~$\lambda_i$,
$p(f)(v) = \prod_{j\neq i}(\lambda_i - \lambda_j)\cdot 0 \cdot v = \vzero$.
Wait---more carefully:
$p(f)(v) = (f-\lambda_1\Id)\cdots(f-\lambda_r\Id)(v)$.  Since
$(f - \lambda_i\Id)(v) = \vzero$, the entire product gives~$\vzero$.
As this holds for every eigenvector and eigenvectors Span~$E$, we get
$p(f) = 0$.  So $\mu_f \mid p$.  Since the roots of $\mu_f$ include
all eigenvalues (by 
ef{prop:prop:minimal-poly-properties}(ii)) and
$\mu_f$ divides a product of distinct linear factors, $\mu_f$ itself
is a product of distinct linear factors.

$(\Leftarrow)$\; Suppose $\mu_f(\lambda)=(\lambda-\lambda_1)\cdots
(\lambda-\lambda_r)$ with distinct $\lambda_i$.  We prove
$E = \bigoplus_{i=1}^r E_{\lambda_i}(f)$ by induction on~$r$.

For $r = 1$: $\mu_f(\lambda) = \lambda - \lambda_1$, so
$f = \lambda_1\Id$ and every nonzero vector is an eigenvector.
Thus $E = E_{\lambda_1}(f)$.

For the inductive step, write
$\mu_f = (\lambda - \lambda_r)\,q(\lambda)$ where
$q(\lambda) = \prod_{i=1}^{r-1}(\lambda - \lambda_i)$.  Since
$\gcd(q, \lambda-\lambda_r) = 1$ (the $\lambda_i$ are distinct),
there exist polynomials $a,b\in\K[\lambda]$ with
$a(\lambda)\,q(\lambda) + b(\lambda)(\lambda-\lambda_r) = 1$.
Substituting $f$:
\[
  a(f)\circ q(f) + b(f)\circ(f - \lambda_r\Id) = \Id.
\]
For any $v\in E$, write $v = a(f)(q(f)(v)) + b(f)((f-\lambda_r\Id)(v))$.
Let $u = q(f)(v)$ and $w = (f-\lambda_r\Id)(v)$.  Then
$(f-\lambda_r\Id)(u) = (f-\lambda_r\Id)(q(f)(v)) = \mu_f(f)(v) = \vzero$,
so $u\in E_{\lambda_r}(f)$.  Similarly,
$q(f)(w) = q(f)(f-\lambda_r\Id)(v) = \mu_f(f)(v) = \vzero$, so
$w\in\Ker(q(f))$.

Now $q(f)$ acts on $F = \Ker(q(f))$, and $q$ annihilates $f\restrict{F}$.
So $\mu_{f|_F}$ divides $q = \prod_{i=1}^{r-1}(\lambda-\lambda_i)$.
By the induction hypothesis, $F = \bigoplus_{i=1}^{r-1}E_{\lambda_i}(f|_F)
\subset \bigoplus_{i=1}^{r-1}E_{\lambda_i}(f)$.

Since $v = a(f)(u) + b(f)(w)$ and $u\in E_{\lambda_r}(f)$,
$w\in F\subset\bigoplus_{i=1}^{r-1}E_{\lambda_i}(f)$, we see that
every $v\in E$ lies in $\sum_{i=1}^r E_{\lambda_i}(f)$.  The sum is
direct by 
ef{prop:prop:eigenspace-direct-sum}.
\end{proof}


% ============================================================
\section{Applications}\label{sec:eigen-applications}
% ============================================================

% -------------------------
\subsection{Computing powers of a matrix}\label{subsec:matrix-powers}
% -------------------------

\begin{proposition}{Powers of a diagonalizable matrix}%
{prop:matrix-power}
\index{matrix power}
If $A = P D \inv{P}$ where $D = \diag(\lambda_1,\ldots,\lambda_n)$,
then for every integer $k\geq 0$,
\[
  A^k = P\,\diag(\lambda_1^k,\ldots,\lambda_n^k)\,\inv{P}.
\]
\end{proposition}

\begin{proof}
By induction, $A^k = (PD\inv{P})^k = PD^k\inv{P}$, since
$(PD\inv{P})(PD\inv{P}) = PD(\inv{P}P)D\inv{P} = PD^2\inv{P}$,
and so on.  Moreover, $D^k = \diag(\lambda_1^k,\ldots,\lambda_n^k)$.
\end{proof}

% -------------------------
\subsection{Linear recurrences: the Fibonacci sequence}%
\label{subsec:fibonacci}
% -------------------------

\begin{example}{Fibonacci sequence via diagonalization}{ex:fibonacci}
\index{Fibonacci sequence}\index{linear recurrence}
The Fibonacci sequence is defined by $F_0 = 0$, $F_1 = 1$, and
$F_{n+2} = F_{n+1} + F_n$.  Set
$V_n = \begin{psmallmatrix} F_{n+1} \\ F_n \end{psmallmatrix}$.  Then
$V_{n+1} = A V_n$ where
$A = \begin{psmallmatrix} 1 & 1 \\ 1 & 0 \end{psmallmatrix}$.

The characteristic polynomial is $\chi_A(\lambda) = \lambda^2 - \lambda - 1$,
with roots $\vphi = \frac{1+\sqrt{5}}{2}$ and
$\psi = \frac{1-\sqrt{5}}{2}$.

Eigenvectors:
$v_1 = \begin{psmallmatrix}\vphi\\1\end{psmallmatrix}$ for $\vphi$
and $v_2 = \begin{psmallmatrix}\psi\\1\end{psmallmatrix}$ for $\psi$.
So $P = \begin{psmallmatrix}\vphi & \psi \\ 1 & 1\end{psmallmatrix}$
and
$A = P\begin{psmallmatrix}\vphi & 0 \\ 0 & \psi\end{psmallmatrix}\inv{P}$.

From $V_n = A^n V_0$ and $V_0 = \begin{psmallmatrix}1\\0\end{psmallmatrix}$,
we extract:
\[
  F_n = \frac{\vphi^n - \psi^n}{\vphi - \psi}
  = \frac{1}{\sqrt{5}}\left[\left(\frac{1+\sqrt{5}}{2}\right)^{\!n}
    - \left(\frac{1-\sqrt{5}}{2}\right)^{\!n}\right].
\]
This is \emph{Binet's formula}\index{Binet's formula}.
\end{example}

% -------------------------
\subsection{Markov chains and steady-state distributions}%
\label{subsec:markov}
% -------------------------

\begin{example}{Steady-state of a Markov chain}{ex:markov}
\index{Markov chain}\index{steady-state distribution}
A weather model has two states: Sunny (S) and Rainy (R), with
transition matrix
\[
  T = \begin{pmatrix} 0.8 & 0.4 \\ 0.2 & 0.6 \end{pmatrix},
\]
where column~$j$ gives the transition probabilities from state~$j$.
If $\vec{p}_n$ is the state vector at time~$n$, then
$\vec{p}_{n+1} = T\vec{p}_n$, so $\vec{p}_n = T^n\vec{p}_0$.

The eigenvalues of~$T$ are $\lambda_1 = 1$ and $\lambda_2 = 0.4$.
The eigenvector for $\lambda_1 = 1$ (normalised to sum to~$1$) is
\[
  \vec{\pi} = \begin{pmatrix} 2/3 \\ 1/3 \end{pmatrix}.
\]
Since $|\lambda_2| < 1$, as $n\to\infty$ the component along the
second eigenvector decays, and
$\vec{p}_n \to \vec{\pi}$ regardless of~$\vec{p}_0$.  Thus the
long-run probability of sunny weather is~$2/3$.
\end{example}

% -------------------------
\subsection{Systems of linear ODEs}\label{subsec:odes}
% -------------------------

\begin{example}{System of ODEs via diagonalization}{ex:odes}
\index{differential equation!system}
Consider the system $\vec{x}'(t) = A\vec{x}(t)$ with
$A = \begin{psmallmatrix} -3 & 1 \\ 1 & -3 \end{psmallmatrix}$.
The eigenvalues are $\lambda_1 = -2$ (eigenvector
$v_1 = \begin{psmallmatrix}1\\1\end{psmallmatrix}$) and
$\lambda_2 = -4$ (eigenvector
$v_2 = \begin{psmallmatrix}1\\-1\end{psmallmatrix}$).

Setting $\vec{x}(t) = P\vec{y}(t)$ where
$P = \begin{psmallmatrix}1 & 1 \\ 1 & -1\end{psmallmatrix}$,
the system decouples to $\vec{y}'(t) = D\vec{y}(t)$ with
$D = \diag(-2,-4)$.  The solution is
\[
  \vec{y}(t) = \begin{pmatrix} c_1 e^{-2t} \\ c_2 e^{-4t} \end{pmatrix},
  \qquad
  \vec{x}(t) = c_1 e^{-2t}\begin{pmatrix}1\\1\end{pmatrix}
    + c_2 e^{-4t}\begin{pmatrix}1\\-1\end{pmatrix},
\]
where $c_1,c_2$ are determined by the initial condition
$\vec{x}(0)$.
\end{example}


% ============================================================
\section{Exercises}\label{sec:eigenvalue-exercises}
% ============================================================

\begin{exercise}{Eigenvalues of a $2\times 2$ matrix \basic}%
{exo:eigen-2x2-basic}
Find the eigenvalues and eigenvectors of
$A = \begin{pmatrix} 5 & 4 \\ 1 & 2 \end{pmatrix}$.
Is $A$ diagonalizable?  If so, find $P$ and $D$ such that
$A = PD\inv{P}$.
\end{exercise}

\begin{exercise}{Eigenvalues of a triangular matrix \basic}%
{exo:triangular-eigen}
Let $A\in\Mat_n(\K)$ be upper triangular.  Show that the eigenvalues
of~$A$ are exactly the diagonal entries $a_{11},a_{22},\ldots,a_{nn}$.
\end{exercise}

\begin{exercise}{Eigenvalues of a projection \basic}{exo:proj-eigen}
Let $p\in\End(E)$ be a projection, i.e.\ $p^2 = p$.  Show that the
only possible eigenvalues of~$p$ are $0$ and~$1$, and that $p$ is
diagonalizable.

\noindent\textit{Hint:} What is the minimal polynomial of~$p$?
\end{exercise}

\begin{exercise}{Eigenvalues and trace/determinant \basic}%
{exo:trace-det-eigen}
Let $A\in\Mat_2(\K)$ with eigenvalues $\lambda_1,\lambda_2$.  Show
that $\tr(A) = \lambda_1 + \lambda_2$ and
$\det(A) = \lambda_1\lambda_2$.  Generalise to $n\times n$ matrices.
\end{exercise}

\begin{exercise}{Eigenvalues of $A^2$, $A^{-1}$, $A + cI$ \intermediate}%
{exo:eigen-transformations}
Let $A\in\Mat_n(\K)$ with eigenvalue $\lambda$ and eigenvector~$v$.
\begin{enumerate}[label=(\alph*)]
  \item Show that $\lambda^2$ is an eigenvalue of $A^2$ with the same
    eigenvector~$v$.  Generalise to $A^k$.
  \item If $A$ is invertible, show that $\lambda^{-1}$ is an eigenvalue
    of~$\inv{A}$.
  \item Show that $\lambda + c$ is an eigenvalue of $A + cI$ for
    any $c\in\K$.
\end{enumerate}
\end{exercise}

\begin{exercise}{Diagonalization of a $3\times 3$ matrix \intermediate}%
{exo:diag-3x3}
Diagonalize the matrix
$A = \begin{pmatrix} 1 & 2 & 0 \\ 0 & 3 & 0 \\ 2 & 1 & -1 \end{pmatrix}$.
\end{exercise}

\begin{exercise}{Non-diagonalizable matrix \intermediate}{exo:non-diag}
Show that $A = \begin{pmatrix} 1 & 1 & 0 \\ 0 & 1 & 1 \\
0 & 0 & 1 \end{pmatrix}$ is not diagonalizable over any field.
What is its minimal polynomial?
\end{exercise}

\begin{exercise}{Cayley--Hamilton application \intermediate}%
{exo:cayley-hamilton-app}
Let $A = \begin{pmatrix} 1 & 2 \\ 3 & 4 \end{pmatrix}$.
\begin{enumerate}[label=(\alph*)]
  \item Find the characteristic polynomial $\chi_A$.
  \item Use the Cayley--Hamilton theorem to express $\inv{A}$ as a
    polynomial in~$A$ (i.e.\ in the form $aI + bA$).
  \item Compute $A^5$ using Cayley--Hamilton to reduce to lower
    powers.
\end{enumerate}
\end{exercise}

\begin{exercise}{Powers via diagonalization \intermediate}{exo:powers}
\index{matrix power!via diagonalization}
Let $A = \begin{pmatrix} 2 & 1 \\ 0 & 3 \end{pmatrix}$.
\begin{enumerate}[label=(\alph*)]
  \item Diagonalize $A$.
  \item Compute $A^{100}$.
\end{enumerate}
\end{exercise}

\begin{exercise}{Simultaneous diagonalization \challenging}%
{exo:simul-diag}
\index{simultaneous diagonalization}
Let $A,B\in\Mat_n(\K)$ be diagonalizable matrices that commute
($AB = BA$).  Show that $A$ and $B$ are \emph{simultaneously
diagonalizable}: there exists an invertible $P$ such that both
$\inv{P}AP$ and $\inv{P}BP$ are diagonal.

\noindent\textit{Hint:} Show that eigenspaces of~$A$ are invariant
under~$B$, then diagonalize the restriction of~$B$ to each eigenspace.
\end{exercise}

\begin{exercise}{Computing the Fibonacci sequence \intermediate}%
{exo:fibonacci}
\begin{enumerate}[label=(\alph*)]
  \item Diagonalize $A = \begin{psmallmatrix} 1 & 1 \\ 1 & 0
    \end{psmallmatrix}$ over~$\R$.
  \item Derive Binet's formula:
    $F_n = \dfrac{1}{\sqrt{5}}\!\left[\left(\dfrac{1+\sqrt{5}}{2}
    \right)^{\!n} - \left(\dfrac{1-\sqrt{5}}{2}\right)^{\!n}\right]$.
  \item Show that $F_n$ is the nearest integer to
    $\dfrac{1}{\sqrt{5}}\left(\dfrac{1+\sqrt{5}}{2}\right)^{\!n}$
    for all $n\geq 0$.
\end{enumerate}
\end{exercise}

\begin{exercise}{Minimal polynomial \challenging}{exo:minimal-poly}
For each of the following matrices, find the minimal polynomial and
determine whether the matrix is diagonalizable.
\begin{enumerate}[label=(\alph*)]
  \item $A = \begin{pmatrix} 3 & 0 & 0 \\ 0 & 3 & 0 \\ 0 & 0 & 5
    \end{pmatrix}$
  \item $B = \begin{pmatrix} 3 & 1 & 0 \\ 0 & 3 & 0 \\ 0 & 0 & 5
    \end{pmatrix}$
  \item $C = \begin{pmatrix} 3 & 1 & 0 \\ 0 & 3 & 1 \\ 0 & 0 & 3
    \end{pmatrix}$
\end{enumerate}
\end{exercise}

\begin{exercise}{Markov chain steady state \intermediate}{exo:markov}
\index{Markov chain!exercise}
A rat in a maze has three rooms.  At each time step, it moves
according to the transition matrix
\[
  T = \begin{pmatrix} 0.5 & 0.25 & 0.25 \\ 0.25 & 0.5 & 0.25 \\
  0.25 & 0.25 & 0.5 \end{pmatrix}.
\]
\begin{enumerate}[label=(\alph*)]
  \item Find the eigenvalues and eigenspaces of~$T$.
  \item Determine the steady-state distribution $\vec{\pi}$ satisfying
    $T\vec{\pi} = \vec{\pi}$, $\sum_i \pi_i = 1$.
  \item Compute $T^n$ explicitly and verify that
    $T^n \to \vec{\pi}\,\transpose{\vec{1}}$ as $n\to\infty$ (where
    $\vec{1} = (1,1,1)^{\mathsf{T}}/3$).
\end{enumerate}
\end{exercise}


% ============================================================
\section{Chapter summary}\label{sec:eigenvalue-summary}
% ============================================================

\begin{itemize}
  \item An \textbf{eigenvalue} $\lambda$ of $f\in\End(E)$ satisfies
    $f(v) = \lambda v$ for some nonzero~$v$ (the \textbf{eigenvector}).
    The \textbf{eigenspace} $E_\lambda(f) = \Ker(f - \lambda\Id)$ is a
    subspace.

  \item The \textbf{characteristic polynomial}
    $\chi_f(\lambda) = \det(f - \lambda\Id)$ has degree~$n$; its roots
    in~$\K$ are exactly the eigenvalues.  The set of eigenvalues is
    the \textbf{spectrum} $\Spec(f)$.

  \item The \textbf{Cayley--Hamilton theorem}: every matrix satisfies
    its own characteristic polynomial, $\chi_A(A) = 0$.

  \item For each eigenvalue~$\lambda$, the \textbf{geometric
    multiplicity} $m_g(\lambda) = \dim E_\lambda$ satisfies
    $1 \leq m_g(\lambda) \leq m_a(\lambda)$ (the \textbf{algebraic
    multiplicity}).

  \item $f$ is \textbf{diagonalizable} if and only if $\chi_f$ splits
    over~$\K$ and $m_g = m_a$ for every eigenvalue.  Equivalently,
    $E = \bigoplus_{\lambda\in\Spec(f)} E_\lambda(f)$.

  \item \textbf{Diagonalization algorithm}: compute $\chi_f$, find
    eigenvalues, compute eigenspaces, check $m_g = m_a$, form the
    change-of-basis matrix~$P$.

  \item Every endomorphism whose characteristic polynomial splits is
    \textbf{trigonalizable}.  Over~$\C$, every endomorphism is
    trigonalizable.

  \item The \textbf{minimal polynomial} $\mu_f$ is the monic
    polynomial of smallest degree annihilating~$f$.  It divides~$\chi_f$
    and shares its roots.  $f$ is diagonalizable if and only if $\mu_f$
    splits into distinct linear factors.

  \item \textbf{Applications}: if $A = PD\inv{P}$ then
    $A^k = PD^k\inv{P}$, which allows efficient computation of matrix
    powers, closed-form solutions of linear recurrences (Fibonacci),
    steady-state analysis of Markov chains, and explicit solutions of
    linear systems of ODEs.
\end{itemize}

% ch6-inner-products.tex — Chapter 6: Inner Product Spaces and Orthogonality
% Bilinear and sesquilinear forms, quadratic forms, inner products,
% norms, Cauchy–Schwarz, orthogonality, Gram–Schmidt, orthonormal bases,
% orthogonal/unitary matrices, adjoint operators, applications.

\chapter{Inner Product Spaces and Orthogonality}\label{ch:inner-products}

% ============================================================
% Motivating introduction
% ============================================================

Up to this point, our study of linear algebra has been purely
\emph{algebraic}: we have manipulated vector spaces, linear maps, and
matrices without ever measuring a length, an angle, or a distance.  Yet
geometry---in the intuitive sense of Euclidean space---depends critically
on such measurements.  How long is a vector?  What angle do two vectors
enclose?  Which vector in a subspace is closest to a given point?

The tool that answers all these questions at once is the \emph{inner
product}\index{inner product}: a function that takes two vectors and
returns a scalar, generalising the familiar dot product on~$\R^n$.  Once
an inner product is fixed, the underlying vector space acquires a rich
geometric structure---norms, distances, angles, orthogonality,
projections---and many powerful theorems become available.

\begin{center}
\begin{tikzpicture}[>=Stealth, scale=1.6]
  % --- Angle between two vectors ---
  \draw[gray!30, very thin, step=0.5] (-0.3,-0.3) grid (3.8,2.8);
  \draw[->] (-0.3,0) -- (3.8,0) node[right]{$x_1$};
  \draw[->] (0,-0.3) -- (0,2.8) node[above]{$x_2$};

  \draw[->, thick, blue!70!black] (0,0) -- (3,1)
    node[right] {$\vec{u}$};
  \draw[->, thick, red!70!black] (0,0) -- (1,2.5)
    node[above] {$\vec{v}$};

  % Angle arc
  \draw[thick, purple!70!black] (0.9,0.3) arc[start angle=18.4,
    end angle=68.2, radius=0.95];
  \node[purple!70!black] at (0.75,0.95) {$\theta$};

  % Annotation
  \node[below, text width=4cm, align=center] at (1.8,-0.6)
    {\small $\cos\theta = \dfrac{\inner{\vec{u},\vec{v}}}%
      {\norm{\vec{u}}\;\norm{\vec{v}}}$};
\end{tikzpicture}
\captionof{figure}{The inner product encodes the angle between two
  vectors.  All of Euclidean geometry follows.}
\label{fig:inner-product-angle}
\end{center}

This chapter develops the theory systematically, starting from the
general notion of bilinear forms, specialising to inner products, and
building up the machinery of orthogonality, projections, and
orthonormal bases.  We then study the linear maps that preserve inner
products (orthogonal and unitary matrices) and those that are
``self-dual'' with respect to the inner product (adjoint operators).

Throughout, $\K$ denotes $\R$ or~$\C$, and $E$ denotes a
finite-dimensional $\K$-vector space unless stated otherwise.  We use
the convention that inner products on~$\C$ are linear in the
\emph{first} argument and conjugate-linear in the second (the
\emph{physics convention}).


% ============================================================
\section{Bilinear forms}\label{sec:bilinear-forms}
% ============================================================

\begin{definition}{Bilinear form}{def:bilinear-form}
\index{bilinear form}
Let $E$ be a vector space over~$\K$.  A \emph{bilinear form} on~$E$ is
a map $\varphi\colon E\times E \to \K$ that is linear in each argument:
\begin{enumerate}[label=(\roman*)]
  \item For all $u,v,w\in E$ and $\lambda\in\K$:
    \[
      \varphi(u+\lambda v,\, w) = \varphi(u,w) + \lambda\,\varphi(v,w).
    \]
  \item For all $u,v,w\in E$ and $\lambda\in\K$:
    \[
      \varphi(u,\, v+\lambda w) = \varphi(u,v) + \lambda\,\varphi(u,w).
    \]
\end{enumerate}
\end{definition}

\begin{definition}{Symmetric and antisymmetric bilinear forms}%
{def:symmetric-bilinear}
\index{bilinear form!symmetric}\index{bilinear form!antisymmetric}
A bilinear form $\varphi$ on~$E$ is called:
\begin{itemize}
  \item \emph{symmetric} if $\varphi(u,v) = \varphi(v,u)$ for all
    $u,v\in E$;
  \item \emph{antisymmetric} (or \emph{skew-symmetric}) if
    $\varphi(u,v) = -\varphi(v,u)$ for all $u,v\in E$.
\end{itemize}
\end{definition}

\begin{definition}{Matrix of a bilinear form}{def:matrix-bilinear}
\index{bilinear form!matrix of}
Let $\cB = (\vece{1},\ldots,\vece{n})$ be a basis of~$E$ and let
$\varphi$ be a bilinear form on~$E$.  The \emph{matrix of~$\varphi$
in the basis~$\cB$} is the matrix $G\in\Mat_n(\K)$ defined by
\[
  G_{ij} \deff \varphi(\vece{i},\vece{j}),
  \qquad 1\leq i,j\leq n.
\]
If $u = \sum_i x_i\,\vece{i}$ and $v = \sum_j y_j\,\vece{j}$, then
\[
  \varphi(u,v) = \transpose{X}\, G\, Y,
\]
where $X = (x_1,\ldots,x_n)^{\mathsf{T}}$ and
$Y = (y_1,\ldots,y_n)^{\mathsf{T}}$ are the coordinate vectors.
\end{definition}

\begin{proposition}{Change of basis for bilinear forms}%
{prop:bilinear-change-basis}
\index{bilinear form!change of basis}
If $P$ is the change-of-basis matrix from $\cB$ to $\cB'$, and $G$ is
the matrix of~$\varphi$ in~$\cB$, then the matrix of~$\varphi$
in~$\cB'$ is
\[
  G' = \transpose{P}\, G\, P.
\]
In particular, $\varphi$ is symmetric if and only if $G$ is symmetric
($G = \transpose{G}$).
\end{proposition}

\begin{proof}
Let $u,v\in E$ with coordinates $X'$ in~$\cB'$ and $X = PX'$ in~$\cB$
(similarly $Y = PY'$).  Then
\[
  \varphi(u,v) = \transpose{X}\, G\, Y
    = \transpose{(PX')}\, G\, (PY')
    = \transpose{X'}\,(\transpose{P}\, G\, P)\, Y'.
\]
Since this holds for all $X', Y'$, we have $G' = \transpose{P}\, G\, P$.
\end{proof}

\begin{definition}{Non-degenerate bilinear form}{def:non-degenerate}
\index{bilinear form!non-degenerate}\index{non-degenerate}
A bilinear form $\varphi$ on~$E$ is \emph{non-degenerate} if
\[
  \bigl(\,\forall v\in E,\;\varphi(u,v) = 0\bigr)
  \;\Longrightarrow\; u = \vzero.
\]
Equivalently, the matrix~$G$ of~$\varphi$ in any basis is invertible.
\end{definition}


% ============================================================
\section{Quadratic forms}\label{sec:quadratic-forms}
% ============================================================

\begin{definition}{Quadratic form}{def:quadratic-form}
\index{quadratic form}
Let $E$ be a $\K$-vector space.  A \emph{quadratic form} on~$E$ is a
map $q\colon E\to\K$ such that
\begin{enumerate}[label=(\roman*)]
  \item $q(\lambda v) = \lambda^2\, q(v)$ for all $\lambda\in\K$,
    $v\in E$;
  \item the map $\varphi\colon E\times E\to\K$ defined by
    \[
      \varphi(u,v) \deff \frac{1}{2}\bigl[q(u+v) - q(u) - q(v)\bigr]
    \]
    is bilinear (this is called the \emph{polar form}\index{polar form}
    of~$q$).
\end{enumerate}
Conversely, every symmetric bilinear form $\varphi$ defines a quadratic
form $q(v) = \varphi(v,v)$.
\end{definition}

\begin{proposition}{Polarization identity}{prop:polarization}
\index{polarization identity}
If $\Char(\K)\neq 2$ and $\varphi$ is a symmetric bilinear form with
associated quadratic form $q(v) = \varphi(v,v)$, then
\[
  \varphi(u,v) = \frac{1}{2}\bigl[q(u+v) - q(u) - q(v)\bigr]
    = \frac{1}{4}\bigl[q(u+v) - q(u-v)\bigr].
\]
Thus, $q$ and $\varphi$ determine each other uniquely.
\end{proposition}

\begin{proof}
Expanding $q(u+v) = \varphi(u+v,\,u+v) = q(u) + 2\varphi(u,v) + q(v)$
gives the first identity.  For the second, note that
$q(u-v) = q(u) - 2\varphi(u,v) + q(v)$, so subtracting yields
$q(u+v) - q(u-v) = 4\varphi(u,v)$.
\end{proof}

\begin{theorem}{Sylvester's law of inertia}{thm:sylvester}
\index{Sylvester's law of inertia}\index{signature}
Let $q$ be a quadratic form on a real vector space~$E$ of
dimension~$n$.  There exists a basis of~$E$ in which the matrix of the
associated bilinear form is diagonal with entries $+1$, $-1$, and~$0$:
\[
  G = \diag(\underbrace{1,\ldots,1}_{p},\;
    \underbrace{-1,\ldots,-1}_{r-p},\;
    \underbrace{0,\ldots,0}_{n-r}).
\]
The integers $p$ (the number of $+1$'s), $r-p$ (the number of $-1$'s),
and $n-r$ (the number of~$0$'s) depend only on~$q$, not on the choice
of basis.  The pair $(p,\,r-p)$ is called the \emph{signature} of~$q$,
and $r = \rank(q)$.
\end{theorem}

\begin{remark}{Positive definiteness from the signature}%
{rem:signature-positive-definite}
A real quadratic form~$q$ is \emph{positive definite}\index{positive
definite} (i.e.\ $q(v)>0$ for all $v\neq\vzero$) if and only if its
signature is $(n,0)$, i.e.\ all signs are~$+1$.
\end{remark}


% ============================================================
\section{Inner products}\label{sec:inner-products}
% ============================================================

\begin{definition}{Inner product (real case)}{def:inner-product-real}
\index{inner product!real}
Let $E$ be a real vector space.  An \emph{inner product} on~$E$ is a
bilinear form $\inner{\cdot,\cdot}\colon E\times E\to\R$ that is
\begin{enumerate}[label=(\roman*)]
  \item \emph{symmetric}: $\inner{u,v} = \inner{v,u}$ for all
    $u,v\in E$;
  \item \emph{positive definite}: $\inner{v,v} > 0$ for all
    $v\neq\vzero$.
\end{enumerate}
A real vector space equipped with an inner product is called a
\emph{Euclidean space}\index{Euclidean space}.
\end{definition}

\begin{definition}{Sesquilinear form}{def:sesquilinear}
\index{sesquilinear form}
Let $E$ be a complex vector space.  A \emph{sesquilinear form} on~$E$
is a map $\varphi\colon E\times E\to\C$ that is linear in the first
argument and conjugate-linear in the second:
\begin{enumerate}[label=(\roman*)]
  \item $\varphi(\alpha u + \beta v,\, w) = \alpha\,\varphi(u,w)
    + \beta\,\varphi(v,w)$;
  \item $\varphi(u,\, \alpha v + \beta w) = \bar{\alpha}\,\varphi(u,v)
    + \bar{\beta}\,\varphi(u,w)$.
\end{enumerate}
A sesquilinear form is \emph{Hermitian}\index{Hermitian form} if
$\varphi(u,v) = \overline{\varphi(v,u)}$ for all $u,v\in E$.
\end{definition}

\begin{definition}{Inner product (complex case)}{def:inner-product-complex}
\index{inner product!complex}\index{Hermitian inner product}
Let $E$ be a complex vector space.  A \emph{Hermitian inner product}
(or simply \emph{inner product}) on~$E$ is a sesquilinear form
$\inner{\cdot,\cdot}\colon E\times E\to\C$ that is
\begin{enumerate}[label=(\roman*)]
  \item \emph{Hermitian symmetric}: $\inner{u,v} =
    \overline{\inner{v,u}}$ for all $u,v\in E$;
  \item \emph{positive definite}: $\inner{v,v} > 0$ for all
    $v\neq\vzero$ (note that $\inner{v,v}\in\R$ by~(i)).
\end{enumerate}
A complex vector space equipped with a Hermitian inner product is called
a \emph{unitary space}\index{unitary space} (or \emph{pre-Hilbert
space}).
\end{definition}

\begin{remark}{Unified notation}{rem:unified-inner-product}
Over~$\R$, a Hermitian inner product reduces to a symmetric bilinear
form (since conjugation is trivial).  We therefore write
$\inner{u,v}$ in both cases and speak simply of an ``inner product
space''\index{inner product space}.
\end{remark}


% ============================================================
\section{Examples}\label{sec:inner-product-examples}
% ============================================================

\begin{example}{Standard dot product on $\R^n$}{ex:dot-product}
\index{dot product}\index{standard inner product!on $\R^n$}
The \emph{standard inner product} on $\R^n$ is
\[
  \inner{x,y} \deff \sum_{i=1}^n x_i\, y_i = \transpose{x}\, y,
  \qquad x,y\in\R^n.
\]
Its matrix in the standard basis is $G = I_n$.
\end{example}

\begin{example}{Hermitian product on $\C^n$}{ex:hermitian-Cn}
\index{standard inner product!on $\C^n$}
The \emph{standard Hermitian inner product} on $\C^n$ is
\[
  \inner{z,w} \deff \sum_{i=1}^n z_i\,\bar{w}_i = \hermit{w}\, z,
  \qquad z,w\in\C^n.
\]
This is linear in the first argument and conjugate-linear in the second.
\end{example}

\begin{example}{$L^2$ inner product on $\mathcal{C}([a,b])$}%
{ex:L2-inner-product}
\index{L2 inner product@$L^2$ inner product}
On the space $\mathcal{C}([a,b],\R)$ of continuous real-valued functions
on~$[a,b]$, the map
\[
  \inner{f,g} \deff \int_a^b f(t)\, g(t)\,\mathrm{d}t
\]
is an inner product.  Symmetry and bilinearity are immediate from the
properties of the integral; positive definiteness follows because
$\inner{f,f} = \int_a^b f(t)^2\,\mathrm{d}t \geq 0$, with equality if
and only if $f = 0$ (by continuity).

The complex analogue on $\mathcal{C}([a,b],\C)$ is
$\inner{f,g} = \int_a^b f(t)\,\overline{g(t)}\,\mathrm{d}t$.
\end{example}

\begin{example}{Weighted inner product}{ex:weighted-inner-product}
\index{inner product!weighted}
Let $w_1,\ldots,w_n > 0$ be positive reals.  Then
\[
  \inner{x,y}_w \deff \sum_{i=1}^n w_i\, x_i\, y_i
\]
defines an inner product on~$\R^n$.  Its matrix in the standard basis is
$G = \diag(w_1,\ldots,w_n)$.
\end{example}

\begin{example}{Frobenius inner product}{ex:frobenius}
\index{Frobenius inner product}
On the space $\Mat_n(\R)$, the map
\[
  \inner{A,B} \deff \tr(\transpose{A}\, B) = \sum_{i,j} a_{ij}\, b_{ij}
\]
defines an inner product (the \emph{Frobenius inner product}).
\end{example}


% ============================================================
\section{Norm and distance}\label{sec:norm-distance}
% ============================================================

\begin{definition}{Norm induced by an inner product}{def:norm}
\index{norm}\index{norm!induced by inner product}
Let $(E,\inner{\cdot,\cdot})$ be an inner product space.  The
\emph{norm} of $v\in E$ is
\[
  \norm{v} \deff \sqrt{\inner{v,v}}.
\]
\end{definition}

\begin{definition}{Distance}{def:distance}
\index{distance}
The \emph{distance} between $u$ and~$v$ in an inner product space is
\[
  d(u,v) \deff \norm{u - v}.
\]
\end{definition}

\begin{proposition}{Basic properties of the norm}{prop:norm-properties}
Let $(E,\inner{\cdot,\cdot})$ be an inner product space.  For all
$u,v\in E$ and $\lambda\in\K$:
\begin{enumerate}[label=(\roman*)]
  \item $\norm{v}\geq 0$, with equality if and only if $v = \vzero$;
  \item $\norm{\lambda v} = \abs{\lambda}\,\norm{v}$ \quad
    (absolute homogeneity);
  \item $\norm{u+v}^2 + \norm{u-v}^2 = 2\norm{u}^2 + 2\norm{v}^2$
    \quad (parallelogram law\index{parallelogram law}).
\end{enumerate}
\end{proposition}

\begin{proof}
Parts (i) and (ii) follow directly from the definition and the
properties of the inner product.  For~(iii), expand using bilinearity:
\[
  \norm{u+v}^2 + \norm{u-v}^2
    = \inner{u+v,u+v} + \inner{u-v,u-v}
    = 2\inner{u,u} + 2\inner{v,v}
    = 2\norm{u}^2 + 2\norm{v}^2. \qedhere
\]
\end{proof}


% ============================================================
\section{The Cauchy--Schwarz inequality}\label{sec:cauchy-schwarz}
% ============================================================

\begin{theorem}{Cauchy--Schwarz inequality}{thm:cauchy-schwarz}
\index{Cauchy--Schwarz inequality}
Let $(E,\inner{\cdot,\cdot})$ be an inner product space over~$\K$.
For all $u,v\in E$,
\[
  \abs{\inner{u,v}} \leq \norm{u}\,\norm{v},
\]
with equality if and only if $u$ and~$v$ are linearly dependent
(i.e.\ one is a scalar multiple of the other).
\end{theorem}

\begin{proof}
If $v = \vzero$, both sides are zero and the result is trivial.  Assume
$v\neq\vzero$.

\medskip\noindent\textbf{Real case ($\K = \R$).}\quad
For any $t\in\R$, the positive definiteness of the inner product gives
\[
  0 \leq \norm{u - tv}^2
    = \inner{u - tv,\, u - tv}
    = \norm{u}^2 - 2t\inner{u,v} + t^2\norm{v}^2.
\]
This is a quadratic polynomial in~$t$ that is non-negative for all~$t$.
A real quadratic $at^2 + bt + c\geq 0$ for all~$t$ requires the
discriminant to satisfy $b^2 - 4ac \leq 0$.  Here
$a = \norm{v}^2$, $b = -2\inner{u,v}$, $c = \norm{u}^2$, so
\[
  4\inner{u,v}^2 - 4\norm{u}^2\norm{v}^2 \leq 0,
\]
which gives $\inner{u,v}^2 \leq \norm{u}^2\norm{v}^2$, and taking
square roots yields $\abs{\inner{u,v}}\leq\norm{u}\,\norm{v}$.

Equality holds if and only if the discriminant is zero, i.e.\ the
minimum of the quadratic is zero, which occurs at
$t_0 = \inner{u,v}/\norm{v}^2$, giving $u = t_0 v$.

\medskip\noindent\textbf{Complex case ($\K = \C$).}\quad
For any $t\in\R$ and any $\alpha\in\C$ with $\abs{\alpha}=1$, consider
\[
  0 \leq \norm{u - t\alpha v}^2
    = \norm{u}^2 - 2t\,\Re(\alpha\,\inner{v,u})
      + t^2\norm{v}^2.
\]
Choose $\alpha$ so that $\alpha\,\inner{v,u} = \abs{\inner{u,v}}$
(this is possible by writing $\inner{u,v} = \abs{\inner{u,v}}\,e^{i\theta}$
and setting $\alpha = e^{-i\theta}$).  Then
$\Re(\alpha\,\inner{v,u}) = \abs{\inner{u,v}}$, and the above becomes
\[
  0 \leq \norm{u}^2 - 2t\abs{\inner{u,v}} + t^2\norm{v}^2.
\]
This is again a non-negative real quadratic in~$t$, and the discriminant
argument gives
\[
  4\abs{\inner{u,v}}^2 - 4\norm{u}^2\norm{v}^2 \leq 0,
\]
i.e.\ $\abs{\inner{u,v}}\leq\norm{u}\,\norm{v}$.  Equality analysis is
identical: it holds if and only if $u = t_0\alpha v$ for the optimal
$t_0$, i.e.\ $u$ and~$v$ are linearly dependent.
\end{proof}

\begin{corollary}{Triangle inequality}{cor:triangle-inequality}
\index{triangle inequality}
For all $u,v\in E$,
\[
  \norm{u+v}\leq\norm{u}+\norm{v}.
\]
In particular, $d(\cdot,\cdot)$ satisfies the triangle inequality and
defines a metric on~$E$.
\end{corollary}

\begin{proof}
We compute
\begin{align*}
  \norm{u+v}^2
    &= \inner{u+v,\,u+v}
    = \norm{u}^2 + 2\Re\inner{u,v} + \norm{v}^2 \\
    &\leq \norm{u}^2 + 2\abs{\inner{u,v}} + \norm{v}^2 \\
    &\leq \norm{u}^2 + 2\norm{u}\,\norm{v} + \norm{v}^2
    = (\norm{u} + \norm{v})^2,
\end{align*}
where the last inequality uses Cauchy--Schwarz.  Taking square roots
gives the result.
\end{proof}

\begin{definition}{Angle between vectors}{def:angle}
\index{angle between vectors}
In a real inner product space, for nonzero $u,v\in E$, the
\emph{angle}~$\theta\in[0,\pi]$ between $u$ and~$v$ is defined by
\[
  \cos\theta \deff \frac{\inner{u,v}}{\norm{u}\,\norm{v}}.
\]
This is well-defined by the Cauchy--Schwarz inequality, which ensures
$\abs{\cos\theta}\leq 1$.
\end{definition}


% ============================================================
\section{Orthogonality}\label{sec:orthogonality}
% ============================================================

\begin{definition}{Orthogonal vectors}{def:orthogonal-vectors}
\index{orthogonal vectors}
Two vectors $u,v$ in an inner product space are \emph{orthogonal},
written $u\perp v$, if $\inner{u,v} = 0$.
\end{definition}

\begin{proposition}{Pythagorean theorem}{prop:pythagorean}
\index{Pythagorean theorem}
If $u\perp v$, then $\norm{u+v}^2 = \norm{u}^2 + \norm{v}^2$.
\end{proposition}

\begin{proof}
$\norm{u+v}^2 = \norm{u}^2 + 2\Re\inner{u,v} + \norm{v}^2
= \norm{u}^2 + \norm{v}^2$, since $\inner{u,v} = 0$.
\end{proof}

\begin{definition}{Orthogonal complement}{def:orthogonal-complement}
\index{orthogonal complement}
Let $F$ be a subset (or subspace) of an inner product space~$E$.  The
\emph{orthogonal complement} of~$F$ is
\[
  F^{\perp} \deff \setcond{v\in E}{\inner{v,w} = 0
    \text{ for all } w\in F}.
\]
\end{definition}

\begin{theorem}{Properties of the orthogonal complement}%
{thm:orth-complement-properties}
\index{orthogonal complement!properties}
Let $(E,\inner{\cdot,\cdot})$ be a finite-dimensional inner product
space and let $F$ be a subspace of~$E$.  Then:
\begin{enumerate}[label=(\roman*)]
  \item $F^{\perp}$ is a subspace of~$E$.
  \item $E = F \oplus F^{\perp}$ (orthogonal direct sum).
  \item $\dim F^{\perp} = \dim E - \dim F$.
  \item $(F^{\perp})^{\perp} = F$.
  \item $F\cap F^{\perp} = \set{\vzero}$.
\end{enumerate}
\end{theorem}

\begin{proof}
\textbf{(i)} $F^{\perp}$ is clearly non-empty ($\vzero\in F^{\perp}$)
and closed under addition and scalar multiplication by linearity of the
inner product.

\textbf{(v)} If $v\in F\cap F^{\perp}$, then $\inner{v,v} = 0$, so
$v = \vzero$ by positive definiteness.

\textbf{(ii)--(iv)} We prove these after establishing orthogonal
projection (
ef{thm:thm:orthogonal-projection}).  Once we know that every
$v\in E$ can be written as $v = p + p^{\perp}$ with $p\in F$ and
$p^{\perp}\in F^{\perp}$, the direct sum decomposition follows
from~(v).  Part~(iii) then follows by counting dimensions, and~(iv)
is immediate from $E = F\oplus F^{\perp}$ and the observation that
$F\subset (F^{\perp})^{\perp}$ together with the dimension count
$\dim(F^{\perp})^{\perp} = n - \dim F^{\perp} = \dim F$.
\end{proof}

\begin{proposition}{Orthogonal complement of sums and intersections}%
{prop:orth-sum-intersection}
Let $F_1,F_2$ be subspaces of a finite-dimensional inner product
space~$E$.  Then:
\begin{enumerate}[label=(\roman*)]
  \item $(F_1 + F_2)^{\perp} = F_1^{\perp}\cap F_2^{\perp}$.
  \item $(F_1\cap F_2)^{\perp} = F_1^{\perp} + F_2^{\perp}$.
\end{enumerate}
\end{proposition}

\begin{proof}
\textbf{(i)} $v\in(F_1+F_2)^{\perp}$ iff $\inner{v,w}=0$ for all
$w\in F_1+F_2$, which (since $F_1,F_2\subset F_1+F_2$) is equivalent
to $v\in F_1^{\perp}$ and $v\in F_2^{\perp}$.

\textbf{(ii)} Take orthogonal complements in (i) and use
$(F^{\perp})^{\perp}=F$.
\end{proof}


% ============================================================
\section{Orthogonal projection}\label{sec:orthogonal-projection}
% ============================================================

\begin{theorem}{Orthogonal projection}{thm:orthogonal-projection}
\index{orthogonal projection}
Let $(E,\inner{\cdot,\cdot})$ be a finite-dimensional inner product
space and let $F$ be a subspace of~$E$.  For every $v\in E$, there
exists a unique vector $p\in F$ such that
\[
  v - p \in F^{\perp}.
\]
The vector~$p$ is called the \emph{orthogonal projection} of~$v$
onto~$F$, denoted $p = \proj_F(v)$.  The map
$\proj_F\colon E\to E$ is a linear map satisfying
$\proj_F^2 = \proj_F$, $\im(\proj_F) = F$, and
$\Ker(\proj_F) = F^{\perp}$.
\end{theorem}

\begin{proof}
\textbf{Existence.}\quad Let $(e_1,\ldots,e_k)$ be an orthonormal basis
of~$F$ (which exists by Gram--Schmidt, 
ef{thm:thm:gram-schmidt}).  Set
\[
  p \deff \sum_{i=1}^k \inner{v,e_i}\, e_i.
\]
Then $p\in F$, and for each basis vector $e_j$ we have
\[
  \inner{v - p,\, e_j}
    = \inner{v,e_j} - \sum_{i=1}^k \inner{v,e_i}\inner{e_i,e_j}
    = \inner{v,e_j} - \inner{v,e_j} = 0.
\]
By linearity, $\inner{v-p,\,w} = 0$ for all $w\in F$, so
$v-p\in F^{\perp}$.

\textbf{Uniqueness.}\quad If $p'$ is another such vector, then
$p - p'\in F$ and $p - p'\in F^{\perp}$, so
$p-p'\in F\cap F^{\perp} = \set{\vzero}$.

\textbf{Linearity} is clear from the formula.  The remaining properties
$\proj_F^2 = \proj_F$, $\im(\proj_F) = F$, and
$\Ker(\proj_F) = F^{\perp}$ follow directly from the construction and
uniqueness.
\end{proof}

\begin{theorem}{Best approximation}{thm:best-approximation}
\index{best approximation}\index{orthogonal projection!best approximation}
Let $F$ be a subspace of a finite-dimensional inner product space~$E$
and let $v\in E$.  The orthogonal projection $p = \proj_F(v)$ is the
unique element of~$F$ closest to~$v$:
\[
  \norm{v - p} < \norm{v - w}
  \qquad\text{for all } w\in F,\; w\neq p.
\]
\end{theorem}

\begin{proof}
Let $w\in F$.  Since $v - p\in F^{\perp}$ and $p - w\in F$, the
Pythagorean theorem gives
\[
  \norm{v - w}^2
    = \norm{(v - p) + (p - w)}^2
    = \norm{v - p}^2 + \norm{p - w}^2
    \geq \norm{v - p}^2,
\]
with equality if and only if $p = w$.
\end{proof}

\begin{center}
\begin{tikzpicture}[>=Stealth, scale=1.5]
  % Subspace F (a plane, drawn as a parallelogram)
  \fill[blue!8] (-0.5,-0.3) -- (3.5,-0.3) -- (4.2,1.2) -- (0.2,1.2)
    -- cycle;
  \node[blue!70!black] at (3.8,0.15) {$F$};

  % Vector v
  \coordinate (v) at (1.8,2.5);
  \coordinate (p) at (1.8,0.5);

  \draw[->, thick, red!70!black] (0,0) -- (v)
    node[above right] {$v$};
  \draw[->, thick, blue!70!black] (0,0) -- (p)
    node[below right] {$p = \proj_F(v)$};

  % Perpendicular line
  \draw[thick, dashed, purple!70!black] (v) -- (p)
    node[midway, right, xshift=3pt] {$v - p \in F^{\perp}$};

  % Right angle mark
  \draw[purple!70!black] (1.8,0.75) -- (2.05,0.75) -- (2.05,0.5);

  % Another point w in F
  \coordinate (w) at (3,0.8);
  \draw[->, thick, gray] (0,0) -- (w) node[right] {$w$};
  \draw[dotted, thick, gray] (v) -- (w)
    node[midway, right, xshift=2pt] {\small $\norm{v-w}\geq\norm{v-p}$};
\end{tikzpicture}
\captionof{figure}{Orthogonal projection onto a subspace~$F$.  The
  projection $p = \proj_F(v)$ is the closest point in~$F$ to~$v$.}
\label{fig:orthogonal-projection}
\end{center}

\begin{proposition}{Projection formula with an orthonormal basis}%
{prop:projection-formula}
\index{orthogonal projection!formula}
If $(e_1,\ldots,e_k)$ is an orthonormal basis of a subspace~$F$, then
\[
  \proj_F(v) = \sum_{i=1}^k \inner{v,e_i}\, e_i.
\]
In particular, if $(e_1,\ldots,e_n)$ is an orthonormal basis of the
full space~$E$, then $v = \sum_{i=1}^n \inner{v,e_i}\, e_i$.
\end{proposition}


% ============================================================
\section{Gram--Schmidt orthogonalization}\label{sec:gram-schmidt}
% ============================================================

\begin{theorem}{Gram--Schmidt process}{thm:gram-schmidt}
\index{Gram--Schmidt process}\index{orthogonalization}
Let $(v_1,\ldots,v_k)$ be a linearly independent family in an inner
product space~$E$.  There exists a unique orthonormal family
$(e_1,\ldots,e_k)$ such that for each $j=1,\ldots,k$:
\begin{enumerate}[label=(\roman*)]
  \item $\Span(e_1,\ldots,e_j) = \Span(v_1,\ldots,v_j)$;
  \item $\inner{v_j,e_j} > 0$ (i.e.\ the ``leading coefficient'' is
    positive).
\end{enumerate}
The family $(e_1,\ldots,e_k)$ is constructed by the
\textbf{Gram--Schmidt algorithm}:
\begin{align*}
  w_1 &\deff v_1, &
    e_1 &\deff \frac{w_1}{\norm{w_1}}, \\
  w_j &\deff v_j - \sum_{i=1}^{j-1} \inner{v_j,e_i}\, e_i, &
    e_j &\deff \frac{w_j}{\norm{w_j}},
    \qquad j = 2,\ldots,k.
\end{align*}
\end{theorem}

\begin{proof}
We proceed by induction on~$j$.

\textbf{Base case ($j=1$).}\quad Set $w_1 = v_1$.  Since $v_1\neq\vzero$
(linear independence), $\norm{w_1}>0$ and $e_1 = w_1/\norm{w_1}$ is
well-defined with $\norm{e_1} = 1$.  Clearly
$\Span(e_1) = \Span(v_1)$ and $\inner{v_1,e_1} = \norm{v_1} > 0$.

\textbf{Inductive step.}\quad Assume that $(e_1,\ldots,e_{j-1})$ is an
orthonormal family with $\Span(e_1,\ldots,e_{j-1}) = \Span(v_1,\ldots,
v_{j-1})$ and $\inner{v_i,e_i}>0$ for $i<j$.  Define
\[
  w_j = v_j - \sum_{i=1}^{j-1}\inner{v_j,e_i}\, e_i.
\]
This is exactly $v_j - \proj_{F_{j-1}}(v_j)$, where
$F_{j-1} = \Span(e_1,\ldots,e_{j-1})$.  Hence $w_j\perp e_i$ for all
$i<j$ (by construction of the orthogonal projection).

We must show $w_j\neq\vzero$.  If $w_j = \vzero$, then
$v_j = \sum_{i=1}^{j-1}\inner{v_j,e_i}\,e_i\in F_{j-1} =
\Span(v_1,\ldots,v_{j-1})$, contradicting the linear independence
of $(v_1,\ldots,v_k)$.

Set $e_j = w_j/\norm{w_j}$.  Then $\norm{e_j} = 1$, and $e_j\perp e_i$
for all $i<j$ since $w_j\perp e_i$.  Moreover,
\[
  \Span(e_1,\ldots,e_j) = \Span(e_1,\ldots,e_{j-1},\,w_j)
    = \Span(e_1,\ldots,e_{j-1},\,v_j)
    = \Span(v_1,\ldots,v_j),
\]
using the induction hypothesis and the definition of~$w_j$.  Finally,
$\inner{v_j,e_j} = \inner{v_j, w_j/\norm{w_j}} = \norm{w_j} > 0$
(since $\inner{v_j,w_j} = \inner{w_j + \sum_i c_i e_i,\, w_j} =
\norm{w_j}^2 > 0$).

\textbf{Uniqueness.}\quad Suppose $(\tilde{e}_1,\ldots,\tilde{e}_k)$
also satisfies conditions (i)--(ii).  By~(i),
$\Span(\tilde{e}_1) = \Span(v_1)$, so $\tilde{e}_1 = \pm e_1$;
condition~(ii) forces $\tilde{e}_1 = e_1$.  Inductively, the orthogonal
complement of $\Span(\tilde{e}_1,\ldots,\tilde{e}_{j-1})$ within
$\Span(v_1,\ldots,v_j)$ is one-dimensional, Spanned by~$w_j$.
Normalisation with a positive leading coefficient forces
$\tilde{e}_j = e_j$.
\end{proof}

\begin{example}{Gram--Schmidt in $\R^3$}{ex:gram-schmidt-R3}
\index{Gram--Schmidt process!example}
Apply Gram--Schmidt to $v_1 = (1,1,0)$, $v_2 = (1,0,1)$,
$v_3 = (0,1,1)$ with the standard inner product.

\textbf{Step 1.}\quad $w_1 = v_1 = (1,1,0)$,\quad
$e_1 = w_1/\norm{w_1} = \frac{1}{\sqrt{2}}(1,1,0)$.

\textbf{Step 2.}\quad
$\inner{v_2,e_1} = \frac{1}{\sqrt{2}}(1\cdot 1 + 0\cdot 1 + 1\cdot 0)
  = \frac{1}{\sqrt{2}}$.

$w_2 = v_2 - \inner{v_2,e_1}\,e_1
  = (1,0,1) - \frac{1}{\sqrt{2}}\cdot\frac{1}{\sqrt{2}}(1,1,0)
  = (1,0,1) - (\tfrac{1}{2},\tfrac{1}{2},0)
  = (\tfrac{1}{2},-\tfrac{1}{2},1)$.

$\norm{w_2} = \sqrt{\tfrac{1}{4}+\tfrac{1}{4}+1}
  = \sqrt{\tfrac{3}{2}} = \frac{\sqrt{6}}{2}$,
\quad so $e_2 = \frac{1}{\sqrt{6}}(1,-1,2)$.

\textbf{Step 3.}\quad
$\inner{v_3,e_1} = \frac{1}{\sqrt{2}}(0+1+0) = \frac{1}{\sqrt{2}}$,\quad
$\inner{v_3,e_2} = \frac{1}{\sqrt{6}}(0-1+2) = \frac{1}{\sqrt{6}}$.

$w_3 = v_3 - \frac{1}{\sqrt{2}}\cdot\frac{1}{\sqrt{2}}(1,1,0)
  - \frac{1}{\sqrt{6}}\cdot\frac{1}{\sqrt{6}}(1,-1,2)
  = (0,1,1) - (\tfrac{1}{2},\tfrac{1}{2},0) - (\tfrac{1}{6},-\tfrac{1}{6},\tfrac{1}{3})
  = (-\tfrac{2}{3},\tfrac{2}{3},\tfrac{2}{3})$.

$\norm{w_3} = \tfrac{2}{3}\sqrt{3}$,\quad so
$e_3 = \frac{1}{\sqrt{3}}(-1,1,1)$.

\medskip
The resulting orthonormal basis is:
\[
  e_1 = \frac{1}{\sqrt{2}}\begin{pmatrix}1\\1\\0\end{pmatrix},\quad
  e_2 = \frac{1}{\sqrt{6}}\begin{pmatrix}1\\-1\\2\end{pmatrix},\quad
  e_3 = \frac{1}{\sqrt{3}}\begin{pmatrix}-1\\1\\1\end{pmatrix}.
\]
\end{example}

\begin{center}
\begin{tikzpicture}[>=Stealth, scale=2.2]
  % Gram-Schmidt illustration in 2D (simpler to draw)
  \draw[gray!30, very thin, step=0.5] (-0.3,-0.3) grid (3.3,2.3);
  \draw[->] (-0.3,0) -- (3.3,0) node[right]{$x_1$};
  \draw[->] (0,-0.3) -- (0,2.3) node[above]{$x_2$};

  % v1 and v2
  \draw[->, thick, blue!50, dashed] (0,0) -- (2,1)
    node[right] {\small $v_1$};
  \draw[->, thick, red!50, dashed] (0,0) -- (1,2)
    node[above] {\small $v_2$};

  % e1 (normalised v1)
  \draw[->, very thick, blue!70!black] (0,0) -- ({2/sqrt(5)},{1/sqrt(5)})
    node[below right] {$e_1$};

  % projection of v2 onto e1
  \coordinate (proj) at ({4/(5*sqrt(5))*2},{4/(5*sqrt(5))*1});
  \draw[dotted, gray] (1,2) -- ({4/5},{2/5});

  % w2 = v2 - proj
  \draw[->, thick, orange!80!black] ({4/5},{2/5}) -- (1,2)
    node[midway, left, xshift=-2pt] {\small $w_2$};

  % e2 (normalised w2)
  \pgfmathsetmacro{\wx}{1 - 4/5}
  \pgfmathsetmacro{\wy}{2 - 2/5}
  \pgfmathsetmacro{\wnorm}{sqrt(\wx*\wx + \wy*\wy)}
  \draw[->, very thick, red!70!black] (0,0) --
    ({\wx/\wnorm},{\wy/\wnorm})
    node[above left] {$e_2$};

  % Right angle
  \draw[gray] ({4/5+0.08},{2/5+0.04}) --
    ({4/5+0.04},{2/5+0.12}) -- ({4/5-0.04},{2/5+0.08});
\end{tikzpicture}
\captionof{figure}{Gram--Schmidt in $\R^2$: $v_2$ is orthogonalised
  against $e_1$ to produce $w_2$, then normalised to give~$e_2$.}
\label{fig:gram-schmidt}
\end{center}


% ============================================================
\section{Orthonormal bases}\label{sec:orthonormal-bases}
% ============================================================

\begin{definition}{Orthogonal and orthonormal families}%
{def:orthonormal-family}
\index{orthogonal family}\index{orthonormal family}\index{orthonormal basis}
A family $(e_1,\ldots,e_k)$ in an inner product space is:
\begin{itemize}
  \item \emph{orthogonal} if $\inner{e_i,e_j} = 0$ for all $i\neq j$;
  \item \emph{orthonormal} if it is orthogonal and $\norm{e_i} = 1$ for
    all~$i$.
\end{itemize}
An orthonormal family that is also a basis of~$E$ is called an
\emph{orthonormal basis} (ONB).
\end{definition}

\begin{proposition}{Orthogonal families are linearly independent}%
{prop:orthogonal-indep}
Every orthogonal family of nonzero vectors is linearly independent.
\end{proposition}

\begin{proof}
Let $(e_1,\ldots,e_k)$ be orthogonal with each $e_i\neq\vzero$.
Suppose $\sum_{i=1}^k\alpha_i\,e_i = \vzero$.  Taking the inner
product with~$e_j$:
\[
  0 = \inner*{\sum_i\alpha_i\,e_i,\, e_j}
    = \sum_i\alpha_i\inner{e_i,e_j}
    = \alpha_j\inner{e_j,e_j}
    = \alpha_j\norm{e_j}^2.
\]
Since $\norm{e_j}>0$, we get $\alpha_j = 0$ for all~$j$.
\end{proof}

\begin{corollary}{Existence of orthonormal bases}{cor:onb-existence}
\index{orthonormal basis!existence}
Every finite-dimensional inner product space admits an orthonormal basis.
\end{corollary}

\begin{proof}
Apply the Gram--Schmidt process (
ef{thm:thm:gram-schmidt}) to any basis
of~$E$.
\end{proof}

\begin{theorem}{Parseval's identity}{thm:parseval}
\index{Parseval's identity}
Let $(e_1,\ldots,e_n)$ be an orthonormal basis of an inner product
space~$E$.  Then for all $u,v\in E$:
\[
  \inner{u,v} = \sum_{i=1}^n \inner{u,e_i}\,\overline{\inner{v,e_i}}.
\]
In particular, taking $v = u$:
\[
  \norm{u}^2 = \sum_{i=1}^n \abs{\inner{u,e_i}}^2.
\]
\end{theorem}

\begin{proof}
Write $u = \sum_{i=1}^n\alpha_i\, e_i$ and $v = \sum_{j=1}^n\beta_j\, e_j$
where $\alpha_i = \inner{u,e_i}$ and $\beta_j = \inner{v,e_j}$.  Then
\[
  \inner{u,v} = \inner*{\sum_i\alpha_i\,e_i,\,\sum_j\beta_j\,e_j}
    = \sum_{i,j}\alpha_i\,\bar{\beta}_j\,\inner{e_i,e_j}
    = \sum_i\alpha_i\,\bar{\beta}_i
    = \sum_i\inner{u,e_i}\,\overline{\inner{v,e_i}}. \qedhere
\]
\end{proof}

\begin{theorem}{Bessel's inequality}{thm:bessel}
\index{Bessel's inequality}
Let $(e_1,\ldots,e_k)$ be an orthonormal family (not necessarily a
basis) in an inner product space~$E$.  Then for all $v\in E$:
\[
  \sum_{i=1}^k \abs{\inner{v,e_i}}^2 \leq \norm{v}^2,
\]
with equality if and only if $v\in\Span(e_1,\ldots,e_k)$.
\end{theorem}

\begin{proof}
Let $p = \sum_{i=1}^k\inner{v,e_i}\,e_i$ be the orthogonal projection
of~$v$ onto $F = \Span(e_1,\ldots,e_k)$.  Then $v = p + (v-p)$ with
$p\perp (v-p)$, so by the Pythagorean theorem:
\[
  \norm{v}^2 = \norm{p}^2 + \norm{v-p}^2 \geq \norm{p}^2
    = \sum_{i=1}^k\abs{\inner{v,e_i}}^2.
\]
Equality holds iff $\norm{v-p} = 0$, i.e.\ $v = p\in F$.
\end{proof}

\begin{definition}{Gram matrix}{def:gram-matrix}
\index{Gram matrix}
Let $(v_1,\ldots,v_k)$ be a family of vectors in an inner product
space.  The \emph{Gram matrix}\index{Gram matrix} is the $k\times k$
matrix $G$ defined by
\[
  G_{ij} \deff \inner{v_i,v_j}.
\]
The family is orthonormal if and only if $G = I_k$.  The family is
linearly independent if and only if $\det(G)\neq 0$.
\end{definition}


% ============================================================
\section{Orthogonal and unitary matrices}\label{sec:orthogonal-unitary}
% ============================================================

\begin{definition}{Orthogonal matrix}{def:orthogonal-matrix}
\index{orthogonal matrix}\index{O(n)@$\Orth(n)$}
A matrix $Q\in\Mat_n(\R)$ is \emph{orthogonal} if
\[
  \transpose{Q}\, Q = I_n,
\]
or equivalently $\inv{Q} = \transpose{Q}$.  The set of all $n\times n$
orthogonal matrices is denoted $\Orth(n)$ and forms a group under
matrix multiplication, called the \emph{orthogonal group}.
\end{definition}

\begin{definition}{Unitary matrix}{def:unitary-matrix}
\index{unitary matrix}\index{U(n)@$U(n)$}
A matrix $U\in\Mat_n(\C)$ is \emph{unitary} if
\[
  \hermit{U}\, U = I_n,
\]
or equivalently $\inv{U} = \hermit{U}$.  The set of all $n\times n$
unitary matrices is the \emph{unitary group}~$U(n)$.
\end{definition}

\begin{proposition}{Characterizations of orthogonal/unitary matrices}%
{prop:orthogonal-characterizations}
\index{orthogonal matrix!characterizations}
Let $Q\in\Mat_n(\R)$.  The following are equivalent:
\begin{enumerate}[label=(\roman*)]
  \item $Q$ is orthogonal.
  \item The columns of~$Q$ form an orthonormal basis of~$\R^n$.
  \item The rows of~$Q$ form an orthonormal basis of~$\R^n$.
  \item $Q$ preserves the inner product:
    $\inner{Qx,Qy} = \inner{x,y}$ for all $x,y$.
  \item $Q$ preserves norms: $\norm{Qx} = \norm{x}$ for all~$x$
    (i.e.\ $Q$ is an \emph{isometry}\index{isometry}).
\end{enumerate}
The analogous statements hold for unitary matrices over~$\C$.
\end{proposition}

\begin{proof}
\textbf{(i)$\Leftrightarrow$(ii).}\quad
$\transpose{Q}Q = I$ means $\inner{q_i,q_j} = \delta_{ij}$, where
$q_1,\ldots,q_n$ are the columns of~$Q$.

\textbf{(i)$\Leftrightarrow$(iii).}\quad
$Q\transpose{Q} = I$ (which follows from $\transpose{Q}Q = I$ since
both are square) means the rows are orthonormal.

\textbf{(i)$\Rightarrow$(iv).}\quad
$\inner{Qx,Qy} = \transpose{(Qx)}Qy = \transpose{x}\transpose{Q}Qy
= \transpose{x}y = \inner{x,y}$.

\textbf{(iv)$\Rightarrow$(v).}\quad Take $y = x$.

\textbf{(v)$\Rightarrow$(i).}\quad
$\norm{Qx}^2 = \norm{x}^2$ for all~$x$ means
$\transpose{x}\transpose{Q}Qx = \transpose{x}x$, so
$\transpose{x}(\transpose{Q}Q - I)x = 0$ for all~$x$.  Since
$\transpose{Q}Q - I$ is symmetric, this implies
$\transpose{Q}Q = I$.
\end{proof}

\begin{proposition}{Determinant of orthogonal matrices}%
{prop:det-orthogonal}
If $Q\in\Orth(n)$, then $\det(Q) = \pm 1$.
\end{proposition}

\begin{proof}
$1 = \det(I) = \det(\transpose{Q}Q) = \det(\transpose{Q})\det(Q) =
\det(Q)^2$, so $\det(Q) = \pm 1$.
\end{proof}

\begin{definition}{Special orthogonal group}{def:SO}
\index{SO(n)@$\SO(n)$}\index{special orthogonal group}
The \emph{special orthogonal group} is
$\SO(n) \deff \setcond{Q\in\Orth(n)}{\det(Q) = 1}$.
Its elements are the orthogonal matrices that preserve orientation
(rotations in the geometric sense).
\end{definition}

\begin{example}{Orthogonal matrices in dimension~$2$}{ex:O2}
\index{orthogonal matrix!in dimension 2}
Every matrix in $\SO(2)$ is a rotation:
\[
  R_\theta = \begin{pmatrix} \cos\theta & -\sin\theta \\
    \sin\theta & \cos\theta \end{pmatrix},
  \qquad \theta\in[0,2\pi).
\]
The matrices in $\Orth(2)\setminus\SO(2)$ are reflections:
\[
  S_\theta = \begin{pmatrix} \cos\theta & \sin\theta \\
    \sin\theta & -\cos\theta \end{pmatrix}.
\]
\end{example}

\begin{example}{Permutation matrices}{ex:permutation-orthogonal}
\index{permutation matrix!orthogonality}
Every permutation matrix (obtained by permuting the columns of~$I_n$)
is orthogonal.  Its determinant is $+1$ for even permutations and $-1$
for odd permutations.
\end{example}


% ============================================================
\section{The adjoint operator}\label{sec:adjoint}
% ============================================================

\begin{definition}{Adjoint of a linear map}{def:adjoint}
\index{adjoint operator}\index{adjoint!of a linear map}
Let $(E,\inner{\cdot,\cdot})$ be a finite-dimensional inner product
space and let $f\in\End(E)$.  The \emph{adjoint} of~$f$ is the unique
linear map $f^*\in\End(E)$ satisfying
\[
  \inner{f(u),v} = \inner{u,f^*(v)}
  \qquad\text{for all } u,v\in E.
\]
\end{definition}

\begin{theorem}{Existence and uniqueness of the adjoint}%
{thm:adjoint-existence}
\index{adjoint operator!existence}
For every $f\in\End(E)$ on a finite-dimensional inner product space, the
adjoint $f^*$ exists and is unique.  If $(e_1,\ldots,e_n)$ is an
orthonormal basis and $A = \Mat_{\cB}(f)$ is the matrix of~$f$ in this
basis, then
\[
  \Mat_{\cB}(f^*) = \hermit{A} \quad
  (\text{i.e.\ } \transpose{\bar{A}} \text{ over } \C, \text{ or }
  \transpose{A} \text{ over } \R).
\]
\end{theorem}

\begin{proof}
\textbf{Uniqueness.}\quad Suppose $g,h\in\End(E)$ both satisfy
$\inner{f(u),v} = \inner{u,g(v)} = \inner{u,h(v)}$ for all $u,v$.
Then $\inner{u,g(v)-h(v)} = 0$ for all~$u$, which by non-degeneracy
gives $g(v) = h(v)$ for all~$v$.

\textbf{Existence.}\quad Let $(e_1,\ldots,e_n)$ be an orthonormal basis
and $A = (a_{ij})$ the matrix of~$f$ in this basis, so
$f(e_j) = \sum_i a_{ij}\, e_i$.  Define $g\in\End(E)$ by the matrix
$B = \hermit{A}$, i.e.\ $b_{ij} = \bar{a}_{ji}$.  Then for any two
basis vectors:
\[
  \inner{f(e_j),e_k} = a_{kj},
  \qquad
  \inner{e_j,g(e_k)} = \overline{b_{jk}} = \overline{\bar{a}_{kj}}
    = a_{kj}.
\]
By linearity, $\inner{f(u),v} = \inner{u,g(v)}$ for all $u,v$, so
$f^* = g$ and $\Mat_{\cB}(f^*) = \hermit{A}$.
\end{proof}

\begin{proposition}{Properties of the adjoint}{prop:adjoint-properties}
\index{adjoint operator!properties}
Let $f,g\in\End(E)$ and $\alpha\in\K$.  Then:
\begin{enumerate}[label=(\roman*)]
  \item $(f+g)^* = f^* + g^*$.
  \item $(\alpha f)^* = \bar{\alpha}\, f^*$.
  \item $(f\circ g)^* = g^*\circ f^*$.
  \item $(f^*)^* = f$.
  \item $\Ker(f^*) = (\im f)^{\perp}$.
  \item $\im(f^*) = (\Ker f)^{\perp}$.
\end{enumerate}
\end{proposition}

\begin{proof}
Properties (i)--(iv) follow from the characterizing identity
$\inner{f(u),v} = \inner{u,f^*(v)}$ and uniqueness of the adjoint.

\textbf{(v).}\quad $v\in\Ker(f^*)$ iff $f^*(v) = \vzero$ iff
$\inner{u,f^*(v)} = 0$ for all~$u$ iff $\inner{f(u),v} = 0$ for all~$u$
iff $v\in(\im f)^{\perp}$.

\textbf{(vi).}\quad Apply (v) to $f^*$ to get
$\Ker(f) = (\im f^*)^{\perp}$, then take orthogonal complements:
$\im(f^*) = (\im f^*)^{\perp\perp} = (\Ker f)^{\perp}$.
\end{proof}

\begin{definition}{Self-adjoint, normal, and skew-adjoint operators}%
{def:special-operators}
\index{self-adjoint operator}\index{normal operator}\index{skew-adjoint operator}
An operator $f\in\End(E)$ is called:
\begin{itemize}
  \item \emph{self-adjoint} (or \emph{Hermitian}\index{Hermitian operator})
    if $f^* = f$;
  \item \emph{skew-adjoint} (or \emph{anti-Hermitian}) if $f^* = -f$;
  \item \emph{normal} if $f^*\circ f = f\circ f^*$.
\end{itemize}
In matrix terms (in an ONB): self-adjoint means $A = \hermit{A}$
(symmetric if $\K = \R$), and normal means $\hermit{A}A = A\hermit{A}$.
\end{definition}

\begin{remark}{Orthogonal/unitary operators via the adjoint}%
{rem:orthogonal-adjoint}
An operator $f\in\End(E)$ satisfies $f^*\circ f = \Id$ (equivalently,
$f$ preserves the inner product) if and only if its matrix in an
orthonormal basis is orthogonal (over~$\R$) or unitary (over~$\C$).
In particular, orthogonal and unitary operators are normal.
\end{remark}


% ============================================================
\section{Applications}\label{sec:inner-product-applications}
% ============================================================

% --- Least squares ---

\subsection{Least squares approximation}\label{subsec:least-squares}
\index{least squares}

A fundamental application of orthogonal projection is the \emph{least
squares} method.  Given a system $Ax = b$ that may have no exact
solution (because $b\notin\im(A)$), we seek the vector $\hat{x}$ that
minimises $\norm{Ax - b}$.

\begin{theorem}{Normal equations}{thm:normal-equations}
\index{normal equations}
Let $A\in\Mat_{m\times n}(\R)$ with $m\geq n$ and $\rank(A) = n$.
The least squares solution of $Ax = b$ is the unique vector $\hat{x}$
satisfying
\[
  \transpose{A}\, A\, \hat{x} = \transpose{A}\, b.
\]
The minimum residual is $\norm{b - A\hat{x}} = \norm{b - \proj_{\im A}(b)}$.
\end{theorem}

\begin{proof}
The vector $A\hat{x}$ must be the orthogonal projection of~$b$
onto~$\im(A)$, i.e.\ $b - A\hat{x}\perp\im(A)$.  This means
$\inner{A_j, b - A\hat{x}} = 0$ for each column~$A_j$ of~$A$, which
is equivalent to $\transpose{A}(b - A\hat{x}) = 0$, i.e.\
$\transpose{A}A\hat{x} = \transpose{A}b$.  Since $\rank(A) = n$,
the matrix $\transpose{A}A$ is invertible (it is the Gram matrix of
the columns of~$A$), giving the unique solution
$\hat{x} = \inv{(\transpose{A}A)}\,\transpose{A}\, b$.
\end{proof}

% --- QR factorization preview ---

\subsection{QR factorization}\label{subsec:QR}
\index{QR factorization}

The Gram--Schmidt process has an elegant matrix interpretation.  If
$A\in\Mat_{m\times n}(\R)$ has linearly independent columns
$a_1,\ldots,a_n$, then Gram--Schmidt produces orthonormal vectors
$q_1,\ldots,q_n$ such that
\[
  \Span(a_1,\ldots,a_j) = \Span(q_1,\ldots,q_j)
  \qquad\text{for each } j.
\]
In matrix form, this means $A = QR$, where
\begin{itemize}
  \item $Q\in\Mat_{m\times n}(\R)$ has orthonormal columns
    $q_1,\ldots,q_n$;
  \item $R\in\Mat_n(\R)$ is upper triangular with positive diagonal
    entries, given by $r_{ij} = \inner{a_j,q_i}$ for $i\leq j$.
\end{itemize}
This is the \emph{QR factorization} (or QR decomposition) of~$A$.  It
is numerically more stable than directly solving the normal equations and
is a cornerstone of numerical linear algebra.

% --- Fourier series preview ---

\subsection{Fourier series}\label{subsec:fourier-preview}
\index{Fourier series}

Consider the inner product space $\mathcal{C}([-\pi,\pi],\R)$ with the
$L^2$ inner product $\inner{f,g} = \int_{-\pi}^{\pi}f(t)\,g(t)\,\mathrm{d}t$.
The family
\[
  \frac{1}{\sqrt{2\pi}},\;
  \frac{\cos t}{\sqrt{\pi}},\;
  \frac{\sin t}{\sqrt{\pi}},\;
  \frac{\cos 2t}{\sqrt{\pi}},\;
  \frac{\sin 2t}{\sqrt{\pi}},\;\ldots
\]
is orthonormal.  The orthogonal projection of a function~$f$ onto the
Span of the first $2N+1$ elements is the \emph{truncated Fourier
series}:
\[
  S_N(f)(t) = \frac{a_0}{2} + \sum_{k=1}^{N}\bigl(a_k\cos kt
    + b_k\sin kt\bigr),
\]
where $a_k = \frac{1}{\pi}\int_{-\pi}^{\pi}f(t)\cos kt\,\mathrm{d}t$
and $b_k = \frac{1}{\pi}\int_{-\pi}^{\pi}f(t)\sin kt\,\mathrm{d}t$
are the \emph{Fourier coefficients}\index{Fourier coefficients}.
By 
ef{thm:thm:best-approximation}, this is the best approximation of~$f$
by trigonometric polynomials of degree~$\leq N$ in the $L^2$ norm.
Bessel's inequality (
ef{thm:thm:bessel}) gives
\[
  \frac{a_0^2}{2} + \sum_{k=1}^{N}(a_k^2 + b_k^2)
    \leq \frac{1}{\pi}\int_{-\pi}^{\pi}f(t)^2\,\mathrm{d}t,
\]
and Parseval's identity (
ef{thm:thm:parseval}) asserts that equality
holds in the limit $N\to\infty$ (for $f\in L^2$).


% ============================================================
\section{Exercises}\label{sec:inner-product-exercises}
% ============================================================

\begin{exercise}{Verifying inner product axioms \basic}%
{exo:verify-inner-product}
\index{inner product!verification}
Determine which of the following are inner products on~$\R^2$:
\begin{enumerate}[label=(\alph*)]
  \item $\inner{x,y} = x_1 y_1 - x_1 y_2 - x_2 y_1 + 4 x_2 y_2$.
  \item $\inner{x,y} = x_1 y_1 + x_2 y_2 - x_1 y_2 - x_2 y_1$.
  \item $\inner{x,y} = 2x_1 y_1 + x_1 y_2 + x_2 y_1 + 2 x_2 y_2$.
\end{enumerate}
For each, write the matrix~$G$ and check symmetry and positive
definiteness.
\end{exercise}

\begin{exercise}{Cauchy--Schwarz for sums \basic}{exo:cs-sums}
\index{Cauchy--Schwarz inequality!for sums}
Using the standard inner product on~$\R^n$, deduce from the
Cauchy--Schwarz inequality that for all $a_1,\ldots,a_n,b_1,\ldots,b_n
\in\R$:
\[
  \Bigl(\sum_{i=1}^n a_i b_i\Bigr)^2
    \leq \Bigl(\sum_{i=1}^n a_i^2\Bigr)
         \Bigl(\sum_{i=1}^n b_i^2\Bigr).
\]
When does equality hold?
\end{exercise}

\begin{exercise}{Cauchy--Schwarz for integrals \basic}{exo:cs-integrals}
\index{Cauchy--Schwarz inequality!for integrals}
Deduce from the Cauchy--Schwarz inequality (applied to the $L^2$ inner
product on $\mathcal{C}([a,b])$) that for continuous $f,g$:
\[
  \Bigl(\int_a^b f(t)\,g(t)\,\mathrm{d}t\Bigr)^2
    \leq \int_a^b f(t)^2\,\mathrm{d}t\;\cdot\;
         \int_a^b g(t)^2\,\mathrm{d}t.
\]
\end{exercise}

\begin{exercise}{Orthogonal complement in $\R^3$ \basic}%
{exo:orth-complement-R3}
Let $F = \Span\bigl((1,1,0),\,(0,1,1)\bigr)\subset\R^3$ with the
standard inner product.  Find $F^{\perp}$ and verify that
$\R^3 = F\oplus F^{\perp}$.
\end{exercise}

\begin{exercise}{Gram--Schmidt in $\R^3$ \basic}{exo:gs-R3}
\index{Gram--Schmidt process!exercise}
Apply the Gram--Schmidt process to the basis
$v_1 = (1,0,1)$, $v_2 = (1,1,0)$, $v_3 = (0,1,1)$ of~$\R^3$.
\end{exercise}

\begin{exercise}{Orthogonal projection computation \basic}%
{exo:proj-computation}
In $\R^3$ with the standard inner product, compute the orthogonal
projection of $v = (1,2,3)$ onto the plane
$F = \setcond{(x_1,x_2,x_3)\in\R^3}{x_1 + x_2 + x_3 = 0}$.
\end{exercise}

\begin{exercise}{Orthogonal matrix properties \intermediate}%
{exo:orthogonal-properties}
\index{orthogonal matrix!exercise}
Let $Q\in\Orth(n)$.
\begin{enumerate}[label=(\alph*)]
  \item Show that $\det(Q) = \pm 1$.
  \item Show that if $\lambda$ is an eigenvalue of~$Q$ (possibly
    complex), then $\abs{\lambda} = 1$.
  \item Show that if $n$ is odd and $Q\in\SO(n)$, then $1$ is an
    eigenvalue of~$Q$ (hence $Q$ has a fixed axis).
\end{enumerate}
\end{exercise}

\begin{exercise}{Adjoint of a matrix \intermediate}{exo:adjoint-matrix}
\index{adjoint operator!exercise}
Let $A = \begin{pmatrix} 1 & 2 \\ 3 & 4 \end{pmatrix}\in\Mat_2(\R)$,
viewed as an endomorphism of $(\R^2,\inner{\cdot,\cdot})$ with the
standard inner product.
\begin{enumerate}[label=(\alph*)]
  \item Find $A^* = \transpose{A}$.
  \item Verify the identity $\inner{Ax,y} = \inner{x,A^* y}$ for
    $x = (1,0)$ and $y = (0,1)$.
  \item Is $A$ self-adjoint?  Is $A$ normal?
\end{enumerate}
\end{exercise}

\begin{exercise}{Least squares line fitting \intermediate}%
{exo:least-squares-line}
\index{least squares!line fitting}
Find the line $y = a + bx$ that best fits (in the least squares sense)
the data points $(0,1)$, $(1,0)$, $(2,3)$, $(3,2)$.

\noindent\textit{Hint:} Set up the system $Ax = b$ with
$A = \begin{psmallmatrix} 1 & 0 \\ 1 & 1 \\ 1 & 2 \\ 1 & 3
\end{psmallmatrix}$, $b = (1,0,3,2)^{\mathsf{T}}$, and solve the
normal equations.
\end{exercise}

\begin{exercise}{QR factorization \intermediate}{exo:qr-factorization}
\index{QR factorization!exercise}
Compute the QR factorization of the matrix
\[
  A = \begin{pmatrix} 1 & 1 \\ 1 & 0 \\ 0 & 1 \end{pmatrix}
\]
by applying the Gram--Schmidt process to the columns of~$A$.
Verify that $A = QR$.
\end{exercise}

\begin{exercise}{Isometries of $\R^2$ \intermediate}{exo:isometries-R2}
\index{isometry!classification in $\R^2$}
Show that every $Q\in\Orth(2)$ is either a rotation
$R_\theta\in\SO(2)$ or a reflection.  Describe the geometric action in
each case.
\end{exercise}

\begin{exercise}{Self-adjoint operators have real eigenvalues \intermediate}%
{exo:self-adjoint-eigenvalues}
\index{self-adjoint operator!eigenvalues}
Let $(E,\inner{\cdot,\cdot})$ be a complex inner product space and
$f\in\End(E)$ self-adjoint ($f^* = f$).
\begin{enumerate}[label=(\alph*)]
  \item Show that all eigenvalues of~$f$ are real.
  \item Show that eigenvectors corresponding to distinct eigenvalues
    are orthogonal.
\end{enumerate}
\end{exercise}

\begin{exercise}{Polarization identity \challenging}{exo:polarization}
\index{polarization identity!exercise}
Let $E$ be a real inner product space.  Prove that the inner product is
completely determined by the norm via the \emph{polarization identity}:
\[
  \inner{u,v} = \frac{1}{4}\bigl(\norm{u+v}^2 - \norm{u-v}^2\bigr).
\]
In the complex case, show the analogous identity:
\[
  \inner{u,v} = \frac{1}{4}\sum_{k=0}^{3} i^k\,\norm{u + i^k v}^2.
\]
\end{exercise}

\begin{exercise}{Orthogonal projection onto a line \challenging}%
{exo:proj-line}
Let $a\in\R^n$ be nonzero.  Show that the matrix of the orthogonal
projection onto $\Span(a)$ is
\[
  P = \frac{a\,\transpose{a}}{\transpose{a}\,a}.
\]
Verify that $P^2 = P$, $\transpose{P} = P$, and $\rank(P) = 1$.
\end{exercise}

\begin{exercise}{Unitary diagonalization of Hermitian matrices \challenging}%
{exo:unitary-diag-hermitian}
\index{Hermitian matrix!diagonalization}
Let $A = \begin{pmatrix} 2 & 1-i \\ 1+i & 3 \end{pmatrix}\in\Mat_2(\C)$.
\begin{enumerate}[label=(\alph*)]
  \item Verify that $A$ is Hermitian.
  \item Find the eigenvalues of~$A$ and show they are real.
  \item Find an orthonormal eigenbasis and write $A = U D\hermit{U}$
    for a unitary matrix~$U$ and diagonal matrix~$D$.
\end{enumerate}
\end{exercise}

\begin{exercise}{Fourier coefficients and best approximation \challenging}%
{exo:fourier-best-approx}
\index{Fourier series!exercise}
Let $f(t) = t$ on $[-\pi,\pi]$.
\begin{enumerate}[label=(\alph*)]
  \item Compute the Fourier coefficients $a_k$ and $b_k$ for
    $k = 0,1,2,3$.
  \item Write the best trigonometric polynomial approximation of
    degree~$\leq 3$ in the $L^2$ norm.
  \item Verify Bessel's inequality for $N = 3$.
\end{enumerate}
\end{exercise}


% ============================================================
\section{Chapter summary}\label{sec:inner-product-summary}
% ============================================================

\begin{itemize}
  \item A \textbf{bilinear form} $\varphi\colon E\times E\to\K$ is
    linear in each argument.  It is \textbf{symmetric} if
    $\varphi(u,v) = \varphi(v,u)$ and \textbf{non-degenerate} if its
    matrix is invertible.  The associated \textbf{quadratic form} is
    $q(v) = \varphi(v,v)$.

  \item \textbf{Sylvester's law of inertia}: every real quadratic form
    can be reduced to diagonal form $\diag(+1,\ldots,+1,-1,\ldots,-1,
    0,\ldots,0)$; the \textbf{signature} $(p,r-p)$ is an invariant.

  \item An \textbf{inner product} is a positive definite symmetric
    bilinear form (real case) or a positive definite Hermitian
    sesquilinear form (complex case).

  \item The \textbf{norm} $\norm{v} = \sqrt{\inner{v,v}}$ satisfies
    the \textbf{Cauchy--Schwarz inequality}
    $\abs{\inner{u,v}}\leq\norm{u}\norm{v}$ and the
    \textbf{triangle inequality}
    $\norm{u+v}\leq\norm{u}+\norm{v}$.

  \item Two vectors are \textbf{orthogonal} if $\inner{u,v} = 0$.  The
    \textbf{orthogonal complement} $F^{\perp}$ satisfies
    $E = F\oplus F^{\perp}$.

  \item The \textbf{orthogonal projection} $\proj_F(v)$ is the unique
    closest point to~$v$ in a subspace~$F$ (\textbf{best approximation
    property}).

  \item The \textbf{Gram--Schmidt process} turns any linearly
    independent family into an orthonormal one while preserving the
    successive Spans.  Every inner product space has an
    \textbf{orthonormal basis}.

  \item \textbf{Parseval's identity}: in an ONB,
    $\inner{u,v} = \sum_i\inner{u,e_i}\overline{\inner{v,e_i}}$.
    \textbf{Bessel's inequality}: partial sums of Fourier coefficients
    are bounded by $\norm{v}^2$.

  \item \textbf{Orthogonal matrices} ($\transpose{Q}Q = I$) preserve
    inner products and norms; they form the group $\Orth(n)$, with the
    subgroup $\SO(n)$ (rotations, $\det = 1$).  Over~$\C$, the
    analogues are \textbf{unitary matrices}.

  \item The \textbf{adjoint} $f^*$ of a linear map satisfies
    $\inner{f(u),v} = \inner{u,f^*(v)}$; in an ONB, its matrix is
    $\hermit{A}$.  Key identities:
    $\Ker(f^*) = (\im f)^{\perp}$ and
    $\im(f^*) = (\Ker f)^{\perp}$.

  \item \textbf{Applications}: \emph{least squares} (solve overdetermined
    systems via $\transpose{A}A\hat{x} = \transpose{A}b$), \emph{QR
    factorization} (matrix form of Gram--Schmidt), and \emph{Fourier
    series} (orthogonal projection onto trigonometric polynomials).
\end{itemize}

% ch7-spectral-theorem.tex — Chapter 7: Spectral Theorem and Symmetric Matrices
% Self-adjoint operators, spectral theorem (real and complex), orthogonal
% diagonalization, normal operators, positive definiteness, SVD, applications.

\chapter{Spectral Theorem and Symmetric Matrices}\label{ch:spectral}

% ============================================================
% Motivating introduction
% ============================================================

In the preceding chapters we studied eigenvalues, eigenvectors, and
diagonalization of general linear maps.  We saw that not every matrix
is diagonalizable, and that when diagonalization succeeds the change-of-basis
matrix need not enjoy any particular geometric property.  In this chapter
we restrict our attention to a distinguished class of operators ---
\emph{self-adjoint}\index{self-adjoint operator} (symmetric or Hermitian)
operators --- for which the theory becomes strikingly clean: every
self-adjoint operator is diagonalizable, all eigenvalues are real, and
one can find an \emph{orthonormal} basis of eigenvectors.

This is the content of the \emph{Spectral Theorem},\index{Spectral Theorem}
one of the most important results in all of linear algebra.  Its applications
pervade mathematics, science, and engineering:
\begin{itemize}
  \item \textbf{Covariance matrices.}  In statistics and data science, the
    covariance matrix of a random vector is symmetric and positive
    semi\-definite; its eigenvectors are the \emph{principal components}, and
    PCA\index{PCA}\index{principal component analysis} amounts to the spectral
    decomposition of this matrix.
  \item \textbf{Hessians and optimization.}  The Hessian of a twice
    differentiable function $f\colon\R^n\to\R$ is a real symmetric matrix;
    its eigenvalues determine whether a critical point is a local minimum,
    maximum, or saddle point.
  \item \textbf{Quadratic forms.}  The classification of quadratic forms
    (ellipses, hyperbolas, etc.)\ relies on the signs of the eigenvalues of
    the associated symmetric matrix.
  \item \textbf{Vibrations and normal modes.}  The natural frequencies of a
    vibrating system are the square roots of the eigenvalues of a symmetric
    matrix, and the normal modes are its eigenvectors.
  \item \textbf{Quantum mechanics.}  Observables are Hermitian operators on a
    Hilbert space; the Spectral Theorem guarantees that measurement outcomes
    (eigenvalues) are real numbers.
\end{itemize}

\begin{center}
\begin{tikzpicture}[>=Stealth, scale=1.2]
  % --- Conic classification via eigenvalues ---
  % Ellipse
  \begin{scope}[shift={(-4,0)}]
    \draw[thick, blue!70!black] (0,0) ellipse (1.5 and 0.8);
    \draw[->, thick, red!70!black] (0,0) -- (1.5,0)
      node[below right] {\small $\lambda_1>0$};
    \draw[->, thick, red!70!black] (0,0) -- (0,0.8)
      node[above right] {\small $\lambda_2>0$};
    \node[below] at (0,-1.2) {\small Ellipse};
  \end{scope}

  % Hyperbola
  \begin{scope}[shift={(0,0)}]
    \draw[thick, blue!70!black, domain=-1.2:1.2, samples=80]
      plot ({cosh(\x)}, {0.7*sinh(\x)});
    \draw[thick, blue!70!black, domain=-1.2:1.2, samples=80]
      plot ({-cosh(\x)}, {0.7*sinh(\x)});
    \draw[->, thick, red!70!black] (0,0) -- (1.3,0)
      node[below right] {\small $\lambda_1>0$};
    \draw[->, thick, green!60!black] (0,0) -- (0,0.9)
      node[above right] {\small $\lambda_2<0$};
    \node[below] at (0,-1.2) {\small Hyperbola};
  \end{scope}

  % Parabola (degenerate: one eigenvalue zero)
  \begin{scope}[shift={(4,0)}]
    \draw[thick, blue!70!black, domain=-1:1, samples=60]
      plot ({1.2*\x*\x}, {\x});
    \draw[->, thick, red!70!black] (0,0) -- (1.3,0)
      node[below right] {\small $\lambda_1>0$};
    \draw[->, thick, gray] (0,0) -- (0,0.9)
      node[above right] {\small $\lambda_2=0$};
    \node[below] at (0.5,-1.2) {\small Parabola};
  \end{scope}
\end{tikzpicture}
\captionof{figure}{Classification of conics by the eigenvalues of the
  associated symmetric matrix.  All eigenvalues positive: ellipse;
  mixed signs: hyperbola; one eigenvalue zero: parabola.}
\label{fig:conic-classification}
\end{center}

Throughout this chapter, $E$ denotes a finite-dimensional inner product
space over $\K = \R$ or~$\C$, with inner product $\inner{{\cdot},{\ \cdot}}$
and induced norm~$\norm{\cdot}$.  We write $\dim E = n$.  Matrices are
identified with their corresponding endomorphisms via a chosen basis
whenever convenient.


% ============================================================
\section{Self-adjoint (symmetric and Hermitian) operators}%
\label{sec:self-adjoint}
% ============================================================

\begin{definition}{Adjoint of a linear map}{def:adjoint-ch7}
\index{adjoint operator}
Let $(E,\inner{\cdot,\cdot})$ be a finite-dimensional inner product
space and $f\in\End(E)$.  The \emph{adjoint} of~$f$ is the unique
linear map $f^*\in\End(E)$ satisfying
\[
  \inner{f(u),\, v} \;=\; \inner{u,\, f^*(v)}
  \qquad\text{for all } u,v\in E.
\]
\end{definition}

\begin{remark}{Existence and uniqueness of the adjoint}{rem:adjoint-exists}
Existence and uniqueness follow from the Riesz representation theorem
(or can be verified directly using an orthonormal basis).  If $\cB$ is
an orthonormal basis and $A = \Mat_\cB(f)$, then
$\Mat_\cB(f^*) = \hermit{A}$, the conjugate transpose.  Over~$\R$
this reduces to $\Mat_\cB(f^*) = \transpose{A}$.
\end{remark}

\begin{proposition}{Properties of the adjoint}{prop:adjoint-properties-ch7}
\index{adjoint operator!properties}
For $f,g\in\End(E)$ and $\alpha\in\K$:
\begin{enumerate}[label=\textup{(\roman*)}]
  \item $(f+g)^* = f^*+g^*$.
  \item $(\alpha f)^* = \bar\alpha\, f^*$.
  \item $(f\circ g)^* = g^*\circ f^*$.
  \item $(f^*)^* = f$.
  \item $f$ is invertible if and only if $f^*$ is, and then
    $(f^*)^{-1}=(f^{-1})^*$.
\end{enumerate}
\end{proposition}

\begin{proof}
All properties follow from the definition and the sesquilinearity
(or bilinearity over~$\R$) of the inner product.  We verify~(iii):
for all $u,v\in E$,
\[
  \inner{(f\circ g)(u),\,v}
  = \inner{f(g(u)),\,v}
  = \inner{g(u),\,f^*(v)}
  = \inner{u,\,g^*(f^*(v))}
  = \inner{u,\,(g^*\circ f^*)(v)},
\]
so $(f\circ g)^* = g^*\circ f^*$ by uniqueness of the adjoint.
The remaining items are proved similarly.
\end{proof}

\begin{definition}{Self-adjoint operator}{def:self-adjoint}
\index{self-adjoint operator}\index{symmetric matrix}\index{Hermitian matrix}
An endomorphism $f\in\End(E)$ is \emph{self-adjoint} if $f^* = f$,
i.e.\ $\inner{f(u),v} = \inner{u,f(v)}$ for all $u,v\in E$.

In matrix terms with respect to an orthonormal basis:
\begin{itemize}
  \item Over $\R$: $A$ is \emph{symmetric},\index{symmetric matrix}
    i.e.\ $\transpose{A} = A$.
  \item Over $\C$: $A$ is \emph{Hermitian},\index{Hermitian matrix}
    i.e.\ $\hermit{A} = A$.
\end{itemize}

We denote the real vector space of $n\times n$ real symmetric matrices by
$\Sym_n(\R)\deff\setcond{A\in\Mat_n(\R)}{\transpose{A}=A}$.
\end{definition}

\begin{example}{Symmetric matrices}{ex:symmetric-examples}
The matrix
$A = \begin{pmatrix} 2 & -1 \\ -1 & 3 \end{pmatrix}$
is symmetric.  The matrix
$B = \begin{pmatrix} 1 & 2+i \\ 2-i & 3 \end{pmatrix}$
is Hermitian (the diagonal entries are real and the off-diagonal entries
are complex conjugates).
\end{example}

\begin{proposition}{Dimension of $\Sym_n(\R)$}{prop:dim-sym}
$\dim\Sym_n(\R) = \dfrac{n(n+1)}{2}$.
\end{proposition}

\begin{proof}
A symmetric matrix is determined by its entries on and above the diagonal.
There are $n$ diagonal entries and $\binom{n}{2}$ entries strictly above
the diagonal, giving $n + \binom{n}{2} = \frac{n(n+1)}{2}$.
\end{proof}


% ============================================================
\section{Eigenvalues of self-adjoint operators are real}%
\label{sec:eigenvalues-real}
% ============================================================

\begin{theorem}{Eigenvalues of self-adjoint operators are real}%
{thm:eigenvalues-real}
\index{eigenvalue!of self-adjoint operator}\index{self-adjoint operator!eigenvalues}
Let $f$ be a self-adjoint operator on an inner product space~$E$
(over $\R$ or~$\C$).  Then every eigenvalue of~$f$ is real.
\end{theorem}

\begin{proof}
Let $\lambda$ be an eigenvalue and $v\neq\vzero$ an eigenvector:
$f(v)=\lambda v$.  Then
\[
  \lambda\,\inner{v,v}
  = \inner{\lambda v,\,v}
  = \inner{f(v),\,v}
  = \inner{v,\,f(v)}
  = \inner{v,\,\lambda v}
  = \bar\lambda\,\inner{v,v}.
\]
Since $v\neq\vzero$, we have $\inner{v,v}>0$, so dividing yields
$\lambda = \bar\lambda$, which means $\lambda\in\R$.
\end{proof}

\begin{remark}{Real eigenvalues over $\R$}{rem:real-eigenvalues-R}
Over~$\R$, one must be slightly careful: a real symmetric matrix might
a priori have no eigenvalues in~$\R$ (since the characteristic polynomial
need not split).  However, the Spectral Theorem will establish that it
does.  The above proof works directly if we allow eigenvalues in~$\C$
and use the standard Hermitian inner product on~$\C^n$.  Concretely:
if $A\in\Sym_n(\R)$ and $\lambda\in\C$ satisfies $Av = \lambda v$ for
some $v\in\C^n\setminus\set{\vzero}$, then $\lambda\in\R$.
\end{remark}


% ============================================================
\section{Orthogonality of eigenvectors for distinct eigenvalues}%
\label{sec:eigenvectors-orthogonal}
% ============================================================

\begin{theorem}{Eigenvectors for distinct eigenvalues are orthogonal}%
{thm:eigenvectors-orthogonal}
\index{eigenvector!orthogonality}\index{self-adjoint operator!orthogonal eigenvectors}
Let $f$ be a self-adjoint operator on~$E$, and let $v_1,v_2$ be
eigenvectors of~$f$ associated with distinct eigenvalues
$\lambda_1\neq\lambda_2$.  Then $\inner{v_1,v_2}=0$.
\end{theorem}

\begin{proof}
We compute $\inner{f(v_1),v_2}$ in two ways:
\[
  \inner{f(v_1),\,v_2}
  = \inner{\lambda_1 v_1,\,v_2}
  = \lambda_1\,\inner{v_1,v_2},
\]
and, using the self-adjointness of~$f$,
\[
  \inner{f(v_1),\,v_2}
  = \inner{v_1,\,f(v_2)}
  = \inner{v_1,\,\lambda_2 v_2}
  = \bar\lambda_2\,\inner{v_1,v_2}
  = \lambda_2\,\inner{v_1,v_2},
\]
where the last equality uses the fact that $\lambda_2\in\R$
(
ef{thm:thm:eigenvalues-real}).  Subtracting:
\[
  (\lambda_1 - \lambda_2)\,\inner{v_1,v_2} = 0.
\]
Since $\lambda_1\neq\lambda_2$, we conclude $\inner{v_1,v_2}=0$.
\end{proof}

\begin{corollary}{Eigenspaces of a self-adjoint operator are mutually
orthogonal}{cor:eigenspaces-orthogonal}
\index{eigenspace!orthogonality}
If $f$ is self-adjoint with distinct eigenvalues
$\lambda_1,\ldots,\lambda_r$, then
\[
  E_{\lambda_i} \perp E_{\lambda_j}
  \qquad\text{for all } i\neq j,
\]
where $E_{\lambda_k} = \Ker(f - \lambda_k\Id)$ is the eigenspace for
$\lambda_k$.
\end{corollary}

\begin{proof}
Every vector in $E_{\lambda_i}$ is an eigenvector for~$\lambda_i$
(or zero), and likewise for~$E_{\lambda_j}$.  The result follows from

ef{thm:thm:eigenvectors-orthogonal} applied to each pair of nonzero vectors.
\end{proof}


% ============================================================
\section{The Spectral Theorem for real symmetric matrices}%
\label{sec:spectral-real}
% ============================================================

We now prove the central result of this chapter.

\begin{lemma}{Invariance of the orthogonal complement}%
{lem:orth-complement-invariant}
\index{orthogonal complement!invariance under self-adjoint maps}
Let $f$ be a self-adjoint operator on~$E$ and let $W\subseteq E$ be a
subspace invariant under~$f$ (i.e.\ $f(W)\subseteq W$).  Then the
orthogonal complement $W^\perp$ is also invariant under~$f$.
\end{lemma}

\begin{proof}
Let $v\in W^\perp$.  We must show $f(v)\in W^\perp$, i.e.\
$\inner{f(v),w}=0$ for all $w\in W$.  Since $f$ is self-adjoint,
\[
  \inner{f(v),\,w} = \inner{v,\,f(w)}.
\]
Now $w\in W$ and $f(W)\subseteq W$, so $f(w)\in W$.  Since
$v\in W^\perp$, we get $\inner{v,f(w)}=0$.  Hence $\inner{f(v),w}=0$.
\end{proof}

\begin{lemma}{A self-adjoint operator over $\R$ has a real eigenvalue}%
{lem:self-adjoint-has-eigenvalue}
\index{self-adjoint operator!existence of eigenvalue}
Let $A\in\Sym_n(\R)$.  Then $A$ has at least one (real) eigenvalue.
\end{lemma}

\begin{proof}
The characteristic polynomial $\chi_A(\lambda) = \det(A - \lambda I_n)$
has degree~$n$ with real coefficients.  Over~$\C$ it has a root
$\lambda_0\in\C$, and by 
ef{thm:thm:eigenvalues-real} (applied in the
Hermitian setting) $\lambda_0$ is real.  Hence $A$ has a real eigenvalue.
\end{proof}

\begin{theorem}{Spectral Theorem for real symmetric matrices}%
{thm:spectral-real}
\index{Spectral Theorem!real symmetric matrices}
Let $A\in\Sym_n(\R)$ be a real symmetric matrix.  Then:
\begin{enumerate}[label=\textup{(\roman*)}]
  \item All eigenvalues of~$A$ are real.
  \item There exists an orthogonal matrix $P\in\Orth(n)$
    (i.e.\ $\transpose{P}P = I_n$) such that
    \[
      \transpose{P}\,A\,P = \diag(\lambda_1,\ldots,\lambda_n),
    \]
    where $\lambda_1,\ldots,\lambda_n$ are the eigenvalues of~$A$
    (counted with multiplicity).
  \item Equivalently, $\R^n$ has an orthonormal basis of eigenvectors
    of~$A$.
\end{enumerate}
\end{theorem}

\begin{proof}
We prove by strong induction on~$n$.

\medskip\noindent\textbf{Base case ($n=1$).}\quad
A $1\times 1$ matrix $A=(a)$ is already diagonal with
eigenvalue~$a\in\R$, and $P=(1)\in\Orth(1)$.

\medskip\noindent\textbf{Inductive step.}\quad
Assume the theorem holds for all real symmetric matrices of size less
than~$n$.  Let $A\in\Sym_n(\R)$.

\medskip\noindent\emph{Step~1: Find a real eigenvalue and unit eigenvector.}\quad
By 
ef{lem:lem:self-adjoint-has-eigenvalue}, $A$ has a real eigenvalue
$\lambda_1\in\R$.  Let $v_1$ be an eigenvector for~$\lambda_1$ with
$\norm{v_1}=1$.

\medskip\noindent\emph{Step~2: Set up the orthogonal complement.}\quad
Let $W = \Span(v_1)$.  Since $A v_1 = \lambda_1 v_1$, the subspace~$W$
is $A$-invariant.  By 
ef{lem:lem:orth-complement-invariant}, $W^\perp$
is also $A$-invariant, and $\dim W^\perp = n-1$.

\medskip\noindent\emph{Step~3: Restrict and apply the induction hypothesis.}\quad
Let $g = f\restrict{W^\perp}$, where $f$ is the linear map
$x\mapsto Ax$.  Since $f(W^\perp)\subseteq W^\perp$, the map~$g$ is a
well-defined endomorphism of~$W^\perp$.  Moreover, $g$ is self-adjoint
on~$W^\perp$ (the inner product is the restriction of the standard one,
and the self-adjointness condition is inherited).

Choose any orthonormal basis of~$W^\perp$ and let $B$ be the matrix
of~$g$ in this basis.  Then $B\in\Sym_{n-1}(\R)$.  By the induction
hypothesis, there exists an orthonormal basis
$(v_2,\ldots,v_n)$ of~$W^\perp$ consisting of eigenvectors of~$g$
(equivalently, of~$f$), with real eigenvalues
$\lambda_2,\ldots,\lambda_n$.

\medskip\noindent\emph{Step~4: Assemble the orthonormal eigenbasis.}\quad
The family $(v_1,v_2,\ldots,v_n)$ is orthonormal:
\begin{itemize}
  \item $v_1\perp v_j$ for $j\geq 2$, since $v_j\in W^\perp$;
  \item $(v_2,\ldots,v_n)$ is orthonormal by the induction hypothesis;
  \item $\norm{v_1}=1$ by construction.
\end{itemize}
Each $v_j$ is an eigenvector of~$A$ with eigenvalue~$\lambda_j$.
Let $P = \begin{pmatrix} v_1 & v_2 & \cdots & v_n \end{pmatrix}$.
Then $P\in\Orth(n)$ and $\transpose{P}AP = \diag(\lambda_1,\ldots,\lambda_n)$.
\end{proof}

\begin{remark}{Spectral decomposition form}{rem:spectral-decomposition}
\index{spectral decomposition}
The Spectral Theorem can be rephrased as follows.  If $A\in\Sym_n(\R)$
has distinct eigenvalues $\mu_1,\ldots,\mu_r$ with orthogonal projections
$P_i$ onto the eigenspace $E_{\mu_i}$, then
\[
  A = \mu_1 P_1 + \mu_2 P_2 + \cdots + \mu_r P_r,
\]
where $P_1 + \cdots + P_r = I_n$ and $P_i P_j = 0$ for $i\neq j$.
This is called the \emph{spectral decomposition}\index{spectral decomposition}
of~$A$.
\end{remark}


% ============================================================
\section{The Spectral Theorem for Hermitian matrices}%
\label{sec:spectral-hermitian}
% ============================================================

The real Spectral Theorem extends naturally to the complex setting.

\begin{theorem}{Spectral Theorem for Hermitian matrices}%
{thm:spectral-hermitian}
\index{Spectral Theorem!Hermitian matrices}
Let $A\in\Mat_n(\C)$ be Hermitian, i.e.\ $\hermit{A}=A$.  Then:
\begin{enumerate}[label=\textup{(\roman*)}]
  \item All eigenvalues of~$A$ are real.
  \item There exists a unitary matrix $U\in\Mat_n(\C)$
    (i.e.\ $\hermit{U}U = I_n$) such that
    \[
      \hermit{U}\,A\,U = \diag(\lambda_1,\ldots,\lambda_n),
    \]
    with $\lambda_1,\ldots,\lambda_n\in\R$.
  \item Equivalently, $\C^n$ has an orthonormal basis of eigenvectors
    of~$A$.
\end{enumerate}
\end{theorem}

\begin{proof}
The proof is identical to the real case: over~$\C$ the characteristic
polynomial always splits, so the base step of the induction is
guaranteed.  The self-adjointness condition $\hermit{A}=A$ ensures that
eigenvalues are real (
ef{thm:thm:eigenvalues-real}), eigenvectors for
distinct eigenvalues are orthogonal
(
ef{thm:thm:eigenvectors-orthogonal}), and the orthogonal complement
of an invariant subspace is invariant
(
ef{lem:lem:orth-complement-invariant}).  The inductive argument carries
over verbatim, with ``orthogonal matrix'' replaced by ``unitary matrix''
and ``transpose'' replaced by ``conjugate transpose.''
\end{proof}


% ============================================================
\section{Orthogonal diagonalization: procedure and examples}%
\label{sec:orth-diag}
% ============================================================

The Spectral Theorem provides both a theoretical guarantee and a
practical algorithm for orthogonally diagonalizing a symmetric matrix.

\begin{remark}{Orthogonal diagonalization procedure}{rem:orth-diag-procedure}
\index{orthogonal diagonalization!procedure}
Given $A\in\Sym_n(\R)$:
\begin{enumerate}
  \item Compute the characteristic polynomial
    $\chi_A(\lambda) = \det(A - \lambda I_n)$.
  \item Find all eigenvalues $\lambda_1,\ldots,\lambda_r$ (they are real).
  \item For each~$\lambda_i$, solve $(A - \lambda_i I_n)X = \vzero$ to
    find a basis of the eigenspace $E_{\lambda_i}$.
  \item Apply the Gram--Schmidt process within each eigenspace to obtain
    an orthonormal basis of~$E_{\lambda_i}$.
    (Eigenvectors from \emph{different} eigenspaces are automatically
    orthogonal, so Gram--Schmidt is only needed within a single
    eigenspace when $\dim E_{\lambda_i}\geq 2$.)
  \item Form $P$ by placing the orthonormal eigenvectors as columns.
    Then $P\in\Orth(n)$ and $\transpose{P}AP = \diag(\lambda_1,\ldots,\lambda_n)$.
\end{enumerate}
\end{remark}

\begin{example}{Orthogonal diagonalization of a $2\times 2$ matrix}%
{ex:orth-diag-2x2}
Diagonalize $A = \begin{pmatrix} 5 & 2 \\ 2 & 2 \end{pmatrix}$
orthogonally.

\emph{Step~1.} $\chi_A(\lambda) = (5-\lambda)(2-\lambda)-4
= \lambda^2-7\lambda+6 = (\lambda-1)(\lambda-6)$.

\emph{Step~2.} Eigenvalues: $\lambda_1=1$, $\lambda_2=6$.

\emph{Step~3.}
For $\lambda_1=1$: $(A-I_2)X=\vzero$ gives
$\begin{pmatrix}4&2\\2&1\end{pmatrix}X=\vzero$, so $2x+y=0$.
Eigenvector: $v_1 = \begin{psmallmatrix}1\\-2\end{psmallmatrix}$.

For $\lambda_2=6$: $(A-6I_2)X=\vzero$ gives
$\begin{pmatrix}-1&2\\2&-4\end{pmatrix}X=\vzero$, so $x=2y$.
Eigenvector: $v_2 = \begin{psmallmatrix}2\\1\end{psmallmatrix}$.

\emph{Step~4.} Normalise:
$\hat v_1 = \frac{1}{\sqrt{5}}\begin{psmallmatrix}1\\-2\end{psmallmatrix}$,
$\hat v_2 = \frac{1}{\sqrt{5}}\begin{psmallmatrix}2\\1\end{psmallmatrix}$.

\emph{Step~5.}
$P = \frac{1}{\sqrt{5}}\begin{pmatrix}1&2\\-2&1\end{pmatrix}$
satisfies $\transpose{P}AP = \begin{pmatrix}1&0\\0&6\end{pmatrix}$.
\end{example}

\begin{example}{Orthogonal diagonalization of a $3\times 3$ matrix}%
{ex:orth-diag-3x3}
Let
$A = \begin{pmatrix} 2 & 1 & 1 \\ 1 & 2 & 1 \\ 1 & 1 & 2 \end{pmatrix}$.
One verifies that $\transpose{A}=A$.

\emph{Characteristic polynomial:}
$\chi_A(\lambda) = -(\lambda-4)(\lambda-1)^2$.

\emph{Eigenvalues:} $\lambda_1=4$ (multiplicity~$1$),
$\lambda_2=1$ (multiplicity~$2$).

\emph{For $\lambda_1=4$:}
$(A-4I_3)X=\vzero$ gives $x_1=x_2=x_3$, so
$E_4 = \Span\!\left(\begin{psmallmatrix}1\\1\\1\end{psmallmatrix}\right)$.
Normalise: $u_1 = \frac{1}{\sqrt{3}}\begin{psmallmatrix}1\\1\\1\end{psmallmatrix}$.

\emph{For $\lambda_2=1$:}
$(A-I_3)X=\vzero$ gives $x_1+x_2+x_3=0$, so
$E_1 = \Span\!\left(
  \begin{psmallmatrix}1\\-1\\0\end{psmallmatrix},\,
  \begin{psmallmatrix}1\\0\\-1\end{psmallmatrix}
\right)$.
Apply Gram--Schmidt:
$w_1 = \begin{psmallmatrix}1\\-1\\0\end{psmallmatrix}$,
$w_2 = \begin{psmallmatrix}1\\0\\-1\end{psmallmatrix}
  - \frac{\inner{w_1,\begin{psmallmatrix}1\\0\\-1\end{psmallmatrix}}}{\inner{w_1,w_1}}\,w_1
= \begin{psmallmatrix}1/2\\1/2\\-1\end{psmallmatrix}$.

Normalise:
$u_2 = \frac{1}{\sqrt{2}}\begin{psmallmatrix}1\\-1\\0\end{psmallmatrix}$,
$u_3 = \frac{1}{\sqrt{6}}\begin{psmallmatrix}1\\1\\-2\end{psmallmatrix}$.

The orthogonal matrix is
$P = \begin{pmatrix}
\frac{1}{\sqrt{3}} & \frac{1}{\sqrt{2}} & \frac{1}{\sqrt{6}} \\[4pt]
\frac{1}{\sqrt{3}} & -\frac{1}{\sqrt{2}} & \frac{1}{\sqrt{6}} \\[4pt]
\frac{1}{\sqrt{3}} & 0 & -\frac{2}{\sqrt{6}}
\end{pmatrix}$,
and $\transpose{P}AP = \diag(4,1,1)$.
\end{example}


% ============================================================
\section{Normal operators}\label{sec:normal}
% ============================================================

Self-adjoint operators are a special case of a broader class.

\begin{definition}{Normal operator}{def:normal}
\index{normal operator}\index{normal matrix}
An operator $f\in\End(E)$ is \emph{normal} if it commutes with its
adjoint:
\[
  f\circ f^* = f^*\circ f.
\]
A matrix $A\in\Mat_n(\C)$ is \emph{normal} if $A\hermit{A} = \hermit{A}A$.
\end{definition}

\begin{example}{Classes of normal operators}{ex:normal-classes}
The following are all normal:
\begin{itemize}
  \item Self-adjoint (symmetric / Hermitian) operators: $f^*=f$.
  \item Skew-adjoint operators: $f^*=-f$ (equivalently,
    $\transpose{A}=-A$ over~$\R$ or $\hermit{A}=-A$ over~$\C$).
  \item Orthogonal / unitary operators: $f^*=\inv{f}$.
\end{itemize}
\end{example}

\begin{proposition}{Characterization of normality}{prop:normal-char}
Let $f\in\End(E)$ where $E$ is a finite-dimensional inner product space
over~$\C$.  The following are equivalent:
\begin{enumerate}[label=\textup{(\roman*)}]
  \item $f$ is normal.
  \item $\norm{f(v)} = \norm{f^*(v)}$ for all $v\in E$.
  \item If $f(v)=\lambda v$, then $f^*(v)=\bar\lambda v$.
    (That is, $v$ is an eigenvector of~$f$ if and only if it is an
    eigenvector of~$f^*$, with conjugate eigenvalue.)
\end{enumerate}
\end{proposition}

\begin{proof}
(i)$\Rightarrow$(ii): $\norm{f(v)}^2 = \inner{f(v),f(v)}
= \inner{f^*f(v),v} = \inner{ff^*(v),v} = \inner{f^*(v),f^*(v)}
= \norm{f^*(v)}^2$.

(ii)$\Rightarrow$(iii): Assume $f(v)=\lambda v$.  Apply~(ii) to
$f-\lambda\Id$ (which is also normal, as one verifies):
$\norm{(f-\lambda\Id)(v)}=\norm{(f-\lambda\Id)^*(v)}
=\norm{(f^*-\bar\lambda\Id)(v)}$.
The left side is~$0$, so $f^*(v)=\bar\lambda v$.

(iii)$\Rightarrow$(i): We show $f^*f = ff^*$ by checking on an
eigenbasis.  Over~$\C$, every normal operator is unitarily
diagonalizable (see the next theorem), so this direction also follows.
Alternatively, one can verify
$\inner{(f^*f-ff^*)(v),v}=0$ for all~$v$ using~(iii) and polarization.
\end{proof}

\begin{theorem}{Spectral Theorem for normal operators over $\C$}%
{thm:spectral-normal}
\index{Spectral Theorem!normal operators}
Let $A\in\Mat_n(\C)$ be normal.  Then there exists a unitary matrix
$U$ such that
\[
  \hermit{U}\,A\,U = \diag(\lambda_1,\ldots,\lambda_n),
\]
where $\lambda_1,\ldots,\lambda_n\in\C$ are the eigenvalues of~$A$.

Conversely, every unitarily diagonalizable matrix is normal.
\end{theorem}

\begin{proof}
\emph{Forward direction.}  By induction on~$n$.  The case $n=1$ is
trivial.  For the inductive step, since $\C$ is algebraically closed,
$A$ has an eigenvalue $\lambda_1$ with unit eigenvector~$v_1$.
Let $W=\Span(v_1)$.

\emph{Claim:} $W^\perp$ is $A$-invariant.  Let $u\in W^\perp$.  We need
$Au\perp v_1$, i.e.\ $\inner{Au,v_1}=0$.  Now
$\inner{Au,v_1} = \inner{u,\hermit{A}v_1} = \inner{u,\bar\lambda_1 v_1}
= \lambda_1\inner{u,v_1} = 0$,
where we used 
ef{prop:prop:normal-char}~(iii): since $Av_1=\lambda_1 v_1$,
we have $\hermit{A}v_1=\bar\lambda_1 v_1$.

The restriction $A\restrict{W^\perp}$ is normal on the $(n-1)$-dimensional
space~$W^\perp$.  By the induction hypothesis it is unitarily
diagonalizable.  Assembling the eigenvectors gives a unitary matrix~$U$
diagonalizing~$A$.

\emph{Converse.}  If $A = U D \hermit{U}$ with $D$ diagonal and $U$
unitary, then $\hermit{A} = U\bar D\hermit{U}$.  Since diagonal
matrices commute, $A\hermit{A} = UD\bar D\hermit{U}
= U\bar D D\hermit{U} = \hermit{A}A$.
\end{proof}

\begin{remark}{Normal operators over $\R$}{rem:normal-real}
\index{normal operator!over $\R$}
Over~$\R$, normal operators need not be orthogonally diagonalizable.
For instance, the rotation matrix
$R_\theta = \begin{psmallmatrix}\cos\theta & -\sin\theta \\
\sin\theta & \cos\theta\end{psmallmatrix}$
is normal (in fact orthogonal) but has no real eigenvalues when
$\theta\notin\set{0,\pi}$.  The real spectral theorem for normal
operators states that such a matrix can be brought to a block-diagonal
form with $1\times 1$ real blocks and $2\times 2$ rotation blocks.
\end{remark}


% ============================================================
\section{Positive definite and positive semidefinite matrices}%
\label{sec:positive-definite}
% ============================================================

\begin{definition}{Positive definite and positive semidefinite matrices}%
{def:positive-definite}
\index{positive definite matrix}\index{positive semidefinite matrix}
A real symmetric matrix $A\in\Sym_n(\R)$ is:
\begin{itemize}
  \item \emph{positive definite}\index{positive definite matrix}
    if $\transpose{x}Ax > 0$ for all $x\in\R^n\setminus\set{\vzero}$;
  \item \emph{positive semidefinite}\index{positive semidefinite matrix}
    if $\transpose{x}Ax \geq 0$ for all $x\in\R^n$;
  \item \emph{negative definite} if $-A$ is positive definite;
  \item \emph{negative semidefinite} if $-A$ is positive semidefinite;
  \item \emph{indefinite} if $\transpose{x}Ax$ takes both positive and
    negative values.
\end{itemize}
\end{definition}

\begin{theorem}{Characterizations of positive definiteness}%
{thm:pd-characterizations}
\index{positive definite matrix!characterizations}
For $A\in\Sym_n(\R)$ the following are equivalent:
\begin{enumerate}[label=\textup{(\roman*)}]
  \item $A$ is positive definite.
  \item All eigenvalues of~$A$ are strictly positive.
  \item All leading principal minors of~$A$ are strictly positive
    (\emph{Sylvester's criterion}\index{Sylvester's criterion}).
  \item There exists an invertible matrix $L\in\Mat_n(\R)$ such that
    $A = \transpose{L}\,L$.
  \item $A$ admits a \emph{Cholesky decomposition}\index{Cholesky decomposition}:
    there exists a unique lower triangular matrix $L$ with strictly positive
    diagonal entries such that $A = L\transpose{L}$.
\end{enumerate}
\end{theorem}

\begin{proof}
(i)$\Rightarrow$(ii): If $\lambda$ is an eigenvalue with eigenvector~$v$,
then $0 < \transpose{v}Av = \lambda\transpose{v}v = \lambda\norm{v}^2$,
so $\lambda > 0$.

(ii)$\Rightarrow$(i): By the Spectral Theorem, $A = P\,D\,\transpose{P}$
with $P$ orthogonal and $D = \diag(\lambda_1,\ldots,\lambda_n)$.  For
$x\neq\vzero$, let $y = \transpose{P}x \neq \vzero$.  Then
$\transpose{x}Ax = \transpose{y}Dy = \sum_{i=1}^n \lambda_i y_i^2 > 0$
since all $\lambda_i>0$ and $y\neq\vzero$.

(i)$\Rightarrow$(iii) (\emph{Sylvester's criterion}): The $k$-th
leading principal minor $\Delta_k$ equals the determinant of the
$k\times k$ upper-left submatrix~$A_k$.  But $A_k$ is itself symmetric,
and for any $x\in\R^k\setminus\set{\vzero}$ we have
$\transpose{x}A_k x = \transpose{\tilde x}A\tilde x > 0$, where
$\tilde x = (x_1,\ldots,x_k,0,\ldots,0)^T$.  Hence $A_k$ is positive
definite, so its eigenvalues are positive, so
$\Delta_k = \det(A_k) > 0$.

(iii)$\Rightarrow$(i): By induction on~$n$.  For $n=1$, $\Delta_1 = a_{11} > 0$,
so $A=(a_{11})$ is positive definite.  For the inductive step, write
$A = \begin{psmallmatrix} A_{n-1} & b \\ \transpose{b} & a_{nn}\end{psmallmatrix}$.
By induction, $A_{n-1}$ is positive definite.  Performing one step of block
Gaussian elimination (valid since $\det(A_{n-1})=\Delta_{n-1}>0$):
\[
  A = \begin{pmatrix} I & 0 \\ \transpose{b}\inv{A_{n-1}} & 1\end{pmatrix}
  \begin{pmatrix} A_{n-1} & 0 \\ 0 & s\end{pmatrix}
  \begin{pmatrix} I & \inv{A_{n-1}}b \\ 0 & 1\end{pmatrix},
\]
where $s = a_{nn} - \transpose{b}\inv{A_{n-1}}b
= \det(A)/\det(A_{n-1}) = \Delta_n/\Delta_{n-1} > 0$.
This expresses $A = \transpose{M}
\begin{psmallmatrix}A_{n-1}&0\\0&s\end{psmallmatrix} M$ with $M$ invertible.
For $x\neq\vzero$, $\transpose{x}Ax = \transpose{(Mx)}
\begin{psmallmatrix}A_{n-1}&0\\0&s\end{psmallmatrix}(Mx) > 0$
since the block-diagonal matrix is positive definite.

(ii)$\Rightarrow$(iv): By the Spectral Theorem, $A = PD\transpose{P}$,
so $A = PD^{1/2}\transpose{(PD^{1/2})} = \transpose{L}L$ where
$L = D^{1/2}\transpose{P}$.

(iv)$\Rightarrow$(i): $\transpose{x}Ax = \transpose{x}\transpose{L}Lx
= \norm{Lx}^2 > 0$ since $L$ is invertible and $x\neq\vzero$.

The existence and uniqueness of the Cholesky factorization~(v) is proved
by an explicit recursive construction using the positive leading
principal minors (we omit the details).
\end{proof}

\begin{example}{Testing positive definiteness}{ex:pd-test}
Consider $A = \begin{pmatrix}2 & -1 \\ -1 & 2\end{pmatrix}$.

\emph{Via eigenvalues:} $\chi_A(\lambda)=(\lambda-1)(\lambda-3)$, so
$\lambda_1=1>0$ and $\lambda_2=3>0$.  Hence $A$ is positive definite.

\emph{Via Sylvester's criterion:} $\Delta_1=2>0$, $\Delta_2=\det(A)=3>0$.
Positive definite.

\emph{Cholesky decomposition:}
$A = \begin{pmatrix}\sqrt{2}&0\\-\frac{1}{\sqrt{2}}&\frac{\sqrt{3}}{\sqrt{2}}\end{pmatrix}
\begin{pmatrix}\sqrt{2}&-\frac{1}{\sqrt{2}}\\0&\frac{\sqrt{3}}{\sqrt{2}}\end{pmatrix}
= L\transpose{L}$.
\end{example}

\begin{proposition}{Characterization of positive semidefiniteness}%
{prop:psd-characterizations}
\index{positive semidefinite matrix!characterizations}
For $A\in\Sym_n(\R)$ the following are equivalent:
\begin{enumerate}[label=\textup{(\roman*)}]
  \item $A$ is positive semidefinite.
  \item All eigenvalues of~$A$ are nonnegative.
  \item There exists a matrix $B\in\Mat_n(\R)$ such that $A=\transpose{B}B$.
  \item All principal minors of~$A$ are nonnegative.
\end{enumerate}
\end{proposition}

\begin{proof}
The proofs are analogous to 
ef{thm:thm:pd-characterizations}, with
strict inequalities relaxed to non-strict ones.  For (ii)$\Rightarrow$(iii),
use the spectral decomposition $A = PD\transpose{P}$ and set
$B = D^{1/2}\transpose{P}$ (noting $D^{1/2}$ is well-defined since all
diagonal entries are nonnegative).
\end{proof}


% ============================================================
\section{Singular Value Decomposition}\label{sec:svd}
% ============================================================

The Spectral Theorem applies to symmetric (or Hermitian) matrices.  The
\emph{Singular Value Decomposition} (SVD)\index{SVD}\index{singular value decomposition}
extends the idea of diagonalization to \emph{arbitrary} matrices,
including rectangular ones.

\begin{definition}{Singular values}{def:singular-values}
\index{singular value}
Let $A\in\Mat_{m\times n}(\R)$.  The \emph{singular values} of~$A$
are the nonnegative square roots of the eigenvalues of the positive
semidefinite matrix $\transpose{A}A\in\Sym_n(\R)$:
\[
  \sigma_i \deff \sqrt{\lambda_i(\transpose{A}A)},
  \qquad i = 1,\ldots,n,
\]
ordered so that $\sigma_1\geq\sigma_2\geq\cdots\geq\sigma_n\geq 0$.
\end{definition}

\begin{theorem}{Singular Value Decomposition}{thm:svd}
\index{SVD!theorem}\index{singular value decomposition!theorem}
Let $A\in\Mat_{m\times n}(\R)$ with $\rank(A)=r$.  Then there exist:
\begin{itemize}
  \item an orthogonal matrix $U\in\Orth(m)$ (left singular vectors),
  \item an orthogonal matrix $V\in\Orth(n)$ (right singular vectors),
  \item a ``diagonal'' matrix $\Sigma\in\Mat_{m\times n}(\R)$ with
    $\Sigma_{ii}=\sigma_i$ for $i=1,\ldots,\min(m,n)$ and all other
    entries zero,
\end{itemize}
such that
\[
  A = U\,\Sigma\,\transpose{V}.
\]
The nonzero singular values $\sigma_1\geq\cdots\geq\sigma_r>0$ are
uniquely determined.  Moreover,
\[
  A = \sum_{i=1}^{r} \sigma_i\, u_i\,\transpose{v_i},
\]
where $u_i$ and~$v_i$ are the $i$-th columns of~$U$ and~$V$.
\end{theorem}

\begin{proof}[Proof sketch]
The matrix $\transpose{A}A\in\Sym_n(\R)$ is positive semidefinite
(
ef{prop:prop:psd-characterizations}).  By the Spectral Theorem, there
exists an orthogonal matrix $V$ such that
$\transpose{V}\transpose{A}A\,V = \diag(\sigma_1^2,\ldots,\sigma_n^2)$
with $\sigma_1\geq\cdots\geq\sigma_r>0=\sigma_{r+1}=\cdots=\sigma_n$.

Define $u_i = \frac{1}{\sigma_i}Av_i$ for $i=1,\ldots,r$.  These are
orthonormal:
\[
  \inner{u_i,u_j}
  = \frac{1}{\sigma_i\sigma_j}\inner{Av_i,Av_j}
  = \frac{1}{\sigma_i\sigma_j}\transpose{v_i}\transpose{A}Av_j
  = \frac{\sigma_j^2}{\sigma_i\sigma_j}\delta_{ij}
  = \delta_{ij}.
\]
Extend $(u_1,\ldots,u_r)$ to an orthonormal basis
$(u_1,\ldots,u_m)$ of~$\R^m$, and set
$U = \begin{pmatrix}u_1&\cdots&u_m\end{pmatrix}$.

By construction, $AV = U\Sigma$ (checking on each column of~$V$), so
$A = U\Sigma\transpose{V}$.
\end{proof}

\begin{center}
\begin{tikzpicture}[>=Stealth, scale=1.6]
  % --- SVD geometry: unit circle -> ellipse ---
  % Left: unit circle in R^n
  \begin{scope}[shift={(-2.5,0)}]
    \draw[thick, blue!50!black] (0,0) circle (1);
    \draw[->, thick, red!70!black] (0,0) -- (0.707,0.707)
      node[above right] {\small $v_1$};
    \draw[->, thick, green!60!black] (0,0) -- (-0.707,0.707)
      node[above left] {\small $v_2$};
    \node[below] at (0,-1.3) {\small Unit circle in $\R^n$};
  \end{scope}

  % Arrow
  \draw[->, thick, dashed, gray] (-0.8,0) -- (0.8,0)
    node[midway, above] {\small $A$};

  % Right: ellipse in R^m
  \begin{scope}[shift={(2.5,0)}]
    \draw[thick, blue!50!black, rotate=30] (0,0) ellipse (1.5 and 0.6);
    \draw[->, thick, red!70!black] (0,0) --
      ({1.5*cos(30)},{1.5*sin(30)})
      node[above right] {\small $\sigma_1 u_1$};
    \draw[->, thick, green!60!black] (0,0) --
      ({-0.6*sin(30)},{0.6*cos(30)})
      node[above left] {\small $\sigma_2 u_2$};
    \node[below] at (0,-1.3) {\small Image ellipse in $\R^m$};
  \end{scope}
\end{tikzpicture}
\captionof{figure}{Geometric interpretation of the SVD\@.  The linear map~$A$
  sends the unit sphere to an ellipsoid.  The right singular vectors~$v_i$
  are mapped to the principal axes $\sigma_i u_i$ of the ellipsoid.}
\label{fig:svd-geometry}
\end{center}

\begin{example}{SVD of a $2\times 2$ matrix}{ex:svd-2x2}
Let $A = \begin{pmatrix}3&0\\0&-2\end{pmatrix}$.  Then
$\transpose{A}A = \begin{pmatrix}9&0\\0&4\end{pmatrix}$, with
eigenvalues $9$ and~$4$, so $\sigma_1=3$, $\sigma_2=2$.

The right singular vectors are $v_1=\vece{1}$, $v_2=\vece{2}$.
The left singular vectors are
$u_1 = \frac{1}{3}Av_1 = \vece{1}$,
$u_2 = \frac{1}{2}Av_2 = -\vece{2}$.

Thus $A = U\Sigma\transpose{V}$ with $U=\begin{psmallmatrix}1&0\\0&-1\end{psmallmatrix}$,
$\Sigma=\begin{psmallmatrix}3&0\\0&2\end{psmallmatrix}$, $V=I_2$.
\end{example}

\begin{proposition}{Properties of the SVD}{prop:svd-properties}
\index{SVD!properties}
Let $A = U\Sigma\transpose{V}$ be the SVD of $A\in\Mat_{m\times n}(\R)$
with $\rank(A)=r$.
\begin{enumerate}[label=\textup{(\roman*)}]
  \item $\rank(A)$ equals the number of nonzero singular values.
  \item $\norm{A}_2 = \sigma_1$ (operator norm) and
    $\norm{A}_F = \sqrt{\sigma_1^2+\cdots+\sigma_r^2}$ (Frobenius norm).
  \item $\im(A) = \Span(u_1,\ldots,u_r)$ and
    $\Ker(A) = \Span(v_{r+1},\ldots,v_n)$.
  \item If $A$ is square and invertible, $\inv{A} = V\inv{\Sigma}\transpose{U}$
    and the condition number is $\kappa(A) = \sigma_1/\sigma_n$.
\end{enumerate}
\end{proposition}

\begin{proof}
(i) follows from $\rank(A) = \rank(\transpose{A}A)$ (the nonzero
eigenvalues of $\transpose{A}A$ are $\sigma_1^2,\ldots,\sigma_r^2$).
(ii)--(iv) follow directly from the decomposition $A=U\Sigma\transpose{V}$.
\end{proof}


% ============================================================
\section{Applications}\label{sec:spectral-applications}
% ============================================================

% ---------- PCA ----------
\subsection{Principal Component Analysis}\label{subsec:pca}
\index{principal component analysis}\index{PCA}

Given data points $x_1,\ldots,x_N\in\R^n$ (centered so that
$\sum x_i = \vzero$), the \emph{sample covariance matrix} is
\[
  S = \frac{1}{N-1}\sum_{i=1}^{N} x_i\transpose{x_i}
  \;\in\;\Sym_n(\R).
\]
Since $S$ is symmetric and positive semidefinite, the Spectral Theorem
gives $S = P\diag(\lambda_1,\ldots,\lambda_n)\transpose{P}$ with
$\lambda_1\geq\cdots\geq\lambda_n\geq 0$.  The columns of~$P$ are the
\emph{principal components}: the first principal component~$p_1$
maximizes the variance $\transpose{p}\,S\,p$ over unit vectors~$p$,
the second maximizes the variance orthogonally to~$p_1$, and so on.

Dimensionality reduction to $k$ components amounts to projecting onto
$\Span(p_1,\ldots,p_k)$, retaining a fraction
$\frac{\lambda_1+\cdots+\lambda_k}{\lambda_1+\cdots+\lambda_n}$
of the total variance.

% ---------- Quadratic form classification ----------
\subsection{Classification of quadratic forms}%
\label{subsec:quadratic-forms}
\index{quadratic form!classification}

A \emph{quadratic form} on~$\R^n$ is a function
$Q(x) = \transpose{x}Ax$ where $A\in\Sym_n(\R)$ (since only the
symmetric part of a matrix contributes to the quadratic form).
By the Spectral Theorem, there is an orthogonal change of variables
$x = Py$ that diagonalizes~$Q$:
\[
  Q(x) = \transpose{y}\diag(\lambda_1,\ldots,\lambda_n)\,y
  = \lambda_1 y_1^2 + \cdots + \lambda_n y_n^2.
\]
The \emph{signature}\index{signature!of a quadratic form} $(p,q)$ ---
where $p$ is the number of positive~$\lambda_i$ and $q$ the number of
negative ones --- classifies the quadratic form up to congruence
(Sylvester's law of inertia\index{Sylvester's law of inertia}).

% ---------- Optimization ----------
\subsection{Second-order optimality conditions}%
\label{subsec:optimization}
\index{Hessian}\index{optimization!second-order conditions}

Let $f\colon\R^n\to\R$ be $C^2$.  At a critical point $x_0$
($\nabla f(x_0)=\vzero$), the Hessian
$H = \bigl(\frac{\partial^2 f}{\partial x_i\partial x_j}(x_0)\bigr)$
is symmetric.  By the Spectral Theorem:
\begin{itemize}
  \item $H$ positive definite $\Rightarrow$ $x_0$ is a strict local minimum.
  \item $H$ negative definite $\Rightarrow$ $x_0$ is a strict local maximum.
  \item $H$ indefinite $\Rightarrow$ $x_0$ is a saddle point.
\end{itemize}

% ---------- Vibrations ----------
\subsection{Vibrations and normal modes}\label{subsec:vibrations}
\index{vibrations}\index{normal modes}

A coupled system of $n$ harmonic oscillators (masses connected by
springs) leads to the equation $M\ddot{x} + Kx = \vzero$, where
$M,K\in\Sym_n(\R)$ are the mass and stiffness matrices (both positive
definite).  Setting $y = M^{1/2}x$ transforms this into
$\ddot{y} + M^{-1/2}KM^{-1/2}y = \vzero$.

The symmetric matrix $\tilde K = M^{-1/2}KM^{-1/2}$ has positive
eigenvalues $\omega_1^2\leq\cdots\leq\omega_n^2$; the~$\omega_i$ are
the \emph{natural frequencies} and the corresponding eigenvectors
(transformed back via $M^{-1/2}$) are the \emph{normal modes} of the
system.


% ============================================================
\section{Exercises}\label{sec:spectral-exercises}
% ============================================================

\begin{exercise}{Eigenvalues of a symmetric matrix \basic}{exo:sym-eigen}
\index{symmetric matrix!eigenvalues}
Let $A = \begin{pmatrix} 4 & 2 \\ 2 & 1 \end{pmatrix}$.
\begin{enumerate}[label=\textup{(\alph*)}]
  \item Find the eigenvalues of~$A$ and verify they are real.
  \item Find an orthonormal basis of eigenvectors.
  \item Write $A = P D \transpose{P}$ with $P$ orthogonal and $D$ diagonal.
\end{enumerate}
\end{exercise}

\begin{exercise}{Orthogonal diagonalization \basic}{exo:orth-diag}
\index{orthogonal diagonalization}
Orthogonally diagonalize
$A = \begin{pmatrix} 6 & -2 \\ -2 & 3 \end{pmatrix}$.
\end{exercise}

\begin{exercise}{A $3\times 3$ symmetric matrix \basic}{exo:3x3-sym}
Let $A = \begin{pmatrix} 1 & 0 & 1 \\ 0 & 1 & 0 \\ 1 & 0 & 1\end{pmatrix}$.
Find the eigenvalues, eigenspaces, and an orthogonal matrix~$P$ such that
$\transpose{P}AP$ is diagonal.
\end{exercise}

\begin{exercise}{Proof practice: eigenvalues of $A^2$ \basic}{exo:eigen-A2}
Let $A\in\Sym_n(\R)$ with eigenvalues $\lambda_1,\ldots,\lambda_n$.
Prove that the eigenvalues of~$A^2$ are $\lambda_1^2,\ldots,\lambda_n^2$.
\end{exercise}

\begin{exercise}{Positive definiteness tests \intermediate}{exo:pd-tests}
\index{positive definite matrix!tests}
Determine whether each matrix is positive definite, positive semidefinite,
or neither:
\begin{enumerate}[label=\textup{(\alph*)}]
  \item $A = \begin{pmatrix} 3 & 1 \\ 1 & 3 \end{pmatrix}$.
  \item $B = \begin{pmatrix} 1 & 2 \\ 2 & 1 \end{pmatrix}$.
  \item $C = \begin{pmatrix} 4 & 2 & 0 \\ 2 & 5 & 3 \\ 0 & 3 & 6
    \end{pmatrix}$.
\end{enumerate}
\end{exercise}

\begin{exercise}{Cholesky decomposition \intermediate}{exo:cholesky}
\index{Cholesky decomposition}
Compute the Cholesky decomposition $A = L\transpose{L}$ for
$A = \begin{pmatrix} 4 & 2 \\ 2 & 5 \end{pmatrix}$.
Verify your answer.
\end{exercise}

\begin{exercise}{SVD computation \intermediate}{exo:svd-compute}
\index{SVD!computation}
Compute the SVD of
$A = \begin{pmatrix} 1 & 1 \\ 0 & 1 \\ 1 & 0 \end{pmatrix}$.
\end{exercise}

\begin{exercise}{Spectral decomposition \intermediate}{exo:spectral-decomp}
\index{spectral decomposition}
Let $A = \begin{pmatrix} 3 & 1 \\ 1 & 3 \end{pmatrix}$.
Write $A$ in spectral decomposition form
$A = \lambda_1 P_1 + \lambda_2 P_2$, where $P_1,P_2$ are the
orthogonal projections onto the eigenspaces.  Verify that
$P_1+P_2 = I_2$ and $P_1P_2 = 0$.
\end{exercise}

\begin{exercise}{Normal but not self-adjoint \intermediate}{exo:normal-not-sa}
\index{normal operator}
Let $A = \begin{pmatrix} 0 & -1 \\ 1 & 0 \end{pmatrix}$.
\begin{enumerate}[label=\textup{(\alph*)}]
  \item Show that $A$ is normal.
  \item Show that $A$ has no real eigenvalues, hence is not orthogonally
    diagonalizable over~$\R$.
  \item Find the eigenvalues over~$\C$ and a unitary matrix~$U$ such
    that $\hermit{U}AU$ is diagonal.
\end{enumerate}
\end{exercise}

\begin{exercise}{Simultaneous diagonalization \challenging}{exo:simult-diag}
\index{simultaneous diagonalization}
Let $A,B\in\Sym_n(\R)$ with $AB=BA$.  Prove that $A$ and~$B$ can be
\emph{simultaneously} orthogonally diagonalized, i.e.\ there exists
$P\in\Orth(n)$ such that both $\transpose{P}AP$ and $\transpose{P}BP$
are diagonal.

\emph{Hint:} Show that each eigenspace of~$A$ is invariant under~$B$,
then apply the Spectral Theorem to the restriction of~$B$ to each
eigenspace.
\end{exercise}

\begin{exercise}{Low-rank approximation via SVD \challenging}{exo:low-rank}
\index{SVD!low-rank approximation}
Let $A\in\Mat_{m\times n}(\R)$ with SVD\
$A = \sum_{i=1}^r \sigma_i\,u_i\transpose{v_i}$.
Define $A_k = \sum_{i=1}^k \sigma_i\,u_i\transpose{v_i}$ for $k\leq r$.
\begin{enumerate}[label=\textup{(\alph*)}]
  \item Show that $\rank(A_k)=k$.
  \item Prove the Eckart--Young theorem\index{Eckart--Young theorem}:
    among all matrices of rank at most~$k$, $A_k$ minimises
    $\norm{A-B}_F$, and the minimum value is
    $\sqrt{\sigma_{k+1}^2+\cdots+\sigma_r^2}$.
\end{enumerate}
\end{exercise}

\begin{exercise}{Rayleigh quotient and eigenvalue bounds \challenging}%
{exo:rayleigh}
\index{Rayleigh quotient}
For $A\in\Sym_n(\R)$, the \emph{Rayleigh quotient}\index{Rayleigh quotient}
is $R(x) = \frac{\transpose{x}Ax}{\transpose{x}x}$ for
$x\neq\vzero$.
\begin{enumerate}[label=\textup{(\alph*)}]
  \item Prove that $\lambda_{\min}(A)\leq R(x)\leq\lambda_{\max}(A)$.
  \item Prove that $\max_{x\neq\vzero}R(x) = \lambda_{\max}(A)$
    and $\min_{x\neq\vzero}R(x) = \lambda_{\min}(A)$, and identify
    the optimisers.
  \item (\emph{Courant--Fischer minimax theorem}\index{Courant--Fischer theorem})
    Prove that the $k$-th largest eigenvalue satisfies
    \[
      \lambda_k = \max_{\dim W = k}\;\min_{x\in W,\,x\neq\vzero}
        R(x)
      = \min_{\dim W = n-k+1}\;\max_{x\in W,\,x\neq\vzero} R(x).
    \]
\end{enumerate}
\end{exercise}


% ============================================================
\section{Chapter summary}\label{sec:spectral-summary}
% ============================================================

\begin{enumerate}
  \item \textbf{Self-adjoint operators.}  An operator~$f$ is self-adjoint
    if $f^*=f$; in an orthonormal basis, this corresponds to symmetric
    ($\transpose{A}=A$) or Hermitian ($\hermit{A}=A$) matrices.

  \item \textbf{Eigenvalues are real.}  Every eigenvalue of a self-adjoint
    operator is a real number.

  \item \textbf{Orthogonal eigenvectors.}  Eigenvectors corresponding to
    distinct eigenvalues of a self-adjoint operator are orthogonal.

  \item \textbf{Spectral Theorem (real).}  Every real symmetric matrix is
    orthogonally diagonalizable: $A = P\diag(\lambda_1,\ldots,\lambda_n)\transpose{P}$
    with $P\in\Orth(n)$.

  \item \textbf{Spectral Theorem (complex).}  Every Hermitian matrix is
    unitarily diagonalizable with real eigenvalues.

  \item \textbf{Normal operators.}  An operator is normal if $ff^*=f^*f$.
    Over~$\C$, normal operators are precisely the unitarily
    diagonalizable ones.

  \item \textbf{Positive definiteness.}  $A\in\Sym_n(\R)$ is positive
    definite if and only if all eigenvalues are positive, if and only if
    all leading principal minors are positive (Sylvester's criterion).

  \item \textbf{SVD.}  Every $m\times n$ matrix admits a decomposition
    $A = U\Sigma\transpose{V}$ with $U,V$ orthogonal and $\Sigma$
    ``diagonal.''  The singular values are the square roots of the
    eigenvalues of $\transpose{A}A$.

  \item \textbf{Applications.}  The Spectral Theorem underlies PCA,
    quadratic form classification, second-order optimality conditions,
    and the analysis of coupled vibrations.
\end{enumerate}

% ch8-jordan.tex — Chapter 8: Jordan Normal Form
% Nilpotent endomorphisms, generalized eigenspaces, Jordan blocks,
% the Jordan normal form theorem, matrix exponential, applications.

\chapter{Jordan Normal Form}\label{ch:jordan}

% ============================================================
% Motivating introduction
% ============================================================

In 
ef{ch:eigenvalues} we saw that a matrix is diagonalizable if and
only if the algebraic and geometric multiplicities of every eigenvalue
coincide.  When they do, one obtains the simplest possible matrix
representation --- a diagonal matrix.  But what happens when
diagonalization \emph{fails}?

Consider the matrix
$A = \begin{psmallmatrix} 2 & 1 \\ 0 & 2 \end{psmallmatrix}$.
Its only eigenvalue is $\lambda = 2$ with algebraic multiplicity~$2$,
yet the eigenspace $E_2(A) = \Span\!\bigl(\begin{psmallmatrix} 1 \\ 0
\end{psmallmatrix}\bigr)$ has dimension~$1$.  There is no basis
of~$\R^2$ consisting of eigenvectors of~$A$, so $A$ is not
diagonalizable.  Nevertheless, $A$ is already in a rather pleasant
form: it is ``almost diagonal,'' differing from $\diag(2,2)$ by a
single off-diagonal entry.  Is this an accident, or can every matrix be
brought into a similarly nice shape?

The answer is provided by the \emph{Jordan normal form}\index{Jordan
normal form}, one of the most powerful results in linear algebra.  It
asserts that over an algebraically closed field (such as~$\C$), every
square matrix is similar to an essentially unique matrix built from
simple building blocks called \emph{Jordan blocks}.  This canonical form
completely determines the matrix up to similarity and provides a
systematic way to handle all the phenomena that arise when
diagonalization fails.

The Jordan form has far-reaching applications: it yields explicit
formulas for matrix powers and matrix exponentials, it explains the
structure of solutions to linear systems of differential equations with
non-diagonalizable coefficient matrices, and it clarifies the
relationship between the characteristic and minimal polynomials.

Throughout this chapter, unless stated otherwise, $\K$ denotes an
algebraically closed field (typically $\K = \C$), $E$ is a
finite-dimensional $\K$-vector space with $\dim E = n$, and
$f \in \End(E)$.  We freely use results from 
ef{ch:eigenvalues},
especially the notions of eigenvalue, eigenspace, characteristic
polynomial, and minimal polynomial.


% ============================================================
\section{Nilpotent endomorphisms}\label{sec:nilpotent}
% ============================================================

The key to understanding non-diagonalizable maps is to understand a
special class of endomorphisms that are ``as far from diagonalizable as
possible'' (aside from the zero map): the \emph{nilpotent}
endomorphisms.

\begin{definition}{Nilpotent endomorphism}{def:nilpotent}
\index{nilpotent!endomorphism}\index{endomorphism!nilpotent}
An endomorphism $f \in \End(E)$ is called \emph{nilpotent} if there
exists an integer $p \geq 1$ such that $f^p = 0$ (the zero
endomorphism).  A matrix $A \in \Mat_n(\K)$ is \emph{nilpotent} if
$A^p = 0$ for some $p \geq 1$.
\end{definition}

\begin{definition}{Index of nilpotency}{def:nilpotent-index}
\index{nilpotent!index of}\index{index of nilpotency}
If $f \in \End(E)$ is nilpotent, the \emph{index of nilpotency}
(or \emph{nilpotency index}) of~$f$ is the smallest positive
integer~$r$ such that $f^r = 0$.  We then have
\[
  f^{r-1} \neq 0 \quad\text{and}\quad f^r = 0.
\]
\end{definition}

\begin{proposition}{Basic properties of nilpotent endomorphisms}%
{prop:nilpotent-properties}
\index{nilpotent!properties}
Let $f \in \End(E)$ be nilpotent with $\dim E = n$.  Then:
\begin{enumerate}[label=(\roman*)]
  \item\label{it:nilp-eigenvalue}
    The only eigenvalue of~$f$ is $\lambda = 0$.
  \item\label{it:nilp-trace}
    $\tr(f) = 0$.
  \item\label{it:nilp-det}
    $\det(f) = 0$; in particular, $f$ is not invertible.
  \item\label{it:nilp-charpoly}
    The characteristic polynomial of~$f$ is $\chi_f(\lambda) = (-\lambda)^n$.
  \item\label{it:nilp-index-bound}
    The index of nilpotency satisfies $r \leq n$.
  \item\label{it:nilp-minpoly}
    The minimal polynomial of~$f$ is $\mu_f(\lambda) = (-\lambda)^r$,
    where $r$ is the index of nilpotency.
\end{enumerate}
\end{proposition}

\begin{proof}
\ref{it:nilp-eigenvalue}\;
Suppose $f(v) = \lambda v$ for some $v \neq 0$.  Then
$f^p(v) = \lambda^p v$, so $f^p = 0$ implies $\lambda^p v = 0$,
hence $\lambda^p = 0$ (since $v \neq 0$), so $\lambda = 0$.

\ref{it:nilp-trace}\ and \ref{it:nilp-det}\;
Since $f$ is trigonalizable (every endomorphism over~$\C$ is, and
the argument extends because all eigenvalues are~$0$), there
exists a basis in which the matrix of~$f$ is upper triangular with
all diagonal entries equal to~$0$.  Hence $\tr(f) = 0$ and
$\det(f) = 0$.

\ref{it:nilp-charpoly}\;
From the triangular form, $\chi_f(\lambda) = (0 - \lambda)^n =
(-\lambda)^n$.

\ref{it:nilp-index-bound}\;
By Cayley--Hamilton, $\chi_f(f) = 0$, i.e.\ $(-1)^n f^n = 0$, so
$f^n = 0$.  The index of nilpotency, being the smallest such
exponent, satisfies $r \leq n$.

\ref{it:nilp-minpoly}\;
Since $\chi_f(\lambda) = (-\lambda)^n$ and the minimal polynomial
divides the characteristic polynomial and has the same roots, we
have $\mu_f(\lambda) = (-\lambda)^s$ for some $1 \leq s \leq n$.
By definition, $\mu_f(f) = 0$ means $f^s = 0$, and $s$ is minimal
with this property.  Hence $s = r$.
\end{proof}

\begin{example}{Nilpotent matrices}{ex:nilpotent-matrices}
\index{nilpotent!examples}
\begin{enumerate}[label=(\alph*)]
  \item The matrix $N = \begin{pmatrix} 0 & 1 \\ 0 & 0 \end{pmatrix}$
    satisfies $N^2 = 0$, so it is nilpotent with index~$2$.

  \item The matrix
    $N = \begin{pmatrix} 0 & 1 & 0 \\ 0 & 0 & 1 \\ 0 & 0 & 0
    \end{pmatrix}$ satisfies $N^2 \neq 0$ but $N^3 = 0$, so it is
    nilpotent with index~$3$.

  \item The zero matrix $0_n$ is nilpotent with index~$1$.

  \item The matrix
    $\begin{pmatrix} 0 & 1 & 0 \\ 0 & 0 & 0 \\ 0 & 0 & 0
    \end{pmatrix}$ is nilpotent with index~$2$.
\end{enumerate}
\end{example}

\begin{proposition}{Kernel filtration of a nilpotent endomorphism}%
{prop:nilpotent-kernel-chain}
\index{nilpotent!kernel filtration}
Let $f \in \End(E)$ be nilpotent with index of nilpotency~$r$.  Then
\[
  \set{\vzero}
  \subsetneq \Ker(f)
  \subsetneq \Ker(f^2)
  \subsetneq \cdots
  \subsetneq \Ker(f^{r-1})
  \subsetneq \Ker(f^r)
  = E.
\]
In particular, each inclusion is strict, and $\dim\Ker(f^k) \geq k$
for all $1 \leq k \leq r$.
\end{proposition}

\begin{proof}
The inclusions $\Ker(f^k) \subseteq \Ker(f^{k+1})$ follow from the
fact that $f^k(v) = 0$ implies $f^{k+1}(v) = f(f^k(v)) = f(0) = 0$.
Since $f^r = 0$, we have $\Ker(f^r) = E$.  It remains to show that
each inclusion is strict.

Suppose for contradiction that $\Ker(f^k) = \Ker(f^{k+1})$ for
some $k < r$.  We claim this forces $\Ker(f^k) = \Ker(f^{k+j})$
for all $j \geq 0$.  Indeed, let $v \in \Ker(f^{k+2})$, so
$f^{k+2}(v) = 0$, meaning $f(v) \in \Ker(f^{k+1}) = \Ker(f^k)$,
hence $f^{k+1}(v) = f^k(f(v)) = 0$, so $v \in \Ker(f^{k+1}) =
\Ker(f^k)$.  By induction, $\Ker(f^k) = \Ker(f^m)$ for all
$m \geq k$.  In particular, $\Ker(f^k) = \Ker(f^r) = E$, so
$f^k = 0$, contradicting the minimality of~$r$ since $k < r$.

Since the inclusions are strict in a space of dimension~$n$, at
each step the dimension increases by at least~$1$, giving
$\dim\Ker(f^k) \geq k$.
\end{proof}


% ============================================================
\section{Invariant subspaces}\label{sec:invariant-subspaces}
% ============================================================

\begin{definition}{Invariant subspace}{def:invariant-subspace}
\index{invariant subspace}
Let $f \in \End(E)$ and let $F \subseteq E$ be a subspace.
We say that $F$ is \emph{$f$-invariant} (or \emph{invariant under
$f$}) if $f(F) \subseteq F$, i.e.\ $f(v) \in F$ for every $v \in F$.
\end{definition}

\begin{remark}{Restriction to an invariant subspace}{rem:restriction-invariant}
If $F$ is $f$-invariant, then $f$ restricts to a well-defined
endomorphism $f\restrict{F} \in \End(F)$.  This simple observation is
the key to decomposing endomorphisms into simpler pieces.
\end{remark}

\begin{proposition}{Examples of invariant subspaces}%
{prop:invariant-examples}
\index{invariant subspace!examples}
Let $f \in \End(E)$.  The following are $f$-invariant subspaces:
\begin{enumerate}[label=(\roman*)]
  \item $\set{\vzero}$ and $E$ (trivial invariant subspaces).
  \item $\Ker(f)$ and $\im(f)$.
  \item More generally, $\Ker(f^k)$ and $\im(f^k)$ for every $k \geq 1$.
  \item Every eigenspace $E_\lambda(f) = \Ker(f - \lambda\Id)$.
  \item $\Ker(P(f))$ for any polynomial $P \in \K[\lambda]$.
\end{enumerate}
\end{proposition}

\begin{proof}
We prove~(v), which implies the others as special cases.  Let
$P \in \K[\lambda]$ and $v \in \Ker(P(f))$, so $P(f)(v) = 0$.
Since $f$ commutes with $P(f)$ (both are polynomials in~$f$), we
have
\[
  P(f)\bigl(f(v)\bigr) = f\bigl(P(f)(v)\bigr) = f(0) = 0,
\]
so $f(v) \in \Ker(P(f))$.
\end{proof}

\begin{proposition}{Block-diagonal structure from invariant decomposition}%
{prop:block-diagonal}
\index{invariant subspace!block-diagonal form}
Let $E = F_1 \oplus F_2 \oplus \cdots \oplus F_s$ where each~$F_i$ is
$f$-invariant.  If $\cB_i$ is a basis of~$F_i$ and
$\cB = \cB_1 \cup \cB_2 \cup \cdots \cup \cB_s$ is the resulting
basis of~$E$, then the matrix of~$f$ in~$\cB$ is block-diagonal:
\[
  \operatorname{Mat}_\cB(f) = \begin{pmatrix}
    A_1 &        &        \\
        & A_2    &        \\
        &        & \ddots \\
        &        &        & A_s
  \end{pmatrix},
\]
where $A_i = \operatorname{Mat}_{\cB_i}(f\restrict{F_i})$.
Moreover, $\chi_f = \chi_{f\restrict{F_1}} \cdots
\chi_{f\restrict{F_s}}$.
\end{proposition}

\begin{proof}
Since each $F_i$ is $f$-invariant, $f$ maps basis vectors of
$\cB_i$ to linear combinations of vectors in $\cB_i$ only.  Hence
the matrix has the stated block-diagonal form.  The factorization
of the characteristic polynomial follows from the multiplicativity
of determinants for block-diagonal matrices.
\end{proof}


% ============================================================
\section{Generalized eigenspaces}\label{sec:generalized-eigenspaces}
% ============================================================

When the ordinary eigenspace $E_\lambda(f) = \Ker(f - \lambda\Id)$ is
too small (i.e.\ $\dim E_\lambda(f) < m_a(\lambda)$), we enlarge it
by looking at the kernels of higher powers of $f - \lambda\Id$.

\begin{definition}{Generalized eigenvector and generalized eigenspace}%
{def:generalized-eigenspace}
\index{generalized eigenvector}\index{generalized eigenspace}
\index{eigenspace!generalized}
Let $f \in \End(E)$ and $\lambda \in \K$.  A nonzero vector
$v \in E$ is a \emph{generalized eigenvector} of~$f$ for the
eigenvalue~$\lambda$ if there exists an integer $k \geq 1$ such that
\[
  (f - \lambda\Id)^k(v) = 0.
\]
The \emph{generalized eigenspace} of~$f$ for~$\lambda$ is
\[
  G_\lambda(f) \deff \bigcup_{k=1}^{\infty} \Ker\bigl(
    (f - \lambda\Id)^k
  \bigr).
\]
\end{definition}

\begin{remark}{Stabilization of the kernel chain}%
{rem:generalized-eigenspace-stabilizes}
Since $\dim E = n$ is finite, the ascending chain
$\Ker(f - \lambda\Id) \subseteq \Ker(f - \lambda\Id)^2 \subseteq
\cdots$ must stabilize at some index $k_0 \leq n$.  Therefore
\[
  G_\lambda(f) = \Ker\bigl((f - \lambda\Id)^{k_0}\bigr)
\]
for some $k_0 \leq n$, and the union in the definition is in fact
achieved at a finite stage.
\end{remark}

\begin{proposition}{Properties of generalized eigenspaces}%
{prop:generalized-eigenspace-properties}
\index{generalized eigenspace!properties}
Let $f \in \End(E)$ and suppose $\chi_f$ splits over~$\K$ as
\[
  \chi_f(\lambda) = \prod_{i=1}^{s} (\lambda_i - \lambda)^{n_i},
\]
where $\lambda_1, \ldots, \lambda_s$ are the distinct eigenvalues
and $n_i = m_a(\lambda_i)$ is the algebraic multiplicity.  Then:
\begin{enumerate}[label=(\roman*)]
  \item\label{it:gen-eig-f-inv}
    Each $G_{\lambda_i}(f)$ is $f$-invariant.
  \item\label{it:gen-eig-dim}
    $\dim G_{\lambda_i}(f) = n_i$.
  \item\label{it:gen-eig-kernel}
    $G_{\lambda_i}(f) = \Ker\bigl((f - \lambda_i\Id)^{n_i}\bigr)$.
  \item\label{it:gen-eig-restriction}
    The restriction $(f - \lambda_i\Id)\restrict{G_{\lambda_i}(f)}$
    is nilpotent.
  \item\label{it:gen-eig-direct-sum}
    The generalized eigenspaces are in direct sum:
    \[
      E = G_{\lambda_1}(f) \oplus G_{\lambda_2}(f) \oplus \cdots
        \oplus G_{\lambda_s}(f).
    \]
\end{enumerate}
\end{proposition}

\begin{proof}
\ref{it:gen-eig-f-inv}\;
This follows from 
ef{prop:prop:invariant-examples}, since
$G_\lambda(f) = \Ker\bigl((f - \lambda\Id)^k\bigr)$ for
sufficiently large~$k$, and the kernel of any polynomial in~$f$ is
$f$-invariant.

\medskip\noindent
We prove \ref{it:gen-eig-dim}, \ref{it:gen-eig-kernel}, and
\ref{it:gen-eig-direct-sum} together using the primary decomposition
theorem.

Write $\chi_f(\lambda) = \prod_{i=1}^s (\lambda_i - \lambda)^{n_i}$.
Define the polynomials $P_i(\lambda) = (\lambda_i - \lambda)^{n_i}$
and $Q_i(\lambda) = \prod_{j \neq i} (\lambda_j - \lambda)^{n_j}$.
Then $\gcd(P_i, Q_i) = 1$ since the~$\lambda_j$ are distinct.
By B\'ezout's identity in~$\K[\lambda]$, there exist polynomials
$U_i, V_i$ such that $U_i P_i + V_i Q_i = 1$.

By Cayley--Hamilton, $\chi_f(f) = 0$, i.e.\ $P_i(f)\,Q_i(f) = 0$
on $G_{\lambda_i}(f)$, and indeed
$\prod_{i=1}^s P_i(f) = 0$ on all of~$E$.

From $U_i(f) P_i(f) + V_i(f) Q_i(f) = \Id$, for any $v \in E$ we
have $v = U_i(f)P_i(f)(v) + V_i(f)Q_i(f)(v)$.  One checks that
$V_i(f)Q_i(f)(v) \in \Ker(P_i(f)) = \Ker\bigl((f -
\lambda_i\Id)^{n_i}\bigr)$ by applying $P_i(f)$ and using
$P_i(f)Q_i(f) = \chi_f(f) = 0$.  Similarly,
$U_i(f)P_i(f)(v) \in \Ker(Q_i(f))$.

Define $\pi_i = V_i(f)Q_i(f)$.  Then $\sum_{i=1}^s \pi_i = \Id$,
each $\pi_i$ maps $E$ into $\Ker(P_i(f)) = G_{\lambda_i}(f)$, and
$\pi_i \pi_j = 0$ for $i \neq j$.  This gives the direct sum
decomposition $E = \bigoplus_{i=1}^s G_{\lambda_i}(f)$.

For the dimensions, since
$\chi_{f\restrict{G_{\lambda_i}(f)}}$ divides~$\chi_f$ and
the only eigenvalue of~$f$ on $G_{\lambda_i}(f)$ is~$\lambda_i$,
we get $\chi_{f\restrict{G_{\lambda_i}(f)}}(\lambda) =
(\lambda_i - \lambda)^{d_i}$ with $d_i = \dim G_{\lambda_i}(f)$.
Since $\chi_f = \prod_i \chi_{f\restrict{G_{\lambda_i}(f)}}$
(
ef{prop:prop:block-diagonal}), comparing degrees gives $d_i = n_i$.

\ref{it:gen-eig-restriction}\;
On $G_{\lambda_i}(f)$, by definition
$(f - \lambda_i\Id)^{n_i}\restrict{G_{\lambda_i}(f)} = 0$,
so the restriction is nilpotent with index at most~$n_i$.
\end{proof}

\begin{theorem}{Primary decomposition theorem}{thm:primary-decomposition}
\index{primary decomposition theorem}\index{decomposition!primary}
Let $f \in \End(E)$ and suppose $\chi_f$ splits over~$\K$ as
$\chi_f(\lambda) = \prod_{i=1}^{s} (\lambda_i - \lambda)^{n_i}$.
Then
\[
  \boxed{E = G_{\lambda_1}(f) \oplus G_{\lambda_2}(f) \oplus \cdots
    \oplus G_{\lambda_s}(f),}
\]
where $\dim G_{\lambda_i}(f) = n_i$ and each $G_{\lambda_i}(f)$ is
$f$-invariant.  Moreover, the restriction $f\restrict{G_{\lambda_i}(f)}$
has the unique eigenvalue~$\lambda_i$.
\end{theorem}

\begin{proof}
This was established in the proof of

ef{prop:prop:generalized-eigenspace-properties}.
\end{proof}


% ============================================================
\section{Jordan blocks and Jordan matrices}\label{sec:jordan-blocks}
% ============================================================

\begin{definition}{Jordan block}{def:jordan-block}
\index{Jordan block}\index{block!Jordan}
For $\lambda \in \K$ and $k \geq 1$, the \emph{Jordan block} of
size~$k$ associated with~$\lambda$ is the $k \times k$ matrix
\[
  J_k(\lambda) \deff
  \begin{pmatrix}
    \lambda & 1       &        &         \\
            & \lambda & 1      &         \\
            &         & \ddots & \ddots  \\
            &         &        & \lambda & 1 \\
            &         &        &         & \lambda
  \end{pmatrix}
  = \lambda I_k + N_k,
\]
where $N_k$ is the $k \times k$ nilpotent matrix with $1$'s on the
superdiagonal and $0$'s elsewhere.
\end{definition}

\begin{remark}{Properties of a Jordan block}{rem:jordan-block-props}
\index{Jordan block!properties}
Let $J = J_k(\lambda)$.
\begin{enumerate}[label=(\roman*)]
  \item $\chi_J(\mu) = (\lambda - \mu)^k$.
  \item The only eigenvalue of~$J$ is~$\lambda$, with algebraic
    multiplicity~$k$ and geometric multiplicity~$1$.
  \item $J - \lambda I_k = N_k$ is nilpotent with index~$k$.
  \item $\mu_J(\mu) = (\lambda - \mu)^k$ (the minimal polynomial
    equals the characteristic polynomial).
\end{enumerate}
\end{remark}

\begin{definition}{Jordan matrix}{def:jordan-matrix}
\index{Jordan matrix}\index{matrix!Jordan}
A \emph{Jordan matrix} is a block-diagonal matrix of the form
\[
  J = \begin{pmatrix}
    J_{k_1}(\lambda_1) &                     &        \\
                        & J_{k_2}(\lambda_2) &        \\
                        &                     & \ddots \\
                        &                     &        & J_{k_m}(\lambda_m)
  \end{pmatrix},
\]
where each $J_{k_i}(\lambda_i)$ is a Jordan block.  The eigenvalues
$\lambda_1, \ldots, \lambda_m$ need not be distinct.
\end{definition}


% ============================================================
\section{The Jordan Normal Form theorem}\label{sec:jordan-theorem}
% ============================================================

\begin{theorem}{Jordan Normal Form}{thm:jordan-normal-form}
\index{Jordan normal form!theorem}
Let $E$ be a finite-dimensional vector space over an algebraically
closed field~$\K$ (e.g.\ $\K = \C$), and let $f \in \End(E)$.  Then
there exists a basis~$\cB$ of~$E$ in which the matrix of~$f$ is a
Jordan matrix:
\[
  \operatorname{Mat}_\cB(f)
  = \begin{pmatrix}
    J_{k_1}(\lambda_1) &                     &        \\
                        & J_{k_2}(\lambda_2) &        \\
                        &                     & \ddots \\
                        &                     &        & J_{k_m}(\lambda_m)
  \end{pmatrix}.
\]
Equivalently, every matrix $A \in \Mat_n(\K)$ is similar to a Jordan
matrix: $A = P J \inv{P}$ for some invertible~$P$.

Moreover, the Jordan matrix is unique up to permutation of the blocks.
\end{theorem}

The proof occupies the next section.  The overall strategy is:
\begin{enumerate}
  \item Use the primary decomposition to reduce to the case where $f$
    has a single eigenvalue.
  \item When $f$ has a single eigenvalue~$\lambda$, write
    $f = \lambda\Id + g$ where $g$ is nilpotent.
  \item Prove that every nilpotent endomorphism has a Jordan form
    (a direct sum of nilpotent Jordan blocks $J_k(0) = N_k$).
\end{enumerate}


% ============================================================
\section{Proof of the Jordan Normal Form theorem}\label{sec:jordan-proof}
% ============================================================

We begin with the nilpotent case, which is the heart of the proof.

\begin{lemma}{Jordan form for nilpotent endomorphisms}%
{lem:jordan-nilpotent}
\index{nilpotent!Jordan form}
Let $g \in \End(E)$ be nilpotent with $\dim E = n$.  Then there
exist positive integers $k_1 \geq k_2 \geq \cdots \geq k_m$ with
$k_1 + k_2 + \cdots + k_m = n$ and a basis of~$E$ in which the
matrix of~$g$ is
\[
  \begin{pmatrix}
    N_{k_1} &         &        \\
            & N_{k_2} &        \\
            &         & \ddots \\
            &         &        & N_{k_m}
  \end{pmatrix},
\]
where $N_k = J_k(0)$ is the $k \times k$ nilpotent Jordan block.
The integers $k_1, \ldots, k_m$ are uniquely determined by~$g$.
\end{lemma}

\begin{proof}
We proceed by strong induction on $n = \dim E$.

\textbf{Base case.}  If $n = 1$, then $g = 0$ and the matrix is the
$1 \times 1$ zero matrix $N_1$, which is already in Jordan form.

\textbf{Inductive step.}  Assume $n \geq 2$ and the result holds for
all spaces of dimension $< n$.  Let $r$ be the index of nilpotency
of~$g$, so $g^r = 0$ and $g^{r-1} \neq 0$.

If $r = 1$, then $g = 0$ and the matrix is the zero matrix, which is
a direct sum of $n$ copies of $N_1$.

Assume $r \geq 2$.  Consider the subspace $F = \im(g)$.  Since
$g \neq 0$, we have $F \neq \set{\vzero}$ and $\dim F < n$ (because
$g$ is not invertible).  Moreover, $F$ is $g$-invariant: for any
$v \in E$, $g(g(v)) = g^2(v) \in \im(g) = F$.

The restriction $g\restrict{F}$ is nilpotent on~$F$, and
$\dim F < n$.  By the induction hypothesis, there exists a basis of~$F$
in which $g\restrict{F}$ has the form
$\diag(N_{k_1-1}, N_{k_2-1}, \ldots, N_{k_p-1}, N_{k_{p+1}}, \ldots)$
(rearranging and shifting indices as needed).

Concretely, there exist \emph{cyclic chains} for~$g$.  Choose a vector
$v_1 \in E$ such that $g^{r-1}(v_1) \neq 0$.  Then the vectors
\[
  v_1,\; g(v_1),\; g^2(v_1),\; \ldots,\; g^{r-1}(v_1)
\]
are linearly independent (if $\sum_{j=0}^{r-1} \alpha_j g^j(v_1) = 0$,
applying $g^{r-1-j}$ for the smallest~$j$ with $\alpha_j \neq 0$ gives
a contradiction).  Write $W_1 = \Span(v_1, g(v_1), \ldots,
g^{r-1}(v_1))$.  In the basis
$(g^{r-1}(v_1), g^{r-2}(v_1), \ldots, g(v_1), v_1)$, the matrix
of $g\restrict{W_1}$ is $N_r$.

We now construct a complement.  Set
$d_j = \dim\Ker(g^j)$ for $0 \leq j \leq r$.  Define
$m_j = 2d_j - d_{j-1} - d_{j+1}$ for $1 \leq j \leq r-1$ and
$m_r = 2d_r - d_{r-1} - d_r = d_r - d_{r-1}$.  We claim that~$m_j$
equals the number of Jordan blocks of size~$j$.

More precisely, we construct the full Jordan basis as follows.
For each $j$ from $r$ down to~$1$, choose vectors $v$ such that
$g^{j-1}(v) \neq 0$ but $g^{j}(v) = 0$, and such that
$v, g(v), \ldots, g^{j-1}(v)$ are linearly independent of the vectors
already chosen.  The number of such chains of length~$j$ is
$m_j = d_j - d_{j-1} - (\text{number of chain elements of length}
> j\text{ lying in }\Ker(g^j)\setminus\Ker(g^{j-1}))$.

To be fully explicit, define $\ell_j = d_j - d_{j-1}$ (with $d_0 = 0$).
One can show that $\ell_1 \geq \ell_2 \geq \cdots \geq \ell_r \geq 1$
and that $\ell_j$ equals the number of Jordan blocks of size $\geq j$.
Hence the number of blocks of size exactly~$j$ is
$\ell_j - \ell_{j+1}$ (where $\ell_{r+1} = 0$).

We verify this by induction.  The subspace $W_1$ Spanned by the first
chain has the property that $g\restrict{W_1}$ corresponds to $N_r$,
and $E = W_1 \oplus W_1'$ where $W_1'$ is a $g$-stable complement.
Such a complement exists: choose a complement $C$ of $\Ker(g^{r-1})$
inside $\Ker(g^r) = E$, so $\dim C = \ell_r$.  For each basis
vector~$c$ of~$C$, the chain $c, g(c), \ldots, g^{r-1}(c)$ Spans
a cyclic subspace isomorphic to $N_r$.  We can choose these chains so
that their union is linearly independent.  After removing these chains,
the remaining space (a suitable complement) has a nilpotent restriction
of strictly smaller dimension, and the induction hypothesis applies.

This procedure produces a basis of~$E$ consisting of cyclic chains,
and in this basis the matrix of~$g$ is block-diagonal with nilpotent
Jordan blocks.

\textbf{Uniqueness.}  The sizes of the Jordan blocks are determined
by the sequence $d_j = \dim\Ker(g^j)$: the number of blocks of
size~$\geq j$ is $d_j - d_{j-1}$, and the number of blocks of size
exactly~$j$ is $(d_j - d_{j-1}) - (d_{j+1} - d_j)$.  Since the
dimensions $d_j$ are invariants of~$g$, the block sizes are uniquely
determined.
\end{proof}

\begin{proof}[Proof of 
ef{thm:thm:jordan-normal-form}]
\index{Jordan normal form!proof}

\textbf{Existence.}  By the primary decomposition theorem
(
ef{thm:thm:primary-decomposition}),
\[
  E = G_{\lambda_1}(f) \oplus G_{\lambda_2}(f) \oplus \cdots
    \oplus G_{\lambda_s}(f),
\]
where $\lambda_1, \ldots, \lambda_s$ are the distinct eigenvalues
of~$f$ and $\dim G_{\lambda_i}(f) = n_i = m_a(\lambda_i)$.

On each $G_{\lambda_i}(f)$, define
$g_i = (f - \lambda_i\Id)\restrict{G_{\lambda_i}(f)}$, which is
nilpotent.  By 
ef{lem:lem:jordan-nilpotent}, there is a basis of
$G_{\lambda_i}(f)$ in which $g_i$ has the form
$\diag(N_{k_1^{(i)}}, \ldots, N_{k_{m_i}^{(i)}})$.

Since $f\restrict{G_{\lambda_i}(f)} = \lambda_i\Id + g_i$, the
matrix of $f\restrict{G_{\lambda_i}(f)}$ in this basis is
\[
  \diag\bigl(J_{k_1^{(i)}}(\lambda_i), \ldots,
    J_{k_{m_i}^{(i)}}(\lambda_i)\bigr).
\]
Concatenating the bases of all generalized eigenspaces gives a basis
of~$E$ in which $f$ has a Jordan matrix.

\textbf{Uniqueness.}  On each generalized eigenspace
$G_{\lambda_i}(f)$, the Jordan block sizes are determined by
the dimensions $\dim\Ker\bigl((f - \lambda_i\Id)^k\bigr)$
by 
ef{lem:lem:jordan-nilpotent}.  Since the eigenvalues, their
algebraic multiplicities, and the kernel dimensions are similarity
invariants, the Jordan form is unique up to the ordering of blocks.
\end{proof}


% ============================================================
\section{Computing the Jordan form}\label{sec:computing-jordan}
% ============================================================

Given a matrix $A \in \Mat_n(\K)$, the following algorithm computes
its Jordan normal form.

\begin{enumerate}[label=\textbf{Step \arabic*:},leftmargin=3cm]
  \item \textbf{Find the eigenvalues.}
    Compute $\chi_A(\lambda) = \det(A - \lambda I)$ and factor it:
    $\chi_A(\lambda) = \prod_{i=1}^s (\lambda_i - \lambda)^{n_i}$.

  \item \textbf{For each eigenvalue~$\lambda_i$, compute the kernel
    dimensions.}
    Compute $d_k^{(i)} = \dim\Ker\bigl((A - \lambda_i I)^k\bigr)$
    for $k = 1, 2, 3, \ldots$ until $d_k^{(i)} = n_i$ (which must
    happen for some $k \leq n_i$).

  \item \textbf{Determine the Jordan block sizes.}
    For each eigenvalue~$\lambda_i$:
    \begin{itemize}
      \item The number of Jordan blocks for $\lambda_i$ is
        $d_1^{(i)} = \dim\Ker(A - \lambda_i I)$ (the geometric
        multiplicity).
      \item The number of blocks of size $\geq k$ is
        $d_k^{(i)} - d_{k-1}^{(i)}$.
      \item The number of blocks of size exactly~$k$ is
        $(d_k^{(i)} - d_{k-1}^{(i)}) - (d_{k+1}^{(i)} - d_k^{(i)})
        = 2d_k^{(i)} - d_{k-1}^{(i)} - d_{k+1}^{(i)}$.
    \end{itemize}

  \item \textbf{Assemble the Jordan matrix.}
    Write $J = \diag\bigl(J_{k_1}(\lambda_1), J_{k_2}(\lambda_2),
    \ldots\bigr)$.
\end{enumerate}

\begin{remark}{Finding the transition matrix}{rem:jordan-transition}
\index{Jordan normal form!transition matrix}
To find the invertible matrix~$P$ such that $A = P J \inv{P}$, one
must compute a \emph{Jordan basis}: for each Jordan block of
size~$k$ associated with eigenvalue~$\lambda$, find a chain of
generalized eigenvectors $v_1, v_2, \ldots, v_k$ satisfying
\[
  (A - \lambda I)v_1 = 0, \quad
  (A - \lambda I)v_2 = v_1, \quad
  \ldots, \quad
  (A - \lambda I)v_k = v_{k-1}.
\]
The columns of~$P$ are then these generalized eigenvectors
$(v_k, v_{k-1}, \ldots, v_1)$ (note the reversed order: $v_k$ comes
first so that $AP = PJ$).
\end{remark}


% ============================================================
\section{Worked examples}\label{sec:jordan-examples}
% ============================================================

\begin{example}{Jordan form of a $3 \times 3$ matrix}{ex:jordan-3x3}
\index{Jordan normal form!example, $3\times 3$}
Let
\[
  A = \begin{pmatrix} 2 & 1 & 0 \\ 0 & 2 & 1 \\ 0 & 0 & 2
  \end{pmatrix}.
\]

\textbf{Step 1.}  $\chi_A(\lambda) = (2 - \lambda)^3$, so $\lambda = 2$
is the only eigenvalue with $n_1 = 3$.

\textbf{Step 2.}  Compute the kernel dimensions.
\[
  A - 2I = \begin{pmatrix} 0 & 1 & 0 \\ 0 & 0 & 1 \\ 0 & 0 & 0
  \end{pmatrix}, \quad
  (A - 2I)^2 = \begin{pmatrix} 0 & 0 & 1 \\ 0 & 0 & 0 \\ 0 & 0 & 0
  \end{pmatrix}, \quad
  (A - 2I)^3 = 0.
\]
So $d_1 = \dim\Ker(A - 2I) = 1$, $d_2 = \dim\Ker(A - 2I)^2 = 2$,
$d_3 = 3$.

\textbf{Step 3.}  Number of blocks of size~$\geq 1$: $d_1 - d_0 = 1$.
Number of blocks of size~$\geq 2$: $d_2 - d_1 = 1$.
Number of blocks of size~$\geq 3$: $d_3 - d_2 = 1$.
So there is exactly one block of size~$3$.

\textbf{Step 4.}  The Jordan form is
\[
  J = J_3(2) = \begin{pmatrix} 2 & 1 & 0 \\ 0 & 2 & 1 \\ 0 & 0 & 2
  \end{pmatrix}.
\]
Indeed, $A$ is already in Jordan form.
\end{example}

\begin{example}{Jordan form of a $4 \times 4$ matrix}{ex:jordan-4x4}
\index{Jordan normal form!example, $4\times 4$}
Let
\[
  A = \begin{pmatrix}
    3 & 1 & 0 & 0 \\
    0 & 3 & 0 & 0 \\
    0 & 0 & 3 & 1 \\
    0 & 0 & 0 & 5
  \end{pmatrix}.
\]

\textbf{Step 1.}  $\chi_A(\lambda) = (3 - \lambda)^3(5 - \lambda)$.
Eigenvalues: $\lambda_1 = 3$ with $n_1 = 3$, and $\lambda_2 = 5$
with $n_2 = 1$.

\textbf{Step 2.}  For $\lambda_1 = 3$:
\[
  A - 3I = \begin{pmatrix}
    0 & 1 & 0 & 0 \\
    0 & 0 & 0 & 0 \\
    0 & 0 & 0 & 1 \\
    0 & 0 & 0 & 2
  \end{pmatrix}, \quad
  (A - 3I)^2 = \begin{pmatrix}
    0 & 0 & 0 & 0 \\
    0 & 0 & 0 & 0 \\
    0 & 0 & 0 & 2 \\
    0 & 0 & 0 & 4
  \end{pmatrix}.
\]
$d_1 = \dim\Ker(A - 3I) = 2$ (the kernel is Spanned by
$e_1$ and $e_3$), $d_2 = \dim\Ker(A-3I)^2 = 2$.

Wait---we need $d_k = \dim\Ker(A-3I)^k$ restricted to the generalized
eigenspace, or equivalently the rank of $(A - 3I)^k$.
Since $d_2 = 2$ and we need $d_k = 3 = n_1$, let us recompute.

Actually, $\rank(A - 3I) = 2$ (two independent rows), so
$d_1 = 4 - 2 = 2$.  For $(A - 3I)^2$: $\rank = 1$, so
$d_2 = 4 - 1 = 3$.

Number of blocks for $\lambda_1 = 3$: $d_1 = 2$ blocks total.
Blocks of size $\geq 2$: $d_2 - d_1 = 1$.
Blocks of size $\geq 3$: $d_3 - d_2 = 0$ (since $d_3 = 3 = d_2$
for the generalized eigenspace dimension).
So: one block of size~$2$ (from $1$ of size $\geq 2$, minus $0$ of
size $\geq 3$), and one block of size~$1$ (from $2$ total, minus $1$
of size~$\geq 2$).

For $\lambda_2 = 5$: one block of size~$1$ (the algebraic multiplicity
is~$1$).

\textbf{Step 3.}  The Jordan form is
\[
  J = \begin{pmatrix}
    3 & 1 & 0 & 0 \\
    0 & 3 & 0 & 0 \\
    0 & 0 & 3 & 0 \\
    0 & 0 & 0 & 5
  \end{pmatrix}
  = \diag\bigl(J_2(3),\; J_1(3),\; J_1(5)\bigr).
\]

\textbf{Jordan basis.}  For the block $J_2(3)$: we need $v_1 \in
\Ker(A - 3I)$ and $v_2$ with $(A - 3I)v_2 = v_1$.  Take
$v_1 = e_1$, then $(A - 3I)v_2 = e_1$ gives $v_2 = e_2$ (since
$(A-3I)e_2 = e_1$).  For $J_1(3)$: take $v_3 = e_3 \in \Ker(A-3I)$.
For $J_1(5)$: solve $(A - 5I)v_4 = 0$; we find
$v_4 = \begin{psmallmatrix} -1\\0\\-1\\2 \end{psmallmatrix}$
(after row reduction).

The transition matrix is $P = (v_2 \mid v_1 \mid v_3 \mid v_4)$:
\[
  P = \begin{pmatrix}
    0 & 1 & 0 & -1 \\
    1 & 0 & 0 & 0 \\
    0 & 0 & 1 & -1 \\
    0 & 0 & 0 & 2
  \end{pmatrix}, \qquad
  A = P\,J\,\inv{P}.
\]
One may verify that $AP = PJ$.
\end{example}

\begin{example}{A non-trivial $4 \times 4$ Jordan form}{ex:jordan-4x4-nontrivial}
Consider
\[
  A = \begin{pmatrix}
    5 & 4 & 2 & 1 \\
    0 & 1 & -1 & -1 \\
    -1 & -1 & 3 & 0 \\
    1 & 1 & -1 & 2
  \end{pmatrix}.
\]
One computes $\chi_A(\lambda) = (\lambda - 1)^2(\lambda - 3)^2$.

For $\lambda_1 = 1$: $\rank(A - I) = 2$, so $d_1 = 4 - 2 = 2 = n_1$.
Since $d_1 = n_1$, the generalized eigenspace equals the eigenspace.
Thus there are two blocks of size~$1$: $J_1(1) \oplus J_1(1)$.

For $\lambda_2 = 3$: $\rank(A - 3I) = 3$, so
$d_1 = 4 - 3 = 1$.  Since $d_1 = 1 < 2 = n_2$, we need to go further:
$\rank(A - 3I)^2 = 2$, so $d_2 = 4 - 2 = 2 = n_2$.
Number of blocks: $d_1 = 1$.  Size $\geq 2$: $d_2 - d_1 = 1$.
So there is one block of size~$2$: $J_2(3)$.

The Jordan form is:
\[
  J = \diag\bigl(J_1(1),\; J_1(1),\; J_2(3)\bigr)
  = \begin{pmatrix}
    1 & 0 & 0 & 0 \\
    0 & 1 & 0 & 0 \\
    0 & 0 & 3 & 1 \\
    0 & 0 & 0 & 3
  \end{pmatrix}.
\]
\end{example}


% ============================================================
\section{Jordan form and the minimal polynomial}%
\label{sec:jordan-minpoly}
% ============================================================

The Jordan form makes the relationship between a matrix and its minimal
polynomial completely transparent.

\begin{theorem}{Minimal polynomial from the Jordan form}%
{thm:jordan-minpoly}
\index{minimal polynomial!and Jordan form}
Let $A \in \Mat_n(\K)$ with Jordan form
$J = \diag\bigl(J_{k_1}(\lambda_1), \ldots, J_{k_m}(\lambda_m)\bigr)$,
where $\lambda_1, \ldots, \lambda_s$ are the distinct eigenvalues.
For each distinct eigenvalue~$\lambda_i$, let $r_i$ denote the
\emph{largest} Jordan block size associated with~$\lambda_i$.  Then
the minimal polynomial of~$A$ is
\[
  \boxed{\mu_A(\lambda) = \prod_{i=1}^{s}
    (\lambda - \lambda_i)^{r_i}.}
\]
\end{theorem}

\begin{proof}
A polynomial $P(\lambda)$ satisfies $P(A) = 0$ if and only if
$P(J) = 0$ (since $A$ and $J$ are similar).  Moreover,
$P(J) = \diag\bigl(P(J_{k_1}(\lambda_1)), \ldots,
P(J_{k_m}(\lambda_m))\bigr)$, so $P(J) = 0$ if and only if
$P(J_{k_j}(\lambda_j)) = 0$ for every block.

For a single Jordan block $J_k(\lambda)$, one computes
\[
  P(J_k(\lambda)) = \begin{pmatrix}
    P(\lambda) & P'(\lambda) & \frac{P''(\lambda)}{2!} & \cdots
      & \frac{P^{(k-1)}(\lambda)}{(k-1)!} \\
    & P(\lambda) & P'(\lambda) & \cdots
      & \frac{P^{(k-2)}(\lambda)}{(k-2)!} \\
    & & \ddots & \ddots & \vdots \\
    & & & P(\lambda) & P'(\lambda) \\
    & & & & P(\lambda)
  \end{pmatrix}.
\]
Hence $P(J_k(\lambda)) = 0$ if and only if
$P(\lambda) = P'(\lambda) = \cdots = P^{(k-1)}(\lambda) = 0$,
which holds if and only if $(\lambda_0 - \lambda)^k \mid P(\lambda)$
where $\lambda_0 = \lambda$ is the eigenvalue.

Therefore, $P(J) = 0$ if and only if $(\lambda_i - x)^{r_i}$ divides
$P(x)$ for every distinct eigenvalue~$\lambda_i$, where $r_i$ is the
largest block for~$\lambda_i$.  The monic polynomial of smallest degree
with this property is $\mu_A(\lambda) = \prod_{i=1}^s
(\lambda - \lambda_i)^{r_i}$.
\end{proof}

\begin{corollary}{Diagonalizability via the Jordan form}%
{cor:diag-jordan}
\index{diagonalizable!and Jordan form}
A matrix is diagonalizable if and only if its Jordan form is diagonal,
which happens if and only if all Jordan blocks have size~$1$.
Equivalently, the minimal polynomial has no repeated roots.
\end{corollary}

\begin{proof}
All blocks have size~$1$ if and only if $r_i = 1$ for each
eigenvalue, if and only if $\mu_A$ has only simple roots, if and only
if $A$ is diagonalizable.
\end{proof}


% ============================================================
\section{Cayley--Hamilton revisited}\label{sec:cayley-hamilton-jordan}
% ============================================================

The Jordan form gives an elegant proof of the Cayley--Hamilton theorem.

\begin{theorem}{Cayley--Hamilton theorem}{thm:cayley-hamilton-jordan}
\index{Cayley--Hamilton theorem!proof via Jordan form}
Let $A \in \Mat_n(\K)$ with $\K$ algebraically closed.  Then
$\chi_A(A) = 0$.
\end{theorem}

\begin{proof}
Let $J = P^{-1}AP$ be the Jordan form.  Then
$\chi_A(A) = P\,\chi_A(J)\,P^{-1} = P\,\chi_J(J)\,P^{-1}$
(since $\chi_A = \chi_J$).

Write $J = \diag(J_{k_1}(\lambda_1), \ldots, J_{k_m}(\lambda_m))$.
Then $\chi_J(\lambda) = \prod_{j=1}^m (\lambda_j - \lambda)^{k_j}$,
so
\[
  \chi_J(J) = \diag\bigl(\chi_J(J_{k_1}(\lambda_1)), \ldots,
    \chi_J(J_{k_m}(\lambda_m))\bigr).
\]
For each block $J_{k_j}(\lambda_j)$, the factor
$(\lambda_j - \lambda)^{k_j}$ divides $\chi_J(\lambda)$, hence
$\chi_J(J_{k_j}(\lambda_j)) = 0$ by the analysis in

ef{thm:thm:jordan-minpoly}.

Therefore $\chi_J(J) = 0$, and so $\chi_A(A) = 0$.
\end{proof}


% ============================================================
\section{Matrix exponential}\label{sec:matrix-exponential}
% ============================================================

One of the most important applications of the Jordan form is the
computation of the \emph{matrix exponential}, which plays a central
role in the theory of linear differential equations.

\begin{definition}{Matrix exponential}{def:matrix-exponential}
\index{matrix exponential}\index{exponential!matrix}
For $A \in \Mat_n(\C)$, the \emph{matrix exponential} of~$A$ is
defined by the power series
\[
  e^A \deff \exp(A) \deff \sum_{k=0}^{\infty} \frac{A^k}{k!}
  = I + A + \frac{A^2}{2!} + \frac{A^3}{3!} + \cdots.
\]
This series converges absolutely for every $A \in \Mat_n(\C)$.
\end{definition}

\begin{proposition}{Properties of the matrix exponential}%
{prop:matrix-exp-properties}
\index{matrix exponential!properties}
\begin{enumerate}[label=(\roman*)]
  \item $e^{0} = I_n$.
  \item If $AB = BA$, then $e^{A+B} = e^A e^B$.
  \item $e^A$ is always invertible, with $\inv{(e^A)} = e^{-A}$.
  \item $e^{PAP^{-1}} = P\,e^A\,P^{-1}$ for any invertible~$P$.
  \item $\det(e^A) = e^{\tr(A)}$.
\end{enumerate}
\end{proposition}

\begin{proof}
(i) is clear.  (ii) follows by expanding the product of power series,
using $AB = BA$ to collect terms as in the scalar case.  (iii) follows
from (ii) with $B = -A$.  (iv) follows from
$(PAP^{-1})^k = PA^kP^{-1}$.

For~(v), if $A = PJP^{-1}$ is in Jordan form, then $e^A = Pe^JP^{-1}$
and $\det(e^A) = \det(e^J)$.  For a Jordan block $J_k(\lambda)$,
the matrix $e^{J_k(\lambda)}$ is upper triangular with $e^\lambda$ on
the diagonal, so $\det(e^{J_k(\lambda)}) = e^{k\lambda}$.
Thus $\det(e^J) = \prod_j e^{k_j\lambda_j} = e^{\sum_j k_j\lambda_j}
= e^{\tr(J)} = e^{\tr(A)}$.
\end{proof}

\begin{proposition}{Exponential of a Jordan block}%
{prop:exp-jordan-block}
\index{matrix exponential!of Jordan block}
For a Jordan block $J_k(\lambda) = \lambda I_k + N_k$, since
$\lambda I_k$ and $N_k$ commute:
\[
  e^{J_k(\lambda)} = e^{\lambda I_k}\,e^{N_k}
  = e^{\lambda}\,e^{N_k}.
\]
Moreover, since $N_k^k = 0$:
\[
  e^{N_k} = I_k + N_k + \frac{N_k^2}{2!} + \cdots
    + \frac{N_k^{k-1}}{(k-1)!}.
\]
Explicitly:
\[
  e^{J_k(\lambda)} = e^{\lambda}
  \begin{pmatrix}
    1      & 1      & \frac{1}{2!}  & \cdots & \frac{1}{(k-1)!} \\
    0      & 1      & 1             & \cdots & \frac{1}{(k-2)!} \\
    \vdots &        & \ddots        & \ddots & \vdots            \\
    0      & \cdots &               & 1      & 1                 \\
    0      & \cdots &               & 0      & 1
  \end{pmatrix}.
\]
More generally, for the \emph{matrix exponential of~$tA$}
(with $t \in \R$):
\[
  e^{tJ_k(\lambda)} = e^{\lambda t}
  \begin{pmatrix}
    1      & t      & \frac{t^2}{2!}  & \cdots & \frac{t^{k-1}}{(k-1)!} \\
    0      & 1      & t               & \cdots & \frac{t^{k-2}}{(k-2)!} \\
    \vdots &        & \ddots          & \ddots & \vdots                  \\
    0      & \cdots &                 & 1      & t                       \\
    0      & \cdots &                 & 0      & 1
  \end{pmatrix}.
\]
\end{proposition}

\begin{proof}
Since $\lambda I_k$ and $N_k$ commute,
$e^{J_k(\lambda)} = e^{\lambda I_k + N_k} = e^{\lambda I_k} e^{N_k}
= e^\lambda I_k \cdot e^{N_k}$.  The entries of $e^{N_k}$ follow
from $(N_k^j)_{pq} = \delta_{p+j, q}$ (the matrix with $1$'s on the
$j$-th superdiagonal), giving
$(e^{N_k})_{pq} = \frac{1}{(q-p)!}$ for $q \geq p$ and $0$
otherwise.  The formula for $e^{tJ_k(\lambda)}$ is analogous, replacing
$N_k$ by $tN_k$.
\end{proof}

\begin{example}{Computing $e^{tA}$}{ex:matrix-exp}
\index{matrix exponential!example}
Let $A = \begin{pmatrix} 2 & 1 \\ 0 & 2 \end{pmatrix} = J_2(2)$.
Then
\[
  e^{tA} = e^{2t}
  \begin{pmatrix} 1 & t \\ 0 & 1 \end{pmatrix}
  = \begin{pmatrix} e^{2t} & te^{2t} \\ 0 & e^{2t} \end{pmatrix}.
\]
\end{example}


% ============================================================
\section{Applications}\label{sec:jordan-applications}
% ============================================================

% ------------------------------------------------------------
\subsection{Systems of linear ODEs}\label{subsec:jordan-ode}
% ------------------------------------------------------------

\index{differential equations!and Jordan form}

Consider the linear system of ordinary differential equations
\begin{equation}\label{eq:ode-system}
  \vec{x}'(t) = A\,\vec{x}(t), \qquad \vec{x}(0) = \vec{x}_0,
\end{equation}
where $A \in \Mat_n(\R)$ (or $\Mat_n(\C)$) and
$\vec{x}(t) \in \R^n$ (or $\C^n$).

The solution is
$\vec{x}(t) = e^{tA}\vec{x}_0$.
When $A$ is diagonalizable, $A = P\diag(\lambda_1,\ldots,\lambda_n)
\inv{P}$ and the solution is a superposition of pure exponentials
$e^{\lambda_i t}$.  When $A$ is \emph{not} diagonalizable, the Jordan
form tells us exactly what happens.

\begin{proposition}{Solution structure via Jordan form}%
{prop:ode-jordan}
\index{differential equations!non-diagonalizable case}
If $A = PJ\inv{P}$ with Jordan form~$J$, then
$e^{tA} = P\,e^{tJ}\,\inv{P}$ and
$\vec{x}(t) = P\,e^{tJ}\,\inv{P}\vec{x}_0$.

For a Jordan block $J_k(\lambda)$, the corresponding components of
the solution involve terms of the form
\[
  t^j e^{\lambda t}, \qquad 0 \leq j \leq k-1.
\]
In particular, non-trivial Jordan blocks (size $> 1$) produce
\emph{polynomial-exponential} solutions (polynomial times
exponential), not pure exponentials.
\end{proposition}

\begin{example}{A $3 \times 3$ system of ODEs}{ex:ode-3x3}
Consider $\vec{x}'(t) = A\vec{x}(t)$ with
\[
  A = \begin{pmatrix} 1 & 1 & 0 \\ 0 & 1 & 1 \\ 0 & 0 & 1
  \end{pmatrix} = J_3(1).
\]
Then
\[
  e^{tA} = e^{t}
  \begin{pmatrix}
    1 & t & \tfrac{t^2}{2} \\
    0 & 1 & t \\
    0 & 0 & 1
  \end{pmatrix}.
\]
With initial condition
$\vec{x}_0 = \begin{psmallmatrix} 1\\0\\1 \end{psmallmatrix}$:
\[
  \vec{x}(t) = e^{t}
  \begin{pmatrix}
    1 + \tfrac{t^2}{2} \\
    t \\
    1
  \end{pmatrix}.
\]
Note the appearance of the polynomial factor $1 + t^2/2$ in the first
component, which is characteristic of a $3\times 3$ Jordan block.
\end{example}

% ------------------------------------------------------------
\subsection{Matrix functions via Jordan form}%
\label{subsec:matrix-functions}
% ------------------------------------------------------------

\index{matrix function}

The Jordan form allows us to define and compute general matrix
functions.  If $h\colon \C \to \C$ is a function that is sufficiently
smooth (or analytic) in a neighbourhood of each eigenvalue, we can
define $h(A)$ as follows.

\begin{definition}{Matrix function via Jordan form}%
{def:matrix-function}
Let $A \in \Mat_n(\C)$ with $A = PJ\inv{P}$ and
$J = \diag(J_{k_1}(\lambda_1), \ldots, J_{k_m}(\lambda_m))$.
Define
\[
  h(A) \deff P\,\diag\bigl(h(J_{k_1}(\lambda_1)), \ldots,
    h(J_{k_m}(\lambda_m))\bigr)\,\inv{P},
\]
where
\[
  h(J_k(\lambda)) \deff
  \begin{pmatrix}
    h(\lambda) & h'(\lambda) & \frac{h''(\lambda)}{2!} & \cdots
      & \frac{h^{(k-1)}(\lambda)}{(k-1)!} \\
    & h(\lambda) & h'(\lambda) & \cdots
      & \frac{h^{(k-2)}(\lambda)}{(k-2)!} \\
    & & \ddots & \ddots & \vdots \\
    & & & h(\lambda) & h'(\lambda) \\
    & & & & h(\lambda)
  \end{pmatrix}.
\]
\end{definition}

\begin{example}{Computing $A^k$ and $\sqrt{A}$}{ex:matrix-functions}
For $A = \begin{psmallmatrix} 4 & 1 \\ 0 & 4 \end{psmallmatrix}
= J_2(4)$ and $h(\lambda) = \lambda^k$:
\[
  A^k = \begin{pmatrix} 4^k & k\cdot 4^{k-1} \\ 0 & 4^k
  \end{pmatrix}.
\]
For $h(\lambda) = \sqrt{\lambda}$ (choosing the principal square root):
\[
  \sqrt{A} = \begin{pmatrix} 2 & \tfrac{1}{4} \\ 0 & 2
  \end{pmatrix},
\]
since $h(4) = 2$ and $h'(4) = \frac{1}{2\sqrt{4}} = \frac{1}{4}$.
\end{example}


% ============================================================
\section{Exercises}\label{sec:jordan-exercises}
% ============================================================

\begin{exercise}{Nilpotent matrices \basic}{exo:nilpotent-basic}
Determine whether each matrix is nilpotent.  If so, find the index of
nilpotency.
\begin{enumerate}[label=(\alph*)]
  \item $A = \begin{pmatrix} 0 & 2 & 1 \\ 0 & 0 & -3 \\ 0 & 0 & 0
    \end{pmatrix}$
  \item $B = \begin{pmatrix} 1 & 1 \\ 0 & 1 \end{pmatrix}$
  \item $C = \begin{pmatrix} 0 & 0 & 0 \\ 1 & 0 & 0 \\ 0 & 1 & 0
    \end{pmatrix}$
\end{enumerate}
\end{exercise}

\begin{exercise}{Properties of nilpotent endomorphisms \basic}%
{exo:nilpotent-properties}
Let $f \in \End(E)$ be nilpotent.  Show that:
\begin{enumerate}[label=(\alph*)]
  \item $\Id + f$ is invertible.
    \emph{Hint:} consider $(\Id + f)(\Id - f + f^2 - \cdots)$.
  \item If $g \in \End(E)$ is nilpotent and $fg = gf$, then $f + g$
    is nilpotent.
\end{enumerate}
\end{exercise}

\begin{exercise}{Jordan form of $2 \times 2$ matrices \basic}%
{exo:jordan-2x2}
Find the Jordan normal form of each matrix over~$\C$:
\begin{enumerate}[label=(\alph*)]
  \item $A = \begin{pmatrix} 3 & 1 \\ 0 & 3 \end{pmatrix}$
  \item $B = \begin{pmatrix} 0 & -1 \\ 1 & 0 \end{pmatrix}$
  \item $C = \begin{pmatrix} 5 & 3 \\ -3 & -1 \end{pmatrix}$
\end{enumerate}
\end{exercise}

\begin{exercise}{Jordan form of a $3 \times 3$ matrix \intermediate}%
{exo:jordan-3x3}
Find the Jordan normal form and a Jordan basis for
\[
  A = \begin{pmatrix} 2 & 1 & 0 \\ 0 & 2 & 0 \\ 0 & 0 & 3
  \end{pmatrix}.
\]
\end{exercise}

\begin{exercise}{A $4 \times 4$ matrix with repeated eigenvalues
\intermediate}{exo:jordan-4x4-repeated}
\index{Jordan normal form!exercise}
Find the Jordan normal form of
\[
  A = \begin{pmatrix}
    2 & 1 & 0 & 0 \\
    0 & 2 & 1 & 0 \\
    0 & 0 & 2 & 0 \\
    0 & 0 & 0 & 2
  \end{pmatrix}.
\]
How many distinct Jordan forms are possible for a $4 \times 4$ matrix
whose only eigenvalue is~$2$?  List them all.
\end{exercise}

\begin{exercise}{Uniqueness via kernel dimensions \intermediate}%
{exo:jordan-kernel-dimensions}
Let $A, B \in \Mat_n(\C)$ be two matrices with the same characteristic
polynomial.  Show that $A$ and $B$ are similar if and only if
$\rank(A - \lambda I)^k = \rank(B - \lambda I)^k$ for every
eigenvalue~$\lambda$ and every $k \geq 1$.
\end{exercise}

\begin{exercise}{Minimal polynomial from Jordan form \intermediate}%
{exo:minpoly-jordan}
For each of the following Jordan matrices, determine the minimal
polynomial:
\begin{enumerate}[label=(\alph*)]
  \item $J = \diag(J_3(2), J_1(2))$
  \item $J = \diag(J_2(1), J_2(3))$
  \item $J = \diag(J_2(0), J_2(0), J_1(0))$
\end{enumerate}
\end{exercise}

\begin{exercise}{Matrix exponential \intermediate}{exo:matrix-exp}
\index{matrix exponential!exercise}
Compute $e^{tA}$ for:
\begin{enumerate}[label=(\alph*)]
  \item $A = \begin{pmatrix} 0 & 1 \\ 0 & 0 \end{pmatrix}$
  \item $A = \begin{pmatrix} 1 & 1 & 0 \\ 0 & 1 & 1 \\ 0 & 0 & 1
    \end{pmatrix}$
  \item $A = \begin{pmatrix} 2 & 1 & 0 & 0 \\ 0 & 2 & 0 & 0 \\
    0 & 0 & 3 & 1 \\ 0 & 0 & 0 & 3 \end{pmatrix}$
\end{enumerate}
\end{exercise}

\begin{exercise}{Solving a non-diagonalizable ODE system
\intermediate}{exo:ode-jordan}
\index{differential equations!exercise}
Solve the initial value problem
\[
  \vec{x}'(t) = \begin{pmatrix} 3 & 1 \\ 0 & 3 \end{pmatrix}
  \vec{x}(t), \qquad
  \vec{x}(0) = \begin{pmatrix} 1 \\ 2 \end{pmatrix}.
\]
\end{exercise}

\begin{exercise}{Cayley--Hamilton via Jordan form \challenging}%
{exo:cayley-hamilton-jordan}
Let $A \in \Mat_n(\C)$ with Jordan form $J$.
\begin{enumerate}[label=(\alph*)]
  \item Compute $\chi_A(J_{k}(\lambda))$ directly, showing each
    entry is zero.
  \item Deduce the Cayley--Hamilton theorem for matrices over~$\C$.
  \item Explain why the result extends to matrices over any field~$\K$.
    \emph{Hint:} use the fact that $\chi_A(A) = 0$ is a polynomial
    identity in the entries of~$A$.
\end{enumerate}
\end{exercise}

\begin{exercise}{Powers of nilpotent-plus-scalar matrices
\challenging}{exo:powers-jordan}
Let $A = \lambda I + N$ where $N$ is nilpotent with index~$r$.
\begin{enumerate}[label=(\alph*)]
  \item Show that for all $m \geq 0$,
    \[
      A^m = \sum_{j=0}^{r-1} \binom{m}{j} \lambda^{m-j} N^j.
    \]
  \item Deduce the entries of $J_k(\lambda)^m$ explicitly.
  \item Use part~(a) to give another proof that if $A$ is nilpotent
    and $\tr(A^k) = 0$ for all $k \geq 1$, then $A$ is nilpotent.
\end{enumerate}
\end{exercise}

\begin{exercise}{Jordan form and commuting matrices \challenging}%
{exo:commuting-jordan}
\index{commuting matrices}
Let $A \in \Mat_n(\C)$.
\begin{enumerate}[label=(\alph*)]
  \item Show that $B$ commutes with $J_k(\lambda)$ if and only if $B$
    is an upper triangular Toeplitz matrix (i.e.\ $B_{ij}$ depends only
    on $j - i$).
  \item Deduce that if $A$ has $n$ distinct eigenvalues, then every
    matrix commuting with~$A$ is a polynomial in~$A$.
  \item Give an example showing this fails when $A$ has repeated
    eigenvalues with multiple Jordan blocks.
\end{enumerate}
\end{exercise}

\begin{exercise}{Logarithm of a matrix \challenging}{exo:matrix-log}
\index{matrix logarithm}
Let $A \in \Mat_n(\C)$ be invertible (so all eigenvalues are nonzero).
\begin{enumerate}[label=(\alph*)]
  \item For a Jordan block $J_k(\lambda)$ with $\lambda \neq 0$,
    show that $J_k(\lambda) = \lambda(I + \lambda^{-1}N_k)$ and use
    the power series $\log(1 + x) = x - \frac{x^2}{2} +
    \frac{x^3}{3} - \cdots$ (which terminates since $N_k$ is
    nilpotent) to define $\log J_k(\lambda)$.
  \item Compute $\log J_k(\lambda)$ explicitly for $k = 2$ and $k = 3$.
  \item Verify that $e^{\log J_k(\lambda)} = J_k(\lambda)$.
\end{enumerate}
\end{exercise}

\begin{exercise}{Jordan form over $\R$ \challenging}{exo:jordan-real}
\index{Jordan normal form!over $\R$}
The matrix
$A = \begin{pmatrix} 0 & -1 & 1 & 0 \\ 1 & 0 & 0 & 1 \\
0 & 0 & 0 & -1 \\ 0 & 0 & 1 & 0 \end{pmatrix}$
has characteristic polynomial $\chi_A(\lambda) = (\lambda^2+1)^2$.
\begin{enumerate}[label=(\alph*)]
  \item Find the Jordan form of~$A$ over~$\C$.
  \item Show that $A$ is not similar over~$\R$ to a Jordan matrix.
  \item Find the \emph{real Jordan form} of~$A$, i.e.\ a block
    matrix with $2\times 2$ ``rotation-scaling'' blocks
    $\begin{psmallmatrix} a & -b \\ b & a \end{psmallmatrix}$
    replacing complex eigenvalues $a \pm bi$.
\end{enumerate}
\end{exercise}


% ============================================================
\section{Chapter summary}\label{sec:jordan-summary}
% ============================================================

\begin{itemize}
  \item \textbf{Nilpotent endomorphisms.}\index{nilpotent!summary}
    An endomorphism~$f$ is nilpotent if $f^r = 0$ for some $r \geq 1$.
    Its only eigenvalue is~$0$, $\chi_f(\lambda) = (-\lambda)^n$, and
    the index of nilpotency satisfies $r \leq n$.

  \item \textbf{Generalized eigenspaces.}\index{generalized eigenspace!summary}
    For each eigenvalue~$\lambda$ of~$f$, the generalized eigenspace
    $G_\lambda(f) = \Ker\bigl((f - \lambda\Id)^{n_\lambda}\bigr)$ has
    dimension equal to the algebraic multiplicity~$n_\lambda$.  The
    primary decomposition theorem gives
    $E = \bigoplus_{\lambda \in \Spec(f)} G_\lambda(f)$.

  \item \textbf{Jordan blocks.}\index{Jordan block!summary}
    $J_k(\lambda) = \lambda I_k + N_k$ is the basic building block.
    It has eigenvalue~$\lambda$ with algebraic multiplicity~$k$ and
    geometric multiplicity~$1$.

  \item \textbf{Jordan Normal Form.}\index{Jordan normal form!summary}
    Over an algebraically closed field, every endomorphism has a
    unique (up to block ordering) Jordan form.  The block sizes are
    determined by the dimensions $\dim\Ker(f - \lambda\Id)^k$.

  \item \textbf{Minimal polynomial.}
    $\mu_f(\lambda) = \prod_{i} (\lambda - \lambda_i)^{r_i}$,
    where $r_i$ is the size of the \emph{largest} Jordan block for
    eigenvalue~$\lambda_i$.  A matrix is diagonalizable iff
    $\mu_f$ has only simple roots.

  \item \textbf{Matrix exponential.}\index{matrix exponential!summary}
    $e^{tJ_k(\lambda)} = e^{\lambda t} \cdot (\text{upper triangular
    polynomial matrix in~$t$})$.  This gives the explicit solution
    $\vec{x}(t) = e^{tA}\vec{x}_0$ to $\vec{x}' = A\vec{x}$, even
    when $A$ is not diagonalizable.  Non-trivial Jordan blocks produce
    solutions of the form $t^j e^{\lambda t}$ (polynomial times
    exponential).

  \item \textbf{Matrix functions.}
    Any sufficiently smooth function~$h$ can be extended to matrices
    via the Jordan form: $h(J_k(\lambda))$ involves
    $h(\lambda), h'(\lambda), \ldots, h^{(k-1)}(\lambda)$.
\end{itemize}


% ============================================================
% Back matter
% ============================================================
\backmatter

\printbibliography[heading=bibintoc,title={Bibliography}]

\printindex

\end{document}
